\documentclass[11pt]{article}
\usepackage[margin=1.1in]{geometry}
\usepackage{amsmath,amssymb,amsthm}
\usepackage{enumitem}
\usepackage[T1]{fontenc}
\usepackage{hyperref}
\hypersetup{colorlinks=true,linkcolor=black,citecolor=black,urlcolor=blue}

\newtheorem{theorem}{Theorem}[section]
\newtheorem{lemma}[theorem]{Lemma}
\newtheorem{proposition}[theorem]{Proposition}
\newtheorem{corollary}[theorem]{Corollary}
\newtheorem{conjecture}[theorem]{Conjecture}
\newtheorem{problem}[theorem]{Problem}
\theoremstyle{definition}
\newtheorem{definition}[theorem]{Definition}
\theoremstyle{remark}
\newtheorem{remark}[theorem]{Remark}
\newtheorem*{measured}{Measured}

\newcommand{\Tr}{\operatorname{Tr}}
\newcommand{\step}{\operatorname{step}}
\newcommand{\ALL}{\mathrm{ALL}}
\newcommand{\N}{\mathbb{N}}
\newcommand{\Z}{\mathbb{Z}}
\newcommand{\F}{\mathbb{F}_2}
\newcommand{\Q}{\mathbb{Q}}
\newcommand{\PROVED}{\textsc{Proved}}
\newcommand{\EXACT}{\textsc{Exact}}
\newcommand{\MEASURED}{\textsc{Measured}}
\newcommand{\KNOWN}{\textsc{Known}}
\newcommand{\OPEN}{\textsc{Open}}

\title{The centre column of Rule~30:\\ excluded periods, decimation of the Duhamel class,\\ and the front of the diagonal defects}
\author{N.~Berzai\thanks{Computational and editorial assistance from Claude (Anthropic); three parallel referee dialogues (Opus~5) contributed Part~III. Assembled 2026-09-16 from the working documents listed in Section~\ref{sec:sources}; revised 2026-09-24; citations of other public work added in October 2026. No claim is made beyond its label.}}
\date{}

\begin{document}
\maketitle

\begin{abstract}
Let $c=(c_t)$ be the centre column of the elementary cellular automaton Rule~30 started from a single black cell. Wolfram's first Prize Problem asks whether $c$ is eventually periodic; the second asks whether its density is $\tfrac12$; the third asks whether computing $c_t$ requires at least $O(t)$ computational effort. Wolfram names an algorithm below the direct $O(t^{2})$ as a partial result; one running in $O(t^{2}/\log t)$ time is already in print, and it is not Problem~3 itself. We collect what this programme established about $c$, each statement labelled \PROVED, \EXACT\ (a finite computation with no fitting), \MEASURED, \KNOWN\ (in the literature) or \OPEN.
Part~I proves that infinitely many period words are impossible as eventual periods of $c$, for every number-like configuration: every necklace all of whose cyclic zero-runs have length at least $19$ (with $1$-runs arbitrary; at least $17$ for every necklace tested), every word with isolated ones and zero-runs of length at least $10$, the words $0^m11$ ($m=6$, $m\ge8$), $1^b0^a$ ($b\ge9$), $1^bv$ ($b\ge29$, $|v|\le12$), and both constants; the method is a finite-state over-approximation of the right half whose reach is characterised exactly. At the boundary of that reach two facts are proved --- a finite block that traps the interface, and a cap of six cells on straight influence into it --- and what is left open there is whether the interface language over $01$ contains an interface that is not eventually periodic: uncountability of that language would suffice, is not necessary, and no enumeration can reach it (positive entropy of that language is \MEASURED\ on random right halves; its uncountability is \OPEN). Some of Part~I is public elsewhere, and no priority is claimed for it. An infinite family of excluded periods, $01^{q}$ for $q=7$ and $q\ge9$ (our $1^b0^a$ with $a=1$), was made public on 10--12 September 2026 in the repository \texttt{cochon123/rule30-prize}; we found it independently and extend it to the other families listed. That neither constant is an eventual period of $c$ is \KNOWN, and the method in kind is in other public work. The sources are in Remark~\ref{ex:prior}.
Part~II writes $c$ as a linear functional of the quadratic source of Rule~30 through the Rule~150 propagator, proves that this class of functionals is closed under the Cartier operators (so the $2$-kernel of $c$ is the image of an explicit orbit), gives the functional in closed form in the variable $\alpha$ with $y=\alpha/(1+\alpha+\alpha^2)$, and shows that only the centre line and the chiral part of the source are visible to $c$; certified computations give no Mahler relation of depth $\le12$ and degree $\le8000$, and a $2$-kernel census that is injective wherever it has been compared, giving $|\mathcal K(c)|\ge138\,672$ pairwise distinct truncations from $2^{21}$ terms; controls show, however, that such a census measures counting room rather than aperiodicity, and the family it can certify from $2^{\nu}$ terms is capped at $2^{\nu-4}-1$ elements (about $n/16$), which closes the census route.
Part~III splits each left diagonal into a periodic background and a finitely supported defect, proves the reset formula, the exact degree formula, monotonicity of the defect front, an exact failure condition, and Theorem~H$^*$, which reduces the conjecture that the periodic texture never crosses the centre to two inequalities on the front, on cancellation depths and on stall runs, whose margins are measured and wide; it also re-derives, in the diagonal frame, the unboundedness of the diagonal periods; that result is \KNOWN, being Gravner--Griffeath's Proposition~3.3 (a left-permutative edge cellular automaton has no periodic solution), stated there for Rule~30, and the argument given here is theirs in different coordinates. Other items of Part~III are not ours either: the split into background and defect and the monotone front are in the public repository \texttt{Dibujaron/rule30} (7--8 September 2026), and the code of Wolfram's 2019 prize post already takes a running maximum; the conjecture itself is a question of that post; and the one-step law behind the reset formula is Rowland's (2006). The citations are in Section~\ref{sec:verdict}.
Part~III also treats the existence of the front limit $\lim M_d/d$, which no theorem in the literature supplies for this edge, and reduces it twice: the near-diagonal inequality it needs follows if $M$ stays within $O(N^{\alpha})$ of its least concave majorant for any $\alpha<1$ --- strictly weaker than the constant we had been pursuing, since $\sum n^{\alpha-2}$ converges throughout that range --- and, separately, $M_N/N$ equals exactly $\sum_{k\ge1}q_k(N)$ at every $N$, where $q_k(N)$ is the frequency of jumps of size at least $k$ below $N$, so the existence of the limit reduces to the convergence of one increment's distribution with a summable dominant.
Part~IV records the collapse pyramids over the column, with an exact identity for the traffic rules and a two-sided criterion, together with measurements showing every bounded-window predictor of $c$ tested ($w\le22$) at chance, and a remark on why that statement must stay \MEASURED.
Part~V asks what is provable about Problem~3. The window transformation monoid of Rule~30 contains $A_5$ and is therefore not solvable, placing the rule on the hard side of the Barrington--Th\'erien and Moore dichotomies. A rank lower bound on the prediction matrix does not reach the column. The rank is exactly $2$ at every $n\le10$ for Rules~$90$, $150$ and~$105$ and about $1.59^{n}$ for Rule~30 on that range (\MEASURED), but Rules~$22$ and~$54$, whose centre columns are respectively eventually $0$ and periodic with period~$4$, have rank over $\F$ of the same order as Rule~30's or larger at every $n$ computed (\EXACT): the rank measures the prediction function over all inputs, not the single-seed column, and that route to Problem~3 is closed. A second control inside our own family --- Rule~$126$, whose centre column agrees for every $t\le2^{14}$ with the Fredholm--Rueppel sequence, which has the perfect linear-complexity profile and whose $t$-th bit is computable in $O(\log t)$ --- bounds what any complexity measurement can establish.
No route to any of the three Prize Problems is claimed. Part~VI states exactly where each route stops, and a separate section states why no measure-level argument can reach Problem~2.
\end{abstract}

\tableofcontents

\section{Conventions and labels}\label{sec:conventions}

Rule~30 acts on configurations $x\in\{0,1\}^{\Z}$ by
\begin{equation}\label{eq:rule30}
a(n+1,i)=a(n,i-1)\oplus\bigl(a(n,i)\vee a(n,i+1)\bigr),\qquad a(0,\cdot)=x,
\end{equation}
equivalently, over $\F$, $u^{t+1}_x=u^t_{x-1}+u^t_x+u^t_{x+1}+u^t_xu^t_{x+1}$. Part~I uses the notation $a(n,i)$ of the exclusion paper; Parts~II--IV use $u^t_x$, with $c_t=u^t_0$ the centre column and $u^0_x=[x=0]$ the single seed. A sequence is \emph{eventually periodic} if it agrees with a periodic sequence from some index on.

Every statement carries one of five labels. \PROVED: a written argument, machine-checked where finite, with the certificate identified. \EXACT: an exact finite computation, reproducible from the script named. \MEASURED: a statistic, a fit, or an out-of-sample test; never used as a hypothesis of a theorem. \KNOWN: in the literature or in other public work, cited; where that work is of 2026 no priority over it is claimed. \OPEN: not settled here. Nothing labelled \MEASURED\ is used in a proof.

The three Prize Problems are: \textbf{P1}, $c$ is not eventually periodic; \textbf{P2}, the density of ones in $c$ tends to $\tfrac12$; \textbf{P3}, computing $c_t$ requires at least $O(t)$ computational effort (Wolfram's wording, with $t$ for his $n$; Section~\ref{sec:p3}). None is settled here. An algorithm below the direct $O(t^{2})$ is the partial result Wolfram asks for, not P3 itself, and it is \KNOWN: $c_t$ is computable in $O(t^{2}/\log t)$ time \cite{NA}.

\section{Two tiers of equivalence, and which one is P1}\label{sec:tiers}

Several distinct conditions on $c$ are easy to conflate, and the routes in this paper divide according
to which tier they aim at, so we set them out once. Write $F(x)=\sum_n c_nx^n\in\F[[x]]$.

\begin{proposition}[\KNOWN]\label{pr:tiers}
The following are equivalent, and each is the negation of \textup{P1}:
\begin{enumerate}[nosep,label=\textup{(\roman*)}]
\item $c$ is eventually periodic;
\item $F$ is \emph{rational} over $\F(x)$;
\item $c$ satisfies a finite linear recurrence with constant coefficients;
\item the subword complexity $p(n)$ is \emph{bounded};
\item the linear complexity of $c$ is finite.
\end{enumerate}
The following are equivalent to each other, and are strictly weaker:
\begin{enumerate}[nosep,label=\textup{(\alph*)}]
\item $c$ is $2$-automatic;
\item $F$ is \emph{algebraic} over $\F(x)$;
\item the $2$-kernel $\mathcal K(c)$ is finite.
\end{enumerate}
The second tier implies $p(n)=O(n)$ and does not imply the first.
\end{proposition}
\begin{proof}
In the first tier, (i)$\Leftrightarrow$(ii)$\Leftrightarrow$(iii) is the standard correspondence between
rational series and constant-coefficient recurrences, together with the observation that over a finite
field a recurrence of order $d$ has a finite state space and so forces a cycle; (i)$\Leftrightarrow$(iv)
is Morse--Hedlund; (v) restates (iii). In the second tier (a)$\Leftrightarrow$(b) is Christol's theorem
\cite{CKMFR} and (a)$\Leftrightarrow$(c) is the Eilenberg characterisation \cite{AS}. The implication to
$p(n)=O(n)$ is Cobham's. Both non-implications are witnessed: Thue--Morse is $2$-automatic and not
eventually periodic, so the second tier does not imply the first; and the Fibonacci word has
$p(n)=n+1$, linear yet not automatic, so linear complexity does not imply the second tier either.
\end{proof}

Two consequences worth keeping in view. First, the weakest hypothesis that settles \textup{P1} is that
$p(n)$ be \emph{unbounded} --- superlinear growth is more than is needed, and neither Christol's theorem
nor any transcendence argument is required for that route, since Morse--Hedlund gives it directly.
Second, and this is why the automaticity machinery of Part~II aims at the second tier rather than the
first: a proof that $c$ is not $2$-automatic gives \textup{P1} as a corollary, because the first tier
implies the second, and the second tier is where the certified computations of
Theorem~\ref{du:certified} have purchase. What no computation can do is establish either tier from a
prefix, by Remark~\ref{du:censuscaveat}: we compute $p(n)=2^n$ exactly for $n\le16$ (\EXACT) and that bounds no
period beyond the range computed.

\part{Excluded period words}

\noindent\emph{Source: \texttt{rule30\_excluded\_words.tex} (2026-09-09), reproduced with its labels prefixed \texttt{ex:}. The only input from the seed is Theorem~\ref{ex:thm:zero}; everything else holds for every number-like configuration. Other public work on excluded periods is credited in Remark~\ref{ex:prior}.}

% Part I body, from rule30_excluded_words.tex (labels prefixed ex:)
\section{Setting and the two-column theorem}\label{ex:setting}

Rule~30 acts on configurations $x\in\{0,1\}^{\Z}$ by
\begin{equation}\label{ex:eq:rule}
a(n+1,i)=a(n,i-1)\oplus\bigl(a(n,i)\vee a(n,i+1)\bigr),\qquad a(0,\cdot)=x .
\end{equation}
A configuration is \emph{number-like} \cite[Def.~2.4]{Kopra} if $x\ne 0^{\Z}$ and $x_i=0$ for all $i<N$, some $N$. The single seed $\delta_0$ is number-like.

For columns $i\le j$ the \emph{trace} is $\Tr_{[i,j]}(x)(t)=(a(t,i),\dots,a(t,j))$. Write $c=\Tr_0$ (the centre column), $u=\Tr_1$, $\ell=\Tr_{-1}$. A sequence $s$ is \emph{eventually periodic} if $s(t+p)=s(t)$ for all $t\ge T$, some $p\ge1$, $T\ge0$.

A word $w\in\{0,1\}^p$ is a \emph{tail} of $c$ if there is $N$ with $c(t)=w[(t-N)\bmod p]$ for all $t\ge N$. If $w$ is a tail so is every rotation of $w$; statements below are about necklaces, and any rotation may be certified.

\begin{theorem}[two-column theorem; \KNOWN: Jen \cite{Jen86,Jen90}; Kopra {\cite[Thm.~3.5]{Kopra}}]\label{ex:thm:zero}
Rule~30 is left permutive---\eqref{ex:eq:rule} is a bijection in $a(n,i-1)$ for fixed $a(n,i),a(n,i+1)$---and left spreading, since $001\mapsto1$. Consequently, for every number-like $x$ and every $i$, the width-$2$ trace $\Tr_{[i,i+1]}(x)$ is not eventually periodic.
\end{theorem}

\begin{proof}[Proof (for completeness; the argument is Jen's)]
Suppose $\Tr_{[i,i+1]}$ is $p$-periodic for $t\ge T$. Solving \eqref{ex:eq:rule} for $a(n,i-1)$,
\begin{equation}\label{ex:eq:left}
a(n,i-1)=a(n+1,i)\oplus\bigl(a(n,i)\vee a(n,i+1)\bigr),
\end{equation}
so $\Tr_{i-1}$ is $p$-periodic for $t\ge T$, and by induction every column to the left of $i$ is $p$-periodic for $t\ge T$. Let $M\ge N$ be the position of the leftmost $1$ of $x$. Since $x_j=0$ for $j<M$ and $001\mapsto1$, $000\mapsto0$, the leftmost $1$ moves one site left per step: the column at site $M-d$ is $0$ for $t<d$ and $1$ at $t=d$. For $d>\max(T+p,\,M-i-1)$, so that the column at site $M-d$ is one of those shown to be $p$-periodic from $T$, this column takes the values $0$ at time $d-p$ and $1$ at time $d$, both $\ge T$, so it is not $p$-periodic there.
\end{proof}

Only the zeros to the left of $N$ are used, which is why the theorem holds for all number-like $x$ and not just for finite-support ones. Kopra \cite[closing remarks]{Kopra} observes that Rule~90 satisfies every hypothesis of Theorem~\ref{ex:thm:zero} and has a periodic centre column from the seed; so no argument at the level of Theorem~\ref{ex:thm:zero} alone can decide the centre column, and everything below uses the specific gate of Rule~30, the OR in \eqref{ex:eq:rule}. The eventually periodic structure of the \emph{left} side of the seed's evolution is described by Rowland \cite{Rowland} and Gravner--Griffeath \cite{GG}; nothing there constrains column~$0$.

\begin{remark}[other public work on excluded periods; no priority is claimed]\label{ex:prior}
Three things in Part~I are public elsewhere.

(a) \emph{An infinite family.} The repository \cite{Cochon} states that the words $01^{q}$ with $q=7$ or $q\ge9$ cannot be eventual periods of the centre column (commit \texttt{0ff6930}, 10 September 2026). Its own review two days later (commit \texttt{adb0139}, 12 September 2026) calls the written argument incomplete and replaces it with a short uniform proof. This is the family $1^b0$ of Theorem~\ref{ex:thm:1b0} together with the word $1^70$ of Remark~\ref{ex:omega}, and it has the same hole at $b=8$. Both dates are before this paper was public (1 October 2026). We found the family independently. Theorems~\ref{ex:thm:zeroruns}, \ref{ex:thm:19}, \ref{ex:thm:1b0a}, \ref{ex:thm:1bv} and~\ref{ex:thm:0m11} go beyond it. Of their content \cite{Cochon} has three pieces: the word $0^{6}11$; words $0^{q}1$ of a single run length, by enumeration and with no uniform proof; and, in the proof of 12 September, those words $1^bv$ in which every zero is isolated and every run of ones has length at least $9$. We found no other excluded family in that repository (every branch searched by keyword, 7--8 October 2026).

(b) \emph{The two constants.} That the centre column of the seed is neither eventually $0$ nor eventually $1$ is \KNOWN; Theorem~\ref{ex:thm:const} states it for number-like configurations. It is in \cite{WRyan} (commit \texttt{8960181}, 22 July 2026); in \cite{Woffinden} (23 August 2026, Propositions~10.3--10.4, with the same two-line proof); in \cite{Condrey} (8 September 2026, for finite configurations; Section~\ref{ex:condrey}); and in \cite{Cochon} and \cite{Dib}.

(c) \emph{The method in kind.} A finite strip beside the column with a free boundary, the recurrent components of its transition graph, and the two-column theorem of Jen and Kopra: this is in \cite{Ikram} (15 May 2026, as a target), in \cite{Pat} (commit \texttt{de9c29d}, 8 September 2026) and in \cite{Cochon} (10 September 2026).
\end{remark}

\section{Left reconstruction and the interface}

\begin{lemma}[\PROVED]\label{ex:lem:left}
$\ell(n)=c(n+1)\oplus\bigl(c(n)\vee u(n)\bigr)$. On a row with $c(n)=1$, $\ell(n)=c(n+1)\oplus1$ does not depend on $u(n)$; on a row with $c(n)=0$, $\ell(n)=c(n+1)\oplus u(n)$.
\end{lemma}
\begin{proof}
\eqref{ex:eq:left} at $i=0$.
\end{proof}

Fix a word $w$ of length $p$ with $Z$ zero positions, and suppose $c$ has tail $w$ from row $N$. The \emph{interface sequence} $u^*$ is the sequence of values $u(n)$ over the rows $n\ge N$ with $c(n)=0$, in order.

\begin{corollary}[\PROVED]\label{ex:cor:interface}
If $x$ is number-like and $c$ has tail $w$, then $u^*$ is not eventually periodic.
\end{corollary}
\begin{proof}
Suppose $u^*$ is eventually periodic with period $q$ (counted in zero-rows). The zero-rows occur at positions periodic in $n$ with period $p$, $Z$ per period, so $u$ restricted to zero-rows is eventually periodic in $n$ with period $pq$, and so is $c$. By Lemma~\ref{ex:lem:left}, $\ell$ is eventually periodic with period $pq$. Then $\Tr_{[-1,0]}=(\ell,c)$ is eventually periodic, contradicting Theorem~\ref{ex:thm:zero} at $i=-1$.
\end{proof}

So to exclude $w$ it suffices to show that every configuration with $c=w^\infty$ (from some row) has an eventually periodic interface. This is a statement about the fibre of the trace map over $w^\infty$ and makes no reference to the seed: the seed enters only through Theorem~\ref{ex:thm:zero}.

\section{The $K$-window automaton}\label{ex:sec:window}

Fix $K\ge1$. The \emph{window} at row $n$ is $W(n)=(a(n,1),\dots,a(n,K))\in\{0,1\}^K$, so $u(n)=W(n)_1$. By \eqref{ex:eq:rule},
\[
W(n+1)=\Phi\bigl(c(n),W(n),f(n)\bigr),\qquad f(n)=a(n,K+1),
\]
where $\Phi(c,w,f)_i=w_{i-1}\oplus(w_i\vee w_{i+1})$ with $w_0=c$ and $w_{K+1}=f$. For $S\subseteq\{0,1\}^K$ define
\begin{align*}
\step(S,c)&=\{\Phi(c,w,f):w\in S,\ f\in\{0,1\}\}\quad\text{(the cell $K+1$ is unknown and left free);}\\
Z_u(S)&=\step(\{w\in S:w_1=u\},0)\quad\text{(a zero-row on which the interface bit is $u$);}\\
O(S)&=\step(S,1)\quad\text{(a one-row).}
\end{align*}

\begin{lemma}[soundness; \PROVED]\label{ex:lem:sound}
For any configuration $x$ and any row $N$, put $S(N)=\{0,1\}^K$ and $S(n+1)=Z_{u(n)}(S(n))$ if $c(n)=0$, $S(n+1)=O(S(n))$ if $c(n)=1$. Then $W(n)\in S(n)$ for all $n\ge N$; in particular $S(n)\ne\emptyset$.
\end{lemma}
\begin{proof}
Induction on $n$. If $W(n)\in S(n)$ then $W(n+1)=\Phi(c(n),W(n),f(n))$; on a zero-row $W(n)_1=u(n)$, so $W(n)$ lies in the selected subset; in either case $W(n+1)\in S(n+1)$.
\end{proof}

The automaton is an over-approximation: it forgets everything right of column $K$ and remembers only a set of possible windows. What it never does is discard a window that the true configuration occupies. That the method is an over-approximation rather than a decision procedure is not a weakness of this construction but a necessity: deciding whether a given temporal sequence of states can be generated in $|s|$ steps --- which is exactly the question Part~I asks of a candidate period word --- is NP-complete, for an automaton Green constructs for the purpose \cite{Green87}. So no complete and efficient test for this question can exist in general, and a sound incomplete criterion with a characterised reach is the right kind of object to look for.

\section{The period graph and the exclusion criterion}\label{ex:sec:graph}

For a word $w$ define the finite directed labelled multigraph $G_w$:
\begin{itemize}[nosep]
\item vertices: the subsets $S\subseteq\{0,1\}^K$ reachable from $\ALL=\{0,1\}^K$ by the edges below;
\item edges: from $S$, for each $u$-word $(u_1,\dots,u_Z)\in\{0,1\}^Z$, apply the letters of $w$ in order, using $Z_{u_j}$ at the $j$-th zero position and $O$ at each one position; if the result $T$ is nonempty put an edge $S\to T$ labelled $(u_1,\dots,u_Z)$. Distinct labels are distinct edges even when their targets coincide.
\end{itemize}

\begin{lemma}[\PROVED]\label{ex:lem:path}
If $x$ is number-like and $c$ has tail $w$ from row $N$ ($N$ chosen so that row $N$ is the start of $w$), then $u^*$, cut into blocks of $Z$, is the label sequence of an infinite path in $G_w$ starting at $\ALL$.
\end{lemma}
\begin{proof}
Lemma~\ref{ex:lem:sound} with the period structure of $w$; every state along the path is nonempty, so every edge exists.
\end{proof}

\begin{lemma}[single-cycle criterion; \PROVED]\label{ex:lem:cycle}
Suppose every strongly connected component $C$ of $G_w$ that contains at least one edge has exactly as many edges (with multiplicity, both ends in $C$) as vertices. Then the label sequence of every infinite path in $G_w$ is eventually periodic.
\end{lemma}
\begin{proof}
$G_w$ is finite, so an infinite path eventually stays inside one strongly connected component $C$, which then contains an edge. Strong connectivity gives every vertex of $C$ at least one outgoing edge inside $C$ (a one-vertex component with an edge has a self-loop); with $\#\text{edges}=\#\text{vertices}$ each vertex has exactly one. So $C$ is a single directed cycle with one label per step, and once the path is in $C$ its labels are periodic with period $|C|$ blocks.
\end{proof}

\begin{theorem}[exclusion criterion; \PROVED]\label{ex:thm:A}
If for some $K$ the graph $G_w$ satisfies the hypothesis of Lemma~\ref{ex:lem:cycle}, then $w$ is not a tail of the centre column of any number-like configuration, in particular not of the single seed.
\end{theorem}
\begin{proof}
Lemmas~\ref{ex:lem:path} and~\ref{ex:lem:cycle} make $u^*$ eventually periodic; Corollary~\ref{ex:cor:interface} forbids it.
\end{proof}

\begin{remark}\label{ex:rem:onesided}
The criterion is one-sided. If $G_w$ fails the hypothesis at some $K$, nothing follows: the over-approximation may branch where the true fibre does not. Only an exact computation on the fibre (Section~\ref{ex:sec:boundary}) decides that a word is genuinely beyond the method. The hypothesis is equivalent to every recurrent component having Perron root exactly~$1$, but the proofs rest on the combinatorial count, which is exact.
\end{remark}

\begin{lemma}[transfer in the length of a run of ones; \PROVED]\label{ex:lem:transfer}
Let $w_b=v\,1^b$. Suppose that in $G_{w_B}$, for every vertex $S$ and every admissible label, the state $T$ reached after $v$ and $B$ one-rows satisfies $O(T)=T$. Then $G_{w_b}=G_{w_B}$ (same vertices, same labelled edges) for every $b\ge B$.
\end{lemma}
\begin{proof}
Induction on $b\ge B$. From any vertex and any label, the state after $v$ and $b$ one-rows is a $T$ with $O(T)=T$, so after $b+1$ one-rows it is $O(T)=T$: same target, same label. Starting from $\ALL$ the two graphs are generated identically, and the $O$-fixed property persists.
\end{proof}

Hence one certificate for $w_B$, together with the $O$-fixed check, certifies all $w_b$ with $b\ge B$.

\begin{remark}[what kind of object the exclusion criterion is]\label{ex:rem:notsync}
The graph $G_w$ is the accessible subset automaton of a \emph{partial nondeterministic} transition
system, and the criterion applied to it is the Perron-root condition of
Remark~\ref{ex:rem:onesided}. It is not synchronisation, for three independent reasons: nothing
collapses to a single state; the empty set is unreachable, since $S\ne\emptyset$ forces some
$Z_u(S)\ne\emptyset$; and the interface bits are tests on the state rather than free input letters, so
although mortal words exist in the free monoid none is realisable, because $u(n)=1$ with $c(n)=0$
forces $u(n+1)=1$. The quantifier differs too: a \v{C}ern\'y-type bound asserts that \emph{some} short
word synchronises, while the exclusion needs a conclusion for \emph{every} word. Where the comparison
is legitimate --- the deterministic window map does synchronise --- the bound is vacuous, the exact
reset thresholds being $3,4,9,11,14,17$ for $K=2,\dots,7$ against the quadratic $(2^K-1)^2$, too weak
by a factor of $949$ at $K=7$.
\end{remark}

\section{Words with a long run of zeros}

The exclusions in this section were obtained by a direct argument on the fibre, recorded because it exhibits the mechanism. An \emph{impulse} is a row with $c=1$ preceded and followed by rows with $c=0$; \emph{zero drive} is a run of rows with $c=0$. Write $v,y,z$ for columns $2,3,4$.

\begin{lemma}[quench; \PROVED]\label{ex:lem:quench}
Let $x$ be any configuration. Suppose $c$ has $\ge10$ consecutive zero-rows, then an impulse at row $t_0$, then zero drive. Then, for as long as the zero drive continues,
\begin{enumerate}[nosep,label=(\roman*)]
\item $u(t)=1$ for all $t\ge t_0+7$;
\item $(v(t),y(t))\ne(0,0)$ for all $t\ge t_0+11$;
\item $(u,v,y)(t)=(1,0,1)$ for all $t\ge t_0+12$.
\end{enumerate}
The state at a row depends only on $c$ at earlier rows, so the conclusions hold at the row of the next impulse as well.
\end{lemma}
\begin{proof}
\emph{Nesting.} A row under zero drive for $a\ge10$ steps is a row under zero drive for exactly $10$ steps started from its own predecessor at step $a-10$; so any fact holding for every configuration with pre-run exactly $10$ holds for pre-run $\ge10$.

\emph{Three finite facts.} With $c=0^{10}\,1\,0^{20}$ and every other cell of row~$0$ free, the following are unsatisfiable (certificates \texttt{drive2.py}, \texttt{freeze.py}; CaDiCaL~1.5.3 and Glucose~4 agree): $u=0$ at $t_0+7$; $(v,y)=(0,0)$ at $t_0+11$; $(v,y)=(0,0)$ at $t_0+12$. Every cell in columns $\le3$ at rows $\le T$ depends only on row-$0$ cells at columns $\le T+3$, so the instance width is sufficient.

\emph{Monotonicity.} Under zero drive $u(t+1)=u(t)\vee v(t)$, so $u=1$ persists; with the first fact, (i).

\emph{Chain.} Suppose $u(t-1)=u(t-2)=1$ and $(v,y)(t)=(0,0)$, with $t-2\ge t_0+7$. From $v(t)=u(t-1)\oplus(v\vee y)(t-1)$ we get $(v\vee y)(t-1)=1$. From $y(t)=v(t-1)\oplus(y\vee z)(t-1)=0$: if $v(t-1)=0$ then $y(t-1)=1$ and $y(t)=1$, contradiction; so $v(t-1)=1$, whence $(v\vee y)(t-2)=0$, i.e.\ $(v,y)(t-2)=(0,0)$. Descending by $2$ from any $t\ge t_0+13$ lands on $t_0+11$ or $t_0+12$, both excluded. This is (ii).

\emph{Freeze.} For $t\ge t_0+12$, $v(t)=1$ would need $(v\vee y)(t-1)=0$ with $t-1\ge t_0+11$, excluded. So $v(t)=0$, and $(v\vee y)(t)=1$ gives $y(t)=1$. With (i), (iii).
\end{proof}

\begin{lemma}[block step; \PROVED]\label{ex:lem:block}
An impulse landing on $(u,v,y)=(1,0,1)$, followed by zero drive, gives $u=0,0,1$ at the next three rows.
\end{lemma}
\begin{proof}
From \eqref{ex:eq:rule}: $u(+1)=1\oplus(1\vee0)=0$, $v(+1)=1\oplus(0\vee1)=0$, $y(+1)=0\oplus(1\vee z)=1$; $u(+2)=u(+1)\vee v(+1)=0$, $v(+2)=u(+1)\oplus(v\vee y)(+1)=1$; $u(+3)=u(+2)\vee v(+2)=1$.
\end{proof}

\begin{theorem}[isolated ones, zero-runs $\ge10$; \PROVED\ from \EXACT\ inputs]\label{ex:thm:zeroruns}
Let $w$ be a necklace in which every $1$ is isolated and every cyclic zero-run has length $\ge11$. Then $w$ is not a tail of the centre column of any number-like configuration. The same holds with zero-runs $\ge10$ (one further finite fact, \texttt{run10.py}), and for the word $0000000001$ (\S\ref{ex:sec:spine}).
\end{theorem}
\begin{proof}
Take any configuration with $c=w^\infty$ from some row. The first impulse is preceded by $\ge10$ zero-rows, so Lemma~\ref{ex:lem:quench} freezes $(u,v,y)$ at $(1,0,1)$ from $12$ rows after it; the next impulse is $\ge12$ rows later and lands on $(1,0,1)$. By Lemma~\ref{ex:lem:block} the interface on the following block reads $0,0,1,1,\dots$ ($u=1$ is then preserved by monotonicity until the next impulse), and the next impulse again finds $(1,0,1)$ by Lemma~\ref{ex:lem:quench} applied afresh. By induction on blocks $u^*$ is eventually periodic with the period of $w$; Corollary~\ref{ex:cor:interface} concludes.
\end{proof}

\subsection{The word $0000000001$}\label{ex:sec:spine}
Zero-run $9$ is one short of Lemma~\ref{ex:lem:quench}. Here the fibre was enumerated exactly: writing $\tau_j$ for the number of leading zeros of $u$ on the $j$-th block of zero-rows, every admissible $\tau$-sequence is, after at most three symbols, an alternating $1212\ldots$ spine that either continues forever or is eventually absorbed into $2222\ldots$, and this is proved by five bounded-window facts ($11$, $2221$, $21221$ never occur; $1221$ only at positions $\le1$; symbols outside $\{1,2\}$ only at positions $\le2$) which reduce, by the observation that the row at time $p$ is itself a phase-$0$ configuration of the fibre, to an enumeration at depth~$6$. Every admissible interface is therefore eventually periodic and Corollary~\ref{ex:cor:interface} applies. Details in \cite[\S2]{Q}.

\begin{theorem}[$K=14$ family certificate, \EXACT; see Section~\ref{ex:sec:repro}]\label{ex:thm:19}
Every necklace all of whose cyclic zero-runs have length $\ge19$, with $1$-runs of arbitrary length, is excluded for every number-like configuration (certificate \texttt{cover14\_19\_48}, Section~\ref{ex:sec:repro}). At $K=14$ every necklace tested with all cyclic zero-runs $\ge17$ is excluded as well, and at $K=16$ so is every necklace tested whose cyclic zero-runs all equal some $Z\ge15$, of the shapes $0^Z1^a$ with $a\le5$ and $0^Z1^a0^Z1^b$ with $a,b\le3$ (\EXACT\ for each such word; Remark~\ref{ex:omega}); the family-level statement at $17$ rests on a product construction we have not rebuilt (see the sweep below) and is \OPEN\ as a family statement. At zero-runs $16$ the necklaces with $1$-run $3$, and with the $1$-run pair $(3,3)$, are not certified at $K=14$ ($0^{16}1^3$ is certified at $K=16$).

Two routes bear on this statement and they stop in different places, which is worth separating. By \emph{forcing} --- exact SAT on the fibre --- the interface bit fails to be forced at blocks $4$ and $5$ already at zero-runs $18$ with the $1$-run pair $(3,3)$, so that route halts at $19$. By the \emph{single-cycle} criterion of Theorem~\ref{ex:thm:A}, applied word by word at $K=14$, every necklace tested is excluded down to $17$ (the family statement at $17$ being \OPEN, as above), and the obstruction at $16$ is the $1$-run shape $1^3$ rather than the zero-run length.
\end{theorem}

The two theorems say: a long run of zeros in the period word quenches the interface. The threshold $10$ (for isolated $1$s) comes from the finite facts of Lemma~\ref{ex:lem:quench}; below a run of $5$ the state $(v,y)=(0,0)$ can be sustained indefinitely by the right context and no quench occurs.

\section{Words with a long run of ones}

The same mechanism operates from the other side: a long run of ones also quenches the interface. These results are certificates of Theorem~\ref{ex:thm:A} at $K=10$, produced by two independent implementations of Sections~\ref{ex:sec:window}--\ref{ex:sec:graph} (\texttt{kauto3.py}, \texttt{s3\_fam1b0a.py}), which agree on every word both cover ($1^b0$ for $b=9,\dots,12$; $0^m11$ for $m=8,\dots,11$), and spot-checked by a third (\texttt{kwin.py}, Appendix) on $1^90$, $1^{12}0$, all rotations of $1^80$ and of $0^m11$ for $m=6,7$, and $0^m11$ for $m=3,4,5,8,9$, with identical verdicts on every word both cover. ``Certified'' means the graph $G_w$ was built and every recurrent component has $\#\text{edges}=\#\text{vertices}$; ``transfer'' means the $O$-fixed check of Lemma~\ref{ex:lem:transfer} at $B=29$.

\begin{theorem}[\EXACT\ certificate of Theorem~\ref{ex:thm:A}; also public in \cite{Cochon}, see Remark~\ref{ex:prior}]\label{ex:thm:1b0}
$1^b0$ is excluded for every $b\ge9$.
\end{theorem}
\begin{proof}[Certificate]
$b=9,\dots,34$ individually ($26$ graphs; for $b=9$ the graph has $3$ states and $3$ edges); transfer at $B=29$ gives all $b\ge29$.
\end{proof}

\begin{theorem}[\EXACT\ certificate of Theorem~\ref{ex:thm:A}]\label{ex:thm:1b0a}
$1^b0^a$ is excluded for every $b\ge9$ and every $a\ge1$.
\end{theorem}
\begin{proof}[Certificate]
For each $a=1,\dots,19$: $b=9,\dots,34$ individually and transfer at $B=29$ ($494$ graphs, $19$ transfer checks). For $a\ge19$, Theorem~\ref{ex:thm:19}.
\end{proof}

\begin{theorem}[\EXACT\ certificate of Theorem~\ref{ex:thm:A}]\label{ex:thm:1bv}
For every word $v$ with $|v|\le12$ that begins and ends with $0$, the word $1^bv$ is excluded for every $b\ge29$.
\end{theorem}
\begin{proof}[Certificate]
$b=29$ with transfer, for all $2048$ such $v$.
\end{proof}

\begin{theorem}[\EXACT\ certificate of Theorem~\ref{ex:thm:A}; the word with $m=6$ is also in \cite{Cochon}]\label{ex:thm:0m11}
$0^m11$ is excluded for $m=6$ and for every $m\ge8$.
\end{theorem}
\begin{proof}[Certificate]
$m=6$ (all eight rotations at $K=10$, e.g.\ $63$ states, $89$ edges for $00000011$; also at $K=12$), $m=8$ ($39$ states, $59$ edges), $m=9,\dots,20$ individually; for $m\ge19$, Theorem~\ref{ex:thm:19}. The word $0^711$ has a branching recurrent component in every rotation at $K=10$ ($6$ vertices, $10$ edges), $K=12$ and $K=14$ ($14$--$16$ vertices, $22$--$26$ edges), and $m\le5$ branch at $K=10$; these are not decided by the method at these $K$.
\end{proof}

\begin{remark}[sharpness]
$1^80$ has a branching recurrent component at $K=10$ ($8$ vertices, $15$ edges) and at $K=12$ ($21$ vertices, $39$ edges), and it is uncertified at every $K$ up to $16$. By Remark~\ref{ex:rem:onesided} this proves nothing about the word, and $b=9$ is not a threshold: $1^70$ is certified at $K=10$ (Remark~\ref{ex:omega}), so $1^80$ is a hole rather than the edge of the family, and the lower bound in Theorems~\ref{ex:thm:1b0}--\ref{ex:thm:1b0a} is only where the whole family is certified. The non-monotonicity in Theorem~\ref{ex:thm:0m11} ($m=6$ certified, $m=7$ not, at every window tried) is a hole of the same kind. It does not mirror the forcing of the interface after mixed zero-runs $0^{m_1}1\,0^{m_2}1$, which is monotone in each run length separately (\MEASURED: SAT, exact per cell, $6\le m_1,m_2\le15$); an earlier reading of that forcing as a non-monotone landscape compared cells differing in both run lengths and is \textsc{withdrawn}. Both holes, $1^80$ and $0^711$, branch in every rotation and the branching grows with $K$; whether they are excluded by some other route is open.
\end{remark}

\begin{theorem}[the constants; \PROVED; \KNOWN\ for the seed, \cite{WRyan,Woffinden,Condrey,Cochon,Dib}, see Remark~\ref{ex:prior}]\label{ex:thm:const}
The centre column of a number-like configuration is neither eventually $0$ nor eventually $1$.
\end{theorem}
\begin{proof}
If $c(t)=0$ for $t\ge T$ then, by \eqref{ex:eq:rule} at $i=1$, $u(t+1)=u(t)\vee v(t)$ for $t\ge T$: $u$ is monotone, hence eventually constant, so $(c,u)$ is eventually periodic, contradicting Theorem~\ref{ex:thm:zero}. If $c(t)=1$ for $t\ge T$ then by Lemma~\ref{ex:lem:left} $\ell(t)=c(t+1)\oplus1=0$ for $t\ge T$, so $(\ell,c)$ is eventually periodic, contradicting Theorem~\ref{ex:thm:zero}.
\end{proof}

\section{The boundary of the method}\label{ex:sec:boundary}

Theorem~\ref{ex:thm:A} uses the seed only through Theorem~\ref{ex:thm:zero}. Its reach is therefore exactly the set of words $w$ for which \emph{every} configuration with centre column $w^\infty$ has an eventually periodic interface. We now show this set excludes the words that matter.

\begin{proposition}[the fibre over a centre column; \PROVED]\label{ex:prop:surj}
For every sequence $c\in\{0,1\}^{\N}$ and every right context $r=(r_1,r_2,\dots)\in\{0,1\}^{\N}$ there is exactly one configuration $x$ with $\Tr_0(x)=c$ and $a(0,i)=r_i$ for all $i\ge1$. In particular every sequence is the centre column of some configuration, and the interface $u^*$ is a function of $(c,r)$.
\end{proposition}
\begin{proof}
Put $a(n,0)=c(n)$ and $a(0,i)=r_i$. The cells $a(n,i)$ with $i\ge1$ are determined forward by \eqref{ex:eq:rule} from row $n-1$ at columns $\ge0$, all of which are already defined; the cells with $i\le-1$ are determined by \eqref{ex:eq:left} from the two columns to their right. The rule \eqref{ex:eq:rule} holds at every column $i\ge1$ by construction and at every column $i\le0$ because \eqref{ex:eq:left} is \eqref{ex:eq:rule} solved for its left input. Uniqueness: every cell has been determined.
\end{proof}

Note that not every pair $(c,u)$ is realised: \eqref{ex:eq:rule} at column~$1$ forces $u(n+1)=\neg c(n)$ whenever $u(n)=1$. The interface is an output of the fibre, not a free input.

Let $L(w)\subseteq\{0,1\}^{\N}$ be the set of interface sequences $u^*$ over all configurations with $c=w^\infty$ from row~$0$, and $L_k(w)$ its set of prefixes of length $k$ (zero-rows). By compactness of $\{0,1\}^{\Z}$ and continuity of the trace, $L_k(w)$ coincides with the set of length-$k$ interface words realisable at finite depth, which is what the enumerations compute.

\begin{proposition}[boundary; \PROVED]\label{ex:prop:boundary}
If $L(w)$ is uncountable---in particular if $|L_k(w)|$ grows exponentially in $k$---then there is a configuration with $c=w^\infty$ whose interface is not eventually periodic, and Theorem~\ref{ex:thm:A} cannot exclude $w$ at any $K$.
\end{proposition}
\begin{proof}
Eventually periodic sequences are countable, so an uncountable $L(w)$ contains one that is not; it is realised by a configuration by definition of $L(w)$. Lemma~\ref{ex:lem:path} uses only Lemma~\ref{ex:lem:sound}, which holds for every configuration, so this configuration's path in $G_w$ has a non-eventually-periodic label sequence, and the hypothesis of Lemma~\ref{ex:lem:cycle} fails at every $K$. If $|L_k(w)|\ge\lambda^k$ for some $\lambda>1$ then $L(w)$, being closed, has positive topological entropy and is uncountable.
\end{proof}

\begin{proposition}[the entropy of the fibre is a limit, not a fit; \PROVED]\label{ex:prop:fekete}
$|L_{m+n}(w)|\le|L_m(w)|\,|L_n(w)|$ for all $m,n\ge1$. Consequently
$h(w):=\lim_k\frac1k\log|L_k(w)|$ exists and equals $\inf_k\frac1k\log|L_k(w)|$.
\end{proposition}
\begin{proof}
The language is factorial, so if $u\in L_{m+n}$ then its length-$m$ prefix lies in $L_m$ and its length-$n$
suffix in $L_n$; the map $u\mapsto(\text{prefix},\text{suffix})$ is injective because $u$ is the
concatenation of the two. Hence $|L_{m+n}|\le|L_m||L_n|$, i.e.\ $k\mapsto\log|L_k|$ is subadditive, and
Fekete's lemma \cite{Fekete} gives the limit and its identification with the infimum.
\end{proof}

\begin{corollary}[the exponential clause is not computationally dischargeable; \PROVED]\label{ex:cor:infimum}
For no word $w$ can a finite enumeration establish the hypothesis $|L_k(w)|\ge\lambda^k$ of
Proposition~\ref{ex:prop:boundary}. By Proposition~\ref{ex:prop:fekete} the entropy is the
\emph{infimum} of $k^{-1}\log|L_k(w)|$, so every computed value is an \emph{upper} bound on $h(w)$ and
none is a lower bound: the infimum over all $k$ may be zero whatever the first thousand terms show.
\end{corollary}

\begin{remark}[what this costs, and where it points]\label{ex:rem:infimum}
The corollary is worth stating because the exponential clause reads like the practical branch of
Proposition~\ref{ex:prop:boundary} and is in fact the unreachable one. It applies uniformly, not just
to the hard case: the comparison necklaces $0000011111$ (nine exact terms, ratio $\to2.0$) and
$0000000101$ (ratio $\approx1.55$) look exponential and are \MEASURED\ so, and by the corollary they
cannot be proved so by extending those enumerations --- the natural next move, and a wasted one. For
$01$ the situation is worse in a second and independent way, the measured ratio being $1.094$ at $k=60$ and
falling, so there the counts do not even point where the hypothesis needs them to; the stream measurement of Remark~\ref{ex:rem:stream} does, but at depths ($k\approx700$) no enumeration reaches.
Consequently uncountability of $L(w)$, for $01$ or for any other background, has to come from a
\emph{construction} rather than from counting: an explicit infinite family, or a self-embedding that
generates one. That is a narrow target but a definite one, and it is the only route to
uncountability that Proposition~\ref{ex:prop:boundary} leaves open. Whether $L(01)$ is uncountable remains \OPEN.
Uncountability is sufficient for the conclusion of Proposition~\ref{ex:prop:boundary}, one interface that
is not eventually periodic, but not necessary: the orbit closure of $z=1\,0\,1\,0^2\,1\,0^3\,1\cdots$
consists of the shifts of $z$, the points $0^a10^{\infty}$ and $0^{\infty}$, so it is a countable closed
shift-invariant set, and $z$ is not eventually periodic. Nor does superpolynomial growth force uncountability: the sequences with
finitely many $1$s in which, if there are $m$ of them, consecutive $1$s are at distance at least $m$ form a
countable closed shift-invariant set in which every point is eventually periodic, with
$\sum_m\binom{k-(m-1)^2}{m}$ words of length $k$, a superpolynomial count that is $3\,334\,829$ at $k=60$ against
$|L_{60}(01)|=103\,220$ (\PROVED; folklore-type examples).
\end{remark}

\begin{remark}[random right halves: the interface of $01$ has positive entropy as measured; \MEASURED, \EXACT\ and \PROVED\ parts as marked]\label{ex:rem:stream}
The counts to $k=60$ are dominated by words anchored at row~$0$. Twenty random right halves driven by $(01)^\infty$ for up to $400\,000$ rows settle, after a transient of a few hundred letters, either into one of the periodic tails already known to be realisable ($(1\,4)^\infty$ and $1^\infty$ in run-length coding; $2$ of $20$) or into a stationary stream of gap letters $2$ and $4$ only, the $2$s isolated and at rate $0.155$--$0.172$ in every seed and every time window, with no coarsening over $200\,000$ letters; the picture is a boundary layer about four cells wide next to the column with ordinary Rule~30 chaos behind it (\MEASURED; ledger 10.238d~D). Every one of $342\,804$ sampled windows of length $36$ of these streams lies in the exact $L_{36}(01)$, rebuilt independently by SAT-driven search with the counts of Section~\ref{sec:record} letter for letter (\EXACT; 10.238d~A, D). The conditional block entropies of the late stream fall slowly, $0.384$ to $0.366$ bits per gap letter at block lengths $10$ to $20$, and extrapolate to about $0.35$ bits per gap letter, $0.07$ per letter of the interface, a growth rate near $1.05$ (\MEASURED; 10.238d~E); with $U=2\,4^{3}$ and $V=2\,4^{5}$, the two commonest late blocks, every one of the $32$ sequences of five blocks is realisable after a free head of $20$ letters and again after $60$, and $116$ of the $128$ sequences of seven blocks finished within the solver's time, all realisable and none failed (\EXACT; 10.238d~F). So $L(01)$ has positive entropy by every measurement here, the growth fits above describe the row-$0$ transient ($2^{0.07k}$ overtakes $\exp(5.91k^{0.28})$ only near $k\approx700$), and the zero-entropy expectation of the earlier text is withdrawn; uncountability remains \OPEN, since what a proof needs is the free choice at every depth. Two facts say why that resists proof (ledger 10.239d). The rightmost cell at which two right-half evolutions with the same boundary differ moves right by exactly one cell per row, forever (\PROVED: the cell to its right reads the difference through its XOR input while its OR inputs agree), so Rule~30 has no localised particle on any background and a free choice cannot be engineered as a glider in a periodic right half; and right halves with random cells on columns $1..N$ only ($N=10$ and $100$) give the same stream in $15$ of $16$ runs, so restricting to number-like right halves changes nothing as far as computed (\MEASURED). Finally, every minimal forbidden space-time block of Rule~30 has width at most $2$, except the twelve blocks of width $3$ and height $2$ that are the rule table itself: a block is globally admissible iff the rule holds at every interior cell and its rightmost two columns form an admissible two-column trace, by left-permutivity on the left and because the right neighbour is read only where the cell is $0$ (\PROVED; ledger 10.236d~B), so no wider excluded block adds a constraint beyond the two-column trace; the Perron roots of the two-column Rauzy graphs descend $3.000$, $2.618$, $2.500$, $2.521$, $2.480$, \dots, $2.342$ at heights $2$--$12$ without a plateau (\EXACT; 10.236d~C). The numbers of two-column traces by height are tabulated to height $17$ in \cite[p.~1087]{NKS} and to $13$ in \cite{Dib}.
\end{remark}

\begin{proposition}[a branching criterion for uncountability; \PROVED]\label{ex:prop:branching}
Call $u\in L(w)$ a \emph{branch word} if $u0$ and $u1$ both lie in $L(w)$. If there is a legal $u$
such that \emph{every} legal word extending $u$ has a branch word among its extensions, then $L(w)$ is
uncountable, and hence by Proposition~\ref{ex:prop:boundary} some configuration with $c=w^{\infty}$ has
an interface that is not eventually periodic.
\end{proposition}
\begin{proof}
Recursively: $u$ has a branch word $u_0$ among its extensions, and $u_00,u_01$ are both legal; each of
those again has a branch word among its extensions. The resulting binary tree of branch words has
$2^{\aleph_0}$ infinite paths, every prefix of every path is legal, and distinct paths give distinct
sequences, so $|L(w)|\ge2^{\aleph_0}$. Note the hypothesis already forces every legal extension of $u$
to extend further, so no separate appeal to the absence of dead ends is needed.
\end{proof}

\begin{remark}[why the criterion cannot be made uniform]\label{ex:rem:branching}
The converse fails --- isolated points may accumulate on the perfect kernel, so an uncountable $L(w)$
need not contain a word all of whose extensions branch; the exact characterisation is the usual one,
that the tree contains a perfect subtree, whose nodes need not be all the extensions of any word. What
is worth recording is that the criterion cannot hold \emph{uniformly} here. If every legal word had a
branch word within $B$ further symbols then $|L_{k+B}|\ge2|L_k|$, so $|L_{mB}|\ge2^m$ and
$h(w)\ge(\log2)/B>0$. For $01$ the measured ratio is $1.094$ at $k=60$ and falling, which read as
$h=0$ would mean that no uniform $B$ exists and the delay before branching is unbounded; the stream
measurement of Remark~\ref{ex:rem:stream} reverses that reading --- the late stream has positive measured
entropy and branches within a bounded number of blocks, and it is the row-$0$ transient that delays. The
criterion may still be of the non-uniform shape that resists proof, since the transient must be crossed
first, which is consistent with the problem's difficulty rather than a defect of the criterion. It also sharpens what to measure: not whether branching occurs, but how the
worst-case delay before it grows with depth. \MEASURED, to depth $130$ on a random sample: a branch
that is forced for the $32$ levels from $28$ to the end of the enumeration resumes branching at $66$
and then branches at every period-$5$ boundary to $121$, so apparent forcing inside a finite window is
a horizon artefact and not isolation.
\end{remark}

\begin{proposition}[a finite block traps the interface; \PROVED]\label{du:prop:trap}
In the right half-plane with column~$0$ pinned to $(01)^{\infty}$, every configuration whose row~$0$
begins with
\[
B=1001100010010100\,0101
\]
in columns $1,\dots,20$ has interface identically $(10)^{\infty}$, whatever the cells beyond column
$20$ are.
\end{proposition}
\begin{proof}
A state is the block of columns $1,\dots,20$ together with the row parity; its successor depends only on
the state and the value of column~$21$, since column~$0$ is the parity. From $(B,0)$, the set $R_t$ of
blocks reachable at row $t$ over \emph{all} input sequences becomes periodic after $21$ rows with period
$8$; the $188$ states on the cycle form a set closed under both inputs, and at every even row column~$1$
takes a single value across $R_t$. Two independent implementations agree on every number
(\texttt{s3\_trap\_repro.py}; \texttt{agent\_reach/proof/p3\_wall\_trap.py}), and $1\,000$ random
tails evolved for $4\,000$ rows all give $(10)^{\infty}$.
\end{proof}

\begin{remark}[what the trap does and does not settle]\label{du:rem:trap}
Three consequences. (i) The push-forward of any product measure on row~$0$ to the interface has an
atom: a random row~$0$ enters this trap with frequency about $0.034$ (\MEASURED, $3\,000$ trials, always
within $100$ rows). So no argument of the form ``no finite cylinder determines the interface'' is
available, and uncountability of $L(01)$, if it holds, must be established despite the atoms. (ii) The
trapped word is nevertheless \emph{not} an isolated point of $L(01)$: $(01)^{k}00$ and $(01)^{k}01$ are
both legal for $k\le60$ (\EXACT), so $(01)^{k}0$ is a branch word in the sense of
Proposition~\ref{ex:prop:branching}, whose hypothesis the trap leaves untouched. (iii) The contrast with
Gravner and Griffeath's Proposition~3.3 \cite{GG} --- a left-permutive edge automaton has no periodic
robust solution against a \emph{zero} edge, Rule~$30$ included --- is that the pinned column here
contains $1$s, and it is the $1$s that build the wall: with a free or all-zero column~$0$ no such block
exists (\EXACT, straight zero diagonals of length $80$ are then realisable). What is not claimed: that
every configuration whose interface no modification of far cells can change is a trap of this kind
(\MEASURED\ only, on every instance found; a configuration shielded only against single far flips need
not be one, the $1^{\infty}$ crystal of Remark~\ref{du:rem:cap}(i) being an example), or anything about
uncountability itself, which remains \OPEN.
\end{remark}

\begin{theorem}[the diagonal cap; \PROVED]\label{du:thm:cap}
In the right half-plane with column~$0$ pinned to $(01)^{\infty}$, $a(t,0)=t\bmod2$, let $g(N)$ be the
largest $l$ for which some configuration has $a(N-k,k)=0$ for $k=1,\dots,l$: the longest run of zeros
ending at column~$1$ on the anti-diagonal $t+i=N$, which is also the longest straight chain of
influence into the interface cell $a(N,1)$, since $\partial\,(l\oplus(c\vee r))/\partial r=1\oplus c$.
Then
\begin{enumerate}[nosep]
\item[(i)] $g(N)=6$ for every even $N\ge58$;
\item[(ii)] $g(N)\le15$ for every even $N\ge30$; with the exact values $g(N)=N$ for even $N\le28$ and
$g(N)=15$ for even $30\le N\le56$ (\EXACT, each even $N\le150$ separately, \texttt{s3\_runcap.py}),
$g$ on the even diagonals is the step function $N,\,15,\,6$;
\item[(iii)] $g(N)=2$ for every odd $N\ge5$, with $g(1)=1$ and $g(3)=3$.
\end{enumerate}
\end{theorem}
\begin{proof}
Three finite facts and a shift. (a) \emph{The triangle.} The cells $a(N-l,l),\dots,a(N-1,1)$ are
computed from row $N-l$, columns $0,\dots,l$, alone, and column~$0$ there is the parity of $N-l$; so
which patterns of columns $1,\dots,l$ produce the run is a brute force over $2^{l}$ cases that depends on
$N$ only through its parity (\texttt{s3\_cap\_verify.py}). For even $N$: the run of $7$ occurs iff row
$N-7$ carries $1111110$ in columns $1,\dots,7$; the run of $6$ iff row $N-6$ carries $000000$ in
columns $1,\dots,6$; the run of $16$ iff row $N-16$ carries $0100100100100110$ or $1001001001001000$ in
columns $1,\dots,16$; in each case these are all the patterns. For odd $N$: the run of $2$ occurs iff
the odd row $N-2$ carries $10$ in columns $1,2$, the run of $3$ iff the even row $N-3$ carries $110$ in
columns $1,\dots,3$, and no pattern of row $N-4$ produces a run of $4$. (b) \emph{The shift.} The right half-plane is an autonomous
one-sided automaton driven by the pinned column, which has period~$2$; so row~$2$ of any configuration
is a legal row~$0$, and a pattern reachable at row $r+2$ is reachable at row $r$. (c) \emph{The
certificate.} On the $200$-row light-cone instance, $1111110$ in columns $1,\dots,7$ is reachable at
row $49$ and unreachable at row $51$, by Cadical153 and by Glucose4 (rows $53$, $61$, $99$ unreachable
as controls; \texttt{s3\_cap\_upper.py}). By (b) it is unreachable at every odd row $\ge51$, so by (a)
no even diagonal $N\ge58$ carries a run of $7$, hence none longer: $g(N)\le6$. Reachability at row $49$
gives the run of $7$ at $N=56$, where the exact table has $g=15$. In the same way both $16$-patterns are
unreachable at row $14$, and one of them is reachable at row $12$ as $g(28)=28$ requires
(\texttt{s3\_cap\_plateau.py}); so no even diagonal $N\ge30$ carries a run of $16$, which with the
exact values is (ii). (d) \emph{The construction.} On the ring of $728$ cells Rule~30 has a cycle of
minimal period $14$ whose cell~$0$ equals the phase parity and whose cells $1,\dots,6$ read $000000$ at
phase $12$ (\texttt{agent\_reach/cap/ring\_cycle.txt}, re-verified cell by cell). Unrolled to the
half-line, column $i$ being cell $i\bmod728$, and started at the even phase $(18-N)\bmod14$, it evolves
as the ring does, its column~$0$ agreeing with the pinned column, and it is at phase $12$ on row $N-6$;
by (a) diagonal $N$ then carries a run of $6$. Direct simulation finds a run of exactly $6$ for every
even $N\in[8,3000]$ and for $N=5\,000$, $10\,000$, $20\,002$, $30\,000$. With (c) this is (i).
(e) \emph{Odd diagonals.} The pattern $110$ in columns $1,\dots,3$ of an even row has no predecessor:
column~$0$ is $1$ on the row above, so $a(\cdot,1)=1$ forces that row's columns $1,2$ to be $0$, then
$a(\cdot,2)=1$ forces its column~$3$ to be $1$, and then $a(\cdot,3)=1$, not $0$. Hence no odd diagonal
$N\ge5$ carries a run of $3$. The crystal of (d) carries $10$ in columns $1,2$ at phase $5$; started at
the even phase $(7-N)\bmod14$ it is at phase $5$ on row $N-2$, so every odd diagonal $N\ge5$ carries a
run of $2$ (exactly $2$ by simulation for every odd $N\in[5,3001]$ and for $N=10\,001$, $30\,001$;
\texttt{s3\_cap\_odd.py}). With $g(1)=1$ and $g(3)=3$ from the free row~$0$, this is (iii).
\end{proof}

\begin{remark}[what the cap says]\label{du:rem:cap}
(i) The witness in (d) is a temporally periodic half-plane of row period $14$ and spatial cycle $728$
whose interface is the gap word $(1\,4)^{\infty}$, the crystal certified in
Section~\ref{sec:record}, and it carries the longest straight chain of influence that remains available
once $N\ge58$. The opposite extreme, a block through which no influence passes, is the trap of
Proposition~\ref{du:prop:trap}, and it is not one of the crystals of Section~\ref{sec:record}: it is a
finite block of $20$ columns whose reachable set cycles with row period $8$ whatever lies beyond
column~$20$. The $1^{\infty}$ crystal (row period $4$, spatial cycle $7$) carries the trap's interface,
$(10)^{\infty}$ up to phase, but is not a wall: single flips at columns $21$--$90$ never reach its interface within
$4\,000$ rows, while random modifications of columns $21$--$120$ (each cell flipped with probability
$\tfrac12$) reach it in $54$ of $300$ trials, against $0$ of $300$ for the trap (\MEASURED,
\texttt{agent\_reach/proof/p13\_cycle7\_shield.py}). (ii) A far-right
perturbation that reaches the interface at a row $N\ge58$ therefore cannot arrive along a straight
chain; it must pause (a vertical step needs the right-neighbour gate to be $0$, as
$\partial\,(l\oplus(c\vee r))/\partial c=1\oplus r$), so its front moves left at speed strictly below
$1$. \MEASURED\ on extracted single-flip fronts: about $0.3$, with pauses and retreats, the bursts into
column~$1$ never exceeding the cap (ledger 10.220--10.222). (iii) The cap is a property of the $1$s in
the pinned column, not of Rule~30's diagonals: with column~$0$ free or all zero, straight zero diagonals
of length $80$ exist (\EXACT), and with column~$0$ pinned to $1^{\infty}$ or $(01)^{\infty}$ they do
not. (iv) We found no statement of this kind in print (ledger 10.224, every source opened): the
Lyapunov and damage-spreading literature for Rule~30 --- Wolfram's measured leftward speed
\cite[p.~949]{NKS}, the defect shapes of Baetens and Gravner \cite{BGrav}, whose Open Problem~9 asks
for a damage-spreading theory for periodic states --- concerns free configurations, and Gravner and
Griffeath's edge automaton \cite{GG} a zero edge. We did not find the cap stated, in the narrow form given here.
\end{remark}

\begin{remark}
\EXACT: for $w=01$ the inequality was checked on all $1770$ ordered pairs with $m+n\le60$ with no violation, which
also fixes the indexing of the count list ($|L_1|=2$, $|L_{44}|=21\,389$, $|L_{60}|=103\,220$; earlier
working notes labelled the same list from $k=2$, which is off by one and is corrected here). The sequence
$\frac1k\log_2|L_k|$ is monotone decreasing from $k=5$ on, through $0.519$ at $k=16$, $0.327$ at $k=44$ and
$0.278$ at $k=60$, so $h(01)\le0.278$ bits. The open question is therefore not whether the entropy exists
but whether this monotone sequence tends to $0$. Every growth law that fits the computed range predicts
that it does, and that conclusion is the one part of the fitting that is robust: for
$k^{1.9}\exp(0.87\sqrt k)$, $\frac1k\log|L_k|\sim1.9\frac{\ln k}{k}+\frac{0.87}{\sqrt k}\to0$, and the
same holds for $\exp(c\ln^2k)$ and for $k^{a}\exp(ck^{\gamma})$ at every $\gamma<1$. Which of them it is
is undetermined (Section~\ref{sec:record}), and no finite depth decides the limit in any case. The submultiplicative
bound itself is Meier and Staffelbach's observation for the analogous count in the finite-ring setting
\cite[eq.~(10)]{MS}, where they use it only to conclude that their attack cost grows at most exponentially.
\end{remark}

\begin{measured}
$|L_k(01)|=2,3,5,8,12,17,25,36$ for $k=1,\dots,8$; $|L_k(0000000101)|$ grows by a factor $1.55$ per period; $|L_k(0000011111)|$ by a factor tending to $2$ (\MEASURED, nine exact terms). For $01$, in run-length coding the alphabet is $\{1,2,3,4\}$, the pairs $(1,2),(2,1),(2,3)$ are forbidden in the recurrent part, and the minimal forbidden words are unbounded \cite[\S11]{Q}, so the language is not of finite type. It has at least $3\,591$ distinct follower sets (\EXACT, by complete enumeration), so any sofic presentation would need more than $3\,591$ states; that it is not sofic at all follows from the growth dichotomy for regular languages --- polynomial or exponential, never intermediate --- and is therefore conditional on the measured intermediate growth, exactly as stated in the list of what is not proved; the stream measurement of Remark~\ref{ex:rem:stream} points instead to exponential growth beyond the row-$0$ transient, under which that inference lapses and only the follower-set bound stands. An earlier draft wrote ``its follower sets are pairwise distinct'', which asserts an infinite statement that no enumeration reaches; that wording is \textsc{withdrawn} in favour of the bound. Exact counts for $01$ reach depth $60$ ($|L_{44}|=21{,}389$, $|L_{45}|=23{,}809$, $|L_{46}|=26{,}482$, $|L_{60}|=103{,}220$, ratio $|L_k|/|L_{k-1}|$ equal to $1.112$ at $k=46$ and $1.094$ at $k=60$, and falling), which is consistent with zero entropy for the words anchored at row~$0$ that dominate these counts (Remark~\ref{ex:rem:stream}); the only bound proved from these counts is $h(01)\le0.278$ bits. Its growth is superpolynomial on that range, and \emph{which} superpolynomial function it is remains \OPEN: no fixed degree fits, since the local exponent $\mathrm d\log|L_k|/\mathrm d\log k$ rises from $3.13$ at $k=12$ to $5.22$ at $k=56$ without levelling, and the phrase ``degree $5$--$6$'' used in earlier working notes reports that exponent \emph{on the window} rather than naming a family; it is withdrawn. Section~\ref{sec:record} gives the out-of-sample table and says exactly which families the data does and does not separate; two gap laws stated in earlier working notes ($444$ never; the $2^a42^b$ rule) were found false on re-enumeration and are withdrawn. Uncountability of $L(01)$ is not proved; positive entropy of $L(01)$ is \MEASURED\ on the stream of Remark~\ref{ex:rem:stream}, not proved. Subject to that, Proposition~\ref{ex:prop:boundary} places $01$, $001$, $0011$ and every word without a run of $\ge9$ ones or $\ge10$ zeros outside the reach of any argument that uses the seed only through Theorem~\ref{ex:thm:zero}.
\end{measured}

\begin{remark}[relation to the cryptanalysis of Rule 30]\label{ex:crypto}
Rule~30 was proposed as a keystream generator, and the attack literature studies exactly the objects of
this section under other names; we record the correspondence because several of our facts are theirs.
Meier and Staffelbach \cite{MS} write the left reconstruction as their equation~(4),
$s_{i-1}(t)=s_i(t+1)\oplus(s_i(t)\vee s_{i+1}(t))$, and observe that ``given the values of the cells in two
adjacent columns, this allows the values of all cells in a triangle to the left of the temporal sequence to
be reconstructed'', which they call \emph{completion backwards} (the same equation and, nearly verbatim, the same sentence are already in Wolfram's CRYPTO~'85 abstract \cite{W85}, which credits the approach to C.~and R.~Feynman); their conclusion that ``knowledge of the
seed is equivalent to knowledge of one of its adjacent sequences'' is the reduction used throughout this
section, and their \emph{right adjacent sequences} are our interface words, so $L_k(w)$ is their count.
Three further items match. Their \emph{completion forwards} statement, that any choice of the sites to the
right has a completion consistent with the given temporal sequence, is the absence of dead ends in $L(w)$.
Their proof that for an all-zero temporal sequence the right adjacent sequence ``can only consist of $k$
digits $0$ followed by $n-k+1$ digits $1$'', so that the count grows only linearly, is the right-hand
companion of the constant-trace fibre classified from the left by Condrey \cite{Condrey}, and predates it by
thirty-five years. And their Table~1 is not merely comparable to our enumeration, it is the same computation: their ``completion forwards'' is a half-line Rule~30 driven on the left by a prescribed boundary column, which is the object we enumerate, and all sixteen of their published means were reproduced exactly here (a first attempt with the wrong boundary convention returned $2.7$ where they report $41.86$, which is how the correspondence was pinned). Their averages are taken over random boundary columns, whereas $L_k(w)$ fixes the column, so the two are the ensemble and the single-word versions of one quantity; the local exponent of their means rises linearly in $\ln N$ with no bend, the same signature our own counts show, so the right-admissible count is not polynomial in either setting. The same literature also has the leftward spreading speed: Meier and Staffelbach note, after Wolfram, that the effect of a bit change propagates to the left with speed roughly $\tfrac14$, and to the right with speed $1$ by the partial linearity of the rule \cite{MS}, and Wolfram's CRYPTO~'85 extended abstract \cite{W85} gives these as Lyapunov exponents, $0.25$ on the left and $1$ on the right; so the leftward speed whose values in \cite{NKS} are discussed in Remark~\ref{fr:slopebar} was published well before \cite{NKS}.  Finally, the $\tfrac34$ that recurs in our measurements --- the rate at which the interface bit
is the complement of the previous centre bit --- is the agreement probability of the best affine approximation
to the Rule~30 gate, $\alpha=\tfrac34$ in Ko\c{c} and Apohan \cite{KA}, and the reason their attacks work is
that the gate lacks first-order correlation immunity and has low nonlinearity \cite{Mariot}. None of that
literature addresses eventual periodicity: its complexity statements are experimental, with no closed-form
exponent.
\end{remark}


The reach of Theorem~\ref{ex:thm:A} is thus, in principle, the set of words whose interface language contains only eventually periodic sequences (Section~\ref{ex:sec:boundary}); zero entropy is necessary for that but not sufficient (Remark~\ref{ex:rem:branching}), and Theorems~\ref{ex:thm:zeroruns}--\ref{ex:thm:const} show that this set contains the families listed there, each of which has a sufficiently long run of one symbol. The words with no long run---those whose interface can be driven freely by the right context---are the ones that carry the difficulty of the Prize Problem, and for them the seed must enter through something other than Theorem~\ref{ex:thm:zero}. The Rauzy graphs of the exact interface languages of ten short residual words ($01$, $0011$, $001$, $011$, $0001$, $0111$, $00011$, $00111$, $01011$, $00101$), aligned to the period, branch at every depth computed, $24$--$26$ rows, with Perron ceilings above $1$ throughout and, for all but $00011$, descending without settling, so none is excluded by the cycle criterion at that depth and each shows the fingerprint of $01$ (\EXACT; ledger 10.236d~E).

\begin{remark}[the criterion is an $\omega$-regular emptiness test, and what that buys]\label{ex:omega}
The exclusion criterion has a standard logical form, and naming it converts a word-at-a-time computation
into a family-at-a-time one.

Read the $K$-window subset automaton as a B\"uchi automaton over the interface alphabet: the criterion
of Lemma~\ref{ex:lem:cycle} asks whether some infinite path has a label sequence that is not eventually
periodic, which is an $\omega$-regular emptiness question. So a \emph{regular family} of candidate
period words can be handled in one finite computation, by taking the product of the subset automaton
with a deterministic safety automaton recognising the family and applying the same
$\#\text{edges}=\#\text{vertices}$ test to the product --- though not literally: whenever the family automaton
must allow a choice of letter, as it must for the zero-run family, the product also branches on that choice, and no
component is then a single cycle (see below). That is the route by which a family threshold of $17$ for Theorem~\ref{ex:thm:19} was
first proposed. On every necklace tested the lowering from $19$ to $17$ holds word by word (\EXACT), but
the family-level product certificate at $17$ has not been reproduced, as recorded below, so the family statement at $17$ is \OPEN.
Lemma~\ref{ex:lem:transfer} is absorbed into it as a one-state family automaton.

A sweep of the other families, and one methodological fact that governs all of them, follows below.
The construction was validated three ways before being relied on: the window map was unit-tested
against direct Rule~30 evolution; the published state and edge counts of Sections~\ref{ex:sec:window}--\ref{ex:sec:graph} were reproduced,
including $|F|=62$ at step $13$; and the family automaton was compared against the word-by-word graph
$G_w$ on all $801$ necklaces of length at most $12$ at $K=8$ and $K=10$, with no disagreement. An
independent implementation of the criterion, working at the letter level rather than one edge per
choice of all interface bits in a period, confirms the boundary directly: at $K=14$ every necklace we
tested with all zero-runs at least $17$ is excluded, while at zero-runs $16$ the words with $1$-run $3$
and with the pair $(3,3)$ are not.

Two things came out of re-running the criterion across the other families of Part~I, and the first is
methodological. \emph{The criterion is not monotone in $K$.} A larger window is a finer
over-approximation, so states that were identified become distinguished, and a component that was a
single cycle can acquire branching; the word $0^{16}1^3$ is not certified at $K=12$ or $K=14$ and is
certified at $K=16$, while $0^{15}1^3$ is certified at $K=14$ but not at $K=12$. Since a certificate at
any single $K$ is a complete proof --- the soundness lemma holds for every $K$ --- the right procedure
is to take the smallest certifying $K$ and never to read a failure at one $K$ as final. Every threshold
below is therefore an upper bound on the true one.

The sweep itself. For the zero-run family, individual necklaces with all cyclic zero-runs equal to $Z$
and $1$-runs up to length $5$, together with the pairs $(a,b)$ for $a,b\le3$, are all certified at
$K=16$ down to $Z=15$; at $Z=14$ the shape $1^3$ and the pair $(3,3)$ resist, and they still resist at
$K=18$. What is independently reproduced and what is not should be separated here. The per-word results above
are ours and were obtained from a second implementation: at $K=16$ every necklace we tested with all
zero-runs at least $15$ is certified, and the obstruction at $14$ survives $K=18$, so it appears
genuine rather than a limit of computation. The \emph{family}-level statement --- one certificate
covering the family with $1$-runs unbounded, which is what would lower the certified threshold from $19$ to $17$ ---
rests on a product construction we have not been able to rebuild. Two reconstructions of it fail: the
single-cycle test cannot apply to the product, because the choice of letter branches and so no
component is a single cycle; and the forcing test on the product's recurrent part leaves a residue
that does not vanish, though it is strikingly small and degrades smoothly, one unforced state at
$Z=18,19,20$, then $7$, $45$ and $3206$ at $Z=17,16,15$ with $K=14$. So our reconstruction differs
from the reported one in a way we have not identified; until the product certificate is rebuilt, the
family statement at $17$ is \OPEN\ and the certified family threshold is the $19$ of \texttt{cover14.py}.
The two recorded holes do not move: $0^711$ and $1^80$ are uncertified at every $K$ up to $16$. Two
smaller facts worth recording because they show the families are not intervals: $1^70$ is certified at
$K=10$ although the published family begins at $b\ge9$, and $0^71$ is certified at $K=12$ while $0^81$
is not, so the isolated-ones family is not monotone in the zero-run length either.

The ceiling of the method is worth stating in the same language, because it explains why the family
route cannot be pushed further. The criterion is expressible as a sentence of monadic second-order
logic over one successor, but its \emph{conclusion} --- ultimate periodicity of the interface sequence
--- is not $\omega$-regular. So the method can quantify over a regular family of candidate periods, and
can never quantify over all periods at once. The gap between the two is exactly the gap between
Part~I's theorems and P1.
\end{remark}

\section{Reproduction}\label{ex:sec:repro}

All certificates are computations with a stated input and a yes/no criterion.
\begin{itemize}[nosep]
\item Lemma~\ref{ex:lem:quench}, finite facts: \texttt{drive2.py}, \texttt{freeze.py}, \texttt{cross.py} (Glucose~4 cross-check); requires \texttt{python-sat}.
\item \S\ref{ex:sec:spine}: \texttt{probe.py} with \texttt{satL.py} (prints $|L|=10,21,31,38,47,54$ at depths $1$--$6$, and checks all five bounded-window facts on every admissible $\tau$-sequence at each depth: $11$, $2221$ and $21221$ never occur, $1221$ only at positions $\le1$, symbols outside $\{1,2\}$ only at positions $\le2$ --- all five hold at every depth to $6$, and the $38$ depth-$4$ sequences are printed in full).
\item Theorems~\ref{ex:thm:1b0}--\ref{ex:thm:0m11}: \texttt{python s3\_fam1b0a.py 10 1 19 7 34}; \texttt{python s3\_thm9.py 10 29 1 12}; \texttt{python kauto3.py <word> 10 <m>}; independent check \texttt{python kwin.py <word> \ldots} (Appendix), which prints the state and edge counts and the failing components if any.
\item Theorem~\ref{ex:thm:19}: family-level certificate \texttt{cover14\_19\_48} ($K=14$, fixed-point coverage over $1$-runs $\le48$), reproduced in full on 2026-09-17 by \texttt{cover14.py 14 19 48} (log \texttt{cover14\_19\_48.out}, $53$ minutes): the zero-run map from the full window set reaches its fixed point at step $28$ with $|F|=556$, the $1$-run map at step $47$, and for every $1$-run length $a=1,\dots,48$ the next zero-run is forced for all $m_1,m_2\ge19$ and every $b=1,\dots,48$, with no undetermined cell in any of the $48$ sweeps. Independently re-verified at $K=14$ by the combinatorial criterion of Section~\ref{ex:sec:graph} (\texttt{K=14 python kwin.py \ldots}) on the necklaces $0^{19}1^b$ for $b=1,2,3,4$ ($42$/$68$, $47$/$72$, $56$/$89$, $44$/$71$ states/edges), $0^{19}\,1\,0^{20}\,11$ ($18$/$44$) and $0^{19}\,111\,0^{19}\,101$ ($32$/$65$): every recurrent component a single cycle.
\end{itemize}
Every word run in more than one implementation received the same verdict in each.

\section{What is not proved here}
\begin{enumerate}[nosep]
\item Theorem~\ref{ex:thm:zero} is cited, with its proof reproduced. Theorem~\ref{ex:thm:19}'s family-level certificate at zero-runs $\ge19$ (uniform in the $1$-run lengths) is reproduced by the code shipped here (\texttt{cover14.py}, Section~\ref{ex:sec:repro}), so it no longer rests on an unpublished working record; the six representative necklaces remain as an independent cross-check by the combinatorial criterion of Section~\ref{ex:sec:graph}. The extension to $17$ is reproduced word by word only (\texttt{s3\_threshold17.py}). It is used above only for tails with a run $\ge19$, which are also covered directly for the computed ranges.
\item The lower bounds in Theorems~\ref{ex:thm:1b0}, \ref{ex:thm:1b0a}, \ref{ex:thm:0m11} are the smallest values at which the whole family is certified; individual smaller words ($1^70$ at $K=10$, $0^71$ at $K=12$) are certified separately, so the families are not intervals, and ``not certified at $K$'' is not ``not excluded''. Since the criterion may be applied to any rotation of a necklace (Lemma~\ref{ex:lem:path}), every uncertified word above was tested in all its rotations.
\item Uncountability of $L(01)$ and its siblings is not proved, and Proposition~\ref{ex:prop:boundary} is stated under exactly that hypothesis --- not entropy --- which suffices for its conclusion but is not necessary for it (Remark~\ref{ex:rem:infimum}). The distinction is not editorial. The proposition's sufficient condition, exponential growth of $|L_k|$, is \emph{unavailable} here rather than merely unproved: the measured ratio is $1.094$ at $k=60$ and falling, which is consistent with zero entropy, so counting cannot deliver the hypothesis and uncountability must be reached by some other route. An earlier version of this item read ``positive entropy of $L(01)$ and its siblings is measured, not proved''; that was \textsc{withdrawn} as contradicting the exact counts recorded above and as misnaming the hypothesis of the proposition it cites. The second objection stands; the first does not, since the stream measurement of Remark~\ref{ex:rem:stream} shows the counts describe the row-$0$ transient: positive entropy of $L(01)$ is now \MEASURED\ (on random right halves, not on the counts, which the row-$0$ transient dominates), the zero-entropy expectation is withdrawn, and uncountability is still not proved.
\item Nothing here bears on the single seed beyond what holds for every number-like configuration, and nothing here excludes any word without a long run.
\end{enumerate}



\section{Note added: the constant fibres and the period-two horizon}\label{ex:condrey}

Condrey \cite{Condrey} (arXiv, September 2026) classifies the two constant-trace fibres exactly: for a prescribed right half the unique left half with identically zero centre is an alternating tail selected by the first one of the right half (a prefix-OR of the right half), and with identically one centre it is one universal checkerboard; every nonzero member of either fibre has infinitely many ones, so no nonzero finite configuration has an eventually constant column, and for support radius $w$ the sharp maximal constant prefix has length $w+2$ with $2^w$ or $2^w-1$ extremisers by parity of $w$. Theorem~\ref{ex:thm:const} above gives the eventual-constancy exclusion for the wider class of number-like configurations by the two-column theorem, and its one-case is attributed by \cite{Condrey} to Jen \cite[Thm.~7a]{Jen86}, a citation we have not been able to verify at source and which should be treated as second-hand; the explicit fibres and the sharp horizons are Condrey's and are not contained here. Other public statements that neither constant is an eventual period, at least two of them earlier than \cite{Condrey}, are listed in Remark~\ref{ex:prior}(b). His Lemma~1 is Proposition~\ref{ex:prop:surj} in its finite form.

Condrey defines $H(p,w)$, the sharp maximal length of an eventually $p$-periodic centre prefix over rows of support radius $w$, notes $H(2,w)\ge w$ from left permutivity, and leaves period two open, which is the target of Part~I. We computed the period-two horizon exactly by exhaustion over all $2^{2w+1}$ rows (\EXACT; prefix lengths, both phases of $01$ counted):
\[
\begin{array}{c|cccccccccc}
w & 1&2&3&4&5&6&7&8&9&10\\ \hline
\text{alternating prefix length} & 7&7&7&7&9&10&10&15&17&17\\
\text{extremisers} & 1&1&1&2&7&11&67&3&1&2
\end{array}
\]
so $H(2,w)=6,6,6,6,8,9,9,14,16,16$ in his index; the same exhaustion reproduces his constant-case table exactly (lengths $3,3,5,5,7,7,9$ and $2,4,4,6,6,8,8$ with counts $2^w-1$ and $2^w$). The period-two horizon runs about six to seven ahead of $w$ with an irregular staircase and no sign of a closed law; the one-sided version of the same quantity (left support bounded, right half free) is the kill instrument of Section~\ref{sec:record}, whose excess over $w$ was measured to grow like $\log w$.

\part{The Duhamel class and its decimation}

\section{The Duhamel formula and the \texorpdfstring{$2$}{2}-kernel}\label{du:prelim}

Write $U_t(z)=\sum_x u^t_x z^x\in\F[z^{\pm1}]$ for row $t$, $\Lambda(z)=z+1+z^{-1}$ for the symbol of Rule~150, and $Q_t(z)=\sum_x u^t_xu^t_{x+1}z^x$ for the census of adjacent pairs. Let $T_m(z)=\Lambda(z)^m$ and $T_m(x)=[z^x]T_m(z)$; $T_m$ is the row $m$ of Rule~150 from the seed, and $T_m(x)=T_m(-x)$.

\begin{lemma}[\KNOWN, elementary]\label{du:kernel}
$T_{2m}(z)=T_m(z^2)$ and $T_{2m+1}(z)=\Lambda(z)\,T_m(z^2)$. Moreover $T_m(0)=1$ for every $m\ge0$.
\end{lemma}
\begin{proof}
In characteristic~$2$, $\Lambda(z)^2=\Lambda(z^2)$, which gives the first two identities. For the third, $T_m(0)=[z^m](1+z+z^2)^m=\sum_{b}\binom{m}{2b}\binom{2b}{b}$ (choose the $b$ factors contributing $z^2$ and the $m-2b$ contributing $z$), and $\binom{2b}{b}$ is even for $b\ge1$ (the binary addition $b+b$ carries), so $T_m(0)\equiv\binom m0=1$.
\end{proof}

\begin{theorem}[Duhamel formula; \PROVED]\label{du:duhamel}
For all $t\ge0$,
\[
U_{t+1}(z)=\Lambda(z)\,U_t(z)+Q_t(z),\qquad\text{hence}\qquad
U_t=\Lambda^tU_0+\sum_{s<t}\Lambda^{t-1-s}Q_s ,
\]
and, reading the constant term with $\delta_x=[x=0]$ and $q^s_x=u^s_xu^s_{x+1}$,
\[
c_t=\langle T_t,\delta\rangle+\sum_{s<t}\langle T_{t-1-s},q^s\rangle
   =1+\sum_{s<t}\sum_xT_{t-1-s}(x)\,u^s_xu^s_{x+1},
\qquad \langle a,b\rangle=\sum_xa_xb_x .
\]
In the notation of Section~\ref{du:decimation}, $c=c[\delta,q]$.
\end{theorem}
\begin{proof}
The rule is $u^{t+1}_x=u^t_{x-1}+u^t_x+u^t_{x+1}+u^t_xu^t_{x+1}$; the three linear terms are $\Lambda U_t$ and the quadratic term is $Q_t$. Iterate. The constant term of $\Lambda^mF$ is $\sum_xT_m(-x)F_x=\langle T_m,F\rangle$ by symmetry of $T_m$, and $\langle T_t,\delta\rangle=T_t(0)=1$ by Lemma~\ref{du:kernel}. Verified numerically for $t\le200$ (\EXACT).
\end{proof}

Rule~30 is thus Rule~150 driven by the source $q$; the same identity with the Rule~90 symbol $z+z^{-1}$ and the source $u^s_x(1+u^s_{x+1})$ (the right-hand ends of black runs) gives $c_t$ as the parity of those sites on the Pascal support (\PROVED, ledger~10.68). Both are the same statement in different linear bases.

\subsection*{Cartier operators, the \texorpdfstring{$2$}{2}-kernel and automaticity (\KNOWN)}
For $G(y)=\sum_tg_ty^t\in\F[[y]]$ put $\Lambda_rG=\sum_tg_{2t+r}y^t$ ($r=0,1$). The \emph{$2$-kernel} of a sequence $g$ is $\mathcal K(g)=\{(g_{2^km+r})_{m\ge0}:k\ge0,\ 0\le r<2^k\}$; its elements are the coefficient sequences of $\Lambda_{r_{k-1}}\cdots\Lambda_{r_0}G$ over binary words. A sequence is $2$-automatic iff its $2$-kernel is finite (Eilenberg \cite{Eilenberg}), iff its generating series is algebraic over $\F(y)$ (Christol \cite{CKMFR}). An eventually periodic sequence is $2$-automatic. Hence \textbf{an infinite $2$-kernel of $c$ implies P1}; this is the automaticity route, and it is strictly stronger than P1.

\begin{proposition}[the uniform-measure two-point function lives on one ray; \PROVED]\label{du:tworay}
Let $F$ be any left-permutive radius-one rule and let $\mu$ be the uniform Bernoulli measure. Then the
two-point function $C(x,t)=\mathbb E_\mu\bigl[(-1)^{u^0_0+u^t_x}\bigr]-\mathbb E_\mu[(-1)^{u^0_0}]
\mathbb E_\mu[(-1)^{u^t_x}]$ vanishes except on the single ray $x=t$. In particular $C(0,t)=0$ for every
$t\ge1$: in the measure in which all of this machinery operates, the centre column is exactly
uncorrelated with the seed.
\end{proposition}
\begin{remark}[the balance itself is \KNOWN]\label{du:rem:balance}
The balance $\Pr_\mu(u^t_0=1)=\tfrac12$ for every $t\ge1$ under every left-permutive rule is not the special case $C(0,t)=0$ but a different statement: it is the one-point identity $\mathbb E_\mu[(-1)^{u^t_0}]=0$, while $C(0,t)=0$ is the two-point identity $\mathbb E_\mu[(-1)^{u^0_0+u^t_0}]=0$. Both follow from $u^t_0=u^0_{-t}\oplus g_t(u^0_{-t+1},\dots,u^0_t)$ with $u^0_{-t}$ uniform and independent of the pair $(u^0_0,g_t)$. For Rule~$30$ the balance is classical: the rule is surjective \cite[Thm.~6.6]{Hedlund}, so by the balance property \cite[Thm.~5.4]{Hedlund} it preserves $\mu$ (Section~\ref{sec:p2}) and $u^t_0$ has the law of $u^0_0$. It is also stated as Theorem~3 of Chan-L\'opez and Mart\'in-Ruiz \cite{CLMR}, but the XOR lemma their proof rests on, that $X\oplus Y$ is a fair bit for every $Y$ even when $Y$ depends on $X$, is false (take $Y=X$); nothing in the balance is new here. What Proposition~\ref{du:tworay} adds is the full two-point function. Its proof uses only that a radius-one rule bijective in its left argument is $l\oplus h(c,r)$ over $\F$, which is all that ``strong'' left permutivity (defined in the proof of Theorem~\ref{du:thm:degree}) means at radius one.
\end{remark}
\begin{proof}
Two lines from strong left permutivity: for $x<t$ the cell $u^t_x$ is, for each fixed assignment of the
other coordinates in its light cone, an exclusive-or with the leftmost coordinate, so summing over that
coordinate cancels the expectation; for $x>t$ the two cells are independent.
\end{proof}

\begin{remark}[why the solvable-circuit analogy must not be made]\label{du:nodual}
The exactly solvable circuit literature has correlations supported only on the light cone, computed by
a transfer operator in the space direction, and the resemblance to Corollary~\ref{du:cor:C} is a trap.
Their confinement is measure-theoretic, following from preservation of the maximal-entropy measure in
both directions together with tracelessness of the observable; ours is an identity over $\F$, valid
orbit by orbit. Proposition~\ref{du:tworay} is the decisive separator: theirs is a statement in the
uniform Bernoulli measure, and in that measure our centre column carries no signal at all. The class is
also empty of informative members --- the elementary rules satisfying the classical hypothesis are the
shifts, whose correlation is frozen, and the linear rules, whose output vanishes --- and the values of
$C$ on the ray $x=t$ for Rule~30 satisfy no linear recurrence of order at most $5$ over $\Q$, where a
dual-unitary qubit circuit forces order at most $3$.

One item in that literature transfers as a target rather than a theorem, and we record why it does not
close either. Influence-solvability in the sense of H\"ubner and Kim \cite{HubnerKim} is characterised by necessary
and sufficient local conditions (their \S III), is stated there for reversible gates with an extension to non-reversible
ones that they call straightforward (\S II.B), and its criterion is a condition on the width-two trace measure; the
deciding question is whether that trace subshift is sofic. Computing its exact language to horizon $11$
--- $38\,616$ realisable traces of $4\,194\,304$ initial windows; the numbers of realisable pairs of adjacent columns are \KNOWN, listed for $t=1,\dots,17$ in \cite[p.~1087]{NKS}, with sampled means to $N=17$ in \cite[Table~1]{MS} and counts to $13$ in \cite{Dib} --- and counting distinct sets of
length-$m$ continuations over contexts of length $k$ gives $4,7,8,11,14,16,18,19$ at $m=3$, then
$4,8,11,18,28,37,47$ and $4,10,16,31,54,89$ at $m=4$ and $m=5$. The flattening of the shallow rows is
not evidence of soficity and cannot be: for fixed $m$ the count is non-decreasing in $k$ and bounded by
$2^{2m}$, so it saturates in \emph{every} subshift, sofic or not, and here the $m=3$ row is still rising
at the deepest context available. What the probe yields is a rising sequence of \emph{lower} bounds on
the number of follower sets --- $7$ at $m=2$, and at least $19$, $47$ and $89$ at $m=3,4,5$ --- and
never an upper bound, while soficity requires finiteness. It is therefore the obstruction of
Remark~\ref{du:censuscaveat} in another guise.
\end{remark}

\begin{remark}[the same object over $\Z_2$, and why the two literatures never meet]\label{du:padic}
Part~II works in $\F[[y]]$, where Christol's theorem lives. The same recursion has a second natural
home, and the comparison explains a gap that is otherwise puzzling.

Read a configuration as a $2$-adic integer and pass to the co-moving frame. Rule~30 becomes
$H(a)=a\oplus(2a\vee4a)$, whose $n$-th output digit depends only on digits $0,\dots,n$ of the input ---
a $1$-Lipschitz map in the sense of \cite{Anashin} --- and the centre column is the \emph{diagonal}
read $c(t)=\mathrm{digit}_t(H^{t}(1))$, which we verified against simulation for $t<3000$. The map is
measure-preserving, being a bijection mod $2^{n}$ for every $n\le16$, and left-permutivity is why: the
opposite orientation $G(a)=4a\oplus(2a\vee a)$ is $1$-Lipschitz but not measure-preserving, its image
mod $2^{16}$ having $18\,869$ of $65\,536$ elements.

Two things follow, and the first corrects a reflex. One is tempted to dismiss ergodic statements here
because the seed is a single point of measure zero; for a $1$-Lipschitz map on $\Z_2$ that is wrong,
since ergodicity is equivalent to transitivity mod $2^{n}$ for every $n$, a statement about every orbit.
The route fails for sharper reasons: $H$ is not ergodic at all, as $H(0)=0$ makes $2^{n}\Z_2$ invariant;
and ergodicity would control each \emph{fixed} digit, that is each fixed diagonal, never
$\mathrm{digit}_t$ at time $t$. The diagonal read is the obstruction, as it is throughout this paper.

The comparison is the useful part. Rule~90 is $(1+y^{2})\cdot a$, natural in $\F[[y]]$ and in $\Z_2$
alike, having neither carries nor an OR; Rule~30's OR is natural only over $\Z_2$, while Christol's
theorem needs characteristic~$2$ and no carries. Each literature holds half of our object and neither
holds it whole, which is why no result crosses between them. For completeness, the $2$-adic complexity
of the centre column is finite exactly when the sequence is eventually periodic, so it restates
Problem~1.
\end{remark}

\section{Certified bounds}\label{du:bounds}

A \emph{Mahler relation of depth $n$ and degree $D$} for $G=\sum c_ty^t$ is $\sum_{i=0}^np_i(y)G(y^{2^i})=0$ with $\deg p_i\le D$, not all zero; $G$ is algebraic iff such a relation exists for some $(n,D)$. Its existence is the vanishing of a null vector of an explicit $\F$-matrix built from $c$, and full rank of a square submatrix on a stated set of rows is a certificate of non-existence for that $(n,D)$.

\begin{theorem}[\EXACT]\label{du:certified}
For the centre column of Rule~30:
\begin{enumerate}[nosep]
\item no Mahler relation of depth $n\le12$ and degree $D\le8000$ exists (\cite[Thm.~9]{Auto}, where each determinant is computed twice, by a row-streaming elimination and by an elimination of the transposed matrix; re-certified by independent code to $D\le4000$ at $n\le12$, with certificates \texttt{ag2\_cert\_n<n>\_D<D>.txt} each re-verified from its file by a transposed elimination, the box $(12,8000)$ itself not being re-run; controls Thue--Morse $(n,D)=(2,4)$, paperfolding $(2,4)$, rational and Rule~90/150 columns all recovered at their known $(n,D)$);
\item the truncated kernel elements are pairwise distinct wherever they have been compared. From $n$ computed terms the maximal window at level $k$ is $W_k=2^{\,n-k}$ simultaneously for every residue (\PROVED), so the census is a clean $(k,W)$ grid; on that grid every computable cell with $W\ge32$ is injective, and all non-injectivity sits in the columns $W\le16$, where it is forced by pigeonhole and is therefore a truncation artefact rather than a kernel coincidence. At $n=21$ the map $(k,r)\mapsto$ element is injective on every pair with $\max(k,k')\le15$, which is $2{,}147{,}385{,}345$ comparisons at window length at least $64$, and the $65\,536$ residues of level $16$ are pairwise distinct on $32$ terms; of the $6{,}442{,}352{,}640$ pairs whose higher level is $16$ exactly one is unresolved, against an independence expectation of $1.5$, and the closest separated pair agrees on $31$ terms and differs at the $32$nd. The sharpest quantity the prefix supports is not the count of a single level but the largest antichain in the prefix trie, two elements being provably distinct exactly when neither prefix extends the other. That quantity has a name: it is Shallit's automaticity function, whose second definition \cite[\S1]{Shallit4} is the maximum number of pairwise $n$-dissimilar sequences in the $n$-truncated $k$-kernel, with dissimilarity meaning that neither is a prefix of the other, and it comes with the same Dilworth optimality remark. So the statement to make is not a bound but an exact value: $A^c_2(2^{20}-1)=80\,660$ and $A^c_2(2^{21}-1)=138\,672$, computed by the algorithm of \cite{Shallit4}, hence at least $138\,672$ states for a least-significant-digit-first automaton and at least $18$ for a most-significant-digit-first one. The order is maximal: $A^s_k(n)=O(n/\log_k n)$ for \emph{every} sequence, and our value sits at $1.389\,n/\log_2 n$, so no computation on a prefix can improve it by more than a constant factor. The figure $1{,}048{,}486$ of \cite{Auto} is $2^{20}-1-89$, that document's own census at $2^{24}$ terms with $89$ residues dropped, a count its method predicts at $96$ under independence; the method is sound but needs $2^{24}$ terms, roughly $46$ hours here, and is not re-derived.
\end{enumerate}
\end{theorem}
\begin{remark}[the census is capped; \PROVED]\label{du:rem:census}
By item~(2), from $2^{\nu}$ terms a level-$k$ kernel element is seen through a window of length $2^{\nu-k}$, and injectivity is certified only at $W\ge32$; so the census can certify at most $\sum_{k\le\nu-5}2^{k}=2^{\nu-4}-1$ elements from $2^{\nu}$ terms, $131\,071$ at $\nu=21$ --- about $n/16$. The antichain bound $138\,672$ of item~(2) is a different object (pairwise incomparable truncations, certified by prefix comparison rather than by window) and is not subject to this cap.
\end{remark}
\begin{remark}[what Christol's theorem gives effectively, and why the conversions are needed]\label{du:christoleff}
Since Part~II converts kernel bounds into exclusions on an algebraic presentation, it matters how much
the classical theorem gives by itself. Reading the original \cite{CKMFR}: the equivalence between
$2$-automaticity and algebraicity is the only clause needing $p$ prime, the equivalence with the finite
$k$-kernel is Eilenberg's \cite{Eilenberg} and is base-agnostic, and $q=p^{s}$-automatic coincides with $p$-automatic. The
effective content, however, is \emph{doubly exponential} --- their \S7 bounds the kernel by
$q^{s(N+1)}$ with Ore height at most $2hq^{d}$ --- so running the theorem backwards from an algebraic
presentation to a state count gives only a doubly logarithmic exclusion. That is why the sharper
conversions of Bridy and of Adamczewski and Yassawi are necessary rather than merely convenient, and it
also locates the binding constraint: the bounds $hd\ge18$, $(h+1)d\ge17$ and $h+d+g\ge19$ are limited
by our census of $138\,672$, not by anything in Christol. One attribution to correct while we are here:
the degree bound $q^{m}-1$ that we took from Bridy is already in \cite[\S6]{CKMFR}.

The same paper's later sections supply a concrete exclusion we can state. For a generalised Morse
sequence $\mu(u)$ the minimal homogeneous Mahler relation has order $2$ and degree $2^{\ell}+2$, which
at $\ell=1$ is $(2,4)$ and reproduces the Thue--Morse relation we use as a control. Theorem
\ref{du:certified}(1) therefore already excludes $c=\mu(u)$ for all $527\,345$ non-trivial generalised
words of length at most $12$ --- a named published family that contains both Thue--Morse and
Rudin--Shapiro. It also exposes a weakness in our control set, which we record rather than repair here:
the non-rational controls we ran reach only degree $82$ --- Thue--Morse and paperfolding sit at $(n,D)=(2,4)$, and the off-centre Rule~$90$ and Rule~$150$ columns at offset $k$, leading zeros stripped, need depth $2$ and degree $D=4k$ and $D=4k+2$ respectively (\EXACT\ at $k=1,2,3,5,7,14,20$, each relation verified on $10^4$ terms) --- while the scan claims to reach degree $8000$, and
\cite{CKMFR} supplies an infinite depth-two ladder at degrees $4,6,10,18,\dots,4098$ with closed-form
minimal polynomials, which is the right family to test against.
\end{remark}

\begin{remark}[a route, with the hypothesis verified as far as we can compute]\label{du:shallitroute}
The same paper supplies the only technique the literature has for proving an automaticity lower bound on
an \emph{explicit} sequence, and it converts our census into a hypothesis rather than an endpoint.
Shallit's separation lemma \cite[Lemmas 3, 12]{Shallit4} needs, for every level $i$ and every pair of
distinct residues $a\neq b$ below $2^i$, a witness $m\le M(2^i)$ with
$c(2^im+a)\neq c(2^im+b)$; if $M$ is polynomially bounded the sequence is not $2$-automatic, and hence
by Christol's theorem $G$ is transcendental over $\F(y)$ and $c$ is not eventually periodic. Our census
\emph{is} that hypothesis, verified to $i=16$ with the constant witness bound $M=32$, far inside the
polynomial regime; with $M=O(\log r)$ the same proof reaches the universal ceiling $\Omega(n/\log n)$.
Two calibrations keep this honest. The witness in every published application comes from an infinite
arithmetic theorem, never from data --- the examples run through the primes, the squarefree numbers, the
Thue--Morse and Sturmian words --- and the best bound proved for any explicit non-automatic sequence is
$n^{2/5}$, some $320$ times below what we measure. And a decision procedure does exist for one class:
Krawczyk and M\"ullner \cite[Thm.~B]{KrawczykMullner} decide automaticity for nonperiodic uniformly recurrent substitutive sequences by
a finite eigenvector check on a return substitution. One correction to how we first read that route:
the missing input is \emph{not} simply ``a substitution for $c$''. Christol, Kamae, Mend\`es France and
Rauzy's substitution clause \cite{CKMFR} is for \emph{uniform} substitutions and is Cobham's theorem, so
exhibiting a uniform substitution would make $c$ automatic --- the opposite of what we want. Any useful
substitution here must be non-uniform, and then it is not Cobham that applies.
\end{remark}

\begin{remark}[what the census does and does not evidence]\label{du:censuscaveat}Controls fix the honest reading. The same census returns exactly $2$ on Thue--Morse and exactly $4$ on the paperfolding sequence, so the implementation is sound; but a periodic sequence of period $12{,}345$, whose $2$-kernel is finite, scores $47{,}086$, which is larger than any figure above, and two independent coin sequences score within $0.24\%$ of Rule~30's value. A kernel census therefore measures how much counting room a prefix allows, not aperiodicity, and a large count is not evidence for an infinite kernel. The limitation is structural rather than a matter of taking more terms, and it has a one-line witness: the sequence that agrees with $c$ on the first $2^{20}$ terms and is zero thereafter is ultimately periodic, hence $2$-automatic with a finite $2$-kernel --- computed at $107\,047$ classes --- and it agrees with $c$ on every term any census of that prefix can read. Since every finite word extends to an automatic sequence, no measured value, growth rate, or proximity
to the universal ceiling can imply non-automaticity. The right name for this is elementary and we use
it rather than a grander one: the ultimately periodic sequences are dense in the Cantor topology on
$\{0,1\}^{\N}$, since $w\cdot0^{\infty}$ lies in the cylinder of $w$. It is not a profinite statement
--- the profinite topology on the free monoid is discrete, and profinite closeness and agreement on a
prefix are incomparable notions. This is \KNOWN, under the name \emph{diversity}: Shallit \cite{Shallit4} proves that almost all binary sequences are maximally diverse, warns that ``it is not so easy to prove that any individual sequence has the maximally diverse property'', and closes by saying he knows no explicit example of a diverse sequence with diversity measure $O(\log n)$ --- which is exactly the wall met here, stated in 1996. The one part that does appear to be unrecorded is the quantitative control: a periodic sequence can score high because for odd period $p$ one has $|\mathcal K|\le p\cdot\mathrm{ord}_p(2)$, which accounts for the $47\,086$ of the period-$12\,345$ control. Moreover the certified family is capped at $2^{\nu-4}-1$ elements from $2^{\nu}$ terms, about $n/16$ (Remark~\ref{du:rem:census}), so no census of a prefix can produce the one-parameter family of distinct kernel elements that the criterion of Allouche, Shallit and Yassawi requires. The census route to P1 is closed; what survives is the exclusion of low-complexity algebraicity by the Mahler certificates.\end{remark}

\noindent Source: \cite{Auto} and the independent re-run \texttt{agent\_certificate.md} (2026-09-16). The positive control set is audited: \texttt{ag2\_controls.py} and its log record the box shapes and recover Thue--Morse at $(n,D)=(2,4)$, paperfolding at $(2,4)$, the rational cases, and the Rule~90 and Rule~150 off-centre columns at their known depth-$2$ relations, each relation verified on $10^4$ terms, so the search does exercise depth $\ge2$ on non-rational algebraic series; the apparent false negative for Rule~150 at offset~$20$ is a bound of the scan, not a failure of the instrument. These remain bounds, not a proof of non-automaticity.

\begin{proposition}[the centre column is one designated ANF coefficient; \PROVED]\label{du:anf}
Let $F_t$ be the centre cell at time $t$ as a Boolean function of the initial row on the window
$[-t,t]$, and write $x_j$ for the variable at position $j$. Then
\begin{enumerate}[nosep]
\item[(a)] $F_t=x_{-t}\oplus g_t(x_{-t+1},\dots,x_{t})$: flipping the leftmost cell always flips the
centre, so the coefficient of the monomial $x_{-t}$ is $1$ for every $t$.
\item[(b)] The Walsh spectrum $W_t(a)=\sum_x(-1)^{F_t(x)+a\cdot x}$ vanishes identically for every mask
$a$ with $a_{-t}=0$. Half the spectrum is zero, and the entire first-order structure is the single
integer $W_t(e_{-t})$.
\item[(c)] $c(t)=F_t(\delta_0)$ is exactly the ANF coefficient of the single monomial $x_0$.
\item[(d)] Adding that one monomial, $F_t\mapsto F_t+x_0$, flips $c(t)$ while permuting the Walsh
spectrum by $W(a)\mapsto W(a\oplus e_0)$, hence leaving the Walsh \emph{multiset}, the algebraic
degree, and every invariant defined from them unchanged.
\end{enumerate}
\end{proposition}
\begin{proof}
(a) Rule~30 is left-permutive, and the flip of the leftmost cell propagates along the left edge of the
light cone to the centre in exactly $t$ steps. (b) Pair $x$ with $x+e_{-t}$: by (a) the exponent
$F_t(x)$ changes while $a\cdot x$ does not when $a_{-t}=0$, so the two terms cancel. (c) $F_t(0)=0$, so
the constant term of the ANF vanishes, and $F_t(\delta_0)$ is the sum of the coefficients of the
monomials supported in $\{0\}$, which is the coefficient of $x_0$ alone. (d) Since $x_0=e_0\cdot x$, the
exponent $F_t+x_0+a\cdot x$ equals $F_t+(a\oplus e_0)\cdot x$. Verified exactly for $t\le11$ in (a)--(c)
and $t\le7$ in (d), with the ANF computed to $t=15$, where $F_{15}$ has $344\,592\,268$ monomials.
\end{proof}

\begin{theorem}[the degree law for the centre cell; \PROVED]\label{du:thm:degree}
Let $F_t$ be as in Proposition~\ref{du:anf}. For every $t\ge2$, $\deg F_t=2t-1$. For $t\ge3$ the
top-degree part of $F_t$ is the single monomial $\prod_{j=-t+2}^{t}x_j$ --- every light-cone variable
except the two leftmost --- and its coefficient is $1$.
\end{theorem}
\begin{proof}
Three ingredients. (i) \emph{Strong left permutivity.} In the ANF of any cell the leftmost light-cone
variable occurs only as the bare singleton: by induction, the new cell's leftmost variable lies in
$L$'s cone alone, $L=x\oplus(\text{terms free of }x)$, and $C$, $R$ and $CR$ never see it.
(ii) \emph{Upper bound and uniqueness}, by induction from $\deg F_2=3$. With $F_{t+1}=L\oplus C\oplus
R\oplus CR$ and $\deg L,C,R\le2t-1$, only $CR$ can exceed $2t-1$. Every monomial of degree $\ge2$ omits
the new leftmost variable $x_{-t-1}$ by (i), so $\deg\le2t+2$. A monomial of degree $2t+2$ or $2t+1$
that contains $x_{-t}$ must take it from $C$, whose factor is then the singleton $x_{-t}$ by (i),
leaving $R$ to supply $2t+1$ or $2t$ variables against $\deg R\le2t-1$. Hence $\deg F_{t+1}\le2t+1$,
and the only candidate of that degree omits $x_{-t-1}$ and $x_{-t}$: the product of the remaining
$2t+1$ variables. (iii) \emph{Non-vanishing.} By M\"obius inversion the candidate's coefficient is the
parity of $\#\{x:x_{-t}=x_{-t+1}=0,\ F_t(x)=1\}$, which has the parity of $R_t/2$ for
$R_t:=\#\{x:x_{-t}=x_{-t+1}=0,\ \partial F_t/\partial x_t=1\}$. Since
$\partial\bigl(l\oplus(c\vee r)\bigr)/\partial r=1\oplus c$, the centre depends on $x_t$ exactly when
the chain of cells $a(k,t-k-1)$, $k=0,\dots,t-1$, down the right edge of the cone all vanish. Each of
those cells is permutive in its own leftmost variable, the odd-position cell $x_{t-2k-1}$, so the
system is triangular: the odd positions are determined and the even ones free, and $x_{-t+1}$ becomes
a function $\psi_t$ of the $t-1$ free interior cells ($x_t$ enters nowhere), with $R_t=2^t-2\,\mathrm{wt}(\psi_t)$. Thus $R_t\equiv2\pmod4$
iff $\bigoplus_x\psi_t(x)=1$. Put $b(k,j)=a(k,t-1-k-j)$, the cell $j$ steps left of the gate diagonal in row $k$ (so $b(0,j)=x_{t-1-j}$ and $b(0,2t-2)=x_{-t+1}$); the rule reads
$b(k+1,j)=b(k,j+2)\oplus\bigl(b(k,j+1)\vee b(k,j)\bigr)$, permutive in $b(k,j+2)$, and the gates are $b(k,0)$.
Solve sideways: put $B(k,0)=0$, $B(k,1)=c_k$ and $B(k,j+2)=B(k+1,j)\oplus\bigl(B(k,j+1)\vee B(k,j)\bigr)$ (two
adjacent columns fix the space-time of a permutive automaton; read in the coordinates $a$, each step is Meier and Staffelbach's completion backwards \cite[eq.~(4)]{MS}). Then $c\mapsto$ row~$0$ is a bijection from
$\F^{t-1}$ onto the row-$0$ words with all gates closed, $B(k,j)$ depends only on $c_k,\dots,c_{k+\lfloor(j-1)/2\rfloor}$,
and $G_m:=B(0,2m-1)=c_{m-1}\oplus g_m$ with $g_m$ free of $c_{m-1}$. Write $F_m:=B(0,2m)$ and $T_m(X):=\bigoplus_{c\in\F^m}X(c)$,
so that $\bigoplus\psi_t=T_{t-1}(F_{t-1})$ (a cube sum is invariant under a bijection of the cube). Expanding the OR,
$F_m=B(1,2m-2)\oplus G_m\oplus F_{m-1}\oplus G_mF_{m-1}$. The first three terms have $T_m=0$: $B(1,2m-2)$ misses $c_0$,
$G_m$ is the bare $c_{m-1}$ plus a term missing $c_{m-1}$, and $F_{m-1}$ misses $c_{m-1}$; and
$T_m\bigl((c_{m-1}\oplus g_m)F_{m-1}\bigr)=T_{m-1}(F_{m-1})$ because $gh\oplus(1\oplus g)h=h$. Hence
$T_m(F_m)=T_{m-1}(F_{m-1})=\dots=T_1(F_1)=T_1(c_0)=1$, so $\bigoplus\psi_t=1$ for every $t\ge2$, $R_t\equiv2\pmod4$ and the coefficient is $1$. A second,
independent proof runs the same telescoping in the original variables through a two-to-one map onto
the $(t-1)$-instance. Both are reproduced by \texttt{s3\_verify\_topcoef.py} and
\texttt{s3\_gate\_fast.py}; the census $R_t$ is exact to $t=22$ ($R_{22}=2\,066\,214$), and a brute
force of the original definition agrees to $t=12$. On provenance: we have not found the law in print, and its ingredients are classical (completion backwards, and the top ANF coefficient as a cube sum). It is a different statement from the Fibonacci degree growth recorded in \cite{CLMR}, which concerns a univariate polynomial encoding the support of a single-seed row and has no free input variables, whereas $F_t$ is a function of all $2t+1$ light-cone cells and its degree is linear in $t$.
\end{proof}

\begin{proposition}[the right-edge influence law; \PROVED]\label{du:prop:influence}
For uniform $x\in\F^{2t+1}$ let $\mathrm{Inf}_j=\Pr[F_t(x)\ne F_t(x\oplus e_j)]$ be the influence of the
light-cone variable $x_j$ on the centre cell $F_t=a(t,0)$. Then $\mathrm{Inf}_{-t}=1$ and
$\mathrm{Inf}_{t}=\mathrm{Inf}_{t-1}=2^{-t}$ for every $t\ge1$. (\EXACT\ alongside, $t\le12$:
$\mathrm{Inf}_{t-2}=(2t-1)2^{-t}$.) This is the exact right-edge form of the sensitivity profile that
Chan-L\'opez and Mart\'in-Ruiz measure by Monte Carlo for $t\le20$ \cite[\S8.2]{CLMR}, whose ``flat left
profile $\approx\tfrac12$'' is not exact. The endpoint $j=-t$ has value $1$; the next two cells have $\mathrm{Inf}_{-t+1}=\mathrm{Inf}_{-t+2}=\tfrac12$ exactly at every odd $t$ (\PROVED, ledger~10.142c2: for odd $t$ the derivative in $x_{-t+1}$ is permutive in $x_{-t+2}$) and at no even $t$ computed (\EXACT, $t\le12$ for $x_{-t+1}$ and $t\le10$ for $x_{-t+2}$; for instance $\mathrm{Inf}_{-t+1}=5/16$ at $t=4$); and the remaining left cells $x_{-t+3},\dots,x_{-1}$ take values between $0.467$ and $0.523$ at $t=11$ (\EXACT).
\end{proposition}
\begin{proof}
$\mathrm{Inf}_{-t}=1$ is left permutivity. For the right edge let $G_j=a(j,t-j-1)$ and $P_j=a(j,t-j)$ be
the two rightmost cells of row $j$ of the cone and $H_j=a(j,t-j-2)$ the next. (a) \emph{Independence.}
The map from row $j$ to $(\text{row }j+1,G_j,P_j)$ is a bijection, since left permutivity recovers row
$j$ from its two rightmost cells and row $j+1$ one cell at a time from the right; so if row $j$ is
uniform then row $j+1$ is uniform and independent of $(G_j,P_j)$, and by induction the pairs $(G_j,P_j)$,
$j<t$, are independent fair bits, with $H_j=P_{j+1}\oplus(G_j\vee P_j)$. (b) \emph{Damage.} A flip of
$x_t$ or of $x_{t-1}$ changes in row $j$ only $G_j$ and $P_j$, by $(\delta G_j,\delta P_j)$ say, and the
rule gives
\[
\delta G_{j+1}=\delta G_j\,(1\oplus H_j),\qquad
\delta P_{j+1}=(G_j\vee P_j)\oplus\bigl((G_j\oplus\delta G_j)\vee(P_j\oplus\delta P_j)\bigr),
\]
with $\mathrm{Inf}=\Pr[\delta P_t=1]$. For $x_t$ the damage starts at $(0,1)$, $\delta G$ stays $0$ and
$\delta P_{j+1}=\delta P_j(1\oplus G_j)$: the derivative $\partial F_t/\partial x_t$ is
$\prod_{k<t}(1\oplus G_k)$, the gate chain of Theorem~\ref{du:thm:degree}, and its probability is
$2^{-t}$ by (a). For $x_{t-1}$ the damage starts at $(1,0)$, and the chain on
$(\delta G_j,\delta P_j,P_j)$ driven by the fresh bits $G_j$ and $P_{j+1}$ has, after $j\ge1$ steps, the
law $\Pr[(1,1,P)]=4^{-j}/2$ and $\Pr[(0,1,P)]=(2^{j}-1)/(2\cdot4^{j})$ for each $P$,
$\Pr[(1,0,1)]=(2^{j}-1)/4^{j}$, $\Pr[(1,0,0)]=0$, the remaining mass dead; the induction is the two
displayed formulas read case by case. Hence $\Pr[\delta P_j=1]=2^{-j}$. The chain reproduces the
brute-force influences for $t\le12$ and gives $2^{-t}$ exactly to $t=60$ (\texttt{s3\_influence.py});
the value was known by enumeration to $t=11$ (ledger 10.141b), and the phenomenon --- the right
sensitivity decays because $l\oplus(c\vee r)$ is free of $r$ when $c=1$ --- is in \cite[\S8.2]{CLMR}.
\end{proof}

\begin{remark}[which rules carry the parity, and the same monomial as the census]\label{du:degreerules}
The argument in (iii) applies verbatim to every left-permutive rule $a'=l\oplus h(c,r)$ with
$h=h_{00}\oplus h_{10}c\oplus h_{01}r\oplus h_{11}cr$, and gives $\bigoplus\psi_t=h_{10}h_{11}^{\,t-2}$,
so $R_t\equiv2\pmod4$ exactly when $h_{10}=1$ and either $h_{11}=1$ or $t=2$: Rules $30$, $75$, $180$
and $225$ at every $t$, and Rules $60$, $105$, $150$ and $195$ at $t=2$ only. \PROVED\ by two
independent derivations and \EXACT\ by brute force of the original count over all sixteen such rules
for $t\le8$ (\texttt{s3\_verify\_general.py}).

The step from $R_t$ to the top coefficient is Rule~$30$'s, and it does not generalise: it runs through
$\partial\bigl(l\oplus(c\vee r)\bigr)/\partial r=1\oplus c$, whereas the general gate derivative is
$h_{01}\oplus h_{11}c$. The coefficient of $\prod_{j=-t+2}^{t}x_j$ in $F_t$ is in fact
$h_{11}\,(1\oplus h_{10}\oplus h_{01})$, independent of $t$ and of $h_{00}$, so the rules reaching the
maximal degree with that monomial are $30$, $120$, $135$ and $225$ --- not the same four, and Rules
$60$, $105$, $150$ and $195$ do not reach it even at $t=2$. \EXACT, all sixteen rules and $t\le8$
(\texttt{s3\_crossrule\_check.py}); no proof of the general formula is claimed here, only of
Rule~$30$'s (Theorem~\ref{du:thm:degree}). An earlier version of this remark carried the $R_t$ list
over to the top coefficient and read ``the top coefficient is $1$ at every $t$ for exactly Rules $30$,
$75$, $180$ and $225$''; that sentence is \textsc{withdrawn}, being false at four rules.

The cleanest instance is Rule~$30=150\oplus cr$: the single monomial $cr$ switches on both
the parity and, by Remark~\ref{du:gate}, the single-seed aperiodicity that Rule~$150$ lacks. That is the
same quadratic coupling of the centre to a neighbour which the census found necessary --- and, the
census is explicit, far from sufficient --- so this is one monomial appearing in two results, not an
explanation of the second by the first.
\end{remark}

\begin{remark}[why this closes the Boolean-invariant route exactly]\label{du:anfobstruction}
Proposition~\ref{du:anf} turns Problem~1 into a statement about one designated coefficient of an
explicit Boolean function, which sounds like progress and is in fact an obstruction, because (d) says
that coefficient is invisible to the standard invariants. Degree, nonlinearity, correlation immunity,
the rank of the quadratic part, and the dimension of the linear space are all functions of the Walsh
multiset or of the monomial support above degree one, and adding $x_0$ changes $c(t)$ while fixing every
one of them. So no such invariant can decide the centre column, however precisely we compute it --- and
we can compute a good deal, and one thing we can now prove: by Theorem~\ref{du:thm:degree} the degree
is exactly $2t-1$, near-maximal in $2t+1$ variables, with a single top monomial for $t\ge3$, and the
linear space is $\{0,e_{-t}\}$ (\EXACT\ to $t=9$; at $t=2$ it has dimension $2$). One correction to how an earlier draft described the open half of that
law. The divisibility $W_t(e_{-t})\equiv0\pmod 8$ is not it. Since $F_t=x_{-t}\oplus g_t$ one has
$W_t(e_{-t})=2^{2t+1}-4\,\mathrm{wt}(g_t)$, so the divisibility says exactly that $\mathrm{wt}(g_t)$ is
even, i.e.\ that the full monomial of $g_t$ is absent --- which is the \emph{upper} bound
$\deg\le2t-1$, the half that was already proved. The lower bound is the non-vanishing of one designated
coefficient, and $W_t(e_{-t})$ does not see it even modulo $16$: the residues run $8,0,0,0,8,0,0,8$ at
$t=2,\dots,9$ while the coefficient is $1$ throughout. This obstruction is worth distinguishing
from the Rule~90 obstruction of Proposition~\ref{p2:rule90}, which says a hypothesis is too weak to
separate two rules. Here the hypothesis class is blind to the quantity itself, on one rule. For the
record, the low-degree layers of the ANF are not independent information either: the degree-one layer
is the mirrored single-seed row, the degree-two layer is the pairwise superposition defect and the
degree-three layer the triple defect, verified on all $615$ pairs and $1\,035$ triples with no
mismatch. And the bias bound available for a quadratic form over $\F$, which is governed by its rank,
does not apply, its hypothesis $\deg=2$ failing here by a factor that grows to $97$ over the range
computed.
\end{remark}

\begin{remark}[the Hankel rank: one object behind two censuses, and why it cannot discriminate]\label{du:hankel}
The rational-series view makes our two censuses one object and supplies the only bound we have about
\emph{weighted} machines. The right Hankel matrix is not the classical $[c_{i+j}]$ --- finite rank there
means a linear recurrence, which over $\F$ is eventual periodicity, so that object restates Problem~1
--- but the matrix indexed by base-$2$ digit words, whose columns \emph{are} the elements of the
$2$-kernel. Hence $\log_2|\mathcal K(c)|\le\operatorname{rank}\le|\mathcal K(c)|$, which is why the two
censuses report $1024$ and $138\,672$ on the same sequence: they are the logarithm and the cardinality
of one set. By the Fliess--Sch\"utzenberger theorem a recognisable series has a minimal linear
representation of dimension its Hankel rank, so every $\F$-weighted automaton computing $c$ from its
base-$2$ digits has at least $1024$ states --- every other bound in this Part concerns deterministic
machines.

It discriminates nothing, on three counts. Only \emph{boundedness} of the rank characterises anything,
so Remark~\ref{du:censuscaveat} applies verbatim. The domination $R_f\le A_f$ makes the rank strictly
weaker than the antichain census, by a factor of $135$ on our data. And the coin control is a theorem:
a random series over $\F$ has expected rank at least $2^{n/2}$, and the split matrix at $n=8,\dots,18$
gives Rule~30 as $16,31,64,127,256,510$ against a coin's $15,30,63,127,255,511$ on the same prefix,
saturating that ceiling and indistinguishable at every $n$.

That is worth stating as the general fact it is. Every instrument in this paper separates Rule~30 from
Rule~90 without difficulty, and not one separates it from a fair coin. The obstruction has been
described as the firewall rule satisfying our hypotheses; the sharper statement is that our sequence is
indistinguishable from a coin to every quantity we can compute, and Rule~90 is the most convenient
witness of that.
\end{remark}

\begin{remark}[$k$-regularity is not a weaker hypothesis, but one derived statistic does separate]\label{du:regular}
It is natural to hope that $k$-regularity, where the $k$-kernel need only span a finitely generated
module rather than be finite, is a weaker hypothesis that our data could refute outright. For $c$ it is
not weaker at all, and for two independent reasons: a finitely generated module over $\F$ is finite, and
a $k$-regular sequence taking finitely many values is $k$-automatic \cite[p.~441]{AS}. The same applies
to the summatory function of $c$. So nothing is gained on the centre column itself.

For the \emph{unbounded} statistics of the orbit the notions do differ, and there the module rank is a
computable invariant that behaves differently on Rule~30 than on the additive rules --- which is more
than the census does. Certified lower bounds on the rank of the $2$-kernel module are $512$ for the row
weight, $351$ for the diagonal preperiods $\tau_d$ and $294$ for the periods $p_d$, against $1$ for
Rule~90 and $2$ for Rule~150. That this separates Rule~30 from the rules that satisfy more of our
hypotheses than it does is the one feature no other instrument in this Part has. We record it as an
instrument rather than a result: the relevant published measure is Shallit's rationality measure, whose
Theorem~9(b) gives $R_f\le A_f$, so the rank census is strictly weaker than the antichain census and
caps at $\sqrt{2N}$ against $N/\log_2 N$. And the decision procedures available --- Krenn and Shallit's
via Jacob's boundedness algorithm, and Krawczyk and M\"ullner's \cite{KrawczykMullner} --- all consume a presentation as input,
which is exactly what these rank bounds say we do not have.
\end{remark}

\begin{proposition}[what the certificates exclude; \PROVED\ from \KNOWN\ bounds]\label{du:convert}
Write $\mathrm{comp}_2(c)$ for the number of states of the minimal $2$-automaton generating $c$, which
equals $|\mathcal K(c)|$ \cite[Cor.~4.1.9, Thm.~6.6.2]{AS}. Adamczewski and Yassawi \cite{AY} record
$\mathrm{comp}_q(a)\le q^{\,d(2H+1)}$ for a series algebraic of degree $d$ and height
$H=\max_i\deg A_i$. Contrapositively, the certificate of Theorem~\ref{du:certified}(1) converts into an
exclusion rather than a mere bound: for every $d\le12$, the series $G$ is not algebraic of degree $d$
with height at most $8000/2^{\,d+1}$. In particular $G$ is not quadratic of height at most $1000$, nor
of degree $4$ with height at most $250$. In the other direction, run backwards against the census of Theorem~\ref{du:certified}(2), Bridy's bound \cite{Bridy} $\mathrm{comp}_q\le(1+o(1))q^{\,h+d+g-1}$ gives $h+d+g\ge19$, the exact form of Rowland, Stipulanti and Yassawi removes the asymptotic caveat and gives $hd\ge18$, and \cite[eq.~5]{AY} gives $(h+1)d\ge17$, for the height, degree and genus of any algebraic presentation of $G$. Over $\F$ a Mahler relation is an Ore relation exactly, the Frobenius being trivial on coefficients.
\end{proposition}


% Part II, from duhamel_decimation_section.tex (labels prefixed du:)
\section{Decimation of the Duhamel class}\label{du:decimation}

The Duhamel formula expresses $c_t$ as a linear functional of the initial row
and of the quadratic source $q^s_x = u^s_x u^s_{x+1}$. In this section we treat
that functional on an arbitrary pair (initial row, source field), show that the
class of such expressions is closed under the Cartier operators, and give a
closed form for the functional. The result is an exact description of the
$2$-kernel of $c$ as the image of an orbit, and a reduction of the algebraicity
question to a series built from the chiral part of the source alone.

\subsection{The class $c[g,f]$}

For a row $g\colon \mathbb Z\to\mathbb F_2$ and a source field
$f=(f^s_x)_{s\ge 0,\,x\in\mathbb Z}$, both of finite support in $x$ for each $s$,
define
\[
  c[g,f]_t \;=\; \langle T_t, g\rangle \;+\; \sum_{s<t}\langle T_{t-1-s}, f^s\rangle,
  \qquad \langle a,b\rangle=\sum_{x} a_x b_x \in \mathbb F_2 .
\]
Theorem~\ref{du:duhamel} states that $c = c[\delta, q]$ with $\delta_x=[x=0]$ and
$q^s_x=u^s_x u^s_{x+1}$. Write $(\Lambda h)_x = h_{x-1}+h_x+h_{x+1}$ for the
action of $\Lambda$ on rows, and $E h$ for the even-site restriction
$(Eh)_x = h_{2x}$.

\begin{theorem}[exact decimation; \PROVED]\label{du:thm:A}
For every $(g,f)$ and every $t\ge 0$,
\[
  c[g,f]_{2t} = c[D_0(g,f)]_t, \qquad c[g,f]_{2t+1} = c[D_1(g,f)]_t ,
\]
where
\begin{align*}
  D_0(g,f) &= \Bigl(E g,\;\; \bigl(E\,(f^{2\sigma+1} + \Lambda f^{2\sigma})\bigr)_{\sigma\ge 0}\Bigr),\\
  D_1(g,f) &= \Bigl(E(\Lambda g + f^{0}),\;\; \bigl(E\,(f^{2\sigma+2} + \Lambda f^{2\sigma+1})\bigr)_{\sigma\ge 0}\Bigr).
\end{align*}
\end{theorem}

\begin{proof}
By Lemma~\ref{du:kernel}, $T_{2m}(z)=T_m(z^2)$ and $T_{2m+1}(z)=\Lambda(z)\,T_m(z^2)$.
A pairing with $T_m(z^2)$ samples the even sites,
$\langle T_m(z^2), h\rangle = \sum_x T_m(x)\,h_{2x} = \langle T_m, Eh\rangle$,
and a pairing with $\Lambda(z)T_m(z^2)$ samples $\Lambda h$ on the even sites,
using the symmetry of $\Lambda$: $\langle \Lambda T_m(z^2), h\rangle = \langle T_m, E\Lambda h\rangle$.

For $c[g,f]_{2t}$ the initial term is $\langle T_{2t}, g\rangle = \langle T_t, Eg\rangle$.
In the source sum split $s$ by parity: $s=2\sigma$ gives $T_{2t-1-2\sigma}=T_{2(t-1-\sigma)+1}=\Lambda T_{t-1-\sigma}(z^2)$,
hence $\langle T_{t-1-\sigma}, E\Lambda f^{2\sigma}\rangle$; and $s=2\sigma+1$ gives
$T_{2(t-1-\sigma)}=T_{t-1-\sigma}(z^2)$, hence $\langle T_{t-1-\sigma}, E f^{2\sigma+1}\rangle$.
Collecting terms with the same $\sigma$ gives $c[D_0(g,f)]_t$.

For $c[g,f]_{2t+1}$ the initial term is $\langle T_{2t+1}, g\rangle = \langle T_t, E\Lambda g\rangle$.
The $s=0$ source term is $\langle T_{2t}, f^0\rangle = \langle T_t, E f^0\rangle$ and joins the
initial row. For $s\ge 1$ write $s=2\sigma+1$ ($T_{2(t-1-\sigma)+1}$, giving $E\Lambda f^{2\sigma+1}$)
and $s=2\sigma+2$ ($T_{2(t-1-\sigma)}$, giving $E f^{2\sigma+2}$). Collecting gives $c[D_1(g,f)]_t$.
\end{proof}

\begin{corollary}[\PROVED]\label{du:cor:orbit}
Let $L(g,f) = (c[g,f]_t)_{t\ge 0}$, a linear map into $\mathbb F_2^{\mathbb N}$, and let
$\mathcal O = \{ D_{w}(\delta,q) : w \in \{0,1\}^* \}$ be the orbit of $(\delta,q)$ under
the words in $D_0, D_1$. Then
\[
  \mathcal K(c) = L(\mathcal O).
\]
In particular $\mathcal K(c)$ is finite if and only if $\mathcal O$ is finite modulo $\ker L$.
\end{corollary}

\begin{proof}
The kernel element $(c_{2^k m + r})_m$ is $\Lambda_{r_{k-1}}\cdots\Lambda_{r_0}G$ with
$r=\sum r_i 2^i$, and Theorem~\ref{du:thm:A} says $\Lambda_{r}\,L(g,f) = L(D_{r}(g,f))$.
\end{proof}

\subsection{Closed form of the functional}

\begin{theorem}[closed form; the resolvent is \KNOWN, Litow--Dumas \cite{LD}]\label{du:thm:B}
Let $\alpha\in y\,\mathbb F_2[[y]]$ be the unique solution of
\[
  y\,\alpha^2 + (1+y)\,\alpha + y = 0, \qquad\text{equivalently}\qquad
  y = \frac{\alpha}{1+\alpha+\alpha^2}.
\]
Then for every $(g,f)$,
\[
  \sum_{t\ge 0} c[g,f]_t\, y^t \;=\; \frac{1}{1+y}\Bigl(\sum_{x} g_x\,\alpha^{|x|}
  \;+\; y\sum_{s,x} f^s_x\, y^{s}\alpha^{|x|}\Bigr).
\]
\end{theorem}

\begin{proof}
Let $\hat g(z)=\sum_x g_x z^x$ and $\hat f^s(z)=\sum_x f^s_x z^x$. Summing the
definition of $c[g,f]$ against $y^t$ and using $T_m=\Lambda^m$,
\[
  \sum_t c[g,f]_t\,y^t = [z^0]\Bigl(\sum_t y^t\Lambda^t\Bigr)\Bigl(\hat g + y\sum_s y^s \hat f^s\Bigr)
  = [z^0]\;\theta\cdot\Bigl(\hat g + y\sum_s y^s \hat f^s\Bigr),
\]
where $\theta=(1+y\Lambda)^{-1}\in \mathbb F_2[z^{\pm 1}][[y]]$ is the unique solution of
$y\theta_{x-1}+(1+y)\theta_x+y\theta_{x+1}=[x=0]$. The ansatz $\theta_x = C\alpha^{|x|}$
satisfies the equation for $x\neq 0$ because $C\alpha^{|x|-1}\bigl(y+(1+y)\alpha+y\alpha^2\bigr)=0$,
and at $x=0$ gives $C\,(2y\alpha + 1+y) = C(1+y)$ in characteristic $2$, so $C=1/(1+y)$.
Since $\alpha = y + O(y^2)$, $\alpha^{|x|}=O(y^{|x|})$ and the ansatz lies in the ring,
hence is the solution. Pairing $\theta$ with $\hat g + y\sum_s y^s\hat f^s$ and taking $[z^0]$
gives the formula.
\end{proof}

\begin{remark}[provenance of Theorem~\ref{du:thm:B}]\label{du:thm:B:prov}
The resolvent is not new. Litow and Dumas \cite{LD} write the generating function of an additive
cellular automaton as $M(x,z)=C_0(z)/(1-xQ(z))$, which is $\theta=(1+y\Lambda)^{-1}$ in the present
notation, and evaluate its coefficients as a sum over the roots of the denominator; for
$Q=z+1+z^{-1}$ that evaluation reduces in characteristic~$2$ to $\alpha^{|x|}/(1+y)$, normalisation
included. Their worked Motzkin example moreover exhibits the quadratic $y\alpha^2+(1+y)\alpha+y=0$ and
hence the element $\alpha$ itself, in print in 1993. What is added here is only the characteristic-$2$
proof by the ansatz together with the jump condition at $x=0$, and the extension of the functional from
an initial row to a pair (initial row, source field), which is what makes Theorem~\ref{du:thm:A} and
Corollary~\ref{du:cor:C} available. The decomposition of Rule~30 as Rule~150 plus the single term $cr$
is likewise \KNOWN\ and stated in Wolfram's own prize announcement \cite{Wolfram}: all of the
complexity is attributed there to that one nonlinear term, ``for without that term, one would have
rule 150''. The reading of Corollary~\ref{du:cor:C} as saying that the centre column is fed by the
asymmetry of the rule is also in Chan-L\'opez and Mart\'in-Ruiz \cite{CLMR}.
\end{remark}


For Rule 30 itself ($g=\delta$, $f=q$), rewriting in the variable $\alpha$ with
$1+y = (1+\alpha)^2/(1+\alpha+\alpha^2)$ gives
\[
  G \;=\; \frac{1+\alpha+\alpha^2}{(1+\alpha)^2} \;+\; \frac{\alpha}{(1+\alpha)^2}\,H,
  \qquad
  H \;=\; \sum_{s,x} q^s_x\,\frac{\alpha^{\,s+|x|}}{(1+\alpha+\alpha^2)^{s}} .
\]
Since $\mathbb F_2(\alpha)/\mathbb F_2(y)$ is a finite extension, $G$ is algebraic over
$\mathbb F_2(y)$ if and only if $H$ is algebraic over $\mathbb F_2(\alpha)$.

\subsection{Only the chiral part of the source is visible}

The weight $\alpha^{|x|}$ in Theorem~\ref{du:thm:B} does not distinguish $x$ from $-x$.
Hence $c[g,f]$ depends on $f$ only through the centre line $f^s_0$ and the
\emph{chiral part}
\[
  \chi^s_x = f^s_x + f^s_{-x}, \qquad x>0 ,
\]
and likewise on $g$ only through $g_0$ and $g_x+g_{-x}$.

\begin{corollary}[\PROVED]\label{du:cor:C}
$\ker L = \{(g,f) : \hat g(\alpha) + y\,\hat f(y,\alpha) = 0\}$ with
$\hat g(\alpha)=\sum_x g_x\alpha^{|x|}$ and $\hat f(y,\alpha)=\sum_{s,x}f^s_x y^s\alpha^{|x|}$.
In particular every mirror-symmetric pair ($g_x=g_{-x}$, $f^s_x=f^s_{-x}$) with
$g_0=0$ and $f^s_0=0$ for all $s$ lies in $\ker L$.
\end{corollary}

\begin{remark}
The source of Rule 30, $q^s_x=u^s_x u^s_{x+1}$, is not mirror-symmetric: the rule
is left-permutive and its orbit from the single cell is not symmetric, so $q$ is not
killed by $\ker L$. For a mirror-symmetric rule started from the single cell the
field is symmetric, $\chi\equiv 0$, and the centre column depends on the source only
through its own centre line $f^s_0$. This is the algebraic form of the observation
that the centre column of Rule 30 is fed entirely by the asymmetry of the rule.
\end{remark}

\begin{remark}[the whole family of single-monomial perturbations; \EXACT]\label{du:family}
Rule~30 and Rule~22 are not isolated cases. Rule~150 has algebraic normal form $l+c+r$, and adding one
monomial gives exactly five rules; Corollary~\ref{du:cor:C} classifies all of them, with no exceptions.
\[
\begin{array}{lcccccc}
\text{added monomial} & \text{none} & 1 & lc & cr & lr & lcr\\
\text{rule} & 150 & 105 & 86 & \mathbf{30} & 54 & 22\\
\text{symmetric in } l\leftrightarrow r & \text{yes} & \text{yes} & \text{no} & \text{no} & \text{yes} & \text{yes}\\
\text{orbit mirror-symmetric} & \text{yes} & \text{yes} & \text{no} & \text{no} & \text{yes} & \text{yes}\\
\text{chiral weight of the source} & 0 & 0 & 16\,784 & 16\,784 & 0 & 0\\
\text{centre density to } t=2^{13} & 1.0000 & 0.5000 & 0.5010 & 0.5010 & 0.5000 & 0.0002\\
\text{centre eventual period} & 1 & 2 & \text{none} & \text{none} & 4 & 1
\end{array}
\]
The unperturbed base case is the sharpest entry and was missing from earlier statements of this table:
the centre column of Rule~150 itself is \emph{identically~$1$}. From a single seed, row~$t$ is
$(1+z+z^2)^t$ over $\F$, so the centre bit is the central trinomial coefficient, which is always odd;
the corresponding statement for Rule~90 is that the centre bit is $\binom{t}{t/2}\bmod2$, and
$\binom{2m}{m}=2\binom{2m-1}{m-1}$ is even for every $m\ge1$, so that column is $1,0,0,\dots$. Both
additive parents therefore have a constant centre column, and the whole contrast of this Part is
between a constant sequence and one with no detected period, across a single monomial.
The three symmetric perturbations give a mirror-symmetric orbit, hence a mirror-symmetric source, hence
chiral weight exactly zero, and in each case the centre column is eventually periodic --- with period
$2$, $4$ and $1$ respectively. The density row is a measurement at a stated depth and nothing more; in
particular Rule~22's centre column contains exactly two ones and then vanishes, so its density is $0$
rather than the small positive number a short run suggests, and the $0.5010$ for Rules~$30$ and~$86$
is a finite-depth figure that drifts (it reads $0.5469$ at $t=2^{6}$ and $0.4810$ at $t=10^{3}$). The two chiral perturbations are Rule~30 and its mirror Rule~86, they
share the same chiral weight as they must, and neither centre column has any period up to $64$. So
within this family, being hard is exactly the same as having a chiral perturbation, which is the content
of Corollary~\ref{du:cor:C} read as a classification rather than as a remark. Rule~54 is a second
exactly solvable nonlinear control for the machinery of this Part, genuinely quadratic and with a
period-$4$ centre column, and Rule~105 a third.
\end{remark}

\begin{remark}[the equivalence does not leave this family; \EXACT]\label{du:family256}
The hedge ``within this family'' in the previous remark is load-bearing, and the control is cheap enough
that it should be stated rather than assumed. Taking all $128$ elementary rules with a quiescent
background ($000\mapsto0$) and running the centre column to $t=2^{12}$, exactly three have no eventual
period with preperiod $\le64$ and period $\le1024$: Rules~$30$ and~$86$, and Rule~$126$. Restricting to
the quiescent rules understates the set, and the full census is given in Remark~\ref{du:gate}; the
conclusion drawn here is unaffected, because the counterexamples below are already inside the quiescent
half. Ninety-four chiral rules are eventually periodic, and one
mirror-symmetric rule is not, so ``chiral'' and ``hard'' are equivalent on the five rules of
Remark~\ref{du:family} and on nothing larger. Two entries deserve naming. Rule~$110$ is chiral,
Turing-universal \cite{Cook}, and its centre column from a single seed is constant~$1$: universality of
the rule says nothing about the column. The same separation is what closes the rank route of
Remark~\ref{p3:rank}: communication rank measures the prediction function over all inputs and not the
column, and Rules~$22$ and~$54$, whose single-seed columns are eventually $0$ and of period $4$, have
rank over $\F$ of the same order as Rule~30's or larger at every $t$ computed.
And $110$ of the $128$ columns are constant within two steps, so the single-seed centre column is a very
thin probe --- it is not that most elementary rules are simple, but that most of them say nothing from
this initial condition. What survives of the classification is therefore a statement about the
Rule~$150$ neighbourhood, not a mechanism; the remark below says why that neighbourhood is nonetheless
the right place to look.
\end{remark}

\begin{remark}[confinement: the period exponent is not a property of the rule, the surjectivity defect is; \EXACT\ and \MEASURED]\label{du:ring}
Confine a rule to a ring of width $N$. Every rule then cycles, by pigeonhole, so periodicity under
confinement discriminates nothing and the content can only be in how the period scales. We report this
experiment together with its own refutation, because the refutation is the useful part.

Measuring the single-seed orbit for $N=10..22$ over all $256$ rules, each tagged by mirror symmetry
\emph{and} by affineness so the two are not confounded, the mean slope of $\log_2(\text{period})$
against $N$ is $0.198$ and $0.171$ for the symmetric and chiral affine rules and $0.071$ and $0.122$ for
the symmetric and chiral nonlinear ones: overlapping ranges, and a sign that reverses between the rows,
so chirality does not act in a consistent direction. On that window the class $\{45,75,89,101\}$ sits at
$0.95$, the class of Rule~30 at $0.45$ --- the two classical laws, a random permutation of $2^N$ states
giving cycles $\sim2^N$ and a random map $\sim2^{N/2}$ --- and everything else below $0.28$. Those
exponents are stable under changing the initial state, with spread $0.07$--$0.22$ against a gap of
$0.51$, and exactly $0$ for the affine rules as a linear map requires.

They are nevertheless \emph{not measurements of anything}, and seed-stability is what misled us: the
instability is in $N$, not in the seed. Extending the same computation to $N=36$ for Rule~30 and $N=52$
for the rules whose periods stay small refutes the window fit outright. The $N\le22$ fit for Rule~30
under-predicts its period at $N=25$ by a factor $46.6$; the fit for Rule~$45$ over-predicts at $N=26$ by
a factor $977$. The fitted exponent is a function of the window and not of the rule: for Rule~30 it
reads $0.447,\,0.718,\,0.329,\,0.587,\,0.511,\,0.059$ and $-0.463$ on the windows
$10..22,\,10..26,\,22..36,\,10..36,\,20..36,\,26..36$ and $30..36$, a range of more than one unit that
changes sign. Even the \emph{ordering} reverses --- Rule~30's period exceeds Rule~$45$'s at
$N=12,16,24,26$ and is smaller by a factor $745$ at $N=19$. The decisive figure is the raw series: Rule~30's
period runs $18\,480\,630$ at $N=35$, then $2\,844$ at $N=36$, then past a cap of $6\times10^{7}$ at
$N=37$, so consecutive widths differ by more than four orders of magnitude. No exponent survives that.
(The $N=36$ value is exactly the $N=18$ value because the orbit does fall into the
spatially $18$-periodic invariant subspace, at $t=193\,613$; Rule~30's ring map is not injective, so such
a collapse is available, and the subspace is shift-invariant hence preserved. It is the \emph{only}
width in $12..40$ where this happens, so it explains this one point and not the scatter: neither $N=35$
nor $N=37$ involves a collapse.) By contrast the rules whose periods
stay small are stable out to $N=52$: Rule~$54$ has period $4$ at every width, exponent exactly $0$ over
$31$ widths, and Rules~$90$, $150$ and~$126$ sit at $0.17$--$0.18$. The instability is specific to the
rules the statistic appeared to single out. We therefore withdraw the exponent and record the route as closed: the ring-period scaling of the
single-seed orbit is not a property of the rule at any width we can reach.

What survives is the quantity the exponents pointed at, and it survives because it involves no fitting.
Writing $\delta_N=1-|{\rm im}|/2^N$ for the surjectivity defect of the width-$N$ ring map: Rule~$45$ has
$\delta_N=2^{-N/2}$ at every even $N$ and $\delta_N=0$ at every odd $N$ --- it is surjective on every odd
ring --- \PROVED\ in Proposition~\ref{du:prop:ringdefect}; Rule~30's defect decays with the exponent
$\log_2\lambda-1=-0.2382$, $\lambda=1.6956$ the real root of $x^{3}-x^{2}-2$, an algebraic constant and
not a fit (the same proposition; the window fits $-0.2366$, $-0.2376$ and $-0.2382$ on $N=8..20$,
$8..28$ and $20..28$ are stable, unlike the period exponent above); Rule~$110$'s image collapses
($+0.061$, \MEASURED\ on $N=8,10,\dots,20$); Rule~$90$ is exactly four-to-one at every even width and
two-to-one at every odd one; and Rule~$150$ is four-to-one exactly at the multiples of three and
bijective otherwise, since $1+z+z^2\mid z^N-1$ iff $3\mid N$. Rules~$45$, $30$ and~$110$ are all chiral
while spanning that whole range.

The sharp form of the conclusion is an identity in the algebraic normal form, and it needs none of the
retracted material. Rule~$30=150+cr$ and Rule~$45=90+1+cr$ carry the \emph{same} chiral monomial over
different additive parents, and differ from each other by the single term $1+c$; their defects differ
not merely in exponent but in kind: Rule~$45$'s is a parity phenomenon in $N$, exactly $2^{-N/2}$ at
even $N$ and absent at odd $N$, while Rule~$30$'s is a genuine algebraic decay at every $N$ with no parity
structure, its Garden-of-Eden counts $1,1,3,5,6,12,22,33,57,101,166,280$ for $N=1,\dots,12$ containing no
power of two past $N=2$. So the chiral monomial does not by itself determine the behaviour even across
the two additive parents this Part is built on. This leaves Corollary~\ref{du:cor:C} untouched,
which is a proved statement about the chiral part of the \emph{source} being the only part visible at
the centre; what it constrains is the reading of Remark~\ref{du:family} as a mechanism, which is now
hedged for a second independent reason.
\end{remark}

\begin{proposition}[ring surjectivity defects from the transfer monoid; \PROVED]\label{du:prop:ringdefect}
For an elementary rule $f$ let $A_b$, $b\in\{0,1\}$, be the Boolean $4\times4$ matrix on pairs of adjacent
cells with $(A_b)_{(p,q),(q,r)}=[f(p,q,r)=b]$. A configuration $y$ on the ring $\Z/N$ has a preimage iff
the Boolean product $A_{y_1}\cdots A_{y_N}$ has a nonzero diagonal entry, so $|{\rm im}_N|$ is the number
of length-$N$ words accepted by the automaton whose states are the finite monoid generated by $A_0,A_1$,
and $|{\rm im}_N|$ satisfies a linear recurrence of order at most the size of that monoid, which is $8$,
$12$, $21$, $38$, $41$, $75$ and $2\,375$ for Rules $90$, $150$, $45$, $18$, $30$, $110$ and $22$.
Consequently:
\begin{enumerate}[nosep]
\item[(a)] Rule~$45$ has exactly $2^{N/2}$ Garden-of-Eden configurations on every even ring and none on
any odd ring: it is surjective on every odd ring.
\item[(b)] Rule~$150$ is four-to-one when $3\mid N$ and bijective otherwise; Rule~$90$ is four-to-one at
even $N$ and two-to-one at odd $N$.
\item[(c)] Rule~$30$'s Garden-of-Eden count $s_N=2^{N}-|{\rm im}_N|$ satisfies the recurrence of
$Q(x)=(x-1)^{7}x^{11}(x^{2}+x+1)^{2}(x^{3}-x^{2}-2)^{5}$, the characteristic polynomial of its automaton
with the factor $(x-2)^{4}$ removed, and $\lim_N N^{-1}\log_2 s_N=\log_2\lambda$ with
$\lambda=1.695620\ldots$ the real root of $x^{3}-x^{2}-2$; its defect exponent is therefore
$\log_2\lambda-1=-0.238186$, an algebraic constant and not a fit.
\end{enumerate}
\end{proposition}
\begin{proof}
The automaton is the transfer-matrix description of preimages on a ring, and the recurrence is
Cayley--Hamilton for its transition matrix $M$, since $|{\rm im}_N|=uM^{N}v$ with $u$ the identity state
and $v$ the accepting states. (a) $2^{N}-|{\rm im}_N|$ satisfies a recurrence of order $\le22$ and
$2^{N/2}[2\mid N]$ one of order $2$, that of $x^{2}-2$; the two agree for $N=1,\dots,60$
(\texttt{s3\_ring\_defect.py}, the automaton counts cross-checked against brute force for $N\le14$),
which exceeds the $24$ consecutive values that force agreement at every later $N$. (b) is linear algebra
over $\F$: the kernel of $x\mapsto x_{i-1}+x_i+x_{i+1}$ on $\Z/N$ has dimension
$\deg\gcd(1+z+z^{2},z^{N}-1)$, which is $2$ iff $3\mid N$, and that of $x\mapsto x_{i-1}+x_{i+1}$ has
dimension $\deg\gcd(1+z^{2},z^{N}-1)$, which is $2$ or $1$ by the parity of $N$. (c) Let $R$ be the
characteristic polynomial of the $41\times41$ matrix $M$, computed exactly, and $R=(x-2)^{4}Q$ with
$Q(2)\ne0$. The sequence $Q(\sigma)s$, $\sigma$ the shift, lies in the four-dimensional solution space of
$(x-2)^{4}$ and vanishes at $N=0,\dots,43$ (\texttt{s3\_ring\_defect2.py}), hence identically; so the
eigenvalue-$2$ component of $s$ is zero and $s$ satisfies $Q$. The roots of $Q$ have moduli $0$, $1$,
$\lambda$ and $(2/\lambda)^{1/2}<\lambda$, and $s$ does not satisfy $Q/(x^{3}-x^{2}-2)^{5}$, so the
$\lambda$-component is nonzero and $s_N=p(N)\lambda^{N}(1+o(1))$ for a nonzero polynomial $p$.
Numerically $s_{N+1}/s_N=1.69562077$ from $N=41$ on, so $p$ is constant to the precision available.
\end{proof}

\begin{remark}[the provable case in this family is the symmetric one; \EXACT]\label{du:rueppel}
Rule~$126$ is $150+lc+cr+lr$, the \emph{full symmetric} quadratic, and it is mirror-symmetric. Its
centre column from a single seed is $1$ exactly at $t=2^{k}-1$, verified for every $t\le2^{14}$, so it
is the Fredholm--Rueppel sequence \cite{Rueppel} as far as computed ($t\le2^{14}$, \EXACT); if the identity holds for all $t$ --- the single-seed pattern of Rule~$126$ is Rule~$90$'s with each cell thickened inward, which we have checked but not written out --- then the column is not eventually periodic, since the gaps double.
The inversion is worth stating plainly. Within the single-monomial neighbourhood of Rule~150, the one
case whose non-periodicity admits a proof is the symmetric one, and the chiral perturbation is precisely
where the proof is unavailable. Nothing here suggests that the chiral monomial \emph{causes}
non-periodicity; what the family shows is that it is the perturbation which removes every handle we
have, which is a statement about our methods and not about the rule.
\end{remark}

\begin{remark}[the full aperiodic census, and the exact necessary condition; \EXACT]\label{du:gate}
Running all $256$ rules rather than the quiescent half --- for a non-quiescent rule the background
blinks, but the centre column is still a well-defined sequence and periodicity is still the right
question --- exactly fifteen have no eventual period at preperiod $\le64$, period $\le1024$ and
$t\le2^{12}$, and two of them are artefacts of the preperiod bound: Rules~$137$ and~$193$, a mirror pair,
are eventually periodic with period $240$ and preperiod $2851$ (\EXACT\ to $t=20\,000$). The other
thirteen show no eventual period at preperiod $\le4000$, period $\le1024$ and $t<8000$:
\[
30,\;45,\;75,\;86,\;89,\;101,\;126,\;129,\;135,\;149,\;161,\;167,\;181 .
\]
The set is consistent with the one symmetry that preserves the experiment: the mirror fixes a symmetric
seed, and every one of the thirteen has its mirror in the set, giving eight mirror classes. The complement
does \emph{not} preserve it, since it sends a single $1$ on a $0$ background to a single $0$ on a $1$
background, and indeed three of the thirteen ($161$, $167$ and~$181$) have a complement outside the set.
Six mirror classes are
complement-closed as well, and they group into three families: Rule~$30$'s $\{30,86,135,149\}$,
Rule~$45$'s $\{45,75,89,101\}$, and Rule~$126$'s $\{126,129\}$. Those three are the rules whose centre
column is aperiodic from a single seed on either background.

The census yields one sharp necessary condition, and it is a condition on a single pair of monomials.
\emph{Every} one of the thirteen has $cr$ or $lc$ in its algebraic normal form --- a quadratic term
coupling the centre to one neighbour --- and no affine rule appears at all. The condition is specific to
that pair rather than to nonlinearity: of the $32$ rules carrying $lr$ but neither $cr$ nor $lc$, none is
aperiodic, and of the $32$ carrying $lcr$ but neither, none is aperiodic. So a quadratic gate on the two
\emph{outer} cells, or a cubic on all three, never suffices; the coupling must reach the centre. This is
the precise form of the informal statement that the chiral monomial is what breaks periodicity, and it
also explains why Rule~54 ($150+lr$) and Rule~$22$ ($150+lcr$) are the tame members of
Remark~\ref{du:family} while $30$ and~$86$ are not.

It is necessary and far from sufficient: $192$ rules carry $cr$ or $lc$ and only thirteen are aperiodic,
and no Boolean expression of at most three literals over the monomial indicators, the light-cone growth
flags and the permutivity flags --- conjunctive or disjunctive --- selects the thirteen exactly. One
tempting half of the informal statement is definitely false: light-cone growth is not necessary, since
Rules~$45$, $101$, $129$ and~$161$ send both $001$ and $100$ to~$0$. What the census
supplies is therefore a clean necessary condition and no sufficient one, which is the honest position:
the gate is required, and what selects thirteen rules out of the $192$ that have it is not visible at the
level of the normal form.
\end{remark}

\begin{remark}[what the literature has solved in this family, and what it has not]\label{du:chiraldisc}
The classification of Remark~\ref{du:family} lines up exactly with the state of the published art, and
the alignment is not a coincidence. Of the five rules, the three symmetric ones have all been solved or
rigorously treated: Rule~105 is affine and falls to \cite{MOW}; Rule~54 is exactly solvable in its
reversible two-step form, with a complete quasiparticle catalogue \cite{BKP}; and Rule~22 has a
rigorous finite-seed theory \cite{GG22} as well as the closed form of \cite{CLMR} that we verified in
Remark~\ref{du:rule22}. Neither chiral one has been solved, and they are Rule~30 and its mirror.
The same split appears one rung down, on Rule~90: the \emph{symmetrised} chiral pair
$90+lc+cr$ is Rule~18, the rule whose defect is the one nonlinear elementary defect with a proved
statistical theorem \cite{EN}, and $90+lc+cr+lcr$ is Rule~146. So the literature has solved every
symmetric member of both families and the symmetrised chiral pair of one of them, and has nothing on a
single chiral monomial. The structural reason is visible in the list of techniques: each of them uses a
reflection symmetry, a sublattice or parity invariant, or a conserved particle number, and one chiral
monomial breaks all three at once.

There is a historical remark that belongs here, because the record turns out to agree with the
classification. Rule~30's first appearance in print is, as far as we can tell, a single unlabelled
panel in Fig.~7 of \cite[p.~608]{W83} --- we checked the panel against the true orbit --- and it is
neither named in the caption nor discussed nor measured anywhere in that paper. The reason is stated on
its p.~603: the analysis is restricted to the $32$ rules called ``legal'', and Rule~30 is excluded
because it is not symmetric under reflection, $110\mapsto0$ against $011\mapsto1$. The body then
records the expectation that the excluded rules ``introduce no qualitatively new features''. So the
exact property that removed Rule~30 from the first systematic survey of these automata is the property
that this remark identifies as the discriminator between the tractable members of its
family and the two intractable ones.

The classification by reflection symmetry itself is published, and we should use its vocabulary rather
than invent one. Wuensche and Lesser's atlas sorts the elementary rules as \emph{symmetric},
\emph{semi-asymmetric} and \emph{full-asymmetric} under reflection, tabulated in
\cite[\S4.3]{Martinez}, and the lists agree with the table above at every entry: $150$, $105$, $54$,
$22$ and $90$ are symmetric, and Rule~30 is semi-asymmetric. The term of art is therefore
\emph{asymmetric}, not chiral, and we keep the word chiral only for the algebraic object of
Corollary~\ref{du:cor:C}, the mirror-antisymmetric part of the source. There is also independent
corroboration from a different direction: Wuensche's four \emph{chain rules}, reached through the
$Z$-parameter and not through any symmetry consideration, are $30$, $45$, $106$ and $154$, every one
semi-asymmetric and not one symmetric. What is ours here is therefore narrower than the observation
that asymmetry matters: it is the single-monomial decomposition of the family, the exact cancellation
of Corollary~\ref{du:cor:C}, and the correlation with what the literature has managed to solve.
\end{remark}

There is a second reason Rule~18 is the wrong precedent for Rule~30, and it is quantitative. The
technique behind every rigorous Rule~18 result is that the nonlinear rule coincides with a linear one
on an invariant set --- the sublattice or invariant-set linearization of Grassberger, Jen and, in its
modern form, Rupe and Crutchfield \cite{RC}. For Rule~30 that set is as small as it can be.

\begin{proposition}[no nontrivial additive subdynamics; \EXACT]\label{du:nosubdyn}
Since $30=150+cr$, Rule~30 and Rule~150 agree on a configuration exactly when it contains no factor
$11$, so the set on which the two rules agree for all time is the largest Rule-30-invariant subshift
inside the golden-mean shift. That set is
\[
\{0^{\Z},\ (01)^{\Z},\ (10)^{\Z}\},
\]
three points, of zero entropy. Moreover the ordered background actually present in the seed's orbit is
not one of them.
\end{proposition}
\begin{proof}[Computation]
The argument is a window exhaustion, with a cyclic sweep as corroboration. Note first that a cyclic
sweep alone would not suffice: the set in question is closed and invariant, and a nonempty subshift need
not contain a periodic point, so excluding periodic points leaves Sturmian possibilities open. What is
needed is a bound on the language.

Take all $2^{23}$ words of length $23$ and iterate, at each step testing only the determined region,
which loses one cell at each end per step, for the factor $11$. The first test alone, at $t=0$, leaves
the $11$-free words, $75\,025=F_{25}$ of them; the second, at $t=1$, leaves six; and those same six
survive all eleven steps. They are $0^{23}$, the two checkerboards, and the three words whose only ones
sit at positions $0$ and $22$, the two cells the determined region never covers at any positive time.
In particular every length-$21$ factor of a configuration in the set is $0^{21}$ or a checkerboard,
which forces the set to be contained in the three configurations listed; and each of the three is
indeed fixed by Rule~30 and $11$-free, the checkerboards because every cell there has exactly one
neighbouring one. Hence equality.

The cyclic sweep agrees and extends the range: for every period $p\le29$, taking all $2^{p}$
configurations, keeping the $11$-free ones ($1\,149\,851$ at $p=29$) and iterating $400$ steps while
discarding any configuration whose image contains $11$, leaves exactly one word for odd $p$ and three
for even $p$ --- $0^{p}$ and, when $2\mid p$, the two checkerboards --- and nothing else at any period.
For the last clause of the statement, the $40$ cells at the light-cone edge are the same word
$1100100011100111100011110100011111101000$ at $t=200$ and at $t=400$, and it contains $11$ nine times,
so the quadratic term fires throughout the ordered region.
\end{proof}

\begin{remark}[the contrast with Rule 90, and why the precedents are on the other side]\label{du:rule90side}
The reason the linearization industry has been productive on Rule~90 and not on Rule~150 is that on
Rule~90 the perturbation is switched off. From the single seed, Rule~90's orbit is supported on
$x\equiv t\pmod 2$, so no two adjacent ones ever occur and no monomial in $\{lc,cr,lcr\}$ ever fires:
rules $18$, $26$, $82$, $146$, $154$, $210$ and $218$ are therefore \emph{identical} to Rule~90 from a
seed, which we verified over $40$ steps, with zero adjacent-one pairs in those rows. On the Rule~150
side the count of adjacent-one pairs in the same $40$ rows is $208$, and rules $30$, $86$, $54$ and
$22$ all diverge from Rule~150 at $t=2$, Rule~150's first row being $111$, which fires every monomial
at once. So the whole class of results that make a nonlinear rule tractable by exhibiting it as a
linear rule on an invariant set lives on the side where the perturbation is off, and
Proposition~\ref{du:nosubdyn} says Rule~30 is not merely on the other side but as far from that side as
the golden-mean shift allows. This raises rather than lowers the standing of Theorem~\ref{du:thm:B}: it
is not a restriction of the dynamics to where the rule is linear, it is a genuine variation-of-constants
formula with a live source.
\end{remark}

\begin{remark}[the literature slot this Part occupies; \KNOWN\ framing]\label{du:fuks}
The decomposition used throughout this Part has a name in the cellular-automaton literature and an
explicit open problem attached to it. Fuk\'s \cite{Fuks25,Fuks23} sets up exactly our situation: write
the local rule as $f=g+b$ with $g$ a rule whose iterate is known in closed form and $b$ a perturbation,
expect a solution of the form $[F^n(x)]_j=[G^n(x)]_j+P(x,n,j)$ by analogy with variation of constants,
and then states flatly that ``we do not know any general method for finding $P(x,n,j)$''
\cite[\S2]{Fuks25}. He carries the programme out for $156=140+b$ with $b$ supported on a single word,
and his appendix lists solution formulae for more than sixty elementary rules --- but always with a
tame $g$. Our Theorem~\ref{du:thm:B} and Corollary~\ref{du:cor:C} are an identification of that $P$ for
$g=$ Rule~150, which is chaotic, and the honest statement of novelty is therefore narrow: the
decomposition and the question are \KNOWN\ and named, and what is ours is the closed form for a chaotic
$g$, together with the chirality obstruction that says which part of $P$ is visible at the centre. His
verification technique --- guess the perturbation term, write it as a density polynomial, induct on $n$
--- is also the natural template for checking any refinement of ours.

Where the topological classification actually places this rule is settled in print, and it is worth
recording because it also ends the enquiry. Sch\"ule and Stoop give the complete K\r{u}rka table for all
$88$ minimal elementary rules \cite{SS12}: the positively expansive class is exactly $\{90,105,150\}$,
and Rule~30 sits one rung below, sensitive but not positively expansive. Their Proposition~13 is the
elementary-rule form of the equivalence between positive expansivity and bipermutivity, and their
Proposition~3 notes that blocking words of length at most $4$ suffice for elementary rules --- the
finite search that the undecidability of the classification in general forbids in the large.

Two consequences. First, no rule is excluded from solvability \emph{by} its topological class. The
criterion sometimes attached to that hierarchy is sufficiency and not necessity: Fuk\'s writes that the
initial value problem is solvable for almost all rules possessing some equicontinuity property, while
noting that simpler methods usually suffice in those cases \cite[\S2]{Fuks25}. Rule~90 is positively
expansive and solvable in closed form, so it is out of scope there rather than a counterexample. What
Fuk\'s does name as the obstruction is defects performing a pseudo-random walk, and that does not reach
us either, since the mechanism assumes a constant defect population while ours grows by
Theorem~\ref{fr:thm:A}.

Second, and this is why we stop here: across that whole table Rule~30 and Rule~90 agree on
sensitivity, surjectivity, transitivity and Devaney chaos, and differ on exactly one label, positive
expansivity. The hierarchy carries one bit of information about the distinction we care about, and that
bit is bipermutivity, which we already use throughout. There is one place where the separation runs our
way instead: an affine rule cannot be globally universal, with Rule~90 named, while Rule~30's status is
open.
\end{remark}

\begin{remark}[a worked nonlinear control: Rule 22]\label{du:rule22}
Rule~22 has algebraic normal form $a+b+c+abc$ and Rule~30 has $a+b+c+bc$, so both are Rule~150 plus a
single monomial: the symmetric cubic in one case, the chiral quadratic in the other. The Duhamel
identity therefore applies verbatim to Rule~22 with the cubic source $q^s_x=u^s_{x-1}u^s_xu^s_{x+1}$,
and Rule~22 is exactly solvable from the single seed: Chan-L\'opez and Mart\'in-Ruiz \cite{CLMR} give
the closed form $Q_{2j}(x)=(x^2+x^{-2})^j$, $Q_{2j+1}(x)=(x^{-1}+1+x)Q_{2j}(x)$ for the row
polynomials, which we have verified against simulation for every $t<129$, together with the support
count $2^{s_2(j)}$ on even rows, $s_2$ the binary digit sum. Since the constant term of
$(x^2+x^{-2})^j$ is $\binom{j}{j/2}\bmod2$, which vanishes for $j\ge1$, the centre column of Rule~22
is $1,1,0,0,\dots$.

This is a complete test of the machinery above, and it isolates what makes Rule~30 hard. \EXACT: the
Duhamel sum reproduces the Rule~22 centre column for all $t<200$; the cubic source is supported
exactly on the odd rows and carries substantial weight, $1251$ ones for $t<200$, yet contributes
nothing beyond $t=1$. Corollary~\ref{du:cor:C} says why. Rule~22 is symmetric, so its single-seed
orbit and hence its source are mirror-symmetric, and the chiral part $\chi$ vanishes identically
(\EXACT: chiral weight $0$ for $t<200$, against $7666$ for Rule~30 over the same range). By
Corollary~\ref{du:cor:C} the centre column then depends on the source only through its own centre
line, which here is the single bit $f^s_0=[s=1]$, and Theorem~\ref{du:thm:B} gives the answer in one
line:
\[
  \sum_t c_t y^t=\frac{1}{1+y}\bigl(\hat g(\alpha)+y\hat f(y,\alpha)\bigr)=\frac{1+y^2}{1+y}=1+y ,
\]
so $c_0=c_1=1$ and $c_t=0$ thereafter, which is the closed form's consequence recovered without the
closed form. The contrast is the content: a nonlinear rule with a source of weight $1251$ can be
completely invisible to its own centre column, and what makes Rule~30's centre column resistant is
not the size of its source but the $7666$ units of \emph{chiral} weight the mirror-symmetric kernel
can see. A symmetric rule is therefore the right reference object for Rule~30 precisely because it is
the $\chi\equiv0$ case of Corollary~\ref{du:cor:C}.
\end{remark}


\begin{remark}[verification]
Theorem~\ref{du:thm:A} was checked for the seed and source $q$ for all $t<60$ in both
parities, and Theorem~\ref{du:thm:B} reproduces $c_t$ for all $t<160$, with $\alpha$
computed to order $160$ as the fixed point of $\alpha = y(1+\alpha+\alpha^2)$
(independent code, 2026-09-16).
\end{remark}

\subsection{What remains}

The question ``does the quadratic source decimate?'' has two readings. Whether the
orbit $\mathcal O$ itself is finite is the weak one: $\ker L$ is large, and
infinitely many distinct sources can give the same kernel element. Whether
$\mathcal O$ is finite modulo $\ker L$ is, by Corollary~\ref{du:cor:orbit}, exactly the
finiteness of $\mathcal K(c)$. Equivalently, by Theorem~\ref{du:thm:B} and
Corollary~\ref{du:cor:C}, $G$ is algebraic over $\mathbb F_2(y)$ if and only if the chiral
series $H$ is algebraic over $\mathbb F_2(\alpha)$. The certified bounds of Section~\ref{du:bounds}
are consistent with $\mathcal K(c)$ being infinite but cannot establish it. The action
of the Cartier operators on the chiral data is not what is missing: $\Lambda$ and $E$ commute with
the mirror $x\mapsto-x$ and preserve the rows with $h_x=h_{-x}$ and $h_0=0$, which are exactly the rows
with zero centre and zero chiral part, so $D_0$ and $D_1$ descend to the half-line data $(g_0,\chi_g;\,f^s_0,\chi^s)$ (\PROVED), and Section~\ref{du:chiral}
writes the descended maps out and checks them (Proposition~\ref{du:chiralcheck}). What remains is to
exhibit an invariant of that recursion that cannot take only finitely many values. We do not
claim such an invariant here.


\begin{remark}[the control the growth measurements were missing; \EXACT]\label{ex:l90}
Every other claim in this paper is checked against Rule~90, and the interface language was not. It
should have been, because the comparison is the sharpest separation we have between our rule and the
linear ones, and it needs no fitting.

For Rule~90 and Rule~150 the interface language is the \emph{full} free monoid: $|L_k|=2^k$ exactly,
verified to $k=40$ by two independent methods, so the entropy is exactly one bit per zero-row. For
Rule~30 it is $2,3,5,8,12,17,25,36,\dots$, reaching $|L_{60}|=103\,220$. Two estimators of the entropy
have to be kept apart here, by the same standing rule that corrected the Rule~150 growth exponent: the
block estimate $k^{-1}\log_2|L_k|$ reads $0.2776$ at $k=60$, but it converges from above and
logarithmically slowly, whereas the local slope $\log_2\bigl(|L_{k+1}|/|L_k|\bigr)$ reads $0.1295$ at $k=60$ ($0.1535$ at $k=46$), a
factor $2.1$ apart on the same data. Both fall, and the local slope is the one to quote. Neither
carries the separation, which needs no estimator at all: Rule~90's entropy is exactly one bit, clearing
both by more than a factor of three. The separation is visible without any asymptotics at all: for a linear rule $|L_k|$ is a power
of two at every depth, and Rule~30's is a power of two at only two of the twenty depths we list, with
$|L_{60}|=103\,220$ lying strictly between $2^{16}$ and $2^{17}$. That is an integrality obstruction at
essentially every depth, and it is exact.

The corollary is worth stating because it inverts the natural reading. Rule~90's fibre algebra
\emph{is} the free algebra on two generators; Rule~30's contains no free subalgebra at all. The rule
with the intermediate growth is the constrained one, not the rich one --- our language is smaller than
the linear rules', not larger, and the difficulty of Rule~30 does not show up here as extra freedom.
\end{remark}

\section{The chiral recursion in practice}\label{du:chiral}

\begin{proposition}[\EXACT]\label{du:chiralcheck}
Let the state be $(g_0,\chi_g;\,f^s_0,\chi^s)$ with $\chi^s(x)=f^s_x+f^s_{-x}$, $x\ge1$. The maps $D_0,D_1$ of Theorem~\ref{du:thm:A} act on this state by
$\chi'^{\sigma}=E\bigl(\chi^{2\sigma+1}+\Lambda\chi^{2\sigma}\bigr)$, $f'^{\sigma}_0=f^{2\sigma+1}_0+f^{2\sigma}_0+\chi^{2\sigma}(1)$ for $D_0$, and the shifted analogues for $D_1$ (with $g'$ absorbing $f^0$), where $(\Lambda\chi)(x)=\chi(x-1)+\chi(x)+\chi(x+1)$ with $\chi(0)=0$ and $(E\chi)(x)=\chi(2x)$. Evaluating Theorem~\ref{du:thm:B} on this state reproduces $c_t$ for $t<64$ and every kernel element $(c_{r+j2^k})_j$ for all $126$ words of length $\le6$ (script \texttt{s3\_chiral.py}).
For depths $1$ to $9$ the orbit is injective on the chiral windows (rows $<16$, $|x|\le16$): all $2^k$ words give distinct windows (depth $10$ exhausts the $16$-row window and is uninformative). For every word to depth~$8$ the chiral data has a nonzero entry within its first four rows.
\end{proposition}

\begin{remark}[what this rules out]\label{du:zerorun}
The zero-run criterion is the elementary observation that if for every $M$ there are $k$ and $r<2^k$ with $c_{r+j2^k}=0$ for all $j<M$ and $c_{r+j2^k}=1$ for some $j\ge M$, then these progressions are nonzero kernel elements of unbounded $y$-adic valuation, so $\mathcal K(c)$ is infinite and $c$ is neither $2$-automatic nor eventually periodic (\PROVED; the clause $j\ge M$ is needed, since an identically zero progression is the zero kernel element). The leading zero run of a kernel element (its $y$-adic valuation) is the only invariant of the orbit that is visible in the data, and it behaves like a coin across residues: over the $4096$ residues $r<2^{12}$ at level $k=12$, read on the first $2^{20}$ terms of $c$, $\Pr[\text{run}\ge m]$ matches $2^{-m}$ to sampling accuracy for every $m\le12$ (\MEASURED), so the criterion's hypothesis holds at the coin's rate as far as it can be tested. Since the chiral data never vanishes near the top, long zero runs arise from cancellation inside the $\alpha$-weighted sum of nonzero data, not from vanishing data. Any invariant that proves the kernel infinite must therefore be a class function of the series $G_w$ itself, not a support or valuation of the chiral arrays. None is known. \OPEN.
\end{remark}

\begin{proposition}[no small linear lift; \EXACT]\label{du:lift}
There is no radius-one linear cellular automaton over $\Z/4$, $\Z/8$, $\mathbb F_4$, $\mathbb F_8$ or $\F^2$ (with $2\times2$ matrix coefficients), no cellwise recolouring to $\{0,1\}$ and no lift of the seed, that reproduces the Rule~30 space-time diagram for $40$ rows (\texttt{s3\_linearlift.py}; $19{,}840$ lifts). \KNOWN\ obstruction: a linear automaton over a finite abelian group from a finite seed has an automatic space-time diagram \cite{AvHPS}, and a cellwise coding preserves automaticity, so any such lift would make $c$ automatic, with an automaton of at least the $138\,672$ states forced by the census above. The implication is in fact an equivalence: over $\mathbb F_q$ in characteristic $p$, a sequence is $p$-automatic if and only if it is a column of a linear cellular automaton with memory from eventually periodic initial conditions (Rowland--Yassawi \cite{RY}, \KNOWN), so the lift search is a search for automaticity in other coordinates. Nesme \cite{Nesme} extends \cite{AvHPS} to commutative monoids and, more to the point here, to $k$-automatic initial configurations, so the conclusion strengthens from ``no lift of the seed'' to ``no $k$-automatic lift of the seed''.
\end{proposition}

\part{The front of the diagonal defects}

\noindent\emph{Source: \texttt{front\_section.tex} (referee dialogue, 2026-09-16/17), labels prefixed \texttt{fr:}; Section~\ref{fr:walk} writes out the unbounded-periods theorem in the diagonal frame; the theorem itself is \KNOWN\ (Gravner--Griffeath \cite{GG}, Proposition~3.3), as recorded in the ledger on 2026-09-17.}

% Part III, from front_section.tex (labels prefixed fr:)
\section{The front of the diagonal defects}\label{fr:sec:front}

Throughout, arithmetic is in $\F$, $u^t_x$ is the Rule~30 evolution
\[
u^{t+1}_x=u^t_{x-1}+u^t_x+u^t_{x+1}+u^t_xu^t_{x+1},\qquad u^0_x=[x=0],
\]
and $c_t=u^t_0$ is the centre column. For a finitely supported
$f\colon\N\to\F$ put $\deg f=\max\{t:f(t)=1\}$ and $\deg 0=-1$.
For $a\in\mathbb{Z}$, $(a)_+=\max(a,0)$.

\subsection{Diagonals, background and defect}

\begin{definition}[Diagonals]\label{fr:def:diag}
For $d\in\mathbb{Z}$ and $t\ge0$ let $D_d(t)=u^t_{d-t}$. Thus $D_0$ is the left edge of the light cone,
and $D_d$ passes through the centre cell $x=0$ at time $t=d$, so $D_d(d)=c_d$.
\end{definition}

\begin{lemma}[\PROVED]\label{fr:lem:diag}
\begin{enumerate}
\item[(a)] $D_d(t+1)=D_{d-2}(t)+D_{d-1}(t)+D_d(t)+D_{d-1}(t)D_d(t)
          =D_{d-2}(t)+\bigl(D_{d-1}(t)\vee D_d(t)\bigr)$.
\item[(b)] (Light cone) $D_d(t)=0$ whenever $d>2t$; $D_d\equiv0$ for $d<0$; $D_d(0)=[d=0]$.
\item[(c)] $D_0\equiv1$ and $D_1(t)=[t\ge1]$.
\end{enumerate}
\end{lemma}
\begin{proof}
(a) Put $x=d-(t+1)$ in the rule; the neighbours $x-1,x,x+1$ are the cells
$(d-2)-t,(d-1)-t,d-t$. (b) $u^t_x=0$ for $|x|>t$ by induction on $t$, and
$d-t>t\iff d>2t$; $d<0$ gives $d-t<-t$. (c) From (a) and (b):
$D_0(t+1)=D_0(t)$, $D_0(0)=1$; $D_1(t+1)=1+D_1(t)+D_1(t)=1$.
\end{proof}

\begin{proposition}[Eventual periodicity and reset formula; \PROVED]\label{fr:prop:period}
Let $d\ge0$.
\begin{enumerate}
\item[(a)] (Reset formula) If $s<t$, $D_d(s)=1$ and $D_d(u)=0$ for $s<u<t$, then
$D_{d+1}(t)=1+\sum_{u=s}^{t-1}D_{d-1}(u)$. If $D_d(u)=0$ for all $u<t$, then
$D_{d+1}(t)=D_{d+1}(0)+\sum_{u<t}D_{d-1}(u)$.
\item[(b)] Every $D_d$ is eventually periodic, with minimal eventual period $p_d$ a power of $2$
(eventual periodicity of the diagonals, with period a power of $2$, is stated by Rowland \cite[\S5]{Rowland} as a consequence of Jen \cite[Thm.~4]{Jen86}, which we have not been able to verify at source, so the attribution is second-hand; the proof below does not use it, and Jen's 1990 paper does not treat diagonals).
Let $\tau_d=\min\{t\ge0: D_d(u+p_d)=D_d(u)\ \forall u\ge t\}$, $T=\max(\tau_{d-1},\tau_d)$,
$q=\max(p_{d-1},p_d)$.
\item[(c)] If $D_d$ is not eventually $0$, then $p_{d+1}\mid q$ and $\tau_{d+1}\le T+p_d$.
\item[(d)] If $D_d$ is eventually $0$, put $w=\sum_{u=T}^{T+p_{d-1}-1}D_{d-1}(u)$. Then
$p_{d+1}\mid p_{d-1}$ if $w=0$, $p_{d+1}=2p_{d-1}$ if $w=1$, and $\tau_{d+1}\le T$.
\end{enumerate}
\end{proposition}
\begin{proof}
(a) By Lemma~\ref{fr:lem:diag}(a), $D_{d+1}(u+1)=D_{d-1}(u)+1$ if $D_d(u)=1$ and
$D_{d+1}(u+1)=D_{d-1}(u)+D_{d+1}(u)$ if $D_d(u)=0$; iterate from $u=s$.
(b)--(d) by induction on $d$, starting from Lemma~\ref{fr:lem:diag}(c) ($p_0=p_1=1$).
(c) Let $s_0=\min\{s\ge T:D_d(s)=1\}\le T+p_d-1$. For $t\ge s_0+1$ let $s(t)$ be the last
$1$ of $D_d$ before $t$; then $s(t)\ge T$, $s(t+q)=s(t)+q$ and, as $p_{d-1}\mid q$,
$\sum_{u=s(t)+q}^{t+q-1}D_{d-1}(u)=\sum_{u=s(t)}^{t-1}D_{d-1}(u)$. By (a),
$D_{d+1}(t+q)=D_{d+1}(t)$ for $t\ge s_0+1$.
(d) For $t\ge T$, $D_{d+1}(t+1)=D_{d+1}(t)+D_{d-1}(t)$, so
$D_{d+1}(t+p_{d-1})=D_{d+1}(t)+w$; if $w=1$ the minimal period divides $2p_{d-1}$ and not
$p_{d-1}$, and is a power of $2$, hence equals $2p_{d-1}$.
\end{proof}

\begin{remark}[status of the reset formula, and of the transient]\label{fr:reset:status}
We have not found the reset formula, or any statement eliminating the recursion in the diagonal index,
in print; Brunnbauer \cite{Brunnbauer} says of his own recursion that ``the recursion
over $n$ cannot be easily removed''. The one-step law that the formula sums is \KNOWN: it is Rowland's \cite[Lemma~3]{Rowland}, and it is also in \cite{Pat} (dated 19 August 2026), in the Lean lemmas of \cite{Dib} and in Nersissian \cite{Nersissian}. What we did not find stated is the summed form. In the frozen case ours collapses to the parity criterion for
period doubling of Wolfram, Rowland and Gravner--Griffeath, so it is a strict generalisation of a
published statement.

One structural fact does have a published explanation: the periods are powers of two because the orbit
closure is an odometer. Coven, Pivato and Yassawi prove that a left-permutive rule with no memory and
local rule $t_0+\psi(t_1,\dots,t_r)$ over $\Z/p$, started from a right-fixed tail with infinite orbit,
has orbit closure topologically conjugate to the $p$-adic odometer with every scale factor $p$
\cite[Thm.~4]{CPY}. Rule~30 as $\sigma\circ F$ satisfies every hypothesis, and the infinite-orbit
condition is Theorem~\ref{fr:unbounded}; instantiating their scale data returns our own doubling
positions. It applies verbatim to Rule~90, so it explains a shared phenomenon rather than separating
the rules, and \emph{where} the doublings occur is not explained by it --- their index is defined as the
doubling position, with no bound on it anywhere.

Of the \emph{transient} there is no theory at all. \cite{GG} contains no occurrence of ``transient'' or
``preperiod'' in its body, Rowland fixes a common starting time and bounds only the period, and the
preperiods exist as data in OEIS A363346 \cite{A363346} ($1\,000$ terms) and, for the $78\,288$ diagonals $0\le d\le78\,287$, in the data file of Wolfram's 2019 prize post \cite{Wolfram19}. The one paper whose title promises otherwise, Jen's on
transience and dislocations, maps certain nonlinear rules onto linear template systems and bounds
collision times within that correspondence, which is valid only while the \emph{number} of dislocations
stays constant; it fails for us twice, since Proposition~\ref{du:nosubdyn} denies any nontrivial set on
which Rule~30 agrees with a linear rule, and Theorem~\ref{fr:thm:B} makes our defect grow. It is
therefore the third instance of one mechanism, with the Rule~18 linearisation and the Rule~22 defect
calculus.

The shape of the gap can be said exactly. For nonlinear rules, every published mechanism we have found for bounding an infinite-lattice
preperiod runs through a \emph{blocking word}, which exists precisely when the rule has an
equicontinuity point --- precisely when the transient is inert. Rule~30 is sensitive and has no blocking
word. Linear rules are handled otherwise, by algebra: for Rule~90 from the seed, Kummer's theorem gives $\tau_d=0$ on every
diagonal (Remark~\ref{fr:nofront}), although Rule~90 is permutive, hence sensitive, and has no blocking word either.
A nonlinear alternative is sometimes credited to Li and Nordahl for Rule~110 \cite{LN}, and does not survive
reading: their paper contains no theorem, its transient is a finite-lattice quantity fitted as
$T_{\mathrm{ave}}\sim N^{1.08}$, and its glider decay is a least-squares fit from uncorrelated random
initial states. So for nonlinear rules the field's only preperiod tool works only where there is nothing to explain, and
the general results in the area --- finite-lattice hardness, and undecidability of ever reaching a
periodic regime --- quantify over rules and say nothing about one fixed rule, exactly as
Remark~\ref{p3:form} distinguishes for Problem~3.
\end{remark}


\begin{definition}[Background and defect]\label{fr:def:bg}
$r_d$ is the unique $p_d$-periodic sequence on $t\ge0$ with $r_d(t)=D_d(t)$ for $t\ge\tau_d$;
the \emph{defect} is $\Delta_d=D_d+r_d$. For $d<0$: $D_d=r_d=\Delta_d=0$, $p_d=1$, $\tau_d=0$.
Put $P_d=\max(p_{d-1},p_d,p_{d+1})$, a common period of $r_{d-1},r_d,r_{d+1}$.
\end{definition}

\begin{lemma}[\PROVED]\label{fr:lem:bgdef}
\begin{enumerate}
\item[(a)] $\Delta_d$ has finite support and $\deg\Delta_d=\tau_d-1$.
\item[(b)] (Background equation) $r_d(t+1)=r_{d-2}(t)+r_{d-1}(t)+r_d(t)+r_{d-1}(t)r_d(t)$ for all $t\ge0$.
\item[(c)] (Defect equation) For $d\ge0$ and all $t\ge0$,
\[
\Delta_{d+1}(t+1)=h_{d+1}(t)+\bigl(1+D_d(t)\bigr)\Delta_{d+1}(t),\qquad
h_{d+1}:=\Delta_{d-1}+(1+r_{d+1})\Delta_d ,
\]
and $\Delta_{d+1}(0)=r_{d+1}(0)$.
\item[(d)] There is no $d\ge0$ with $r_d\equiv r_{d+1}\equiv0$.
\end{enumerate}
\end{lemma}
\begin{proof}
(a) If $\tau_d\ge1$ and $\Delta_d(\tau_d-1)=0$ then
$D_d(\tau_d-1)=r_d(\tau_d-1)=r_d(\tau_d-1+p_d)=D_d(\tau_d-1+p_d)$, contradicting minimality of $\tau_d$.
(b) The identity holds for $t\ge\max(\tau_{d-2},\tau_{d-1},\tau_d)$ by Lemma~\ref{fr:lem:diag}(a); both sides are
periodic with a common period (powers of $2$), so it holds for all $t$.
(c) Subtract (b) at $d+1$ from Lemma~\ref{fr:lem:diag}(a) at $d+1$:
$\Delta_{d+1}(t+1)=\Delta_{d-1}+\Delta_d+(1+D_d)D_{d+1}+(1+r_d)r_{d+1}$, and
$(1+D_d)D_{d+1}+(1+r_d)r_{d+1}=(1+D_d)\Delta_{d+1}+r_{d+1}\Delta_d$.
$D_{d+1}(0)=0$ for $d\ge0$.
(d) If $r_d\equiv r_{d+1}\equiv0$, (b) at $d+1$ gives $r_{d-1}\equiv0$; descending gives
$r_0\equiv r_1\equiv0$, contradicting Lemma~\ref{fr:lem:diag}(c) ($r_0\equiv r_1\equiv1$).
\end{proof}

\begin{definition}[Front quantities]\label{fr:def:front}
For $d\ge0$:
\begin{gather*}
M_d=\max(\deg\Delta_{d-1},\deg\Delta_d),\qquad S_d=M_d-d,\qquad H_{d+1}=\deg h_{d+1},\\
\kappa_{d+1}=M_d-H_{d+1},\qquad \beta_{d+1}=\Delta_{d+1}(M_d+1).
\end{gather*}
$M_d+1=\max(\tau_{d-1},\tau_d)$ is the \emph{front}; $S_d$ the \emph{slack};
$\kappa_{d+1}\ge0$ (since $h_{d+1}(t)=0$ for $t>M_d$) the \emph{cancellation depth}.
If $\deg\Delta_{d-1}\le\deg\Delta_d$ put $\mu_{d+1}=(\deg\Delta_{d-2}-H_{d+1})_+$, otherwise $\mu_{d+1}=0$.
\end{definition}

\subsection{Exact degree formula}

\begin{theorem}[A: exact degree formula; \PROVED]\label{fr:thm:A}
Let $d\ge0$, $H=H_{d+1}$, $\pi=\Delta_{d+1}(H+1)$ and $\rho=\min\{t\ge H+1:D_d(t)=1\}\in\N\cup\{\infty\}$.
\begin{enumerate}
\item[(a)] $\Delta_{d+1}(t)=\sum_{u=s(t)}^{t-1}h_{d+1}(u)$, where $s(t)=\max\{u<t:D_d(u)=1\}$;
if there is no such $u$, $\Delta_{d+1}(t)=r_{d+1}(0)+\sum_{u<t}h_{d+1}(u)$.
In particular $\pi$ is the parity of $h_{d+1}$ over $[s(H+1),H]$, plus $r_{d+1}(0)$ when $D_d$ has no $1$ below $H+1$.
\item[(b)] If $\pi=0$ then $\deg\Delta_{d+1}=H$, and $D_d(H)=0$ when $H\ge0$.
\item[(c)] If $\pi=1$ then $\rho<\infty$ and $\deg\Delta_{d+1}=\rho$.
\item[(d)] $h_{d+1}(t)=\Delta_{d+1}(t+1)+(1+D_d(t))\Delta_{d+1}(t)$; hence
$H_{d+1}\le\deg\Delta_{d+1}$, with equality iff $\Delta_{d+1}=0$ or $D_d(\deg\Delta_{d+1})=0$.
\end{enumerate}
\end{theorem}
\begin{proof}
(a) Iterate Lemma~\ref{fr:lem:bgdef}(c): at a $1$ of $D_d$ the accumulator is overwritten by
$h_{d+1}(s)$, between $1$s it adds $h_{d+1}$.
(b),(c) For $t\ge H+1$, $\Delta_{d+1}(t+1)=(1+D_d(t))\Delta_{d+1}(t)$.
If $\pi=1$, $\Delta_{d+1}=1$ on $[H+1,\rho]$ and $\Delta_{d+1}(\rho+1)=0$; $\rho=\infty$ contradicts
Lemma~\ref{fr:lem:bgdef}(a). If $\pi=0$, $\Delta_{d+1}\equiv0$ on $[H+1,\infty)$; if $H\ge0$,
$0=\Delta_{d+1}(H+1)=1+(1+D_d(H))\Delta_{d+1}(H)$ forces $\Delta_{d+1}(H)=1$, $D_d(H)=0$.
(d) Rearrange Lemma~\ref{fr:lem:bgdef}(c); at $t_0=\deg\Delta_{d+1}$ it gives $h_{d+1}(t_0)=1+D_d(t_0)$.
\end{proof}

\begin{corollary}[Modes; \PROVED]\label{fr:cor:modes}
$\beta_{d+1}=1\iff(\pi=1\text{ and }\rho>M_d)$. If $\beta_{d+1}=0$, exactly one of:
\textup{(E)} $\pi=0$ and $\deg\Delta_{d+1}=M_d-\kappa_{d+1}$;
\textup{(R)} $\pi=1$, $\rho\le M_d$ and $\deg\Delta_{d+1}=\rho\in[M_d-\kappa_{d+1}+1,M_d]$.
In all cases $\deg\Delta_{d+1}\ge M_d-\kappa_{d+1}$. If $\deg\Delta_{d-1}>\deg\Delta_d$ then $\kappa_{d+1}=0$.
\end{corollary}
\begin{proof}
$h_{d+1}=0$ on $(H,M_d]$, so $\Delta_{d+1}(M_d+1)=\pi\prod_{H<u\le M_d}(1+D_d(u))=\pi[\rho>M_d]$.
The rest is Theorem~\ref{fr:thm:A}(b),(c). If $\deg\Delta_{d-1}>\deg\Delta_d$ then
$h_{d+1}(M_d)=\Delta_{d-1}(M_d)=1$.
\end{proof}

\subsection{The front is monotone}

\begin{theorem}[B; \PROVED; (a), (c) and (d) were public earlier as statements, see the paragraph after the proof]\label{fr:thm:B}
Let $d\ge0$.
\begin{enumerate}
\item[(a)] $M_d\le M_{d+1}$.
\item[(b)] If $\beta_{d+1}=0$ then $M_{d+1}=M_d$.
\item[(c)] If $\beta_{d+1}=1$ then $r_d\not\equiv0$ and
$\deg\Delta_{d+1}=M_{d+1}=\min\{t\ge M_d+1:r_d(t)=1\}$; hence the jump
$J_{d+1}:=M_{d+1}-M_d\in[1,p_d]$.
\item[(d)] If $r_d\equiv0$ then $\beta_{d+1}=0$.
\end{enumerate}
\end{theorem}
\begin{proof}
(a) By Lemma~\ref{fr:lem:bgdef}(c),
$\Delta_{d-1}(t)=\Delta_{d+1}(t+1)+(1+r_{d+1}(t))\Delta_d(t)+(1+D_d(t))\Delta_{d+1}(t)$,
which vanishes for $t>M_{d+1}$; so $\deg\Delta_{d-1}\le M_{d+1}$ and $M_d\le M_{d+1}$.
(b) $\Delta_{d+1}(M_d+1)=0$ and $h_{d+1}=0$ beyond $M_d$ give $\deg\Delta_{d+1}\le M_d$, so $M_{d+1}\le M_d$.
(c) For $t\ge M_d+1$: $h_{d+1}(t)=0$ and $D_d(t)=r_d(t)$, so
$\Delta_{d+1}(t+1)=(1+r_d(t))\Delta_{d+1}(t)$ with $\Delta_{d+1}(M_d+1)=1$. Finite support forces a
$1$ of $r_d$ in $[M_d+1,M_d+p_d]$, and $\deg\Delta_{d+1}>M_d\ge\deg\Delta_d$.
(d) Contrapositive of (c).
\end{proof}

\noindent\emph{Provenance of Theorem~\ref{fr:thm:B}.} Parts (a), (c) and (d) are \KNOWN\ as statements. The records structure of the front ($M_d+1=\max_{j\le d}\tau_j$, which is (a) with Definition~\ref{fr:def:front}), its monotonicity, the absence of a move at a frozen diagonal, and the split of a diagonal into background and defect (Definition~\ref{fr:def:bg}) are public: the code of Wolfram's 2019 prize post \cite{Wolfram19} takes a running maximum, and they are in \cite{Dib} (commits \texttt{0a1cda2} of 7 September 2026, \texttt{a8dbbc0} of 8 September 2026, \texttt{94e15f8} and \texttt{9c853e4}). The landing rule (c), that a move ends at the first black cell of the stripe, is in \cite{Dib} (commit \texttt{a8dbbc0}), where it is checked on $60\,065$ of $60\,065$ diagonals. The source of Part~III, which has this theorem, is dated 16--17 September 2026. We learned of \cite{Dib} on 1 October 2026 (ledger 10.165c2) and of this overlap on 8 October 2026 (ledger 10.167c2).

\subsection{The exact failure condition}

\begin{theorem}[C; \PROVED]\label{fr:thm:C}
Let $d\ge0$.
\begin{enumerate}
\item[(a)] $\tau_{d+1}\le d\iff\bigl(\kappa_{d+1}\ge S_d+1\text{ and }\Delta_{d+1}(d)=0\bigr)$.
\item[(b)] If $\tau_d\ge d$ and $\beta_{d+1}=1$ then $\tau_{d+1}\ge d+1$.
\item[(c)] If $\kappa_{d+1}\le S_d$ then $\tau_{d+1}\ge d+1$.
\end{enumerate}
\end{theorem}
\begin{proof}
(a) $\tau_{d+1}\le d$ iff $\Delta_{d+1}(t)=0$ for all $t\ge d$. By Theorem~\ref{fr:thm:A}(d) this implies
$h_{d+1}(t)=0$ for $t\ge d$; conversely, if $\Delta_{d+1}(d)=0$ and $h_{d+1}=0$ on $[d,\infty)$,
Lemma~\ref{fr:lem:bgdef}(c) propagates $\Delta_{d+1}(t)=0$. Finally $h_{d+1}=0$ on $[d,\infty)$ iff
$H_{d+1}\le d-1$ iff $\kappa_{d+1}\ge M_d-d+1$.
(b) $M_d\ge\deg\Delta_d=\tau_d-1\ge d-1$, and by Theorem~\ref{fr:thm:B}(c)
$\tau_{d+1}=M_{d+1}+1\ge M_d+2$.
(c) Immediate from (a).
\end{proof}

\subsection{The cancellation window}

\begin{lemma}[F: window slaving; \PROVED]\label{fr:lem:F}
Let $d\ge1$ with $\deg\Delta_{d-1}\le\deg\Delta_d=:b$, and $W=(H_{d+1},b]$. Put
$\omega_d=1+r_{d-1}+r_d(1+r_{d+1})$, a $P_d$-periodic sequence. For $t\in W$:
\[
\text{(F1)}\quad \Delta_{d-1}(t)=(1+r_{d+1}(t))\,\Delta_d(t),\qquad
\text{(F2)}\quad \Delta_d(t+1)=\Delta_{d-2}(t)+\omega_d(t)\,\Delta_d(t).
\]
\end{lemma}
\begin{proof}
(F1) is $h_{d+1}(t)=0$. Lemma~\ref{fr:lem:bgdef}(c) at diagonal $d$ reads
$\Delta_d(t+1)=\Delta_{d-2}(t)+(1+r_d)\Delta_{d-1}(t)+(1+r_{d-1}+\Delta_{d-1}(t))\Delta_d(t)$.
Insert (F1) and use $\Delta_d^2=\Delta_d$: the coefficient of $\Delta_d(t)$ becomes
$(1+r_d)(1+r_{d+1})+1+r_{d-1}+(1+r_{d+1})=1+r_{d-1}+r_d(1+r_{d+1})$.
\end{proof}

\begin{lemma}[G; \PROVED]\label{fr:lem:G}
Under the hypotheses of Lemma~\ref{fr:lem:F} let $m=\deg\Delta_{d-2}$ and $x=\max(m,H_{d+1})$.
Then $x\le b$, and if $x<b$:
$\Delta_d=1$ on $(x,b]$, $\omega_d=1$ on $[x+1,b-1]$, $\omega_d(b)=0$, and
$\Delta_{d-1}=1+r_{d+1}$ on $(x,b]$. Consequently $b-x\le P_d$ and, for every $d\ge1$,
\[
\kappa_{d+1}\le P_d+\mu_{d+1}.
\]
\end{lemma}
\begin{proof}
$m\le M_{d-1}\le M_d=b$ by Theorem~\ref{fr:thm:B}(a), and $H_{d+1}\le M_d$.
For $t\in(x,b]\subseteq W$, $\Delta_{d-2}(t)=0$, so (F2) gives $\Delta_d(t+1)=\omega_d(t)\Delta_d(t)$.
At $t=b$: $\Delta_d(b)=1$, $\Delta_d(b+1)=0$, so $\omega_d(b)=0$. Descending from $t=b-1$ to $x+1$,
$\Delta_d(t+1)=1$ forces $\omega_d(t)=\Delta_d(t)=1$. The last claim is (F1).
Since $\omega_d$ is $P_d$-periodic and $\omega_d(b)=0$, a run of $1$s of $\omega_d$ has length at most
$P_d-1$, so $b-1-x\le P_d-1$. Then $\kappa_{d+1}=(b-x)+(x-H_{d+1})\le P_d+\mu_{d+1}$.
If $\deg\Delta_{d-1}>\deg\Delta_d$, $\kappa_{d+1}=0$ by Corollary~\ref{fr:cor:modes}.
\end{proof}

\begin{remark}[Level two; identity \PROVED, bound \OPEN]\label{fr:rem:level2}
For $t,t+1\in W$, substituting (F1), (F2) into Lemma~\ref{fr:lem:bgdef}(c) at diagonal $d-1$ gives,
with $y=\Delta_d$,
\[
\Delta_{d-3}(t)=\xi(t)\,y(t+1)+\zeta(t)\,y(t)+(1+r_{d+1}(t))\,y(t)y(t+1),
\]
$\xi(t)=r_{d+1}(t+1)+r_{d-1}(t)$,
$\zeta(t)=(1+r_{d-2}(t))(1+r_{d+1}(t))+\omega_d(t)(r_{d-1}(t)+r_{d+1}(t))$.
Where $\Delta_{d-3}(t)=0$ this is a transition relation on $(y(t),t\bmod P)$. Since
$\Delta\equiv0$ solves Lemma~\ref{fr:lem:bgdef}(c), every back-solved relation is homogeneous in $y$ and the
transition $0\to0$ is always admissible: no level of the telescoping yields a finite bound on
$\mu_{d+1}$ by itself. A bound on $\mu_{d+1}$ is \OPEN.
\end{remark}

\subsection{Plateaus, stall runs and the row reading}

\begin{definition}[Jumps and stall runs]\label{fr:def:jumps}
Let $d_1<d_2<\cdots$ enumerate $\{d\ge1:\beta_d=1\}$. For each $i$:
the front $\sigma_i=M_{d_i}+1$; the stall run $k_i=d_{i+1}-d_i-1\in\N\cup\{\infty\}$
($\infty$ if $d_i$ is the last jump); the jump $J_i=M_{d_i}-M_{d_i-1}$.
The stall steps of plateau $i$ are the $d$ with $d_i\le d\le d_{i+1}-2$ (those with $\beta_{d+1}=0$).
Put $P_i=\max\{P_d:d_i\le d\le d_{i+1}-1\}$ and $\mu_i=\max\{\mu_{d+1}:d_i\le d\le d_{i+1}-2\}$
($\mu_i=0$ if $k_i=0$).
\end{definition}

\begin{lemma}[Row reading; \PROVED]\label{fr:lem:row}
For each $i$:
\begin{enumerate}
\item[(a)] $\sigma_i=\tau_{d_i}$ and $M_j=\sigma_i-1$ for $d_i\le j<d_{i+1}$;
\item[(b)] $\Delta_j(\sigma_i)=0$ for $d_i\le j<d_{i+1}$, and $\Delta_{d_{i+1}}(\sigma_i)=1$ if $k_i<\infty$;
\item[(c)] $k_i=\min\{j>d_i:\Delta_j(\sigma_i)=1\}-d_i-1$, i.e.\ in cells $x=j-\sigma_i$ of row $\sigma_i$:
$k_i+1$ is the distance from the front cell $d_i-\sigma_i$ to the first cell $x>d_i-\sigma_i$ with
$u^{\sigma_i}_x\ne r_{x+\sigma_i}(\sigma_i)$.
\end{enumerate}
\end{lemma}
\begin{proof}
(a) Theorem~\ref{fr:thm:B}(c) gives $\deg\Delta_{d_i}=M_{d_i}$; Theorem~\ref{fr:thm:B}(b) keeps
$M_j$ constant on the plateau. (b) $\deg\Delta_j\le M_j=\sigma_i-1$; and
$1=\beta_{d_{i+1}}=\Delta_{d_{i+1}}(M_{d_{i+1}-1}+1)=\Delta_{d_{i+1}}(\sigma_i)$.
(c) Combine (b) with $D_j(\sigma_i)=u^{\sigma_i}_{j-\sigma_i}$.
\end{proof}

\begin{lemma}[Lattice path; \PROVED]\label{fr:lem:path}
For each $i$: $S_d=\sigma_i-1-d$ for $d_i\le d<d_{i+1}$;
$\sigma_{i+1}=\sigma_i+J_{i+1}$ and $S_{d_{i+1}}=S_{d_i}-k_i+J_{i+1}-1$ if $k_i<\infty$, with
$J_{i+1}=\min\{t\ge\sigma_i:r_{d_{i+1}-1}(t)=1\}-\sigma_i+1\in[1,p_{d_{i+1}-1}]$;
and if $1\le k_i<\infty$, $\min\{S_d:d_i\le d\le d_{i+1}-2\}=\sigma_i-d_i-k_i$.
\end{lemma}
\begin{proof}
Lemma~\ref{fr:lem:row}(a) and Theorem~\ref{fr:thm:B}(c) at $d=d_{i+1}-1$; the minimum is at
$d=d_i+k_i-1$.
\end{proof}

\subsection{Reduction of Conjecture H to stall runs}

\begin{conjecture}[H; \OPEN; the statement was public earlier, see below]\label{fr:conj:H}
$\tau_d\ge d$ for every $d\ge20$.
\end{conjecture}
\noindent\EXACT: the conjecture holds for $20\le d\le4\times10^{6}$ ($\tau_d$ per diagonal, with a certificate; independent
codes agree to $2\times10^{6}$). The threshold $20$ is sharp: $\tau_d<d$ exactly for $d\in\{3,\dots,15,17,19\}$ (E1, and the
same set to $4\times10^{6}$ by E10). An exact simulation of the seed, run by two independent simulators, extends both statements to $d=4\times10^{7}$ (ledger 10.163c2). The statement is \KNOWN: Wolfram's 2019 prize post has it as a question \cite{Wolfram}, and it is stated in \cite{Pat} (commit \texttt{08786d1}, on its main branch on 28 August 2026) and in \cite{Dib}. None of them has a proof. Shallower checks are public: Wolfram's 2019 data file \cite{Wolfram19} reaches $d=78\,287$ and gives exactly the set $\{3,\dots,15,17,19\}$, and \cite{Dib} checks to $d=200\,000$. Theorem~\ref{fr:thm:Hstar} reduces the conjecture to two inequalities on the front; their measured margins are in Remark~\ref{fr:Cmargin}.

\begin{remark}[the column as the sum of two diagonals; \PROVED, \EXACT\ and \MEASURED\ parts as marked]\label{fr:middle}
The centre column is the diagonal of the frame, $c(t)=D_t(t)$, so with $D_t=r_t+\Delta_t$,
\[
c(t)=b_B(t)\oplus\Delta_t(t),\qquad b_B(t):=r_t(t),
\]
the background column of Remark~\ref{fr:frozencolumn} XOR the defect of diagonal $t$ at its own row (\PROVED, an identity; ledger 10.234d~A; the background/defect split is also in \cite{Dib}, see the paragraph after Theorem~\ref{fr:thm:B}). Since $\tau_t>t$ for every $20\le t\le8\times10^{5}$ (\EXACT, ledger 10.242d~A14, strengthening E7's $\tau_d\ge d$; $\tau_{19}\le19$), on that range the second summand is a live transient bit, about $0.32\,t$ rows below the front of its diagonal; its persistence is, up to one row, Conjecture~H (ledger 10.242d~A14). On rows $1\,000$--$1\,130\,000$ the bit $\Delta_t(t)$ has density $0.50015$ (coin standard deviation $0.00047$), geometric runs, and conditional frequencies within $2.1$ standard deviations of $\tfrac12$ against $b_B$, against the wall's phase modulo $16$, $32$ and $64$, against the local context of $r_t$, against the period, and, to $8\times10^{5}$, against the state of the edge at diagonal $t$ --- move or wait, depth below the front, mode (\MEASURED; ledger 10.234d~B, 10.242d~A15). Nothing the left half's order knows is seen to leak into the column's bits at this length. The same entry argues that the edge speed equals the leftward damage speed of a single flipped cell ($0.2429$ against Wolfram's $0.2428$, Remark~\ref{fr:slopebar}); the argument uses the universality of the front, which is \MEASURED\ (Theorem~\ref{fr:universaldata} and the discussion following it), so that lemma is conditional on it and is not a theorem (ledger 10.234d~C).
\end{remark}

\begin{theorem}[H$^*$; \PROVED]\label{fr:thm:Hstar}
Let $d_{i_0}\ge20$ be a jump diagonal and assume \textup{(Base)} $\tau_d\ge d$ for $20\le d\le d_{i_0}$.
\begin{enumerate}
\item[(U)] If $k_i\le\sigma_i-d_i-P_i-\mu_i$ for every $i\ge i_0$ with $k_i\ge1$, then $\tau_d\ge d$ for all $d\ge20$.
\item[(C)] If $\kappa_{d+1}\le P_d-1$ for all $d\ge d_{i_0}$, and $k_i\le\sigma_i-d_i-P_i+1$ for every
$i\ge i_0$ with $k_i\ge1$, then $\tau_d\ge d$ for all $d\ge20$.
\end{enumerate}
\end{theorem}
\begin{proof}
Induction on $d\ge d_{i_0}$; assume $\tau_d\ge d$. If $\beta_{d+1}=1$, Theorem~\ref{fr:thm:C}(b) gives
$\tau_{d+1}\ge d+1$. If $\beta_{d+1}=0$, then $d+1$ lies strictly inside a plateau $i\ge i_0$, so $k_i\ge1$ and
$d$ is a stall step; the hypothesis excludes $k_i=\infty$. By Lemma~\ref{fr:lem:path},
$S_d\ge\sigma_i-d_i-k_i$. Under (U) this is $\ge P_i+\mu_i\ge P_d+\mu_{d+1}\ge\kappa_{d+1}$
(Lemma~\ref{fr:lem:G}); under (C) it is $\ge P_i-1\ge P_d-1\ge\kappa_{d+1}$.
Theorem~\ref{fr:thm:C}(c) gives $\tau_{d+1}\ge d+1$.
\end{proof}

\begin{remark}[the two halves of (C), and their measured margins; \EXACT\ and \MEASURED\ parts as marked, used in no proof]\label{fr:Cmargin}
Condition (C) has two halves. The first, $\kappa_{d+1}\le P_d-1$, is \EXACT\ on $[21,8\times10^{5}]$ (E2, E7; ledger 10.156c2). The second asks $k_i\le\sigma_i-d_i-P_i+1$, which on the seed's data allows stall runs of about $0.32\,d_i-P_i$, far above the $P_i-1$ one would read off if the two halves were conflated (ledger 10.242d~A13). The coin-driven edge machine of Remark~\ref{fr:edgeautomaton} gives the tail of the stall-run distribution to $10^{-10}$ per cycle over $4\times10^{10}$ diagonals: geometric with ratio $0.433$ at both periods, the longest stalls seen (each run's start-up stall of $43$--$62$ excluded) being $24$ at period $32$ and $26$ at period $64$, against the empirical $k_i\le P_i-1$ of E3, that is $31$ and $63$, which is \EXACT\ to $8\times10^{5}$ and not a proven bound; extrapolated, a stall of $32$ once per $2.7\times10^{12}$ diagonals and one of $64$ once per $10^{24}$ (\MEASURED; ledger 10.233d~H). That is a margin and not a proof; the machine can say nothing about a stall that the universal front forbids or allows beyond what its coin sees, and nothing in this remark enters Theorem~\ref{fr:thm:Hstar}.
\end{remark}

\begin{corollary}[Light cone and necessity; \PROVED]\label{fr:cor:lightcone}
Let $1\le k_i<\infty$.
\begin{enumerate}
\item[(a)] $k_i\le2\sigma_i-d_i+2P_i-1$.
\item[(b)] If $\tau_j\ge j$ for all $j\in(d_i,d_{i+1})$, then $k_i\le\sigma_i-d_i$.
\end{enumerate}
\end{corollary}
\begin{proof}
(a) Let $j^*=d_i+k_i$. By Lemma~\ref{fr:lem:bgdef}(d) some $j\in\{j^*-1,j^*\}\subseteq[d_i,d_{i+1})$ has
$r_j\not\equiv0$, so $r_j(t)=1$ for some $t\in[\sigma_i,\sigma_i+p_j-1]$. As
$\deg\Delta_j\le\sigma_i-1$ (Lemma~\ref{fr:lem:row}), $D_j(t)=1$, hence $j\le2t$ by
Lemma~\ref{fr:lem:diag}(b). Thus $d_i+k_i-1\le2\sigma_i+2P_i-2$.
(b) $\tau_{j^*}=\deg\Delta_{j^*}+1\le\sigma_i$ and $\tau_{j^*}\ge j^*=d_i+k_i$.
\end{proof}

\begin{remark}[the shape of this argument is a recognised technique, and the literature prices it]\label{fr:template}
Two published proofs have exactly the shape of Theorem~\ref{fr:thm:Hstar}, and it is worth saying which,
because it converts our saturation computations from an ad hoc device into an instance of a method.

The first is Baetens and Gravner \cite[Thm.~4.2]{BGrav}. They model defect accumulation along a
trajectory as a branching walk and prove that if there is a finite set $M$ with $|M|\ge2$ and a time
$t_M$ such that a single defect at the origin has spread over all of $M$ by time $t_M$ \emph{for
arbitrary} initial configuration, then the defect dynamics is expansive with an explicit lower bound on
the Lyapunov profile. Rule~30 is one of the four rules their theorem handles, but only on its rightward, permutive side:
they apply it with $M=\{1,3\}$ and $t_M=3$, so the bound lives on velocities $\alpha\in[\tfrac13,1]$, and the left edge
$-0.41$ of their defect shape for Rule~30 is an empirical estimate. We have found nothing in print that bounds Rule~30's
leftward damage, defect or front speed from below, and the left side is where our front lies. Two features matter for
us: it is one of the very few theorems in this literature that applies to Rule~30 with \emph{no}
hypothesis on an initial measure, and its proof is a finite verification over all backgrounds in a
window followed by an induction --- which is precisely the form of the hypotheses (U) and (C) above and
of the saturation computations that support them.

The second is Gravner and Griffeath's treatment of Rule~22 \cite{GG22}, the symmetric sibling of our
rule. Their Lemma~3.3 defines an irregular block, gives an explicit succession rule for how the
irregular blocks evolve, and concludes that their number is nonincreasing in time; Theorem~1 then
follows. That is the template a proof of Conjecture~H would follow: a finite list of defect types, a
closed succession rule, a monotone count, all deterministic and all along one orbit. It does not
transfer as a theorem, and the two reasons are worth recording because they are exactly the difficulty.
Their argument needs the nonlinear rule to \emph{be} Rule~90 on the background, which
Proposition~\ref{du:nosubdyn} excludes for Rule~30; and their defects shrink by one cell from each end
per step whereas ours grow, by Theorem~\ref{fr:thm:B}, so the monotone quantity runs the wrong way and
the list of defect types is not finite.

Finally, the general theory says why no borrowed theorem will close the gap. Pivato's kinematic
classification \cite{PivatoDefect} sorts a defect between two subshifts $L$ and $R$ by their entropies.
When both vanish the defect is a finite automaton; when both are positive it is Turing equivalent; our
case is $h(L)=0<h(R)$, the asymmetric pushdown regime, in which he shows the analogous questions ---
including whether the defect eventually stops moving --- are equivalent to the halting problem. His
framework also assumes throughout that defects stay bounded in width, and he singles out the unbounded
case as probably undecidable in general, citing \cite{Sutner}. Our $\Delta_d$ is exactly that case, and
Theorem~\ref{fr:thm:A} is an exact formula for its growth, which is more than the framework has; but the
price is that any proof of Conjecture~H must be specific to Rule~30, and this is the citation for that.
\end{remark}

\begin{remark}[what precision the settling constant actually has; \MEASURED]\label{fr:slopebar}
Every number in this paper that is keyed to $\tau_d$ inherits an uncertainty we had been quoting too
optimistically, and the correction matters for the remark that follows, so we state it here.

Two facts about the sequence $\tau_d$. First, least-squares error bars on it are meaningless. Fitting
$\tau_d=\rho d+c$ on disjoint blocks of the exact data to $d=90\,000$ gives block-to-block standard
deviations of $0.0158$ at block width $10^{4}$ and $0.0077$ at width $3\times10^{4}$, against
within-block standard errors of about $10^{-4}$: the residuals are strongly autocorrelated, because
$\tau_d$ moves in plateaus, and a regression bar understates the spread by roughly two orders of
magnitude. Over $d\in[8000,90\,000]$ the smoothed ratio $\tau_d/d$ is flat at $1.335$ to within
$\pm0.001$, with no drift visible inside that range at all.

Second, the drift is real but lives further out, beyond $90\,000$, and on the exact data to $4\times10^{6}$ none is
visible past about $10^{6}$. The two published constants quoted below are both from \cite{NKS}, but the difference-pattern
constant $0.2428$ has an earlier, two-digit version: Wolfram's CRYPTO~'85 extended abstract \cite{W85} gives the left Lyapunov exponent of this rule as $0.25$, with $1$ on the right, and Meier and
Staffelbach \cite[p.~191]{MS} report that a bit change propagates to the left with speed roughly $\tfrac14$, citing
Wolfram's 1986 paper in \emph{Adv.\ Appl.\ Math.}, which we have not seen. Neither gives the range measured. We read Wolfram's
1983 survey \cite{W83} page by page and it contains no measurement of this boundary speed, nor any
quantity convertible to one --- its only difference-pattern data, Fig.~20, is for Rules~90 and $126$
over some $220$ steps, and the quantity plotted there is a width times an interior density, which for
the symmetric Rule~126 simply recovers the light cone at density $\tfrac12$. So no third range is
available from the literature against which to test the drift. Two public code repositories also give values: \cite{Dib} has a front speed of $0.243$ over $10^{6}$ rows, and \cite{Pat} gives about $1.34\,d$ for the settling row of diagonal $d$. Our own exact values of $M_d/d$, now per diagonal to
$d=4\times10^{6}$ (with a certificate; independent codes agree to $2\times10^{6}$), are $1.336660$ at $d=5\times10^{4}$, then
$1.334278,\,1.332570,\,1.329870,\,1.326183,\,1.324145,\,1.321118,\,1.320079$ at
$d=9\times10^{4},10^{5},2\times10^{5},3\times10^{5},4\times10^{5},6\times10^{5},8\times10^{5}$, and $1.320763$ at
$d=4\times10^{6}$, above the value at $8\times10^{5}$: the ratio is not monotone. The bin medians of
$M_d/d$ from the word-granular frame of E9 fall from $1.3370$ on the bin at $d\approx9.5\times10^{4}$ to $1.3202$ on
$[7\times10^{5},10^{6}]$ and then sit at $1.3209$--$1.3212$ on the six bins covering $[10^{6},4\times10^{6}]$ (spread
$0.0003$), and the per-diagonal medians on those six bins reproduce them to $\pm0.0001$; the linear fit over
$d\ge3\times10^{6}$ is $1.32039\,d+1830$, and the local slope shows no trend from $10^{5}$ to $4\times10^{6}$ (mean $1.32046$
over blocks of $10^{5}$ diagonals, block standard deviation $0.005$, drift $0.00005\pm0.00073$ per $10^{6}$). So the bin
medians are locked to four digits on $[10^{6},4\times10^{6}]$; whether the ratio converges is \OPEN, and these data give no
bound on the limit: the fitted slope $1.32039$ ($v=0.24265$) is a local rate of advance and not a bound. Two earlier
versions of this remark gave upper bounds, $1.3204$ and then, under continued monotone decrease, $1.3209$ (front speed at
most $0.2429$); both are \textsc{withdrawn}, the second because the ratio is not monotone.
Extrapolation does not do better, and that can be shown rather than
assumed. Fitting bin medians on $d\le2\times10^{5}$ to $\rho_\infty+Cd^{-\alpha}$, to
$\rho_\infty+C/d$, to $\rho_\infty+C/\log d$, and by Aitken acceleration on the dyadic subsequence,
returns $1.25$--$1.34$, $1.332$--$1.338$, $1.29$--$1.41$ and $1.340$ respectively; the two tight estimates
($1.332$--$1.338$ and $1.340$) overshoot the bin median $1.3209$ measured at $d=4\times10^{6}$, and the two loose ones merely
contain it. That measurement is a held-out test point which no method predicts, so the approach to the limit is of none of
those forms over the range we hold. This is unsurprising once the input is examined: the bin medians are not
monotone there, rising to $1.3395$ at $d=32\,768$ before falling, and extrapolating a sequence that has
just turned over is exactly the case in which extrapolation misleads. The two published constants are $0.2428$ \cite[p.~949]{NKS} and $0.252$ \cite[p.~871]{NKS}. The locked bin value
$1-1/1.3209=0.2429$ agrees with the first to three digits and is inconsistent with the second at $3.7\%$; an
earlier version of this remark spoke of three numbers agreeing to three digits, which is \textsc{withdrawn}, and the
agreement with $0.2428$ is a coincidence of ranges until a limit is proved to exist. What the data establish is only that
the local slope gives a speed near $0.2427$ on $[10^{5},4\times10^{6}]$; that the limiting speed is below $\tfrac14$ rests on the
conditional bracket of Remark~\ref{fr:dyadic}, and nothing finer is known. Two numerical
coincidences in the older literature should not be mistaken for our quantities: the decay constant
$\lambda\approx\tfrac43$ of \cite{W83} governs a triangle density and is unrelated to the settling
constant near $\tfrac43$ that we refuted as an asymptote, and the ``effective density'' $0.25$ there is
not the mean-field $0.25$ of \cite[p.~949]{NKS}.
\end{remark}

\begin{remark}[the front is a min-plus first-passage law, and not a finite one]\label{fr:tropical}
The front has an exact tropical form, which names the difficulty precisely and then supplies a
negative result.

Theorem~\ref{fr:thm:B} can be written as a first-passage law in the min-plus sense: the front advances
from $M_d$ to the next position at which the background $r_d$ permits it,
\[
M_{d+1}=M_d+\beta_{d+1}\cdot\mathrm{gap}_{r_d}\bigl(M_d \bmod p_d\bigr),
\]
an identity we verified on all $90\,001$ diagonals, with every one of the $49\,121$ jumps lying in
$[1,p_d]$ and mean increment $1.334289$ --- the slope \emph{on that range}, which is the figure to
compare against $M_{90\,000}/90\,000=1.334278$ and not against $1.324145$ at
$d=4\times10^{5}$, since the ratio falls further beyond that range (Remark~\ref{fr:slopebar}). The min-plus part --- the gap
function of a periodic background --- is exactly the kind of object max-plus spectral theory handles.
The gate $\beta_{d+1}$ is not: by Theorem~\ref{fr:thm:A}(a) it is a parity of $h_{d+1}$ on an interval
and is not a function of the background at all. That single factor is what makes the front a cocycle
rather than a matrix power, and it is the whole difficulty of this Part in one line. In the same
vocabulary, Lemma~\ref{fr:lem:path} with Theorem~\ref{fr:thm:Hstar} is a Lindley reflected walk against
the barrier $P_d+\mu_{d+1}$.

Now the negative result, and it is the strongest unconditional statement we have about the front. For an
irreducible max-plus matrix, the Cyclicity Theorem \cite[Thms.~3.109, 3.112]{BCOQ} gives
$A^{t+\gamma}=\lambda^{\gamma}\otimes A^{t}$ for all large $t$, so every entry is eventually affine plus
periodic and its deviation from the line $\lambda t+c$ is \emph{bounded}, at any dimension. Measure that
deviation for our front: the Chebyshev band half-width
$W(D)=\min_a\tfrac12\bigl(\max_{d\le D}(M_d-ad)-\min_{d\le D}(M_d-ad)\bigr)$ takes the values
$20.6,\,29.8,\,56.7,\,109.0,\,175.4,\,216.7$ at $D=10^{3},3\times10^{3},10^{4},3\times10^{4},
6\times10^{4},9\times10^{4}$, a fitted growth of $D^{0.54}$ over these six points ($0.63$ over the last three), while the three natural bounded controls
are flat across the same two orders of magnitude: a Beatty sequence of irrational slope gives $0.50$
throughout, a rational slope $0.33$, and an actual irreducible max-plus matrix power $3.00$. A random
walk with increments in $\{1,2\}$ gives $6.5,29.4,98.2$, so our growth is of that order. So the front is
not the max-plus power of any fixed finite irreducible matrix; the growth of $W$ is \MEASURED, and the
implication from unbounded $W$ is \PROVED.

Two consequences for the dichotomy stated in the next remark. First, it is incomplete as we stated it:
a Beatty sequence with irrational slope has discrepancy below $1$, so irrationality by itself does not
produce a growing band, and there is a third branch --- irrational slope with \emph{unbounded}
discrepancy --- which is the branch the growing band points to (\MEASURED). Second, our own appeal to the growth of the
transfer matrix was the wrong mechanism. Blondel, Gaubert and Tsitsiklis \cite{BGT} show that at
\emph{fixed} finite dimension the max-plus Lyapunov exponent is already NP-hard to approximate and that
deciding whether it vanishes is undecidable, so irrationality does not require an unbounded system: the
dividing line is autonomy, not dimension. Our $2^{k}$ growth is a strengthening of the obstruction, not
its source. Consistently with that, the jump increments do not grow with the period, their mean being
flat at $2.44$ across $p=4,\dots,32$.
\end{remark}

\begin{remark}[the front speed: a dichotomy, and what each branch costs]\label{fr:rational}
Proved front velocities for cellular automata are rational with one class of exception, and the
exception sits on our side of the question. Pivato's Theorem~2.1 \cite{PivatoDefect} explains the rule:
a defect of bounded width between $\sigma$-periodic backgrounds has a finite internal state driven by
periodic input, so its average velocity is rational by construction. The concrete method agrees ---
Gravner and Griffeath's $(a,\tau)$-blockers \cite[Props.~5.1, 5.2]{GG} return
$2,1,\tfrac43,\tfrac34,\tfrac12$ for the range-two edge automata, each by a finite verification, and
their remark that the blocker bound is independent of the rest of the environment makes it the one
published method built for an ordered solution competing against a chaotic one.

The exception is the multiplication automaton: multiply by $p$ in base $pq$, a radius-one rule on a
full shift, whose digit string from a quiescent background grows by $\log_{pq}p$ per step, irrational.
This is Kopra's Theorem~5.5 \cite{Kopra2020}, and Hochman realises every non-negative real as a
Lyapunov exponent on a subshift. Note which side it falls on: that automaton's exponent as a supremum
over configurations is rational, while the exponent \emph{along the single orbit} is irrational, and our
front is the single-orbit quantity. The right general statement is a published question rather than a
theorem: Cyr, Franks and Kra ask whether any automaton on a full shift has an irrational Lyapunov
exponent \cite[Question~5.7]{CFK}.

Blockers do not apply to Rule~30 as it stands, since Theorem~\ref{fr:unbounded} denies it a periodic
solution at the light-cone edge. That obstruction has a published way around it: Proposition~3.3
forbids only \emph{one-sided} solutions, and Gravner and Liu define a weak robust periodic solution ---
doubly periodic, preserved on the left half-line, advancing into any environment --- with a decidable
criterion and the bound $v\ge a/T_a$ \cite{GL23}. That is the object we had described informally as the
ordered texture invading the chaotic bulk, and it separates the rules: Rules~90 and $150$ have none,
their fronts retreating because they are bipermutive, while Rule~30's checkerboard advances against
every environment in a sample of $912$.

So the dichotomy, and its price. A universal $(a,\tau)$-blocker certifies $a/\tau$ and costs $2^{a}$,
both bounds quantifying over all environments of length $a$, so the numerator of the true speed is the
price of the certificate. At the interval $0.242653\pm0.00017$ used when this was computed (an error bar that
Remark~\ref{fr:slopebar} does not support; no bound on the limit is now claimed), no rational of numerator below $25$ lies in it, the
first admissible fraction being $25/103$; the published velocities were certified at $a\le4$. And the
price rises as the measurement sharpens --- $33/136$ at a third of the error bar, $223/919$ at a
hundredth. A small numerator is the only case the published method reaches, and it is the case the
measurement excluded at that error bar. Given Remark~\ref{fr:slopebar} the number $25$ is conditional on an error bar we
do not trust; what is unconditional is that any interval tight enough to identify the speed excludes
every numerator a blocker search could feasibly reach.

\end{remark}

\begin{remark}[the existence of the front limit: one route dead, one live; \MEASURED]\label{fr:exists}
The existence of the limit, as opposed to its value, is cheaper, and one route is dead while another is
live. Dead: the increment $a(d,e)=M_{d+e}-M_d$ is a telescoping difference, hence an exact cocycle,
verified on $20\,000$ random triples with no violation, so every two-parameter subadditivity hypothesis
--- Kingman's or Liggett's --- holds with equality and yields nothing, and Fekete applied to it is a
tautology. The only non-vacuous superstructure quantifies over restart states,
$G^{+}(e)=\sup_d(M_{d+e}-M_d)$, which is subadditive; but $G^{+}(1)=\sup_d J_d$ and the maximal jump
saturates the diagonal period, the ratio $\max J/p$ being $1.00,1.00,1.00,0.94$ at $p=1,4,8,16$, so with
periods unbounded that constant is $+\infty$. Note that this is the same sup-over-environments
structure a blocker quantifies, so the blocker route dies here a second time, independently of its
cost. Kingman cannot be repaired for a deterministic array either: its stationarity hypothesis becomes
equality of increments, forcing $J_d$ constant.

Live, and sharper than we first posed it. The de Bruijn--Erd\H{o}s theorem \cite{dBE}, in F\"uredi and
Ruzsa's form \cite{FR}, needs only near-diagonal almost-subadditivity: if
$M_{d+e}\le M_d+M_e+f(d+e)$ for all $d\le e\le 2d$ with $\sum_nf(n)/n^{2}<\infty$, then $\lim M_d/d$
exists. It quantifies over no restart states and needs no randomness, monotonicity or stationarity,
which is why it survives where the routes above fail. On the exact orbit the inequality holds with $f$
constant: the maximum of $M_{d+e}-M_d-M_e$ over the $13\,333\,533\,333$ band pairs with
$d+e\le400\,000$ is $326$, attained at $(d,e)=(10\,391,16\,064)$, and it is the same at each of
$32\,000$, $90\,000$, $120\,000$, $150\,000$, $180\,000$, $200\,000$ and $400\,000$ with the same
maximising pair throughout; $\tau$ is exact to $d=400\,000$, all $400\,001$ diagonals settled, agreeing
with the independent $d\le90\,000$ computation at all $89\,999$ diagonals settled in both.

A running maximum that never moves is also what one would see if the extremal region were simply never
revisited, so what is reported below is the maximum restricted to each window of $d+e$, which is the only
form in which the figure carries information. The extremum is localised in
$d+e\in[25\,000,50\,000]$ and is never approached again: over the ten billion band pairs with
$d+e\ge200\,000$ --- three times the entire range available before this computation --- the largest
deficit is $94$. The decline is also not mysterious. $M_d/d$ falls across the computed range,
through $1.336660,\,1.332570,\,1.329870,\,1.326183,\,1.324145$ at
$d=5\times10^{4},10^{5},2\times10^{5},3\times10^{5},4\times10^{5}$, and a falling ratio favours
$M_{d+e}<M_d+M_e$ at large scales, though it forces nothing pointwise, the ratio not being monotone (below, and
Remark~\ref{fr:slopebar}); what the constant measures is the small-scale
roughness of $\delta$, and the window table is the evidence that the roughness does not grow.

Examining where that maximum lives improves the target. First, the deficit is exactly the deviation
combination: writing $\delta_d=M_d-\rho d$, we have $M_{d+e}-M_d-M_e=\delta_{d+e}-\delta_d-\delta_e$,
and at the extremal pair this reads $384.3-(-48.7)-107.0=326$. So the inequality is a statement about
the \emph{roughness} of $\delta$ and not its size --- a smooth concave $\delta$ gives a negative
deficit, and only local excursions give a positive one. That the front is smoother than generic noise
of the same amplitude can be checked: a synthetic sequence with our slope and a random walk of matched
amplitude gives band maximum $629$, and one with our measured exponent $1126$, in both cases with the
maximising pair drifting to the far end as the range grows, where ours is fixed.

Second, and this is the useful part, the positive deficits are confined. Taken by window of $d+e$ the
maximum runs
\[
\begin{array}{lrrrrrrrrrr}
d+e/10^{3} & 20\text{--}30 & 30\text{--}40 & 40\text{--}55 & 55\text{--}70 & 70\text{--}90 & 90\text{--}110 & 110\text{--}130 & 130\text{--}150 & 150\text{--}175 & 175\text{--}200\\
\max & +326 & +285 & +26 & +227 & +233 & -54 & -40 & -323 & -503 & -429
\end{array}
\]
A negative maximum means every band pair in that window satisfies $M_{d+e}<M_d+M_e$ outright. On the
range available when that table was computed, $d+e\le200\,000$, every window above $90\,000$ was
negative, and we concluded that the front is plainly subadditive on the band beyond that point with a
margin that grows, so that the target could be stated without an error term as $M_{d+e}\le M_d+M_e$ for
$d+e\ge N_0$. Doubling the range refutes both halves of that, and the refutation is recorded here
because the fallback it triggers was registered in advance. Continuing the same windows, now $20\,000$
wide, to $d+e=390\,000$:
\[
\begin{array}{lrrrrrrrr}
d+e/10^{3} & 190\text{--}210 & 210\text{--}230 & \mathbf{230\text{--}250} & 250\text{--}270 & 270\text{--}290 & 290\text{--}310 & 310\text{--}330 & 370\text{--}390\\
\max & -167 & -93 & \mathbf{+94} & -40 & -571 & -770 & -872 & -1374
\end{array}
\]
The positive deficits \emph{do} recur: exactly one window above $90\,000$ is positive, at
$d+e\in[230\,000,250\,000]$ with maximum $+94$ attained at $(101\,905,146\,203)$. So plain subadditivity
does not hold on the tail $d+e\ge90\,000$, and the margin does not grow monotonically; quoting the
smallest $N_0$ that works on the range computed so far would repeat the error the doubling just exposed.
What survives is the almost-subadditive form with $f$ constant, which is what \cite{FR} needs and which
the recurrence at $+94$ does not threaten, since $94$ is far below the global $326$. That is now the
target, and the clean form is withdrawn. Two cautions, and the second can now be quantified. The crossover
sits at $87\,866$, the last doubling position in the seed's lineage (E8), so the positive region is the
one preceding it, and we do not know whether that is causal. The doublings do not stop there, so
positive deficits may return at the next one --- but the next one is much further away than ``beyond our
range'' suggests. Reading the doubling positions off the exact periods, the first diagonal carrying
period $p$ occurs at
\[
\begin{array}{lrrrrr}
p & 2 & 4 & 8 & 16 & 32\\
\text{first } d & 3 & 8 & 29 & 400 & 87\,867\\
\log_2(\text{first }d) & 1.585 & 3.000 & 4.858 & 8.644 & 16.423
\end{array}
\]
and the successive differences of that last row are $1.415,\,1.858,\,3.786,\,7.779$, whose ratios are
$1.31,\,2.04,\,2.06$. The three usable ratios are consistent with a \emph{doubly exponential} onset in $\log_2 p$ after a short
transient; that is a reading of three ratios and not a law. An earlier version of this remark extrapolated the two stable ratios to a first period-$64$ diagonal near
$2^{32.4}\approx5\times10^{9}$, and said that no computation we could contemplate would reach the next doubling. Both statements are \textsc{withdrawn}.
The label recursion reaches that doubling (Remark~\ref{fr:conjM:refuted}; public earlier in \cite{Dib,Pat}): the first period-$64$ diagonal is $2{,}107{,}985{,}255$ on one branch of the fork at $1{,}420{,}878{,}968$ and not before $9{,}958{,}700{,}505$ on the other, and which branch the seed takes is \OPEN. What stands is that
every exact computation of the front in this paper ends far below that fork, inside the
period-$32$ regime that begins at $87\,867$. The positive deficit we do find above $90\,000$, at $d+e\in[230\,000,250\,000]$, is
accordingly \emph{not} at a doubling, so whatever produces it is not the doubling structure. In the same range,
$399\,251$ of the $400\,001$ diagonals carry period $16$ or $32$, and the mean of $M_d/d$ continues to
fall \emph{within} the period-$32$ class alone --- $1.3321,\,1.3295,\,1.3268,\,1.3250$ over
$d\in[88\,000,150\,000),[150\,000,250\,000),[250\,000,320\,000),[320\,000,400\,000]$ --- but this does not
separate the drift from the period: $M_d/d$ is a running average over all diagonals and carries the range below $87\,867$
with it, while the local slope shows no trend from $10^{5}$ to $4\times10^{6}$ (Remark~\ref{fr:slopebar}). Proving either form of the inequality is the open problem we
would most like solved.

Exact monotonicity of the ratio would suffice and is false. If $\rho_n=M_n/n$ were non-increasing then
$\rho_{d+e}\le\min(\rho_d,\rho_e)$ and hence $M_{d+e}=(d+e)\rho_{d+e}\le d\rho_d+e\rho_e=M_d+M_e$
immediately. But $\rho_{n+1}-\rho_n=(J_n-\rho_n)/(n+1)$ exactly, where $J_n=M_{n+1}-M_n$, so $\rho$
rises whenever a jump exceeds the running ratio, that is whenever $J_n\ge2$ --- which happens at
$39.9\%$ of steps, and at every scale we measured, from $41.7\%$ in $[2^{10},2^{11})$ to $39.8\%$ in
$[2^{17},2\times10^{5}]$. What does decay is the size of each rise, bounded by $(p_n-\rho_n)/(n+1)$
and measured falling from $1.27\times10^{-2}$ to $1.01\times10^{-4}$ across those scales, in proportion
to $1/n$ as that bound predicts.

One further route is priced rather than closed. Baetens and Gravner's transfer-matrix theory
\cite[Thms.~6.1, 6.2]{BGrav} gives the asymptotic shape and Lyapunov profile of a defect on a doubly
periodic background as a Legendre transform, with one state per phase of the background --- so the
matrix has size proportional to the period, which is $2^{k}$ in our regime index. That is the same wall
our own transfer computation hit at $K=18$: the literature's best deterministic front theorem and our
computational ceiling are one ceiling.
\end{remark}

\begin{remark}[the target reduced to one index; \PROVED\ reduction, \MEASURED\ hypothesis]\label{fr:concave}
The band inequality quantifies over pairs. It can be replaced by a statement about a single index, and
the exchange also shows that the constant form of $f$ we have been pursuing is stronger than needed.

Let $G$ be the least concave majorant of $M$ on $[0,N]$ and $h=M-G\le0$. For concave $G$ with $G(0)=c$,
concavity at $x$ and at $y$ against the endpoints $0$ and $x+y$ gives $G(x)+G(y)\ge c+G(x+y)$; here
$c=M_0=-1$. Hence
\[
M_{d+e}-M_d-M_e=\bigl[G_{d+e}-G_d-G_e\bigr]+\bigl[h_{d+e}-h_d-h_e\bigr]
\ \le\ 1+|h_d|+|h_e|\ \le\ 1+2\max_{m\le d+e}|h_m| ,
\]
using $h\le0$ throughout. So $f(n)=1+2\max_{m\le n}|h_m|$ is admissible, and since $\sum n^{\alpha-2}$
converges for every $\alpha<1$:
\begin{quote}
\emph{If $M$ stays within $O(N^{\alpha})$ of its least concave majorant for some $\alpha<1$, then
$\lim M_d/d$ exists.}
\end{quote}
The reduction is \PROVED. The hypothesis is \MEASURED: recomputing the majorant on $[0,N]$ for
$N=12\,500,\,25\,000,\,50\,000,\,10^{5},\,2\times10^{5},\,4\times10^{5}$ gives
$\max|h|=115.9,\,190.7,\,241.7,\,367.5,\,557.6,\,762.8$. So $\max|h|$ grows across the computed range (\MEASURED), and a
constant $f$ does not come from this route --- but the ratio per doubling is
$1.646,\,1.267,\,1.520,\,1.517,\,1.368$, every one below the value $2$ that $\alpha=1$ would require.
That is the point of this formulation: the hypothesis is a \emph{threshold} question rather than an
exponent estimate, so it does not rest on fitting $\alpha$ --- an exercise withdrawn elsewhere in this
paper (Remark~\ref{du:ring}). The largest observed ratio is $1.65$ against a threshold of $2$, and the
growth is consistent with the $D^{0.54}$ Chebyshev band above, which is the same quantity measured
against a line instead of against the majorant.

Two limits, stated so the reduction is not oversold. The bound is lossy in the large: its maximum over
band pairs is $1525.7$, at $(176\,617,176\,617)$, against a measured band maximum of $326$. It is
nevertheless tight where it matters --- the twelve largest deficits all lie in one neighbourhood near
$(10\,390,16\,064)$ and sit $38$ to $41$ below their bound. And the reduction proves nothing on its own:
it converts a statement quantified over $13$ billion pairs into a statement about one sequence, which is
progress in form and not in content. The exact-monotonicity route above is the $\alpha=0$ case of the
same picture, and it is false; this asks only for $\alpha<1$.

Three caveats belong with this, and the first prevents the reduction from being read as more than it
is. \emph{The $o(d)$ form of the hypothesis is equivalent to the conclusion, not merely sufficient for
it.} If $M_d/d\to L$, fix $\varepsilon>0$, choose $D$ with $M_d\le(L+\varepsilon)d$ for $d\ge D$ and put
$C=\max_{d\le D}M_d$; then $A(t)=(L+\varepsilon)t+C$ is a concave majorant of $M$, so $G\le A$ and
$\limsup G_d/d\le L+\varepsilon$ for every $\varepsilon$, while $G\ge M$ gives the reverse. Hence
$\lim G_d/d=L$ and $h_d/d\to0$. So a hypothesis of the form $|h_d|=o(d)$ is a restatement and cannot be
a target; what is stated above, $|h_d|=O(N^{\alpha})$ with $\alpha<1$, is strictly stronger than the
conclusion and is a genuine one. Second, the measured quantity is weaker than it looks: $\max|h|$ on
$[0,N]$ is \emph{non-decreasing in $N$ by construction}, since for $N<N'$ the restriction of $G^{(N')}$
to $[0,N]$ is a concave majorant there and so dominates $G^{(N)}$; and the object in the theorem is the
deficit below the full-horizon majorant $G^{(\infty)}\ge G^{(N)}$, so every finite-$N$ figure is a
\emph{lower} bound on a quantity that must be bounded above. Third, $\sup_dM_d/d<\infty$ is
\MEASURED\ and not proved --- the only proved recursion is
$\tau_{d+1}\le\max(\tau_{d-1},\tau_d)+p_d$ with $p_d$ unbounded --- and without it $G\equiv+\infty$ and
the majorant has no object. The measured value is $\sup_{d\le8\times10^{5}}M_d/d=1.4588$, at $d=85$.
\end{remark}

\begin{remark}[a criterion strictly between the two; \PROVED\ reduction, \MEASURED\ hypothesis]\label{fr:dyadic}
The gap left by the previous remark --- weaker than $O(N^{\alpha})$, stronger than the tautology --- is
filled by a dyadic criterion on the slope itself. Put $s_d=M_d/d$ and
\[
W_k=\max\{\,s_b-s_a\;:\;2^{k}\le a<b\le2^{k+1}\,\},
\]
the largest \emph{upward} move of the slope inside one dyadic window. Then
$\sum_kW_k<\infty$ implies $\lim M_d/d$ exists, with $\lim\le s_a+\sum_{i\ge k}W_i$ for every
$a\ge2^{k}$: given $2^{j}\le a<b\le2^{k+1}$, insert the intervening powers of two, note each consecutive
pair lies inside a single window, and sum, so $\limsup s-\liminf s\le\sum_{i\ge j}W_i\to0$. It is weaker
than the $O(N^{\alpha})$ hypothesis, since $|h|=O(d^{\alpha})$ gives $W(D)=O(D^{\alpha-1})$, summable
exactly when $\alpha<1$; and stronger than the conclusion, since convergence does not imply
summability.

Measured on the exact front to $d=8\times10^{5}$:
\[
\begin{array}{lccccccccc}
k & 10 & 11 & 12 & 13 & 14 & 15 & 16 & 17 & 18\\
W_k & 5.13\text{e--}2 & 2.91\text{e--}2 & 2.31\text{e--}2 & 1.80\text{e--}2 & 8.47\text{e--}3 &
3.77\text{e--}3 & 5.28\text{e--}3 & 3.18\text{e--}3 & 8.25\text{e--}4
\end{array}
\]
with partial sums reaching $0.1431$ and a fit $W_k\sim2^{-0.666k}$; extrapolating the tail beyond
$k=18$ gives $0.0023$ and the conditional bracket $\lim M_d/d\le1.3201+0.0023=1.3224$. Two features
recommend this over $\max|h|$. The windows are \emph{disjoint}, so the nine numbers are independent
measurements rather than a nested sequence, and correspondingly $W_k$ is \emph{not} monotone --- $k=16$
exceeds $k=15$. And it is the threshold, summability, that is being claimed: three estimators of the
same underlying deviation now read $0.54$ from $\max|h|$, $0.666$ from $W_k$ and $0.54$--$0.63$ from the
Chebyshev band, by window, and that spread is exactly why the exponent itself is not quoted.
\end{remark}

\begin{remark}[two routes closed, with their citations; \KNOWN\ theorems, \MEASURED\ inputs]\label{fr:closed2}
Unique ergodicity is unavailable, and not for a technical reason. Exactness of the cocycle makes the
question Birkhoff's rather than Kingman's, and Oxtoby's theorem gives uniform convergence at
\emph{every} point of a uniquely ergodic system --- but three independent obstructions apply. The
observable is not continuous, since $\tau_d$ is a preperiod and no finite window determines it; it is
not bounded, the largest jump to $d=8\times10^{5}$ being $19$ at $d=746\,573$, and
$\max_{d\le D}J_d\sim\log_2D$ because the jump tail halves; and the system is not uniquely ergodic in
any case, because the defect-free pair is itself an orbit in the closure carrying an invariant measure
on which the front observable vanishes identically. The structure that makes the front well defined ---
a periodic background --- is what supplies the second measure.

Second, the max-plus exclusion can be strengthened at no cost. Kohlberg \cite{Kohlberg} shows a
nonexpansive piecewise-linear self-map of $\mathbb R^m$ has a unique invariant half-line, so $f^n(x)-n\alpha$
is \emph{bounded}; any Lipschitz readout of such a model therefore has a bounded band. Our band grows through the computed range (\MEASURED); if it is unbounded, this upgrades the exclusion from ``no autonomous finite max-plus model with
ultimate period at most $44\,000$'' to \emph{no finite-dimensional sup-norm-nonexpansive
piecewise-affine model of the front at any dimension}, conditionally on that unboundedness.

One identity is free and worth recording. Since $M$ is non-decreasing, $J_d\ge0$, and
$d(s_d-s_{d-1})\ge-\sup_dM_d/d$, which is Landau's slowly-decreasing condition; so by the
Hardy--Landau Tauberian theorem for $(C,1)$, $\lim M_d/d$ exists \emph{if and only if} the Ces\`aro
average $D^{-1}\sum_{d\le D}M_d/d$ converges. The logarithmic average is not enough.
\end{remark}

\begin{remark}[the slope as a series over the jump tail; \PROVED\ identity, \MEASURED\ tail law]\label{fr:jumpseries}
There is a second reduction, to a statement about the distribution of a \emph{single} increment. Write
$J_n=M_{n+1}-M_n$ and $q_k(N)=\#\{n<N:J_n\ge k\}/N$. For a non-negative integer variable
$\mathbb E[J]=\sum_{k\ge1}\Pr[J\ge k]$, and $M_N/N$ is the empirical mean of the jumps, so
\[
\frac{M_N}{N}\;=\;\sum_{k\ge1}q_k(N)
\]
exactly, at every $N$ --- verified numerically to six decimals at
$N=1.25\times10^{4},\dots,4\times10^{5}$. By Tannery's theorem, if the $q_k(N)$ are dominated uniformly
in $N$ by a summable sequence and each converges as $N\to\infty$, then $\lim M_N/N$ exists and equals
$\sum_k\lim_Nq_k(N)$. So the front speed exists as soon as the \emph{jump distribution} converges with a
controlled tail.

The domination is measured and is unusually clean: the tail halves. On the fresh block
$n\in[2\times10^{5},4\times10^{5}]$, the ratios $q_{k+1}/q_k$ for $k=2,\dots,11$ read
\[
0.478,\;0.498,\;0.500,\;0.491,\;0.487,\;0.486,\;0.487,\;0.495,\;0.491,\;0.492,
\]
i.e.\ $q_k\asymp2^{-k}$ across a full decade of $k$, and the truncation error behaves as a tail of
ratio just below $\tfrac12$ requires, $\sum_{k>K}q_k$ lying within $3\%$ below $2q_{K+1}$ at every $1\le K\le8$. A uniform
bound $q_k(N)\le C2^{-k}$ is therefore the missing hypothesis of the domination half, and it is the
part of this route that looks provable rather than merely measured. We note but do \emph{not} claim a
connection to the halving of survivors in Section~\ref{sec:record}: that is a different object, and
sharing the factor $2$ is not evidence.

Two features of the distribution deserve recording. It is not monotone --- $\Pr[J=0]=0.455$,
$\Pr[J=1]=0.147$, $\Pr[J=2]=0.207$, so a jump of $2$ is commoner than a jump of $1$, which the
definition $M_d+1=\max(\tau_{d-1},\tau_d)$ makes unsurprising but which any model of the increment must
reproduce. And the convergence half is genuinely open. Measured on disjoint blocks rather than
cumulatively, every $q_k$ drifts downward over the first $4\times10^{5}$ diagonals, by $-0.0027$, $-0.0031$, $-0.0035$, $-0.0048$, $-0.0032$,
$-0.0026$, $-0.0014$, $-0.0006$ for $k=1,\dots,8$ between the first block and the last, and no single $k$ carries
the drift. That drift is a transient of the first $2$--$3\times10^{5}$ diagonals: past them no $q_k$ with
$k\le8$ shows a detectable trend on disjoint blocks of $10^{5}$ diagonals up to $d=2\times10^{6}$, and $q_1$ shows none over
the $39$ such blocks of $[10^{5},4\times10^{6})$, with mean $0.5430$ and block s.d.\ $0.0014$ (\MEASURED).
Flat blocks on a finite range are not convergence: a slow continuing drift, of the form $c/\log d$ with
small $c$, would not be resolved by blocks of this size. So this reduction, like the concave one, changes the shape
of the problem without weakening it: what must be proved is now the convergence of the density of a
local event rather than a property of the front.

One route is closed by this data. If $\delta_d=M_d-\rho d$ carried an explicit-formula structure ---
a main term plus oscillatory corrections indexed by some spectrum, the shape a $2$-adic odometer would
produce and the natural guess given that the diagonal periods are powers of two --- the corrections
would appear as discrete peaks in $\log_2d$. They do not. The spectrum of $\delta$ against $\log_2 d$
over $d\in[10^{3},4\times10^{5}]$ puts $99.81\%$ of its power in the lowest $200$ bins, and its largest
components sit at $1,2,4,5$ times the reciprocal of the window length, which is the signature of a
smooth aperiodic wander rather than of any oscillation. There is no log-periodic structure to expand.

What the tail law does buy is a compression. Since $q_k$ is geometric from $k=2$ with ratio $r$, the
identity closes on three numbers,
\[
\rho \;=\; q_1+q_2+q_2\frac{r}{1-r},
\]
and on the fresh block $d\in[2\times10^{5},4\times10^{5}]$ this reproduces the measured slope to
$5\times10^{-4}$ relative --- $q_1=0.5445$, $q_2=0.3970$, $r=0.4875$ give $1.319136$ against $1.318420$.
It is important not to overstate what that is: $\rho=\sum_kq_k$ is an identity, so the agreement measures
how nearly geometric the tail is and does not \emph{derive} $\rho$ from anything. Its value is that it
reduces the convergence question from a whole distribution to three sequences: $\lim M_d/d$ exists as
soon as $q_1,q_2$ and $r$ converge. Past the transient $q_1,\dots,q_8$ show no detectable drift to
$2\times10^{6}$ (\MEASURED), which, as above, is not convergence.

The same three numbers separate rules, and stably. Running the construction on every rule --- the
diagonal frame transfers verbatim, since $D_d(t+1)=f(D_{d-2},D_{d-1},D_d)$ is the same $f$ --- and
recomputing at $D_{\max}=2\,000,\,6\,000,\,12\,000$ gives $r=0.5000,0.5116,0.5135$ for Rule~30 against
$0.6757,0.6271,0.6445$ for Rule~$110$, with Rules~$135$ and~$137$ tracking their partners to three
decimals at every size. So $r$ is a stable invariant of the rule and not of the window, which the front
exponents of Remark~\ref{du:ring} were not. Two exclusions are needed: Rules~$22$, $126$ and~$182$
report jumps of order $10^{4}$ against a computation of depth $1.1\times10^{4}$, which are horizon artefacts --- Rule~$22$'s
diagonals are purely periodic from the seed with periods doubling without bound (Remark~\ref{fr:nofront}), so past the
horizon the detector reports the horizon --- and their ratios are meaningless; the mirror rules $86,149,101,169$ do not resolve at all,
because the frame $D_d(t)=u^t_{d-t}$ tracks the left front specifically and for a mirror rule the
corresponding front is on the other side at the other speed; and Rule~$45$ has the trivial light-cone front $M_d=2d+1$
(Remark~\ref{fr:nofront}) and resolves nothing for the same reason. That last is a limitation of the frame, not
a fact about those rules.
\end{remark}

\subsection{Computations}

All statements in this subsection are finite computations (\EXACT) unless labelled otherwise. $D_d(t)$ was computed from
Lemma~\ref{fr:lem:diag}(a) for $0\le d\le90\,000$, $0\le t<132\,000$. By
Proposition~\ref{fr:prop:period}(c),(d), $\tau_{d+1}\le\max(\tau_{d-1},\tau_d)+p_d$ and $p_{d+1}\le2\max(p_{d-1},p_d)$.
In the range, $\max_d\tau_d+2\cdot64<132\,000-512$ and $p_d\le32$. Periods were read on a terminal window
of length $512$, and preperiods as the last disagreement with the shift $p_d$. So $p_d$, $\tau_d$, $r_d$ and $\Delta_d$
are determined by data inside the horizon. Part of what follows is \KNOWN\ data. The data file of Wolfram's 2019 prize post \cite{Wolfram19} gives the settling time and the period of the $78\,288$ diagonals $0\le d\le78\,287$, and it agrees with ours at every one (ledger 10.167c2). OEIS A363346 \cite{A363346} has $1\,000$ terms. And \cite{Pat} (commit \texttt{ae8d14d}, dated 19 August 2026) has exact periods for $d<10^{6}$ by direct simulation, with the same seven frozen diagonals and period $32$ at $87\,867$.

\begin{enumerate}
\item[\textbf{E1}] (\EXACT) For $1\le d\le90\,000$: $\tau_d<d$ iff $d\in\{3,\dots,15,17,19\}$;
$\min_{20\le d\le90\,000}(\tau_d-d)=1$, attained exactly at $d\in\{23,29\}$;
$\tau_d/d\ge1.267$ for $1000\le d\le90\,000$ (minimum at $d=1026$).
First occurrences of $p_d=2,4,8,16,32$ are at $d=3,8,29,400,87\,867$.
$r_d\equiv0$ exactly for $d\in\{2,7,28,399,53\,207,58\,286,87\,866\}$.
\item[\textbf{E2}] (\EXACT) $\kappa_{d+1}\le P_d-1$ for $21\le d+1\le90\,000$, and two later exact computations
(one of them the third frame described in E7) extend it without exception to $21\le d+1\le8\times10^{5}$.
Below $90\,000$ the maxima of $\kappa$ for $P_d=4,8,16,32$ are $3,7,14,17$ and $\mu_{d+1}\le11$; to
$8\times10^{5}$ the maxima are $\kappa_{d+1}=19$ and $\mu_{d+1}=15$, both in the regime $P_d=32$. And
$\kappa_{d+1}\le P_d$ whenever $\mu_{d+1}=0$, as Lemma~\ref{fr:lem:G} requires.
\item[\textbf{E3}] (\EXACT) $k_i\le P_i-1$ for every plateau with $20\le d_i$ and $d_{i+1}<90\,000$.
The maxima for $P_i=4,8,16,32$ are $1,5,11,6$; $\max k_i=11$, first at $d_i=4913$.
\item[\textbf{E4}] (\EXACT) Among plateaus with $20\le d_i$, $d_{i+1}<90\,000$ and $k_i\ge1$, the inequalities of
(U) and (C) fail exactly for $d_i\in\{20,22,28,32,34\}$. For $d_i\ge36$ they hold, and
$\min(\sigma_i-d_i-P_i+1-k_i)$ over $d_i\ge D$ equals $2,28,258,3299,16\,812$ for
$D=36,100,1000,10^4,5\cdot10^4$.
\item[\textbf{E5}] (\EXACT) $d=36$ is a jump diagonal, and (Base) of Theorem~\ref{fr:thm:Hstar} holds with $d_{i_0}=36$ (E1).
\item[\textbf{E6}] (\EXACT, independent reproduction) E1 was recomputed from the definition of $\tau_d$ in a frame using
neither a row array nor a windowed readout: iterate the diagonal recursion of Lemma~\ref{fr:lem:diag}(a) over $d$, and for
each candidate power of two $p$ record $\max\{u: D_d(u+p)\neq D_d(u)\}$, so that $\tau_d$ is that index plus one exactly, a
period being accepted only when $\tau_d$ falls at least $1024$ below the horizon. Over $1\le d\le90\,000$ with $T=132\,000$
every diagonal settles and all five statements of E1 are reproduced exactly, including the first period $32$ at $d=87\,867$
and the eventually-zero set $\{2,7,28,399,53\,207,58\,286,87\,866\}$. Their frozen partners, the eventually-\emph{one}
diagonals, are $\{0,1,4,9,30,401,53\,209,58\,288,87\,868\}$ (that an eventually-one diagonal stands two places after each frozen one is \KNOWN: \cite{Rowland}; \cite{Dib}, commit \texttt{59a9c26}); note $D_7\equiv0$ identically, so $\tau_7=0$. This frame
(\texttt{s3\_tau\_convention.py}) reproduces E1 only. E3, the (C) half of E4 and E5 were reproduced separately
from the last-violation array of a second exact per-diagonal run, using no code or data of the paper's frame
($49\,110$ plateaus, $0$ discrepancies), and E2 by the two later computations named in E2; the (U) half of
E4, which needs $\mu_i$ and hence the defect sequences themselves, has not been independently reproduced. Note also that $P_d=\max(p_{d-1},p_d,p_{d+1})$ is
flat across a doubling of $p$, the increase of $P$ sitting one diagonal earlier: $P$ reaches $32$ at $d=87\,866$ while $p$
reaches $32$ at $d=87\,867$.
\item[\textbf{E8}] (\EXACT, from archived public data) Gravner's Stage~I label tree for this rule at period $16$, $2\,159\,030$ nodes to generation $894\,234$, has been public since 2012; parsing it in full reproduces those counts and identifies the seed's lineage, whose branch generations are exactly $2,7,28,399,53\,207,58\,286$ with leaf $87\,866$. So the list in E7 is archived data, not a new computation, and the $8$ in the published list $2,8,28,399$ of \cite[\S8]{GG} is a typo, their own archive giving $7$. The tree has $19$ branchings and $20$ leaves, $16$ dying by period doubling, which are their sixteen succession lines; the branching at $72\,575$ lies on the sibling lineage, not the seed's. Two consequences for us: Conjecture~\ref{fr:conjM} is \emph{true} at $P=16$, the seed's exit at $87\,866$ being the minimum of the sixteen, while remaining refuted at $P=32$; and Proposition~\ref{fr:gaps} survives a much larger test, since across all twenty archived lineages the minimum gap is $5$ and the only gaps at most $30$ anywhere are one $5$ and one $21$, with none of the sixteen excluded values occurring. No period-$32$ tree is archived, so the death generation $65\,154\,361$ quoted in \cite[\S8]{GG} is not reproducible from published files.
\item[\textbf{E7}] (\EXACT) The eventually-zero diagonals below $d=400\,000$ are exactly $2,7,28,399,53\,207,58\,286,87\,866$: the set is unchanged from the paper's earlier range and there is no eighth such diagonal and no period $64$ below $400\,000$ (bit-packed frame, $T=608\,192$, $108$ s, with $p_d$ read twice on disjoint windows and cross-checked against the exact per-diagonal frame with no mismatch). Their six gaps are $5,21,371,5079,29\,580,52\,808$. A previous version of this item quoted a gap histogram with $47$ gaps of $2$; those figures belong to a different population, the $349$ diagonals of $[400,87\,867)$ with $p_d\mid8$, and are withdrawn here --- for the eventually-zero diagonals a gap of $2$ is impossible by Proposition~\ref{fr:gaps} below, and the apparent odd-gap deficit in that other population does not recur at matched Poisson power: in the corresponding population of the period-$32$ regime, the $664$ diagonals of $[87\,867,4\times10^{7}]$ with $p_d<32$, the gap window $[3,977]$, whose Poisson expectation matches that of $[3,10]$ among the $349$, holds $5$ odd gaps against $5.36$ expected (\MEASURED).
A third implementation of the frame, written from Lemma~\ref{fr:lem:diag}(a) alone (rows bit-packed over
$d$, $T=1\,120\,000$, every diagonal settled at least $63\,936$ rows below the horizon), reproduces E1, E2, E3, the (C) margins of E4 and E7 and extends them to
$d=8\times10^{5}$: $\kappa_{d+1}\le P_d-1$ and $k_i\le P_i-1$ hold throughout, the longest stall run is $13$,
first at $d_i=184\,979$ (the maxima $11$ of E3 and $13$ here are also in \cite{Dib}), the (C) margin $\min(\sigma_i-d_i-P_i+1-k_i)$ over $d_i\ge10^{5}$ and over
$d_i\ge4\times10^{5}$ is $33\,230$ and $129\,626$, and there is no eighth eventually-zero diagonal and no
period $64$ below $8\times10^{5}$. The hypotheses of (C) are therefore verified for plateaus ending before
$8\times10^{5}$ (E2, E4 and this extension), and those of (U) only for plateaus ending before $90\,000$
(E4). Hence the computation reproduces $\tau_d\ge d$ on $[20,8\times10^{5}]$ through Theorem~\ref{fr:thm:Hstar}(C) and proves nothing beyond; E10 checks $\tau_d\ge d$ directly to $4\times10^{6}$. Conjecture H itself is \OPEN, and its target is stated as an open question in print by Wolfram \cite{Wolfram}. Gravner and Griffeath's regularity-expansion machinery cannot reach it: it requires the locked region to be recognised by a finite automaton, which unbounded periods forbid, there is no final Stage~I tree for this rule by their own \S8, and they diagnose the same failure for the mirror rule in the words ``no local method from the present paper applies''.
\item[\textbf{E9}] (\MEASURED, word-granular) The bin medians of Remark~\ref{fr:slopebar}, and its linear fit over $d\ge3\times10^{6}$, come from a bit-packed frame (\texttt{s3\_frame\_bitpack.py 4000000 5500000}, $9\,179$~s) holding $64$ diagonals per machine word and recording, per word, the \emph{maximum} last violation over its $64$ diagonals, de-biased by the constant $84$ (validated against the exact per-diagonal frame at $d\le9.5\times10^{4}$: offset min $49$, max $135$, mean $84.8$; de-biased bin medians $1.3370/1.3353$ against the exact $1.3371/1.3353$). It yields bin medians and the slope robustly and no per-diagonal statement: no bound of the form $\tau_d/d\ge X$ for every $d\le4\times10^{6}$ can be read off it, so it does not extend E1's per-$d$ bound; E10 does. Period $32$ settles for all $61\,129$ words of $64$ diagonals in $[87\,808,4\times10^{6}]$, so there is no sixth doubling below $4\times10^{6}$ (the lower bound of \cite{NKS} for it is $2.1\times10^{9}$).
\item[\textbf{E10}] (\EXACT, with a certificate) $\tau_d$ and $p_d$ per diagonal for $0\le d\le4\times10^{6}$ (\texttt{run4Mv2\_tau.npy}, a bit-parallel frame restricted to a band of rows). The certificate: the last block of rows computed contains no violation, and a violation at $(t+1,d)$ needs one at $(t,d')$ for some $d'\in\{d-2,d-1,d\}$, so none occurs later. The file agrees with an earlier exact per-diagonal file to $8\times10^{5}$ and with an independently written code to $2\times10^{6}$ except at $d=0$ and $d=7$, which never violate and which this file records as $-1$ (the other two files give $\tau_0=\tau_7=0$, since $D_0\equiv1$ and $D_7\equiv0$). $\tau_d<d$ exactly for $d\in\{3,\dots,15,17,19\}$, so $\tau_d\ge d$ on $[20,4\times10^{6}]$; $\min\tau_d/d$ over $[1000,4\times10^{6}]$ is still $1.2671$, at $d=1026$, and over $[800\,001,4\times10^{6}]$ it is $1.3194$, at $d=900\,864$; there is no period $64$ below $4\times10^{6}$; and the largest jump is $24$, at $d=1\,680\,044$.
\end{enumerate}


\begin{remark}[the front is the one object here that Rule 90 does not have]\label{fr:nofront}
Throughout this paper the sharpest obstruction to any general argument has been Rule~90, which
satisfies more hypotheses than Rule~30 does and fails every conclusion we want
(Proposition~\ref{p2:rule90}, and the width table of Section~\ref{sec:p2}). The front is the exception,
and it is worth saying so explicitly because it tells us where to spend effort. For Rule~90 from the
single seed, $\tau_d=0$ for every $d$: the diagonals are purely periodic from the start, with no
transient, which follows from Kummer's theorem on the binomial coefficients and which we verified
against direct simulation with no mismatch. So $\Delta_d\equiv0$, $M_d\equiv-1$, and Rule~90 has no
front, no slack and no defect. The same holds for Rules~150, $22$ and $106$ (\EXACT\ on the computed range; Rule~$22$'s diagonal periods grow without bound). Every statement in this
Part is therefore about a structure that the firewall rule does not possess, so no Rule~90 argument can
be run against it --- unlike our automaticity census and every density statement, where it can. The
price is that Rule~90 also cannot serve as a control here: among the rules with a front, the ones we can
compute ($45$ with $M_d=2d+1$, and $60$ with $M_d=d$) satisfy the band inequality with a non-positive
constant precisely because their periods are bounded, so nothing available tests the one feature that
matters, $p_d\to\infty$.
\end{remark}

\section{The background walk and unbounded periods}\label{fr:walk}

Let $r_d$, $p_d$, $\tau_d$ be as in Definition~\ref{fr:def:bg}, with $r_d=0$ for $d<0$. For $P$ a power of $2$ let $V_P$ be the set of $P$-periodic sequences $\N\to\F$, a finite set of size $2^P$.

\begin{definition}
For sequences $b,c\colon\N\to\F$ define $\pi(b,c)=(a,b)$ with
\[
a(t)=c(t+1)+b(t)+\bigl(1+b(t)\bigr)c(t)\qquad(t\ge0).
\]
$\pi$ maps $V_P^2$ to $V_P^2$ for every $P$. Put $v_d=(r_{d-1},r_d)$ for $d\ge0$.
\end{definition}

\begin{lemma}[\PROVED]\label{fr:backward}
$\pi(v_{d+1})=v_d$ for all $d\ge0$. Moreover $v_0=(0,1)$, $\pi(v_0)=(0,0)$ and $\pi(0,0)=(0,0)$.
\end{lemma}
\begin{proof}
The background equation, Lemma~\ref{fr:lem:bgdef}(b) at index $d+1$, reads
$r_{d+1}(t+1)=r_{d-1}(t)+r_d(t)+r_{d+1}(t)+r_d(t)r_{d+1}(t)$ for all $t\ge0$; solving for $r_{d-1}(t)$ gives
$r_{d-1}(t)=r_{d+1}(t+1)+r_d(t)+(1+r_d(t))r_{d+1}(t)$, i.e.\ $\pi(r_d,r_{d+1})=(r_{d-1},r_d)$.
By Lemma~\ref{fr:lem:diag}(c), $r_0\equiv1$, and $r_{-1}\equiv0$ by convention, so $v_0=(0,1)$. Then $\pi(0,1)=(1+0+1,\,0)=(0,0)$ and $\pi(0,0)=(0,0)$.
\end{proof}

\begin{theorem}[unbounded diagonal periods; \KNOWN: Gravner--Griffeath \cite{GG}, Prop.~3.3]\label{fr:unbounded}
The pairs $v_d$, $d\ge0$, are pairwise distinct. Consequently $\sup_dp_d=\infty$: the minimal eventual periods of the left diagonals of Rule~30 from the seed are unbounded.
\end{theorem}
\begin{proof}
Suppose $v_{d+m}=v_d$ with $m\ge1$. Applying $\pi^d$ and using Lemma~\ref{fr:backward} repeatedly, $\pi^d(v_{d+m})=v_m$ and $\pi^d(v_d)=v_0$, so $v_m=v_0$; then $\pi^m(v_0)=\pi^m(v_m)=v_0$. But $\pi(v_0)=(0,0)$ is a fixed point of $\pi$ different from $v_0$, so $\pi^m(v_0)=(0,0)\ne v_0$ for every $m\ge1$, a contradiction.
If the periods were bounded, then, the $p_d$ being powers of $2$ (Proposition~\ref{fr:prop:period}(b)), all $r_d$ would lie in $V_P$ for $P=\max_dp_d$, and the infinitely many distinct $v_d$ would lie in the finite set $V_P^2$.
\end{proof}

\begin{remark}[provenance]
Gravner--Griffeath \cite{GG}, Proposition~3.3, proves that a left-permutative edge cellular automaton
has no periodic solution, by exactly this argument: left permutativity gives every label a unique
predecessor, a repetition on the ancestral line forces the whole ancestral line to carry the zero label,
and that contradicts the root label~$1$. Their labels are our eventual cycles $r_d$, their root is
$r_0\equiv1$, and their unique predecessor is our $\pi$. Rule~30 is named there as an instance. Jen
\cite{Jen90}, Proposition~3, is the related statement. The present section is therefore a translation
into the diagonal frame, not a new theorem; it is kept because Proposition~\ref{fr:universalwall}
and Theorem~\ref{fr:universaldata} are stated in that frame. Read for minimal periods, the fact that $\pi$ maps $V_P^2$ to itself, with Lemma~\ref{fr:backward}, gives the pair-period law: $\max(p_{d-1},p_d)$ never decreases with $d$, so the largest period up to diagonal $d$ is $\max(p_{d-1},p_d)$. In the second form it is Lemma~2 of Nersissian's arXiv:2609.25077 \cite{Nersissian}.
\end{remark}

\begin{remark}
The theorem says nothing about how fast the periods grow. \EXACT: the running maximum of $p_d$ first reaches $2,4,8,16,32$ at $d=3,8,29,400,87\,867$ --- these depths, and a sixth at ``$2{,}107{,}985{,}255$ or more'', are already tabulated in \cite[p.~871]{NKS}, and the per-diagonal periods and preperiods are OEIS A363345 and A363346 \cite{A363346} and, to $d=78\,287$, the data file of Wolfram's 2019 prize post \cite{Wolfram19}, which agrees with ours at every diagonal; Brunnbauer's \cite{Brunnbauer} independently computed table of $410$ pairs agrees with ours on every period, with his offset equal to $\tau_d-\lceil d/2\rceil$ throughout, and $r_d\equiv0$ exactly at $d\in\{2,7,28,399,53\,207,58\,286,87\,866\}$ for $d<90\,000$ (E1 in Section~\ref{fr:sec:front}). Rowland's parity law \cite{Rowland} (the period doubles across a frozen diagonal iff the cycle two diagonals to the left has an odd number of ones) is Proposition~\ref{fr:prop:period}(d) and was verified at all seven frozen diagonals. For $P\le16$ the forward tree of $(0,1)$ under $\pi^{-1}$ was enumerated exactly: $|F_1|=3$, $|F_2|=13$, $|F_4|=97$, $|F_8|=3065$, $|F_{16}|=34\,541\,081$ vertices; the seed's path exits $F_{16}$ at $d=87\,866$ along the shortest maximal branch (\EXACT; at $P=32$ it is not the shortest, Remark~\ref{fr:conjM:refuted}).
\EXACT\ (ledger 10.231d~I, 10.232d~B, 10.243d; the companion note \cite{WN}; public earlier in \cite{Dib,Pat}, as stated below): the label recursion of Lemma~\ref{fr:backward}, run forward in C with no simulation of the rule, gives the eighth frozen diagonal of the seed at $d^{*}=1{,}420{,}878{,}968$, a fork of even weight (the label $r_{d^{*}-1}$ has period $32$ and weight $20$), so its two candidate successors are different walls. On the branch whose label begins \texttt{6cbda353} the running maximum reaches $64$ at exactly $d=2{,}107{,}985{,}255$, Wolfram's number to the digit; on the other branch (\texttt{93425cac}) the next frozen diagonal is $3{,}340{,}408{,}059$, a second fork, and on every continuation of that branch the period does not reach $64$ before $9{,}958{,}700{,}505$ (\EXACT\ over all branches to $2\times10^{10}$; ledger 10.160c2). This is consistent with reading ``or more'' as the fork; why it was written is not recorded. These integers are not new here. The eighth frozen diagonal and the first period-$64$ onsets on its two branches ($2{,}107{,}985{,}255$ and, on the other side, $9{,}958{,}700{,}505$) were, to our knowledge, first made public on 8 September 2026 in two independent code repositories: \cite{Dib} (commit \texttt{94833d2}) and \cite{Pat} (commit \texttt{04f30f6}, from a run of about 30 August 2026). The first gives the frozen diagonal $1{,}420{,}878{,}968$ with its two complementary continuations, the onset $2{,}107{,}985{,}255$ on one of them, which it matches to \cite[p.~871]{NKS}, and the frozen diagonals $3{,}340{,}408{,}059$ and $4{,}989{,}445{,}007$ on the other. The second gives the branch point $1{,}420{,}878{,}969$ with the continuations \texttt{93425cac} and \texttt{6cbda353}, both onsets, and every branch up to $1.2\times10^{10}$. Of the two, only \cite{Pat} has the onset $9{,}958{,}700{,}505$. We obtained the same integers independently on 24 September 2026 (ledger 10.160c2) and learned of these repositories on 1 October 2026 (ledger 10.165c2). We found the numbers in no paper, preprint or OEIS entry. Beyond $1.2\times10^{10}$ the branch tree was followed here to $2\times10^{10}$ (ledger 10.160c2); that part is in neither repository. Neither repository determines the seed's branch: \cite{Pat} records it as unresolved, and \cite{Dib} infers it from the match with \cite{NKS} and lists it as not checked. Beyond the doubling, branch \texttt{6cbda353} has no frozen diagonal up to $10^{12}$, so its period-$64$ regime lasts at least to $10^{12}$: \EXACT\ on the second machine's arithmetic, with $13$ of the run's $8\,492$ steps of $10^{8}$ diagonals recomputed by two independently written programs and the remaining $8\,479$ resting on the second machine alone (ledger 10.243d; a step-by-step recheck of the whole log is running). Which branch the seed itself takes is the subject of Remark~\ref{fr:branchrule}.
\end{remark}

\subsection{The left half as a driven machine, and the universal wall}\label{fr:machine}

The cells $x\le-1$ depend only on cells $x\le0$, so the left half of any Rule~30 evolution is Rule~30 on the half-line $x\le-1$ driven by column~$0$ as a boundary input, from whatever initial left state it has. Call this the \emph{left machine}; its boundary enters through the OR at $x=-1$ and its output returns to the column through the XOR. Symmetrically the right half is a machine driven by the same column, receiving it through the XOR and returning through the OR. The column is $c_{t+1}=L_t+(c_t\vee R_t)$ with $L,R$ the two outputs. Everything in this Part is a statement about the left machine, and Lemma~\ref{fr:backward} and Theorem~\ref{fr:unbounded} use nothing about the boundary: aligning diagonals to the trajectory of the leftmost $1$ (which moves one cell left per step and is all ones, since $001\mapsto1$ and $000\mapsto0$), $D_0\equiv1$ and $D_{-1}\equiv0$ hold for every left machine, so its background pairs form a backward orbit of $\pi$ from $(0,1)$ and its diagonal periods are unbounded.

\begin{proposition}[universal background; \KNOWN\ in a stronger form, Gravner--Griffeath {\cite[\S8]{GG}}]\label{fr:universalwall}
For a power of two $P$ let $F_P$ be the tree of backward orbits of $\pi$ from $(0,1)$ inside $V_P^2$. The left machine's background pairs $v_d$, in edge-aligned indexing, follow a maximal path of $F_{p}$ for each period regime $p$ until the period doubles. The enumerations $|F_1|=3$, $|F_2|=13$, $|F_4|=97$, $|F_8|=3065$ (heights $2,7,28,399$) show that for $P\le8$ the maximal path of $F_P$ is unique up to time translation. Hence, for every boundary sequence and every finite initial left state, the periodic background of the left half agrees with the seed's on every diagonal of edge-aligned index below $400$, up to a single time translation, with the same periods and the same doubling positions $3,8,29,400$.
\end{proposition}
\begin{proof}
Lemma~\ref{fr:backward} holds for every left machine after edge alignment, so $(v_d)$ is a backward orbit of $\pi$ from $(0,1)$; while the periods stay $\le P$ it lies in $F_P$, and by Theorem~\ref{fr:unbounded} it cannot terminate. Within each regime it is therefore a path of that regime's tree that the orbit never leaves early; that it is \emph{maximal} in the tree needs one more step --- the sibling of the orbit's next element, which differs from it by the solution $e$ of $e(0)=1$, $e(t+1)=(1+b(t))e(t)$ (both preimages of $(a,b)$ under $\pi$ obey the same first-order recursion), must not lie in $V_P^2$ --- and that holds whenever $b=r_{d+1}\not\equiv0$, since then $e$ is finitely supported and nonzero, hence not periodic; at the frozen diagonals, where $b\equiv0$, we do not check it here: there the statement is \KNOWN\ from Gravner and Griffeath's label tree to $d=53\,207$ \cite[\S8]{GG} (Remark~\ref{fr:universalwall:prov}) and \EXACT\ to $12\,000$ by Theorem~\ref{fr:universaldata}, and nothing below rests on more (ledger 10.242d~A7). Uniqueness of the maximal path of $F_P$ for $P\le8$ is the finite enumeration, and the time translation is the freedom in choosing the phase of a periodic pattern.
\end{proof}

\begin{remark}[provenance, and the upgrade it supplies]\label{fr:universalwall:prov}
This is \KNOWN\ and the published statement is stronger than the one proved above. Gravner and Griffeath
\cite[\S8]{GG}, writing about this rule, record that ``the first genuine branching occurs for period 16 at
generation 53{,}207'' and conclude that ``no matter what the seed, there is only one limit cycle on
$[0,53207]$'', adding that ``all seeds have only 16 distinct limit cycles on the leftmost 65 million
cells''; Rowland \cite[\S5]{Rowland} states the same reduction, that unless the previous column is
eventually white the period of a column is independent of the seed. Both cover every seed to $d=53\,207$,
where the proposition above covers $d<400$ with an enumeration and the accompanying computation reaches
$12\,000$.

Their argument also removes the enumeration from ours. They prove that at a period doubling ``the two
resulting successors are each other's rotations by $\pi$'' and that ``the only label repetition is among
the two siblings at a period doubling branching'': their label tree at period $\pi$ is our $F_P$, so
uniqueness of the maximal path up to time translation holds for \emph{every} $P$, not merely $P\le8$, and
the finite enumerations $|F_1|=3,\dots,|F_8|=3065$ become illustrations rather than the proof. Their
$\pi=16$ tree, $2{,}159{,}030$ labels over $894{,}234$ generations, matches our $|F_{16}|=34\,541\,081$
once labels are counted with their $16$ phases, the small deficit being the short-period labels. One
bookkeeping note in the other direction: their list of death generations reads $2,8,28,399$ where ours
reads $2,7,28,399$, and $7$ is correct here, confirmed four independent ways.
\end{remark}


\begin{remark}[the branch rule, the phase congruence, and the branch at the eighth frozen diagonal; \PROVED, \EXACT\ and \MEASURED\ parts as marked]\label{fr:branchrule}
Session~d of the ledger (10.231d~I, 10.237d, as corrected by 10.242d~A) makes the mechanism behind Conjecture~\ref{fr:universal} explicit.
(a) \emph{The branch rule} (\PROVED; ledger 10.231d~I(ii); the one-cell branch step is \KNOWN, being Rowland's \cite[Lemma~3]{Rowland} in kind, and a related check was public earlier: in \cite{Pat}, log of 9 September 2026, at the seven frozen diagonals below $10^{6}$ the phase and one boundary bit of the transient select the branch). Let $d$ be a frozen diagonal, so $D_d(t)=0$ for $t\ge\tau_d$ and $D_d(\tau_d-1)=1$, and let $c^{0}$ be the candidate label vanishing at row~$0$. The frame recursion gives $D_{d+1}(\tau_d)=D_{d-1}(\tau_d-1)\oplus1$ and, for $t\ge\tau_d$, $D_{d+1}(t+1)=D_{d-1}(t)\oplus D_{d+1}(t)$, which both candidates satisfy; writing $D_{d-1}=r_{d-1}+\Delta_{d-1}$, the defect $\Delta_{d+1}$ is eventually the constant $1+c^{0}(\tau_d-1)+\sum_{t\ge\tau_d-1}\Delta_{d-1}(t)$. Hence $M_{d+1}=M_d$ (Theorem~\ref{fr:thm:B}(d) again), the label taken is the candidate with $r_{d+1}(\tau_d-1)=1+(\text{parity of }\Delta_{d-1}\text{ on rows }\ge\tau_d-1)$, and nothing below row $\tau_d-1$ has any part in the branch: it is decided by the phase of the front against the wall and by the tail of one defect. Checked cell by cell at $399$, $53\,207$, $58\,286$ and $87\,866$ (\EXACT).
(b) \emph{The phase congruence} (\PROVED\ with an \EXACT\ step; ledger 10.231d~I(iv$'$), 10.242d~A2). Let a left machine whose front stays left of its boundary column have, beyond some diagonal, the seed's wall up to a translation $s_w$ and the seed's front up to a translation $s_b$. Then $s_b\equiv s_w$ modulo the period on every stretch containing $20\,000$--$87\,866$ or $90\,000$--$800\,000$: the landing rule of Theorem~\ref{fr:thm:B}(c), transported to the machine, says that the seed's wall obeys the landing rule shifted by $\delta=s_b-s_w$ at every jump diagonal, and the census of all $37\,086$ jump diagonals of the period-$16$ regime and all $385\,784$ of the period-$32$ regime finds the shifted rule satisfied for $\delta\equiv0$ only (the best other residue holds at $26.1\%$ of them). The six left machines of the discussion following Theorem~\ref{fr:universaldata} confirm it in all twelve cases measured. With (a) this reduces the frozen-diagonal half of Conjecture~\ref{fr:universal} to the universality of the front. That universality fails for some finite starting words: the starting state of \cite{Lania} takes the other branch after $53\,207$ (Conjecture~\ref{fr:universal}), so its front is not the seed's up to translation there. For the boundaries and words measured it holds (\MEASURED).
(c) \emph{The advance criterion} (\PROVED; ledger 10.237d~A). At a frozen diagonal $d$ with $m=M_{d-1}$ the twin labels $r_{d-1}=r_{d-2}=:r$ make the frame recursion read $D_d(t+1)=r(t)\oplus(r(t)\vee D_d(t))$ for $t\ge m+1$, so the front advances at $d$ exactly when $D_d(m+1)=1$, that is, when the frozen diagonal is black on the row after the previous front, and then $M_d=\min\{t\ge m+1: r(t)=1\}$.
(d) \emph{The seed's branch at $d^{*}$} (ledger 10.237d~D, 10.242d~A8--A10; \cite[\S6]{WN}; the fork itself is in \cite{Dib,Pat}, which leave the seed's branch undetermined). Given the front's path over the three diagonals before $d^{*}=1{,}420{,}878{,}968$ (rows $M_{d^{*}-3},M_{d^{*}-2},M_{d^{*}-1}\equiv29,31,31\pmod{32}$: a move of two at $d^{*}-2$, a stall at $d^{*}-1$), five label cells force the branch \texttt{6cbda353}, the one on which the period doubles at Wolfram's number (\PROVED\ given the path). The path itself is \OPEN: it is shared by all eight coin-driven edge machines whose front histories were recorded across the fork, and all twenty-eight such machines take that branch (\MEASURED). These are model verdicts and are not calibrated evidence: checked against the exact simulation of the seed on $[4\times10^{6},4\times10^{7}]$, the same coin-driven machine with windows up to $W=256$ was unanimous and wrong at $d=20\,238\,347$, from $50\,000$ and from $150\,000$ diagonals back (ledger 10.164c2). A census over every front history consistent with the landing rule finds both verdicts at every frozen diagonal checked, so the wall alone decides at none of them (\EXACT; that the labels do not fix the branch is \KNOWN\ in substance, \cite{Rowland,Pat,Lania}); the seed's own advances at $28$ and $58\,286$ are forced by the labels at no depth up to $32$ and are decided by the transient. Every branch of the seed's wall is therefore a finite computation from the labels and the front's phase history, the phase history being the one measured input.
\end{remark}

\begin{theorem}[\EXACT]\label{fr:universaldata}
With the left machine simulated for $20\,000$ rows and diagonals of edge-aligned index in $[20,12\,000)$, all $11\,980$ locked diagonals carry the seed's background pattern up to one global time translation per run, with identical periods and doubling positions, for each of: the seed column with a random initial left state of width $40$ and of width $100$; a random-coin boundary with empty state, a second coin, and a coin with a random state of width $100$; the boundaries $(01)^\infty$ and $(0011)^\infty$ with empty state. The translations are $5,9,13,0,6,8,5 \pmod{16}$; the $(01)^\infty$ wall is the seed's wall translated by half a period. The range covers the whole period-$8$ regime and the period-$16$ regime to $12\,000$ (it ends at $53\,207$). Scripts \texttt{s3\_leftdrive5.py}, \texttt{s3\_leftdrive7.py}.
\end{theorem}

Two further facts from the same runs (\EXACT\ and \MEASURED): the lock slope $\tau_d/d$ of the left
machine takes the same value for every boundary tried in those runs, periodic or random --- the value itself, and why the data give no bound
on its limit, are in Remark~\ref{fr:slopebar}; and two different initial left states under the same boundary disagree at $x=-1$ on half of all rows forever, so the left machine never forgets its initial state.
The slope is not the same for \emph{every} boundary. Driven from the empty left state by the background
column $b_B(t)=r_t(t)$, the left machine is the background field $r_{x+t}(t)$ itself, which obeys the rule by
Lemma~\ref{fr:lem:bgdef}(b) (this field as a Rule~30 evolution of its own is \KNOWN, \cite{Dib}; see Remark~\ref{fr:frozencolumn}): it has no transient and its front stays at the boundary, with
$\tau_d-d\in[-15,0]$ off the frozen diagonals and $M_d/d=0.99999$ at $d=2\times10^{5}$ (\EXACT\ to
$2\times10^{5}$), so its lock slope is $1$. The front's slope is therefore claimed to be a property of the
machine, and not of the column, only for boundaries that differ from $b_B$ at a positive rate; every such
boundary tried ignited a front with the seed's slope (\MEASURED), and that every boundary which ignites a
front has that slope is \OPEN. For the ignited machines the agreement is stronger than equal slopes: six
left machines, driven by $(01)^\infty$, $(0011)^\infty$, a coin and $b\equiv1$, and by the seed column from
random initial left states of widths $40$ and $100$, have the seed's settling times, after edge alignment,
up to one constant translation per machine for every $d>4\,484$ (\EXACT\ to $2\times10^{5}$; in a run to
$8\times10^{5}$ with the seed column for $t<5\,000$ and $b_B$ afterwards, the translation was reset in
episodes where the boundary's influence reached the front, each time by a multiple of the period $32$). Consequently, for the boundaries and initial states tried, the wall is one fixed object that the column does not mark, and every trace of the boundary and of the initial state lives in the transient $\Delta$. That the stripes are the same for every input tried is \KNOWN\ (\cite[\S8]{GG}; \cite{Dib}, commit \texttt{47f8ab5}; \cite{Lania}). It is not true of every finite starting word: the starting state of \cite{Lania} changes the wall after $53\,207$ (Conjecture~\ref{fr:universal}).

\begin{remark}[the edge automaton: rule, wall and one input bit; \PROVED, \MEASURED\ and \EXACT\ parts as marked]\label{fr:edgeautomaton}
The frame recursion $D_{d+1}(t+1)=D_{d-1}(t)\oplus(D_d(t)\vee D_{d+1}(t))$ shows that diagonals $d-1$ and $d$ on rows $[M_d-W,\,M_d+65]$ together with the single cell $D_{d+1}(M_d-W)$ determine diagonal $d+1$ on those rows, and since $M_{d+1}\le M_d+p_d\le M_d+64$ by Theorem~\ref{fr:thm:B}(c) they determine $M_{d+1}$. Hence the front is generated by the wall and exactly one input bit per diagonal, the cell $W$ rows below the front on the new diagonal, for any $W\ge0$ (\PROVED; ledger 10.231d~A; run with the true bits the machine reproduces the seed's $M_d$ with zero deviations). With that bit a fair coin --- and equally with any biased coin or a constant --- the machine reproduces every proved law and every registered frequency of the seed's edge (the near-laws are the exception, below): the move probability $0.5432$, the mean move $2.432$, the share of cycles with two or more waits $0.1082$ and the tight-move fraction $0.3544$ change by at most $2\times10^{-4}$ as the input law runs from constant $0$ to constant $1$ (\MEASURED; ledger 10.231d~H, 10.233d~D). Its speed on the seed's period-$32$ wall is $1.321294\pm0.000019$ over $2.84\times10^{10}$ diagonals (twenty coins) and $1.321438\pm0.000007$ on the period-$64$ wall of branch \texttt{6cbda353} beyond the doubling; within a block of $10^{8}$ diagonals twenty coins agree to $\pm0.0001$ while the block means have standard deviation $0.00013$, seven times the coins' share, so the speed of a stretch of wall is a property of that stretch, the coin only realising it (\MEASURED; ledger 10.233d~F, 10.242d~A16). These are numbers of the coin model on the wall and not a limit of the seed's own front, whose precision is what Remark~\ref{fr:slopebar} says it is. A SAT encoding of a window two rows below the front, with the frame recursion, the definition of the front and Theorem~\ref{fr:thm:B}(c) as its only constraints, proves $82$ of the $86$ deterministic two-event contexts and $278$ of the $279$ three-event contexts of the edge's event table (\EXACT; ledger 10.232d~C); the remaining patterns are broken by the coin machine at $0.15$--$1.2\%$ per case, but $142$ of $147$ exceptions fall while its front is off the seed's $32$-lattice, in states the true dynamics never visits, so they are laws of the universal front rather than of rule and wall (\EXACT, with the mechanism; ledger 10.233d~A).
\end{remark}

\begin{proposition}[gaps between eventually-zero diagonals; \PROVED\ modulo the finite enumerations named in the proof, which are \EXACT]\label{fr:gaps}
Let $g\ge1$ and suppose $r_d\equiv0$ and $r_{d+g}\equiv0$. Then $g\notin\{1,2,3,4,6,7,9,10,11,13,15,17,18,19,20,22\}$ for any
$d$ and any periods; for $g=5$ the only solution is $d=2$ and for $g=21$ the only solution is $d=7$. Hence the only pairs of
eventually-zero diagonals at distance at most $22$ are $\{2,7\}$ and $\{7,28\}$, with the four distances $8,12,14,16$ settled
only for joint period at most $16$.
\end{proposition}
\begin{proof}[Proof sketch]
With $r_d\equiv0$ the background equation at index $d+2$ becomes a closed affine recursion, and the equations at
$d+2,\dots,d+g$ together with $r_{d+g}\equiv0$ force a deterministic map on $\mathbb F_2^{\,g-1}$. Its recurrent subgraph is
functional for every odd $g\le21$, so the periodic solutions form a complete finite list, which is then cut by three
independent facts: Lemma~\ref{fr:lem:bgdef}(d); that a joint period must be a power of two, which is what removes
$g=11,15,17,19$, whose surviving cycles have lengths $11,20,12,5$; and the $\pi$-descent of Lemma~\ref{fr:backward}, which
must reach $v_0=(0,1)$. Two independent calibrations agree: brute force over all $(x,0)$ in $V_L^2$ for $L\le16$ and $g\le40$
returns exactly the pairs at distances $5$, $21$ and $26$, namely $d=2$, $d=7$ and $d=2$, with no false positive or negative,
and the state machine at $g=21$ independently returns $d=7$.
\end{proof}

\begin{remark}[the orbit is a necessary ingredient]\label{fr:gaps:orbit}
Gap $3$ falls to the algebra alone: the only consistent pairs are $(0,0)$ and $(1,1)$ and Lemma~\ref{fr:lem:bgdef}(d) kills
both. Gap $7$ does not. There a family survives every background equation, namely $r_{d+2}\equiv1$, $r_{d+4}=(10)^\infty$,
$r_{d+5}=r_{d+6}=(0110)^\infty$ and $r_{d-1}=(1100)^\infty$, arising from the unique $4$-cycle of an eight-state map; it dies
only because its $\pi$-descent falls into a $28$-cycle instead of reaching $v_0$. So the background equations alone do not
suffice and the backward orbit is needed. The surviving cycle length is non-dyadic exactly at $g=11,15,17,19$ and dyadic at
$g=5,7,13,21$; a lemma forcing that dichotomy would settle every $g$ at once, and is \OPEN.
\end{remark}

\begin{remark}[the fork law; twin law \PROVED\ and public earlier in \cite{Pat,Dib}, law \MEASURED, rate conditional]\label{fr:forklaw}
A frozen diagonal occurs exactly when the two labels before it coincide, $r_{d-1}=r_d$: with $a=b$ the recursion of Lemma~\ref{fr:backward}, $c(t+1)=a(t)\oplus b(t)\oplus(\neg b(t)\wedge c(t))$, reads $c(t+1)=\neg b(t)\wedge c(t)$, and $c$ dies at the first black cell of $b$ (\PROVED; ledger 10.235d~A). This twin criterion is \KNOWN: it is Lemma~A of \cite{Pat} (commit \texttt{d612567} of 19 August 2026, on its main branch on 28 August 2026), with proof, and it is in \cite{Dib} (commit \texttt{59a9c26}, 8 September 2026) and in Lean there. The wall is thus the orbit of a deterministic map on pairs of $P$-bit patterns, $4^{P}$ states of which the $2^{P}$ twin pairs are the forks; doubling the period squares the state space and only doubles the fork set, which is the shape of the gaps $5$, $21$, $371$, $\approx3\times10^{4}$, $\approx10^{9}$ between forks at periods $2,4,8,16,32$. Measured on the seed (ledger 10.235d~B, D): the Hamming distance between consecutive labels over one period is binomial$(P,\tfrac12)$ to statistical error on all $1{,}420{,}788{,}969$ pairs from $90\,000$ to the fork ($\chi^{2}=29.7$ on $31$ classes, every class within two standard deviations, the one fork being the distance-$0$ count against $0.33$ expected) and on $10^{10}$ pairs past the doubling at period $64$ on branch \texttt{6cbda353} (the seed's only if the seed takes that branch), down to distance $8$ and up to distance $56$ on the two tails ($\chi^{2}=51.7$ on $47$ classes) (\MEASURED). A fork is the distance-$0$ class of this law. \emph{If} the law holds down to distance $0$ --- the single fork at period $32$ is the only direct evidence --- then forks come once per $2^{P}$ diagonals, $1.8\times10^{19}$ at period $64$, and, on branch \texttt{6cbda353}, the next doubling of the edge lies near $10^{19}$--$10^{20}$ (conditional; ledger 10.242d~A11). The fair-coin model behind this and the $2^{-P}$ rate are \KNOWN\ (\cite[\S8]{GG}; \cite{Pat}, dated 19 August 2026; \cite{Dib}, 8 September 2026). What is measured here is the distance histogram at periods $32$ and $64$. The law is the wall's own pseudo-randomness statement, of the same kind as P2 for the column; what is proved is the mechanism and Proposition~\ref{fr:gaps}.
\end{remark}

\begin{conjecture}[M; \textsc{refuted} in print, see the remark following]\label{fr:conjM}
For every $P$ the left machine's background follows the shortest maximal path of $F_P$; equivalently, the background of the left half of Rule~30 is unique up to time translation for every boundary and every finite initial state. Verified through $d=12\,000$ as above, and at the branch $d=53\,208$ for $58$ initial words (Conjecture~\ref{fr:universal}).
\end{conjecture}

\begin{remark}[Conjecture~M fails]\label{fr:conjM:refuted}
Two published numbers refute it. Wolfram \cite[p.~871]{NKS} records the doubling depths
$3,8,29,400,87\,867$ and then $64$ at ``$2{,}107{,}985{,}255$ or more'' (read as a fork in \cite{WN}; the fork and the onset on one branch were public earlier, in \cite{Dib,Pat}); Gravner and Griffeath
\cite[\S8]{GG}, continuing the sixteen period-$32$ succession lines, report that ``the earliest stopping
event was a death at generation $65{,}154{,}361$''. The seed's line is one of those sixteen; it reaches the fork at $1{,}420{,}878{,}968$ still at period $32$,
and on branch \texttt{6cbda353} its period-$32$ regime ends exactly at $2{,}107{,}985{,}254$, after $2.1\times10^{9}$
diagonals, with no further frozen diagonal before $10^{12}$, while on branch \texttt{93425cac} the next frozen diagonal, $3{,}340{,}408{,}059$, is a second fork and on no continuation does the regime end before
$9{,}958{,}700{,}504$ (\EXACT, ledger 10.160c2, 10.232d~B, 10.242d~A11, 10.243d; the two forks and the statement for branch \texttt{6cbda353} are also in \cite{Dib,Pat}, and the onset $9{,}958{,}700{,}505$ on the other side is in \cite{Pat}; the range to $10^{12}$ is in neither); either way it is not the shortest maximal path of $F_{32}$: the coincidence
observed at $P=16$, where the seed exits along the shortest branch, does not persist. What survives is
the weaker and now fully attributed statement of Proposition~\ref{fr:universalwall}.
\end{remark}


\begin{conjecture}[left-half universality; \textsc{refuted} by a finite starting word, see the end of the statement]\label{fr:universal}
For every finite initial word, the eventual cycles of the left diagonals, the frozen diagonals and the doubling positions coincide with the seed's, and the left half of the evolution is the seed's left half up to a time delay. Verified for $58$ words (lengths $1$--$40$, $38$ of them random) through the branch at $d=53\,208$; Rowland's ``other possibility'' at that branch is not realised by any of them. The author of a public gist of 21 August 2026 \cite{Lania} reports the same for more than $2^{17}$ random starting states. The gist gives a constructed starting state of $53\,209$ bits that, its author reports, does realise the other possibility (the gist numbers the diagonal $53\,209$). We checked that state (ledger 10.235). Evolved under Rule~30 from a single row, its left diagonals agree with the seed's up to one translation of rows at every diagonal from $52\,800$ to $53\,199$, both have the frozen diagonal $53\,207$, and from $53\,208$ on the labels differ, with no translation of rows matching them on $[53\,300,53\,500]$ (\EXACT). So the conjecture is false as stated: a finite starting word can take the other branch at the first genuine fork. It holds for the $58$ words above and the measured boundaries; which finite words take the seed's branch is \OPEN.
\end{conjecture}

\section{Aperiodic traces inside the ordered region}\label{fr:walltrace}

Kopra \cite{Kopra} places the two-column theorem in a general frame: a cellular automaton is
\emph{left expansive} with dimensions $(h,d,w)$ if the contents of a rectangle of those dimensions
determine the cell to its left, and \emph{rapidly} left expansive if in addition it is left spreading
with speed $s$ and $s<1/h$. An $(m,n)$ left-permutive rule is left expansive with dimensions
$(0,1,m+n)$, and an elementary rule is left spreading exactly when $001\mapsto1$, with speed $1$.
Rule~30 is therefore rapidly left expansive with $w=2$, which gives the two-column theorem, and the
reason the frame stops there is visible in the inequality: $s=1$ forces $h=0$, and at $h=0$ the width
is $m+n=2$. No choice of parameters reaches $w=1$, which is Problem~1.

The gap can be stated more sharply than that, and doing so isolates what Problem~1 asserts. The
width-one trace is always a factor of the width-two trace, by forgetting the second column. For the
$3/2$ automaton --- Kopra's motivating example, and the one other system where a width-one trace has
been proved aperiodic --- that factor map is a \emph{conjugacy}, because there one column determines the
column to its left \cite[Prop.~2.7]{Kopra2}, so all trace widths carry the same information and
aperiodicity descends from width two to width one at no cost. For Rule~30 the same map is a proper
factor map. Problem~1 is exactly the assertion that this particular proper factor map does not collapse
periodicity. Whether Rule~30 has the property that makes the other case easy is a question we can
settle negatively: there is no local map of any support radius from $4$ to $10$ recovering the column to
the left of a single column, at any depth up to $20$, so no inverse of radius $\le10$ exists (\EXACT); a conjugacy would need one of some finite radius, and
this route to Problem~1 is closed as far as computed.

Width~$1$ is not hopeless in every direction, though. Our diagonals are width-one traces along slope
$1$ and they \emph{are} eventually periodic. So eventual periodicity of a width-one trace depends on
its slope, and it is natural to ask where the transition is. For a trace along slope $a$, meaning the
cells $x=\lfloor-at\rfloor+\text{const}$, the cell at time $t$ lies on diagonal $d=(1-a)t+\text{const}$,
so slope $1$ is exactly the case that stays on one diagonal and every $a<1$ drifts across infinitely
many. The ordered region lies between the light-cone edge, which advances at speed $1$, and the front,
which advances at about $0.2427$, so a trace with $0.2427<a<1$ eventually lies \emph{inside} the locked
wall.

\begin{measured}
Traces along slope $a$ from offset $400$, tested to $2\times10^{5}$ rows for any period up to $4096$:
periodic only at $a=1$ (period $16$, preperiod $498$); no period at
$a=\tfrac34,\tfrac12,\tfrac13,\tfrac3{10},\tfrac14,\tfrac15,\tfrac16,\tfrac18,0$. No critical slope at
the front speed is visible: slopes well inside the ordered region show no period up to $4096$ in $2\times10^{5}$
rows (\MEASURED).
\end{measured}

The slope-$\tfrac12$ case deserves to be stated on its own, because it relocates the difficulty of
Problem~1.

\begin{proposition}[a coin-like trace made of background, with no period up to $20\,000$; \EXACT]\label{fr:wallcoin}
Take the width-one trace of slope $\tfrac12$ from offset $400$. Of its first $1.5\times10^{5}$ cells,
$99.01\%$ are locked, that is, they lie at or beyond the preperiod of the diagonal carrying them, and
every diagonal it meets in that range has period exactly $16$. On the locked part the trace is
reproduced by the background alone, $r_d(t)=R_d[(t-\tau_d)\bmod16]$, in $118\,049$ cells with no
disagreement. Nevertheless it has no period up to $20\,000$ over a window of $1.14\times10^{5}$ rows,
and its density of ones is $0.50063$.
\end{proposition}

So the ordered region is not the simple part of the picture. It contains a width-one trace that shows
no period up to $20\,000$ and is coin-like to every test we can apply, although every diagonal it is built from has period
$16$ and the whole background is, by the results above, seven branch bits together with a deterministic
map. The aperiodicity cannot come from the chaotic interior, since the interior contributes $1\%$ of the
cells and none after the first few thousand rows; it comes from the phase relationship \emph{between}
consecutive diagonals.

That relationship is exactly what the backward orbit controls, which is what makes the following a
well-posed problem rather than a restatement of Problem~1. If the slope-$\tfrac12$ trace had period
$p=2m$, then for all large $t$
\[
  R_{\,d(t)+m}\bigl[(t+2m)\bmod16\bigr]=R_{\,d(t)}\bigl[t\bmod16\bigr],
  \qquad d(t)=400+\lceil t/2\rceil ,
\]
a condition on the sequence $d\mapsto R_d$ alone. The full patterns never coincide: at $m=1000$ there
are $0$ coincidences $R_{d+m}=R_d$ among $58\,433$ diagonals, as Theorem~\ref{fr:unbounded} requires,
since a repeated pattern pair cannot occur. But the trace needs agreement at only one phase per
diagonal, and that behaves like a coin, $2486$ agreements in $5000$ at $m=1000$.

\begin{problem}\label{fr:wallproblem}
Prove that the slope-$\tfrac12$ trace of Rule~30 from the single seed is not eventually periodic.
\end{problem}

Unlike Problem~1 this asks for aperiodicity of a sequence every term of which is given by an explicit
finite description, and the obstruction is a single-phase coincidence condition on an orbit we can
enumerate. We do not know whether it is easier; we record it because it is the only width-one
aperiodicity question about Rule~30 we have found in which all of the data is explicit. The background column $b_B$ of the next remark is its slope-one analogue, explicit in the same sense, with no period $\le10^{7}$ found on the last $10^{8}$ of its $10^{9}$ terms (\EXACT; ledger 10.233d~I).

\begin{remark}[the background column to $10^{9}$; public earlier in \cite{Dib}; the column \EXACT, the tests \MEASURED]\label{fr:frozencolumn}
The whole background, read on the diagonal, is the column $b_B(d)=r_d(d)$: the centre column the stripes would make if the wedge reached it. This object is \KNOWN: the background as a Rule~30 evolution of its own, and its centre column, are in \cite{Dib} (commits \texttt{0a1cda2}, \texttt{47f8ab5}, \texttt{b7b9195} and \texttt{980d553}, 7--8 September 2026), to $10^{9}$ terms and in Lean. The label recursion of Lemma~\ref{fr:backward}, run forward in C, writes it directly, $10^{9}$ bits in $35$ seconds, exact to the fork at $1{,}420{,}878{,}968$, on branch \texttt{6cbda353} exact to the doubling at $2{,}107{,}985{,}254$ (both also in \cite{Dib,Pat}), and known up to a global complement beyond it (ledger 10.233d~I; it agrees with the frame's column on all $189\,999$ diagonals $90\,001$--$279\,999$). On $d=90\,001,\dots,10^{9}$: density $0.500029$, every $10^{8}$-block within $0.4999$--$0.5001$; no period $p\le10^{7}$ on the last $10^{8}$ bits; run lengths within $0.1\%$ of $N/2^{L}$ for $L\le12$ and within $1.5\%$ to $L=18$; blocks of $8$, $12$ and $16$ bits each within one standard deviation of a fair coin; the best guess of the next bit from the last $k\le16$ bits scores what a coin scores in sample; autocorrelation within $2.1$ standard deviations at $26$ lags with one exception, lag $64$ at $-6.2$ standard deviations, a trace of the period-$32$ stripes two periods apart (\MEASURED; ledger 10.233d~I as corrected by 10.242d~A3). Its agreement with the real centre column on rows $90\,001$--$1\,130\,000$ is $0.49973$, chance (ledger 10.241d~\#9). A column with no chaos in its construction that every coin test but one (lag $64$) fails to distinguish from a coin, and whose eventual periodicity is the same kind of question as Problem~1, with the same answer so far.
\end{remark}

\part{The column as an object}

\noindent\emph{Source: \texttt{collapse\_section.tex} (2026-09-16), labels prefixed \texttt{co:}; Section~\ref{co:split} is new.}

% Part IV, from collapse_section.tex (labels prefixed co:)
\section{Collapse pyramids over the centre column}

Let $c=(c_0,c_1,\dots)$ be the centre column of Rule~30 from the single black cell. For an elementary
rule $R$ and $N\ge 0$ write the first $2N+1$ bits as a row and apply $R$ upward, each new cell being
$R$ of the three cells beneath it and no assumption being made outside the row. The row loses one cell
at each end per step and closes after $N$ steps to a single cell, the \emph{apex}
\[
  a^R_N \;=\; \bigl(R^N(c_0c_1\cdots c_{2N})\bigr)_0 .
\]
Equivalently $a^R_N$ is the leftmost cell of $R^N$ applied to the whole column; it is a function of
$c_0,\dots,c_{2N}$ only. The apex sequence $(a^R_N)_{N\ge 0}$ is thus a derived sequence of the column,
one bit per two column bits. This section records what these sequences are and what they can and cannot
say about $c$.

\subsection{Inheritance}

\begin{proposition}[downstream only; \PROVED]\label{co:prop:inherit}
If $c$ is eventually periodic then $(a^R_N)$ is eventually periodic for every $R$. \PROVED
\end{proposition}
\begin{proof}
Let $c$ have period $p$ from index $t_0$, and read each row from its left edge: $X_m[j]:=(R^m(c))_{m+j}$, so that
$X_{m+1}[j]=R(X_m[j],X_m[j+1],X_m[j+2])$ and the apex is $a^R_N=X_N[0]$. If $X_m$ is $p$-periodic from relative position
$t_0$ then so is $X_{m+1}$, and $X_0=c$ is; hence every $X_m$ is determined by its first $t_0+p$ cells, a finite set of
states, so $m\mapsto X_m$ is eventually periodic and so is $a^R_N=X_N[0]$. (An earlier version read the apex off the
periodic part of $R^N(c)$; the apex lies in the transient block, which is why the row must be read from its edge.)
\end{proof}

Hence no apex sequence can be used to prove periodicity of $c$, and any apex sequence that is
provably not eventually periodic proves Problem~1. The converse inheritance fails: for the identity rule
$a_N=c_N$, while for the constant rules $a_N$ is constant.

\subsection{Two exact apexes: the traffic rules}

For a row $r=(r_0,\dots,r_{n-1})$ let $H(k)=\sum_{j<k}(2r_j-1)$, $0\le k\le n$, be its balance walk, so
$H(k)\equiv k \pmod 2$ and $r_k = [H(k+1)>H(k)]$. Rule~184 is $r'_k = (r_{k-1}\wedge \neg r_k)\vee(r_k\wedge r_{k+1})$.

\begin{lemma}[height form of Rule 184; \PROVED]\label{co:lem:height}
Let $r'=R_{184}(r)$ be the interior row $(r'_1,\dots,r'_{n-2})$ and let
$H'(k)=\sum_{j=1}^{k}(2r'_j-1)$, $0\le k\le n-2$, be its balance walk. Then
\[
  H'(k) \;=\; \max\bigl(H(k),H(k+2)\bigr) - C, \qquad C=\max\bigl(H(0),H(2)\bigr)\in\{0,2\},
\]
for $0\le k\le n-2$. \PROVED
\end{lemma}
\begin{proof}
Both sides vanish at $k=0$ and both are walks with steps $\pm1$, so it suffices to compare increments.
Write $G(k)=\max(H(k),H(k+2))$. Since $H(k+2)=H(k)+(2r_k-1)+(2r_{k+1}-1)$, the value of
$G(k+1)-G(k)$ is determined by $(r_k,r_{k+1},r_{k+2})$, and checking the eight cases gives
$G(k+1)-G(k)=+1$ exactly for $(r_k,r_{k+1},r_{k+2})\in\{100,\,101,\,011,\,111\}$, which is exactly the set
on which Rule~184 outputs $1$, i.e. $r'_{k+1}=1$. Hence $G(k+1)-G(k)=2r'_{k+1}-1=H'(k+1)-H'(k)$.
(Checked exactly on random rows of length $2001$.)
\end{proof}

\begin{theorem}[record parity; \PROVED]\label{co:thm:record}
Let $S_k=\sum_{j<k}(2c_j-1)$ be the balance walk of the centre column and
$M_N=\max_{0\le k\le 2N+1}S_k$, $m_N=\min_{0\le k\le 2N+1}S_k$. Then
\[
  a^{184}_N \;=\; M_N \bmod 2, \qquad a^{226}_N \;=\; 1-\bigl(|m_N| \bmod 2\bigr). \quad\PROVED
\]
\end{theorem}
\begin{proof}
Iterating Lemma~\ref{co:lem:height}, after $N$ steps $H_N(k)=\max_{0\le j\le N}H(k+2j)-C_N$ for a constant $C_N$. The apex is the
single remaining cell, whose value is $[H_N(1)>H_N(0)]$, i.e. whether the maximum of $H$ over odd
positions in $[0,2N+1]$ exceeds the maximum over even positions. Since $H(k)\equiv k\pmod 2$ the two maxima
differ by exactly $1$, and the odd one is larger exactly when the overall maximum $M_N$ is odd.
Rule~226 is the complement conjugate of Rule~184, $R_{226}(x)=1-R_{184}(1-x)$ (equivalently its mirror image); complementing
the row sends $H$ to $-H$ and the maximum to the minimum.
\end{proof}

\begin{remark}[provenance]
Rule~184 is the ultradiscrete Burgers equation \cite{NT}, which is the theorem-bearing form of the
statement that follows. The height-function (max-plus) representation of Rule~184 is standard in the traffic and ballistic
annihilation literature (Belitsky--Ferrari \cite{BF}; Krug--Spohn \cite{KS}). Theorem~\ref{co:thm:record} is its immediate
corollary for the collapse geometry and is recorded here for the criterion below, not as a new result.
\end{remark}

\subsection{A two-sided criterion}

\begin{corollary}[two-sided criterion; \PROVED]\label{co:cor:twosided}
If both apex sequences $(a^{184}_N)$ and $(a^{226}_N)$ change value infinitely often, then $c$ is not
eventually periodic.
\end{corollary}
\begin{proof}
$a^{184}_N$ changes only when $M_N$ increases, $a^{226}_N$ only when $m_N$ decreases. If $c$ is
eventually periodic with period $p$ and $\delta$ the net balance over one period, then for $\delta=0$ the
walk is bounded and both records are eventually constant, and for $\delta\ne 0$ the walk drifts, so one
record is eventually constant.
\end{proof}

The hypothesis is that the balance walk of $c$ is unbounded above and below; it is neither implied by nor implies
Problem~2 (a bounded walk has density $\tfrac12$; a one-sided drift of order $\sqrt k$ has density $\tfrac12$), and it is not implied by Problem~1. \EXACT to $k=200{,}001$: the walk's
record high is $419$, first reached at $k=24{,}989$ (last new high at $N=12{,}494$), and its record low
$-257$ at $k=172{,}735$ ($N=86{,}367$); the Rule~184 apex has changed $196$ times, the Rule~226 apex
$127$ times. Whether either changes again is \OPEN, and a long freeze is ordinary behaviour at the
$\sqrt{k}$ scale.

\begin{proposition}[why the predictor measurements cannot be upgraded; \PROVED]\label{co:windowbound}
Suppose the word frequencies $d(u)=\lim_N N^{-1}\#\{i<N:c_{i}\cdots c_{i+|u|-1}=u\}$ exist, and for
$w\ge1$ let $A_w=\sum_{|u|=w}\max\bigl(d(u0),d(u1)\bigr)$ be the limiting accuracy of the best
window-$w$ predictor. Then
\[
A_w=\tfrac12+\tfrac12\sum_{|u|=w}\bigl|d(u0)-d(u1)\bigr| ,
\]
and consequently $A_w=\tfrac12$ for a \emph{single} $w\ge1$ already forces $d(1)=\tfrac12$, which is P2.
\end{proposition}
\begin{proof}
$\max(a,b)=\tfrac12(a+b)+\tfrac12|a-b|$ and $\sum_{|u|=w}(d(u0)+d(u1))=1$ give the identity. If the sum
vanishes then $d(u0)=d(u1)=\tfrac12 d(u)$ for every $u$, so
$d(1)=\sum_{|u|=w}d(u1)=\tfrac12\sum_{|u|=w}d(u)=\tfrac12$. The identity was checked against the direct
computation of $A_w$ on $2^{18}$ terms for $w=1,\dots,14$, agreeing to $10^{-12}$ at every $w$.
\end{proof}

\begin{remark}[what this says about the rest of this Part]\label{co:windowbound:read}
The measurements of this Part say that no bounded-window predictor of the centre column beats a coin,
and we have been careful to label them \MEASURED. Proposition~\ref{co:windowbound} explains why they
must stay that way: proving the statement in its natural limiting form, at even one window length,
would prove P2. So the gap between our measurements and a theorem here is not a gap in effort or in
sample size --- it is exactly Problem~2, and any programme that proposes to close it by computing
harder is mistaken about what it is computing. Across all $w$ the same statement is base-$2$
normality.

This also fixes the status of the predictor question in the other direction. The learning-theoretic
machinery that looks as though it should apply does not: mistake bounds and the Littlestone dimension
are min-max statements quantified \emph{over} sequences --- the dimension of the window-$w$ class is
$2^{w}$, whatever sequence is presented --- and regret bounds are relative to a comparator class, so
neither says anything about one fixed sequence. The one branch of that literature which does face a
single sequence, the finite-state predictability of Feder, Merhav and Gutman \cite{FMG}, delivers a universal
algorithm and never a lower bound, a gap those authors state themselves.
\end{remark}

\subsection{What the pyramids cannot do}

Every apex, and every cell of every collapse pyramid, is a function of a bounded window of the column.
A single computation therefore bounds the whole family. Let $P_w$ be the majority predictor of $c_n$ from
$(c_{n-w},\dots,c_{n-1})$ fitted on the first $2^{19}$ bits and evaluated on the second $2^{19}$.
\MEASURED\ (out of sample; \EXACT\ as counts on that split): its accuracy lies in $0.4973$--$0.5011$ for every $w\le 22$ (standard error $\le0.002$), with conditional empirical entropy
$1.000$ bit per symbol through $w=14$. Consequently:
\begin{itemize}
\item no rule $R$, height $k$ and offset $d$ gives $a^R$ or any edge cell of the pyramid a correlation with
$c_{2N+1+d}$ beyond the chance maximum of the number of tests ($256$ rules, $40$ heights, $9$ offsets:
best in-sample $3.2\sigma$; the right-edge family: best $4.6\sigma$ in-sample, $-0.1\sigma$ out of sample);
\item the pyramid's own centre column matches $c$ at no position and shift beyond chance ($6.5\times10^{7}$
tests, best $5.2\sigma$, expected $5.3\sigma$);
\item rules applied to the locked left wall at time $t$ correlate with the causally reachable centre bits
$c_{\lfloor 3t/4\rfloor+d}$ at the chance level out of sample.
\end{itemize}
These are measurements, not theorems; the theorem-level statement is Proposition~\ref{co:prop:inherit}:
the collapse pyramids are downstream of the column, and information does not flow back.


\section{Where the information enters the column}\label{co:split}

\begin{proposition}[parent-instance split; \PROVED]\label{co:parents}
Let $\ell_t=u^t_{-1}$ and $u_t=u^t_1$ be the columns adjacent to the centre. Then:
\begin{enumerate}[nosep]
\item[(a)] $\ell_t$ is a function of $c_0,\dots,c_{t-1}$ alone: the half-plane $x\le-1$ evolves by Rule~30 with column~$0$ as its boundary input.
\item[(b)] If $c_t=1$ then $c_{t+1}=1+\ell_t$; if $c_t=0$ then $c_{t+1}=\ell_t+u_t$.
\item[(c)] $u_{t}=1+c_{t-1}$ exactly when $u^{t-1}_1\vee u^{t-1}_2=1$.
\end{enumerate}
\end{proposition}
\begin{proof}
(a) A cell at $x\le-1$ and time $t$ depends on cells at $x-1,x,x+1\le0$ at time $t-1$; induct. (b) is the rule at $x=0$. (c) is the rule at $x=1$, where $c_{t-1}$ enters by XOR.
\end{proof}

So after every black centre cell the next bit is determined by the column's own past, through an unbounded-memory computation; after every white centre cell one bit enters from the right, and that bit is the complement of the previous centre bit three times in four (\MEASURED: $0.7499$ on $16\,384$ rows). Parent instances of $c_{t+1}$ are equidistributed over the eight triples (\EXACT, counts $1953$--$2107$ of $16\,383$). A bounded-context model on the eight-colour parent sequence predicts the next colour at $0.37$--$0.44$ against $0.125$, entirely from (b) and (c), and the next centre bit at $0.48$--$0.50$ (\MEASURED, train/test halves). Combined with Theorem~\ref{ex:thm:zero}, P1 is equivalent to: \emph{if $c$ were eventually periodic, the bits $u_t$ at the white times $t$ would have to be eventually periodic}. Part~I settles this family by family; the general case is \OPEN.

\part{Computational record, and Problem~3}

\section{Computational record}\label{sec:record}

Each item names its script in the author's working folder (not all of them are public); the ledger section is in brackets.

\subsection*{The two-sided kill instrument (\EXACT, then \MEASURED)}
Fix the tail $w=01$. Let a survivor at $(N,T)$ be a configuration with $a(0,i)=0$ for $i<-N$ whose centre follows $w$ on rows $0..T-1$; $T_{\mathrm{death}}(N)$ is the first $T$ with no survivor and $f(N)=T_{\mathrm{death}}(N)-N$. \PROVED\ (one line): $f$ bounded implies $w$ is not a tail of $c$. \EXACT: a survivor at $(N,T)$ is an admissible interface whose reconstructed row~$0$ vanishes on $(N,T]$; minimal unsatisfiable cores use nearly every centre row but only the $7$--$8$ deepest vacuum cells; the kill depth $m(T)$ (the least number of consecutive forced zeros that no admissible interface can carry at depth $T$) was computed to $T=120$. \MEASURED: $m(T)\approx1.7$--$2\log_2T$; the record $f=18$ from $N=19$ to $120$ broke to $f=20$ at $N=130,140$. Locality test: forcing the centre only on rows $[t_0,T_{\mathrm{death}})$ keeps the instance impossible only for $t_0\le6$. The bounded-$f$ route is not supported by the data: $f$ has already increased once and no finite record can establish boundedness (\MEASURED; whether $f$ is bounded is \OPEN), and the kill is global [10.47--10.50].

\subsection*{The left machine alone, for the tail $01$ (\EXACT; ledger 10.96--10.106)}
At a black centre cell the OR is already $1$, so a $01$ tail forces $u^t_{-1}=1$ at every odd $t$: a condition on the left machine only. With column $0$ alternating from $t=0$, column $-1$ forced at odd times and free at even times, left-permutivity reconstructs the left half, and number-likeness with support $N$ requires the row-$0$ cells $P_k=u^0_{-k}$ to vanish for $k>N$. The odd equations $P_{2m+1}=0$ fix the later free bits (each free bit enters its anti-diagonal by XOR), the even ones are checks, and the candidates are the $2^{\lfloor(N-1)/2\rfloor+1}$ initial free-bit strings. Exhaustive enumeration (validated against SAT at $N\le30$) gives the first depth with no survivor, $T_{\mathrm{death}}(N)=8,18,30,32,38,44,54,66,78,88,98$ for $N=4,8,12,16,20,24,30,32,40,44,50$ (all even $N\le50$ computed; full table in ledger 10.106). The death depth is not an independent law but the unicity distance implied by the halving: against $N+2n_{\mathrm{free}}$ with $n_{\mathrm{free}}=\lfloor(N-1)/2\rfloor+1$ the residuals lie in $[-8,+2]$ with mean $-3.7$ (over $16\le N\le50$; over all even $N\le50$ they lie in $[-8,+6]$ with mean $-2.5$) and vanish at $N=16$ and $N=44$, so quoting a ratio $T_{\mathrm{death}}/N$ as if it were a constant overstates the content (an outside proposal $T_{\mathrm{death}}=N\,p/(p-z)$, exactly $2N$ for $01$, is refuted: the ratio is not constant). Survivors halve at every check (the checks are exact coins conditional on all earlier ones, $m\le22$), and the longest-lived candidates form classes riding a stripe zone in which column $-1$ runs in phase with the centre, the $\ldots0101\ldots$ fixed point of Rule~30. Every such champion at $N\ge8$ is right-inadmissible within $18$ interface symbols, most by the forbidden factor $11$: the right half's language removes exactly the striped candidates, which is why the two-sided instrument above dies with delay $\sim\log N$ while the left machine alone needs $\sim N$. We do not claim that the surviving family is polynomially small: the count of right-admissible completions has the same non-polynomial signature in the ensemble setting of \cite[Table~1]{MS} as $|L_k|$ does here. \PROVED\ alongside: $T_{\mathrm{death}}(N)\ge T_{\mathrm{death}}(N-2)$; an eventually periodic left output is never number-like; survival is invariant under the two-row shift; the row-$0$ polynomials satisfy $P_{k+2}=\sigma(P_k)+(P_{k+1}\vee P_k)+(S_k\vee S_{k-1})$ with $\sigma$ the index shift and $S_j=P_{j+1}+(P_j\vee P_{j-1})$, an exact recursion whose state is the anti-diagonal pair; every check has odd weight. Two ingredients of this instrument are \KNOWN\ and we claim neither: the forcing at a black centre cell, that $u^t_{-1}$ is the complement of the next centre bit, is stated one paragraph after equation~(4) of \cite{MS}, repeated in \cite[p.~1087]{NKS} and proved as Proposition~3.4 of \cite{Spencer}; and the count of free bits, about $N/2$ for a balanced word, is Spencer's bound of $n-j-1$ coins for a word with $j$ ones \cite[\S3.2]{Spencer}, also given as $2^{n/2}$ in \cite{Leon}. What is new here is the use of the row-$0$ equations as checks, the exhaustion of the candidate set and the measurement of a death depth; Meier and Staffelbach explicitly turned away from the left branch, since for key recovery the left restrictions are an obstacle rather than a tool, and Spencer proposed the left programme without developing it. The target, no candidate surviving forever, is \OPEN. The obstruction is that every check is a coin conditional on everything computable, and it is sharpened by a reformulation that is already in print: left-permutivity makes the map from the row-$0$ cells at $x=-1,\dots,-T$ to the centre bits at $t=1,\dots,T$ a bijection, which is in substance Meier and Staffelbach's completion backwards \cite[eq.~(4)]{MS} (\KNOWN; the two-column triangle reconstruction is already in \cite{W85}, credited there to C.~and R.~Feynman), and of which \cite[Lemma~1]{Condrey} is the injectivity half. The reconstruction is therefore unique, there is no ensemble over which to equidistribute, and the question is properly: for the single string that bijection returns, is the density of ones bounded below? That is why no measure-level theorem can reach it.

\subsection*{The fibre over the alternating column (\EXACT)}
$|L_k(01)|$ for $k=1..44$: $2$, $3$, $5$, $8$, $12$, $17$, $25$, $36$, $50$, $68$, $91$, $119$, $156$, $199$, $251$, $316$, $393$, $487$, $596$, $721$, $875$, $1054$, $1255$, $1493$, $1780$, $2111$, $2483$, $2904$, $3378$, $3908$, $4502$, $5153$, $5875$, $6664$, $7541$, $8534$, $9649$, $10876$, $12231$, $13730$, $15384$, $17196$, $19197$, $21389$, extended by an independent encoding to $|L_{45..60}|=23809$, $26482$, $29422$, $32656$, $36193$, $40047$, $44280$, $48878$, $53894$, $59333$, $65240$, $71672$, $78632$, $86169$, $94361$, $103220$; the local exponent rises without a bend from $4.3$ at $k=30$ through $5$ at $k=50$ to $5.34$ at $k=60$, its slope in $\ln k$ itself rising. Out of sample on $k=45..60$, fitted on $k\in[K_0,44]$ and with $K_0$ swept over $8,12,16,20,24$ so that no conclusion rests on one window, and with every family given the same number of parameters, exactly two families are excluded and they are excluded by orders of magnitude: a free-degree polynomial $Ck^{d}$ misses by $21.8$--$43\%$ and the pure exponential $C\lambda^{k}$ by $84$--$402\%$. Nothing else is separated. The quasi-polynomial $\exp(c\ln^2k)$ and the stretched exponentials $k^{a}\exp(ck^{\gamma})$ for every $\gamma$ tried in $[0.25,0.70]$ all sit between $1.48\%$ and $12.90\%$, and their \emph{ranking} reverses with the window --- $\gamma=0.25$ is the best of them at $K_0=8$, $\gamma=0.50$ at $K_0=12$, $\gamma=0.70$ at $K_0=16,20,24$ --- so an earlier phrasing that put the quasi-polynomial ``at errors of $5\%$ or more'' and left ``only a stretched exponential'' surviving is withdrawn: it is true at $K_0=12$ ($5.13\%$) and false at $K_0=8$, where the quasi-polynomial at $2.63\%$ beats $\gamma=0.50$ at $6.70\%$. Note also that the excluded pure exponential is excluded \emph{as a fit}, which says nothing about $h$: the proved bounds are upper bounds --- $h\le0.278$ from the counts (Proposition~\ref{ex:prop:fekete}) and the sharper block-automaton bound $h\le0.132$ quoted below --- and $h>0$ is not excluded. The exponent is not determined, and
we withdraw three stronger phrasings used earlier. It is not true that every pure $\exp(ck^\gamma)$ with
$\gamma\ge0.3$ is excluded: $\exp(5.91k^{0.28})$ predicts $k=46..60$ to $1.31\%$ with one parameter
\emph{fewer} than $k^{1.9}\exp(0.87\sqrt k)$ at $1.38\%$, and beats it on the longer window, $0.60\%$
against $1.08\%$. Nor, and this is the third withdrawal, does that comparison exclude $\gamma\ge0.35$,
because it is exactly a two-parameter form beating a three-parameter one; see the end of this remark.
Nor is $k^{1.9}\exp(0.87\sqrt k)$ stable across
fitting windows: moving the left endpoint gives $(1.10,1.10)$ at $0.45\%$ and $(0.98,1.14)$ at
$0.25\%$, so the fit lives on a degenerate ridge and no value of $c$ is pinned. What matters is that
$\gamma$ is exactly the quantity that decides between mechanisms --- a partition mechanism requires
$\gamma=\tfrac12$ exactly, a hierarchical one requires $\log_\ell r=\tfrac12$ --- and by Cassaigne's
Theorem~3, which realises $\log p(n)\sim n^\alpha$ for \emph{every} $\alpha\in(0,1)$ with a uniformly
recurrent word, no theorem will pin it for us. So $\gamma$ is \OPEN\ --- and, correcting a narrower reading of $[0.25,0.35)$ given earlier, with no interval pinned at all. That reading set a two-parameter pure $\exp(ck^\gamma)$ against a three-parameter $k^{a}\exp(ck^\gamma)$, and the exclusion does not survive equalising the count: at three parameters throughout, $k^{1.91}\exp(0.866\sqrt k)$ predicts $46..60$ from $12..45$ to $1.19\%$, \emph{better} than the pure $\gamma=0.28$ form's $1.60\%$ on the same window, and the held-out optimum over a grid $\gamma\in[0.15,0.85]$ drifts monotonically with the fit window, taking the values $0.150,0.150,0.250,0.525,0.700,0.750$ as the window ends at $k=34,38,42,46,50,54$, while the in-sample optimum on sliding windows of $20$ visits both ends of that grid. An estimate that sweeps its whole grid as the window moves is not an estimate, so what the data supports is superpolynomial growth with $\gamma$ unconstrained, and the exclusions that follow are the ones that do not go through $\gamma$. A cleaner exclusion is available than the follower-set census, which exhibited $3\,591$ distinct
follower sets: a monomial algebra whose basis language is \emph{regular} has growth either polynomial
or exponential, and no intermediate growth is possible (Ufnarovskii; stated in this form by
\cite{BM10}). So intermediate growth alone excludes regularity, hence soficity, with no bound on
follower sets needed. One scope note, since the graph form of the same statement assumes a finite
forbidden set and our language is not of finite type: it is the automata-algebra form, which needs only
regularity, that applies. The inference is conditional on the growth being intermediate, which is
\MEASURED\ on $k\le60$ only; the stream of Remark~\ref{ex:rem:stream} puts the asymptotic growth near $2^{0.07k}$, and if that holds the exclusions of sofic, finite-type and context-free languages here lapse, while the follower-set bound $3\,591$ does not. In the same direction: a sofic language is regular, and the
factor-counting function of a regular language is exactly a polynomial times an exponential on each
congruence class for large $n$ (Berstel; see \cite[Thm.~V.3]{FS} and \cite{SYZS}), so
\emph{superpolynomial together with subexponential} growth already excludes sofic, of finite type and
group shift, with no bound on follower sets needed. The context-free case is stronger than we had it:
by Bridson and Gilman \cite[Thm.~1.1]{BG} a context-free language has growth eventually at least $r^n$
or else bounded, so superpolynomial growth excludes \emph{all} context-free languages with no
unambiguity hypothesis, which answers a question of Flajolet (see also \cite{Incitti}). Topological rank two is excluded as well
\cite[Cor.~6.6]{DDMP}. What is \emph{not} excluded is the indexed, ET0L and EDT0L levels, and that is
where the one published language with partition growth sits. One correction to how we have described
$L(01)$: it is an \emph{anchored} language, not the language of a subshift --- at $k=59$ some
$24.4\%$ of its words occur only at row $0$ --- so the absence of dead ends holds on the right only,
while factoriality and submultiplicativity were verified at all sixty depths and all $1770$ ordered pairs with $m+n\le60$.
These exclusions are therefore conditional exactly on the growth law, and on nothing else (and the growth law is now in question: positive entropy is \MEASURED\ on random right halves, Remark~\ref{ex:rem:stream}); block-automaton entropy bounds reach $h\le\log_2 1.0958=0.132$ bits (proved in the ledger, not reproduced in this paper) [10.16, 10.31, 10.54, 10.105, 10.117]. \PROVED\ by hand: constraint (A), no $11$, gaps at most $4$, factoriality, no dead ends. \PROVED\ finite, on the realisability of constant gap words by temporally periodic half-planes: $1^\infty$, $2^\infty$, $4^\infty$ and $(1\,4)^\infty$ are realised, at row periods $4,6,10$ and $14$ and spatial cycles $7,84,155,728$, each certified and re-verified by independent forward simulation, the period-$14$ cycle length $728$ being pinned between two independent searches (a direct sweep finds no cycle of length $\le700$, a wide solution scan finds the repeat at exactly $728$); $3^\infty$, $(1\,3)^\infty$, $(2\,3)^\infty$, $(2\,4)^\infty$, $(3\,4)^\infty$, $(2\,4\,4)^\infty$, $(1\,1\,3\,4)^\infty$ and $(2\,4^4)^\infty$ are not. Of the $3\,779$ primitive necklaces of width $\le24$, four are unbounded, $3\,767$ are proved bounded and eight are \OPEN, the sharpest being $(2\,4\,4\,4)^\infty$ at row period $36$. A conjectured row period $26$ is \textsc{refuted}. Zero entropy and the growth law remain \OPEN, and the literature now says three separate things about
them. First, our shape is realised: Cassaigne \cite[Thm.~1]{Cass02} exhibits an explicit binary word
with $2^{\lceil\sqrt n\rceil}\le p(n)\le n^2\,2^{\lceil\sqrt n\rceil}$, which is our form exactly,
$\gamma=\tfrac12$ with a positive polynomial prefactor, Mitchell \cite{Mitchell} gives a second such subshift with the polynomial prefactor explicit, and Greenfeld, Moreira and Zelmanov
\cite{GMZ} realise the entire subexponential range by subshifts, even strictly ergodic ones. So no
general theorem forces subexponential growth and none forbids our shape; in particular there is no
polynomial-or-exponential dichotomy to refute it.

Second, there \emph{is} a usable sufficient condition, contrary to what we said in earlier notes. By
Donoso, Durand, Maass and Petite \cite[Cor.~6.5]{DDMP} an $S$-adic subshift whose alphabet growth is
negligible against its word lengths --- in particular one of finite topological rank --- has zero
entropy, and finite rank does not bound complexity, so it is compatible with any subexponential
growth. That is a constructive target rather than a hope, and it is the route worth building. Two
cautions attach to it: word entropy is vacuous for subexponential functions, and the monotone decrease
of $\tfrac1k\log_2|L_k|$ is not evidence for the limit being zero, since $3\cdot2^k$ decreases
monotonically to $1$; the bound $h\le0.278$ from the counts stands as an infimum, but no computation at finite depth
can decide $h$.

Third, the mechanism. The partition heuristic for the $\exp(c\sqrt k)$ form is half right. Structurally
it holds: the ratio of $|L_k|$ to the number of domain-length multisets is $4.91$ at $k=60$ with
increments $+0.39,+0.16,+0.07$, admissible orderings are nearly forced, only $3.6\times10^{-6}$ of the
combinatorially available ones occurring, and the multiset is a genuine partition with $5.56$ distinct
part sizes per word growing like $0.72\sqrt k$ against the random-partition value $0.780\sqrt k$. These
parts are \emph{domain lengths} in the configuration and not the run lengths of the interface word; the
latter take at most the four values $\{1,2,3,4\}$ at every depth to $60$, so the two part sets in this
paper are different objects and neither figure bounds the other. The magnitudes confirm it: partitions
of $60$ into parts $\le4$ number exactly $1\,906$ against $|L_{60}|=103\,220$, a factor $54.2$, where the
domain-length ratio quoted just above is $4.91$. As an
explanation it fails, and the decisive reason is not the one we first gave. The fitted constant sits on
a degenerate ridge where every cell agrees to $1.3$--$1.7\%$ and sparse-part models at
$\gamma\approx0.25$ fit equally well, so ``$c=0.87$ against Meinardus' $2.607$'' proves little on its
own. What does refute it is the measured \emph{positive density} of the part set: every slot of type
other than $3$ is realised contiguously to the space cap, which forces $\gamma=\tfrac12$ \emph{together
with} a prefactor exponent at most $\tfrac14$, and it is the pair that fails. Fitting $k^{a}\exp(c\sqrt
k)$ on $k\in[K_0,44]$ and testing on $45..60$, with $a$ held at $\tfrac14$, misses by $13.5$--$46\%$ as
$K_0$ runs over $8,12,16,20,24$ (and by $16$--$55\%$ with $a=0$); releasing $a$ brings the error down to
$1.5$--$11.8\%$ but only by demanding $a=1.49,1.89,2.36,2.89,3.05$, an order of magnitude above the
$\tfrac14$ the density argument allows. So what is refuted is the \emph{conjunction}. The value
$\gamma=\tfrac12$ on its own is not refuted --- with a free prefactor it is the best-fitting exponent at
the longest windows --- and nothing in the withdrawal of the interval $[0.25,0.35)$ above revives
sorting, whose refutation runs through the bounded part set and the prefactor, not through $\gamma$. The
structural half also undercuts the explanatory half, since a correspondence bijective up to a bounded
factor cannot itself generate growth.

One line of reasoning we pursued has to be withdrawn. We argued that since free segmentation is never
intermediate, the exponent must be forced by the ordering constraint, and that identifying the
mechanism would pin $\gamma$. It will not, and the obstruction is a theorem rather than a gap. The
universal realisation constructions --- Bell and Zelmanov's, and the subshift realisations of
\cite{GMZ} --- all run on a single mechanism, a gap threshold tending to infinity, and the shape
$An^{a}\exp(cn^{\gamma})$ satisfies their conditions for \emph{every} $\gamma\in(0,1)$; our own measured
table satisfies them with room to spare, the largest ratio $|L_m|/|L_n|^2$ being $0.889$ over all $595$
pairs. So no theorem of the form ``exponent $\gamma$ requires mechanism $X$'' can exist. Nor are the
exponents quantised: Bergman's gap remains the only known gap in the space of growth functions, and
sorting realises every exponent too --- over a part set of density exponent $r$ it gives
$\gamma=r/(r+1)$ by Meinardus \cite{Meinardus}. The consequence for us is that $\gamma$ separates branching from
everything else and nothing finer; the quantity that would discriminate sorting is the \emph{density of
the part set}, not the exponent. The sorting mechanism does have a published witness --- Grigorchuk and
Mach\`i's \cite{GMachi} language $\{ab^{i_1}\cdots ab^{i_k}:i_1\le\cdots\le i_k\}$, whose growth is
exactly Euler's partition function --- and several classes remain excluded by our own measurements:
deterministic substitution, capped at $\Theta(n^2)$ by Pansiot \cite{Pansiot}, Toeplitz, and primitive
random substitution. Non-primitive random substitution is not on that list, and the density half of that route is dead in any case:
the frequency theory for random substitutions is primitive-only and almost-every-only, and a
non-primitive subshift carries uncountably many ergodic measures, so even an exact substitution would
not by itself give a density for our particular word. The mechanism is also partly settled: the switching law between domains is bounded-memory, a function of the outgoing type, the run length truncated at $8$ and the incoming type, with every threshold at most $7$, yet it is \emph{not} a function of type, phase and type alone ($1^2\,4$ is realised while $c\,1^2\,4$ is realised for no letter $c$), and $L(01)$ is not left-extendable ($100101\in L_6$ but $0100101\notin L_7$). The domain periods do \emph{not} grow with $k$; what grows is the number of domains needed under any fixed period cap. So neither branch of the earlier dichotomy applies, and since a regular language grows either polynomially or exponentially and never in between, the measured intermediate growth is inconsistent with any finite state space --- conditional on that measurement, which is why the growth law itself is still \OPEN.

\subsection*{Every bounded-window predictor of the column (\MEASURED)}
See Section~\ref{co:split} (``What the pyramids cannot do''), which carries the numbers: accuracy $0.4973$--$0.5011$ for every $w\le22$ on the held-out half, entropy $1.000$ bit through $w=14$, and the same for all $256$ rules applied to the column and to the locked left wall [10.84--10.87].

\subsection*{Other coin readings (\MEASURED)}
Density $0.4996$ to $8000$ rows and $0.5006$ to $2^{20}$; firing rate of the OR gate $0.750$; subword complexity $2^n$ to $n=16$; partial-sum Hurst exponent $0.506$; exponential sums coin-like at all rational frequencies tested; lag-$1$ autocorrelation along the diagonals through the centre $-0.504$ (the wire) and $0$ for the centre column itself; the Rule~30/Rule~86 overlay has four colours at density $\tfrac14$ each with an isotropic agreement field [10.73, 10.77, 10.78]. Smooth (multilinear, spin and wave) extensions of the rule collapse to $\tfrac12$ [10.70].
The single-column Hurst figure above is noisy, because fitting a power law to one realisation of a
random walk is dominated by wherever that path sits near zero; averaging across columns fixes it. Taking
the $201$ columns $x\in[-100,100]$ to $T=2^{20}$ (column $0$ validated bit-for-bit against the stored
centre column on all $1\,048\,576$ rows) and writing $D_x(T)$ for the excess of blacks over whites,
corrected for the $|x|$ rows in which the column lies outside the light cone, the root-mean-square of
$D$ across columns tracks the fair-coin value $\sqrt T/2$ at every scale: the ratio reads
$1.105,\,1.034,\,0.964,\,0.888,\,0.954,\,0.927,\,0.963,\,0.971,\,1.039,\,1.004,\,0.939$ at
$T=2^{10},\dots,2^{20}$, and the fitted exponent is $0.4932$ over the whole range and $0.4974$ over the
last five checkpoints. The largest $|D|$ across the $201$ columns stays between $2.5$ and $3.5$ standard
deviations throughout, against the $2.9$ expected for that many independent walks, and the pooled density
over all columns is $0.499999$.

Two things follow. First, the columns exhibit square-root cancellation with constant $1$, which is the
random-model value --- the same statement that, for the primes, is the \emph{conclusion} of the Riemann
Hypothesis rather than a route to it. Achieving the expected size of cancellation is what a proof would
have to deliver, so measuring it settles nothing; P2 needs only $D_0(T)=o(T)$, which is far weaker, and
no finite computation reaches either. Second, and this is the new negative: \emph{there is no chiral
signature in the column densities}. The root-mean-square of $D$ over the hundred left columns against the
hundred right ones gives ratios $0.913,\,0.928,\,0.828,\,0.974,\,0.951,\,1.021$ at the last six scales,
and the mean of $D$ is $+65.2$ on the left and $-24.6$ on the right against a standard error of $51.2$,
so neither is significant. Rule~30's chirality is plain in the front geometry --- the left edge advances
at speed $1$ and the right at about $0.2427$ --- and is invisible in the first-order statistics of the
columns. That matches what Remarks~\ref{du:family256} and~\ref{du:ring} found from two other directions:
chirality organises the geometry and does not organise the statistics.

\subsection*{The rule chart (\EXACT; criterion \KNOWN)}
From the seed, the aperiodic centre columns among the $256$ rules number thirteen (Remark~\ref{du:gate}); excluding the five degenerate ones ($126,129,161,167,181$; Remark~\ref{du:rueppel} for the type) leaves exactly $30,45,75,86,89,101,135,149$: the one-wire rules $l\oplus g(c,r)$ with gate weight $3$. Rule~135 from the seed is Rule~30's orbit complemented and shifted by $(1,1)$ (\EXACT, $2.2\times10^6$ cells); $75$ is $45$ from the word $11$ [10.36, 10.41].

\begin{remark}[the one hypothesis Rule~90 fails, and why it still does not separate the rules]\label{p2:hotspot}
Throughout this paper Rule~90 has satisfied every structural hypothesis we could state while failing
every conclusion we want, which is what closes route after route. The hot-spot bound is the first
exception, and it is worth recording precisely because the exception is instructive rather than useful.
Rule~90's centre column is $1,0,0,0,\dots$, so the all-zero word of length $m$ occupies almost the whole
prefix and $C_m=2^m$ exactly: it fails the bound as badly as any sequence can, while Rule~30 tracks a
coin. So here at last is a computable quantity that separates them.

It does not give P2, and the reason is structural. Bailey and Misiurewicz's theorem is an
\emph{equivalence}: the hot-spot bound holds if and only if the sequence is normal. So the bound is
not a route to the conclusion, it is the conclusion in other words, and exhibiting a quantity that
separates the two rules is exactly as hard as P2 itself. What the published proofs actually verify the
bound with is self-similarity --- for the Stoneham constants, the orbit is trapped on a grid of about
$n$ points --- and that is the ingredient we lack. The correction to record is that the missing
ingredient is in the \emph{verification} and not in the theorem: nothing in the hypothesis of the
hot-spot lemma requires a finite $k$-kernel, a digit formula or a functional equation, and an earlier
draft of ours said otherwise.

One further negative closes the obvious alternative. One might look for a third family in the
fractal-geometry literature \cite{Bandt}, where a finite automaton recognising an address relation yields a
self-similar quotient space. It is not a third family. The implication runs the wrong way --- the
automaton is the hypothesis and the geometry the conclusion --- and where that literature does produce
frequencies, for the Rauzy tiles of a primitive substitution, the vector of subtile measures \emph{is}
the Perron eigenvector of the incidence matrix. That is family~\textup{(i)} in geometric coordinates,
with the same Perron datum doing the same work.
\end{remark}

\section{What Problem 3 asks, and what is provable about it}\label{sec:p3}

Problem~3 asks whether the $n$-th cell of the centre column can be computed with substantially less
than $O(n)$ effort. We have said least about it, so we set out what the relevant literature settles,
what it forbids, and why the one route it appears to leave open does not reach the centre column. Nothing in this section is ours except the
computation reported at the end.

\begin{remark}[what the decision procedures for automatic sequences can and cannot do]\label{p3:walnut}
B\"uchi's theorem identifies the sets definable in $\langle\N,+,V_2\rangle$ with those whose base-$2$
representations are automaton-recognised, and Walnut \cite{Walnut} mechanises the resulting decision procedure,
eventual periodicity included. It settles Rule~90 outright, whose entire space-time diagram from the
seed is the definable set $\{(t,i):((t+i)/2)\wedge((t-i)/2)=0\}$ by Kummer's theorem, and it cannot be
aimed at Rule~30 at all: its input would be an automaton generating the centre column, and the
existence of that automaton is Problem~1 by Christol's theorem. Every decision procedure we have found
consumes the presentation we are trying to establish. On the difficulty of the statement itself,
eventual periodicity of a computable sequence is $\Sigma^0_2$-complete, so non-periodicity is
$\Pi^0_2$ and each pair (period, preperiod) is finitely refutable; what is missing is uniformity in the
pair, which is what Part~I supplies for the families it covers.
\end{remark}

\begin{remark}[the formalisation, and a category error to avoid]\label{p3:form}
Problem~3 is not the statement that predicting Rule~30 is $\mathsf{P}$-complete, and the two are
formally independent in both directions. $\mathsf{P}$-completeness bounds parallel \emph{depth}, while
Problem~3 bounds total \emph{work}, and \textsc{Circuit Value} is $\mathsf{P}$-complete and computable
in linear time; conversely there are cellular automata with $\mathsf{P}$-complete prediction and
communication complexity $O(n^{1/k})$. The inputs differ as well: the universality results give the
time parameter in unary, so their input length is $\Theta(t)$, whereas Problem~3 gives $n$ in binary,
input length $\Theta(\log n)$. The correct frame is therefore non-uniform --- circuit size for one
fixed sequence --- and not $\mathsf{P}$ versus $\mathsf{NC}$.
\end{remark}

\begin{remark}[status: open in both directions, and citably so]\label{p3:status}
Rule~110 is Turing-universal \cite{Cook}, and prediction for it is logspace-complete for $\mathsf{P}$ \cite[Thm.~1]{NW}. Rule~30 is one of
$33$ elementary rules left unresolved by the only systematic complexity classification of all $256$
\cite[\S4.1]{Meunier}, and Moore and Pnin \cite{Moore98}, whose method predicts the motion of non-colliding defects in Rule~18, write that they have been
unsuccessful in applying it to rules $22$, $30$ and $110$. So there is no lower bound of any kind for
Rule~30 and no $O(n^{2-\epsilon})$ algorithm either. On the upper side, Wolfram's request for any bound below
$O(n^{2})$ --- a ``nice partial result'' in his words \cite{Wolfram}, not Problem~3 itself --- has in fact been met in print: there is an $O(n^{2}/\log n)$-time algorithm, at the cost of
$O(n^{2}/(\log n)^{3})$ space \cite{NA}. No sublinear-space
method is known, and the generic space-efficient simulations are worse here than the trivial bound.
\end{remark}

\begin{remark}[the barrier ladder]\label{p3:barriers}
It is worth knowing which forms of Problem~3 are unreachable rather than merely unproved. A circuit
lower bound of $\Omega(n)$ is \emph{false}: every Boolean function of $n$ variables has circuits of size
$O(2^{n}/n)$, so the $n$-th bit of any sequence has circuits of size $O(n/\log n)$, and that ceiling
applies to us. A bound of $n^{\epsilon}$ would, through Impagliazzo and Wigderson \cite{IW}, give
$\mathsf{BPP}=\mathsf{P}$, and lies behind the natural-proofs barrier. Even $(\log n)^{\omega(1)}$ would
give $\mathsf{E}\not\subset\mathsf{P}/\mathrm{poly}$ and hence $\mathsf{EXP}\ne\mathsf{MA}$. Two
reformulations are foreclosed rather than hard: Levin's $Kt$ of the prefix is $O(\log N)$, since a short
program plus the time to run it describes the sequence, so no incompressibility statement is available
there; and $KT$ complexity is the circuit question relabelled. Relativisation and algebrisation do not apply here, there being no class separation in the statement.

One distinction belongs here, because it is easy to conflate two different obstacles. In restricted
models, lower bounds for \emph{explicit} objects do exist and are not barrier-bound: Nisan's bound for
algebraic branching programs is the Fliess--Sch\"utzenberger theorem in another notation, and it gives
exponential lower bounds for explicit polynomials over any field, $\F$ included. What that argument
consumes is a witness uniform in the degree, which comes from a formula for the object and never from
its initial segment. So in the weighted or branching-program setting our obstacle is not the
circuit-lower-bound barrier of the table above but the prefix obstruction of
Remark~\ref{du:censuscaveat}, and the two should not be quoted as one.
\end{remark}

\begin{remark}[the propagator carries no hardness; the source carries all of it]\label{p3:rigidity}
This is where Part~II's decomposition earns its keep, in the negative direction. One might hope the
Rule~150 propagator is hard to compute, and matrix rigidity is the tool for that. It is not: the
Pascal-mod-$2$ matrix is exactly the Kronecker power $\left(\begin{smallmatrix}1&1\\0&1\end{smallmatrix}\right)^{\otimes m}$,
which is the canonical object of Alman's non-rigidity theorem \cite[Thm.~1.8]{Alman} and is therefore
not Valiant-rigid. Independently and more simply, $\Lambda^{t}=\prod_{j\in\mathrm{bits}(t)}\Lambda(z^{2^{j}})$
by Frobenius, a product of at most $\log_2 N$ three-sparse factors, giving an $\F$-circuit of size
$O(N\log N)$ and depth $O(\log N)$. And the two linear rules settle it from the other side: Rule~90's
centre column is $1,0,0,0,\dots$ and Rule~150's is identically $1$, both computable in constant time.
So all of the difficulty of Problem~3 sits in the quadratic source, exactly where
Corollary~\ref{du:cor:C} says the chirality lives, and none of it in the kernel.
\end{remark}

\begin{remark}[the centre column is not a verifiable delay function]\label{p3:vdf}
Problem~3 has the shape of a cryptographic assumption, so it is worth saying why the centre column is
not a candidate for the one deployed primitive of that kind. A verifiable delay function needs a
challenge with min-entropy; a single fixed sequence has none, so the preprocessing phase of the
standard definition wins trivially. Moreover the sequentiality of iterated squaring is \emph{trapdoor}
hardness, broken in groups of known order, and no unconditional sequentiality bound exists for any
explicit map.
\end{remark}

\begin{proposition}[the window monoid is non-solvable; \EXACT]\label{p3:monoid}
For each pair $(c,f)\in\{0,1\}^2$ the map $\Phi(c,\cdot,f)$ of Section~\ref{ex:setting} is a
transformation of the $2^{K}$ windows; let $T_K$ be the monoid they generate. Then $|T_2|=113$ and
$|T_3|=41\,197$, neither is a group, both have minimum rank $1$, and their largest group
$\mathcal H$-classes below full rank are $\mathrm{Sym}(3)$ and $\mathrm{Sym}(4)$. At $K=4$ the largest
such class is $\mathrm{Sym}(5)$: a group of order $120$ acting faithfully on a five-element image, with
element orders $\{1,2,3,4,5,6\}$. Hence $T_4$ contains $A_5$ and is \emph{not solvable}, and its
Krohn--Rhodes complexity is at least $1$. By contrast the corresponding monoid for Rule~90 is a solvable
\emph{group} at even $K$, of order $8$ at $K=2$ and $96$ at $K=4$.
\end{proposition}
\begin{proof}[Computation]
$T_2$ and $T_3$ were enumerated in full by closure under left multiplication. At $K=4$ the enumeration
was capped at $6\times10^{6}$ elements, which suffices: the elements of a group $\mathcal H$-class are
those whose restriction to their own image is a bijection, and any set of them found within the monoid
generates a subgroup, since a monoid is closed under products and products within a group
$\mathcal H$-class stay in it. Closing the $67$ rank-five elements found under composition gives a group
of order exactly $120$ on five points, which is $\mathrm{Sym}(5)$.
\end{proof}

\begin{remark}[where this puts Rule 30 in Moore's classification]\label{p3:moore}
This is the sharpest algebraic statement we know about the rule, and it is the same invariant that
settled the nearest comparable case. Moore's parallel-prediction theorem \cite{Moore98} places a cellular automaton in
$\mathsf{ACC}^1$ when the blocked local rule's algebra is a \emph{solvable} group, and that is exactly
how single-defect motion in Rule~18 was shown to be computationally reducible; he conjectures that
automata built on non-solvable groups such as $A_5$ are $\mathsf{P}$-complete. Our corpus previously
recorded only that Rule~30's algebra is not a group. Proposition~\ref{p3:monoid} replaces that with a
positive statement: it contains $A_5$. So Rule~30 sits on the side of that dichotomy which Moore
conjectures hard, while Rule~90 --- which satisfies more of our hypotheses than Rule~30 throughout this
paper --- sits on the side where his theorem applies and gives a polylogarithmic-depth algorithm. The
conjecture is a conjecture, so this is evidence and not a bound.

There is, however, a \emph{theorem} at exactly the level our data reaches, and it is the sharper thing
to cite. Barrington and Th\'erien \cite{BT} prove a dichotomy for the word problem of a finite monoid:
if the monoid contains a non-solvable group the word problem is $\mathsf{NC}^1$-complete under
$\mathsf{AC}^0$ reductions, and if the monoid is solvable the word problem lies in $\mathsf{ACC}^0$.
Proposition~\ref{p3:monoid} is precisely the checked hypothesis of the first branch, so the word
problem of our window monoid is $\mathsf{NC}^1$-complete, unconditionally, while Rule~90's sits in
$\mathsf{ACC}^0$. Three limits belong with it and we state them rather than let the result be
overread: $\mathsf{NC}^1$ is a small class and completeness in it is far from the superlinear work
bound Problem~3 asks about; the statement is about the monoid $T_K$, not about the centre column; and
$\mathsf{AC}^0$ reductions are coarse. What it does is replace a conjecture with a theorem on the one
question where we have the hypothesis in hand.

The Rule~90 side of the comparison can be strengthened from measurement to proof. Its four window
generators are affine over $\F$ with a common linear part, so every subgroup of the monoid they
generate is metabelian, for every $K$ --- which is why the orders $8$, $96$, $896$ we computed at
$K=2,4,6$ are not a small-$K$ accident. Rule~30's generators are affine at no $K$.
\end{remark}

\begin{remark}[communication rank: a measurement, and the control it fails; \MEASURED]\label{p3:rank}
What remains is communication complexity. For a cellular automaton that is intrinsically universal, the
deterministic communication complexity of prediction is $\Omega(n)$ \cite[Thm.~3.1]{Meunier}, and the
standard lower bound on that is the logarithm of the rank of the prediction matrix. Taking the maximum
over all cuts, rules $90$, $150$ and $105$ give rank exactly $2$ at every $n\le10$ --- a clean instance
of the filter that has governed this paper --- while Rule~30 gives rank about $1.59^{n}$, a
$\log_2$-rank slope of $0.668$. Three caveats belong with it: the range is only $n\le10$, log-rank is a
lower bound on communication and not an upper one, and Rule~110's own slope declines over the same
range. Note also that Rule~110's intrinsic universality is itself open --- Open Problem~5 of \cite{KariSurvey}, with Ollinger recording that Cook's construction cannot be used to settle it --- so the hypothesis of the theorem above is not known to hold for any elementary rule. Meunier's results put this more sharply than mere openness: intrinsic universality is \emph{refuted} for $55$ of the $88$ minimal rules, and $33$ remain open, including both Rule~110 and Rule~30. An earlier version of this remark called a bound
$\mathrm{rank}\ge2^{\delta n}$ for Rule~30 the first hardness theorem about the centre column itself, and the only such target not barred in principle; that is \textsc{withdrawn}, because the statistic fails our own control (Remark~\ref{p3:rueppel}). Rules~$22$ and~$54$ have trivially computable centre columns --- $1,1,0,0,\dots$ (Remark~\ref{du:rule22}) and periodic with period~$4$ (Remark~\ref{du:family}) --- and their communication rank over $\F$ is of the same order as Rule~30's or larger at every $n$ computed (\EXACT). Taking the maximum over cuts at $n=6,7,8,9$, Rule~30 gives $19,26,44,79$, Rule~22 gives $36,58,94,147$ and Rule~54 gives $21,31,46,69$ (below Rule~30 only at $n=9$); at the centre cut, $n=10,11,12$, Rule~30 gives $53,79,104$, Rule~22 gives $231,355,555$ and Rule~54 gives $100,145,208$. Rule~$126$ exceeds Rule~30 at each of these $n$ too, but its column is identified only for $t\le2^{14}$ (Remark~\ref{du:rueppel}), so the argument does not rest on it. The measurement above is of real rank, which for a $0/1$ matrix is at least the rank over $\F$ (an odd minor is non-zero); for Rule~30 at the centre cut the two agree at $n=10,11$, so there the real rank fails the control in the same way. The reason is that the rank measures the Boolean function $F_n$ of Proposition~\ref{du:anf} over all $2^{2n+1}$ inputs, as Meunier's communication complexity of prediction does, and not the single-seed column: a rank lower bound would be a theorem about predicting Rule~30 and not about the centre column, just as the $\mathsf{NC}^1$-completeness above concerns the window monoid and not the column. Over $\F$ the rank of the prediction matrix equals the rank of the matrix of ANF coefficients $a_{S,T}$ of $F_n$ split across the cut, since the monomial-evaluation matrices are unitriangular in the subset order (\PROVED; ledger 10.230d~A). Computed to $n=14$ (\EXACT; 10.230d~B): the maximum over cuts sits left of centre and moves left with $n$, and at a fixed block of $n_L$ leftmost variables the rank saturates in $n$ at about $2^{n_L}/3$ ($87$, $172$, $343$, $684$ for $n_L=8,\dots,11$), so the slope $0.668$ is a reading on a range where the cut is still moving and not a constant of the rule; whether the slope tends to $1$ is \OPEN.
\end{remark}

\begin{remark}[a control on inferring hardness from a complexity measure; \EXACT]\label{p3:rueppel}
The caveats above are not hypothetical, and the witness is inside the family of Part~II. The centre
column of Rule~$126$ is the Fredholm--Rueppel sequence for every $t\le2^{14}$ (Remark~\ref{du:rueppel}), which covers every bit measured below. Berlekamp--Massey on
its first $2048$ bits returns linear complexity $1024$, exactly $n/2$ --- the perfect-profile value --- and exactly what a random
control returns on the same length; Rule~30's centre column returns $1025$ and Rule~150's returns $1$.
So on this measure Rule~30 is indistinguishable both from a random sequence and from Rule~$126$ ---
whose $t$-th bit on that range is $1$ if and only if $t+1$ is a power of two, and is therefore computable in
$O(\log t)$ time from $t$ alone. A perfect linear-complexity profile is compatible with constant-depth
computability. We claim nothing new in that observation: this sequence is the standard textbook
witness for a perfect linear complexity profile, and its primary source is Rueppel \cite{Rueppel}; although we have not located an earlier source for
it in our own corpus, it should be assumed known and is used here only as a control. The consequence for Problem~3 is a restriction on
what any such measurement can establish: a complexity statistic that Rule~$126$ also maximises cannot,
on its own, be evidence that Rule~30's column is hard to compute. The rank measurement above is of
this kind as well, with Rules~$22$ and~$54$ in the role of Rule~$126$: rules $90$, $150$ and $105$ give rank exactly $2$ at every $n\le10$, which separates the affine rules, but Rules~$22$ and~$54$, whose columns are eventually $0$ and periodic with period~$4$, have communication rank of the same order as Rule~30's or larger at every $n$ computed (Remark~\ref{p3:rank}), so the statistic does not separate the
tractable cases. Rules~$22$, $54$ and~$126$ should be run against any future candidate before it is
reported, and they are the controls we now use.
\end{remark}

\part{Closed routes, verdict, and the measure-level obstruction}

\section{Closed routes}\label{sec:closed}

Most of the work behind this paper was the elimination of routes, and the result is more useful as a
table than as prose. Each row names a route, the reason it fails for this object, and where the reason
lives. ``Rule~90'' as a reason means the firewall of Proposition~\ref{p2:rule90}: the rule satisfies
the hypothesis and fails the conclusion. ``Prefix'' means the obstruction of
Remark~\ref{du:censuscaveat}: every finite word extends to an automatic sequence, so no statistic of a
prefix can decide the question.

\medskip
\noindent\begin{tabular}{@{}p{0.31\textwidth}p{0.45\textwidth}l@{}}
\hline
\textbf{Route} & \textbf{Why it fails here} & \textbf{Where} \\
\hline
Measure-level randomisation & Rule~90 satisfies strictly more and has density $0$ & \ref{p2:rule90}\\
Density at width $w$, then descend & Rule~90's width-$w$ trace has density $0$ at every $w$ and is aperiodic for $w\ge2$ & \S\ref{sec:p2}\\
Spread above $\tfrac12$ & Rule~90 saturates the threshold exactly & \S\ref{sec:p2}\\
Substitutive route to a density & needs linear recurrence; ours is $R(13)\ge105\,901$ & \ref{p2:twofamilies}\\
Self-similarity route to a density & hypothesis consistent, no known verification & \ref{p2:hotspot}\\
Fractal address-relation route & family (i) in geometric coordinates; implication runs backwards & \ref{p2:hotspot}\\
Kernel census as aperiodicity & prefix; a padded prefix is automatic with $107\,047$ classes & \ref{du:censuscaveat}\\
$k$-regularity as a weaker hypothesis & over $\F$ it coincides with automaticity & \ref{du:regular}\\
Hankel rank & prefix, dominated by the antichain census, and a coin saturates it & \ref{du:hankel}\\
Boolean spectral invariants & the designated coefficient is invisible to the Walsh multiset & \ref{du:anfobstruction}\\
Sublattice linearisation & the invariant set is three points & \ref{du:nosubdyn}\\
Conservation laws, solitons, ultradiscrete-integrable systems & no additive invariant of order $1$ or $2$, exhaustively & \ref{fr:closedfamilies}\\
Integrable ``Rule 54/150'' & those are reversible two-step automata, not these maps & \ref{fr:closedfamilies}\\
Mixing for the uniform measure & holds, and says nothing about either problem & \ref{fr:closedfamilies}\\
Solvable-circuit duality & needs bipermutivity; and the class is empty of informative members & \ref{du:nodual}\\
$2$-adic ergodicity & not ergodic, and the obstruction is the diagonal read & \ref{du:padic}\\
Operator-algebraic reformulation & equivalent by definition; the invariant is the suffix count & \S\ref{sec:verdict}\\
\hline
\end{tabular}

\noindent\begin{tabular}{@{}p{0.31\textwidth}p{0.45\textwidth}l@{}}
\hline
\textbf{Route} & \textbf{Why it fails here} & \textbf{Where} \\
\hline
Coupling, Kingman, Fekete & the increment is an exact cocycle & \ref{fr:exists}\\
$(a,\tau)$-blockers & cost $2^{a}$ at numerator $a$, published only to $a\le4$; sup constant infinite & \ref{fr:rational}\\
Front as a finite max-plus system & band width grows; every bounded control is flat & \ref{fr:tropical}\\
Bounded-window prediction & provable at one width would be P2 itself & \ref{co:windowbound}\\
Learning-theoretic mistake bounds & quantified over sequences, not about one & \ref{co:windowbound:read}\\
Decision procedures for automaticity & each consumes the presentation we lack & \ref{p3:walnut}\\
Matrix rigidity of the propagator & the kernel is the canonical non-rigid object & \ref{p3:rigidity}\\
Verifiable-delay-function framing & the challenge has zero min-entropy & \ref{p3:vdf}\\
Communication rank of the prediction matrix & Rules~$22$ and~$54$ (centre columns eventually $0$ and of period $4$) and Rule~$126$ have rank of the same order as Rule~30's or larger at every $n$ computed: the rank measures $F_n$ over all inputs, not the column & \ref{p3:rank}\\
\v{C}ern\'y-type bounds & wrong quantifier, and vacuous by a factor of $949$ & \ref{ex:rem:notsync}\\
Topological classification & one bit of discriminating power, and it is bipermutivity & \ref{du:fuks}\\
Preperiod theory & for nonlinear rules the only mechanism we have found is a blocking word, which needs inertness & \ref{fr:reset:status}\\
\hline
\end{tabular}
\medskip

Two entries deserve emphasis because they are the sharpest form of the obstruction. The first is the
observation recorded in Remark~\ref{du:hankel} (\MEASURED): every instrument tried here separates Rule~30 from
Rule~90, and none tried separates it from a fair coin. The second is that the one
hypothesis Rule~90 does fail, the hot-spot bound, is equivalent to the conclusion we want, so
exhibiting a separating quantity there is exactly as hard as P2.

\section{Verdict}\label{sec:verdict}

\begin{enumerate}[nosep]
\item[\PROVED] Part~I: three infinite families and two thresholds of excluded period words, for every number-like configuration; the reach of the method characterised (Proposition~\ref{ex:prop:boundary}). At that boundary: the exponential clause cannot be discharged by enumeration (Corollary~\ref{ex:cor:infimum}); the branching criterion for uncountability (Proposition~\ref{ex:prop:branching}); a finite row-$0$ block that traps the interface (Proposition~\ref{du:prop:trap}); and the diagonal cap, $g(N)=6$ on every even diagonal $N\ge58$ and $g(N)=2$ on every odd $N\ge5$ (Theorem~\ref{du:thm:cap}). Of Part~I, the family $1^b0$ is also public in \cite{Cochon} (10--12 September 2026), the exclusion of the two constants is \KNOWN\ for the seed, and the method in kind is in other public work; the sources are in Remark~\ref{ex:prior}. Part~II: the Duhamel formula, exact decimation of its class, the closed form in $\alpha$, the chiral-only corollary, the centre-cell degree law $\deg F_t=2t-1$ with its unique top monomial (Theorem~\ref{du:thm:degree}), the right-edge influence law $\mathrm{Inf}_{t}=\mathrm{Inf}_{t-1}=2^{-t}$ (Proposition~\ref{du:prop:influence}), and the ring surjectivity defects of Rules $45$, $150$, $90$ and $30$ from the transfer monoid, Rule~$30$'s exponent being $\log_2\lambda-1$ with $\lambda$ the real root of $x^{3}-x^{2}-2$ (Proposition~\ref{du:prop:ringdefect}). Part~III: reset formula, degree formula, monotone front, exact failure condition, window slaving, Theorem~H$^*$ and the light-cone corollary. Of these the monotone front, with the split into background and defect, is also in \cite{Dib} (commits \texttt{0a1cda2} and \texttt{a8dbbc0}, 7--8 September 2026), and the code of Wolfram's 2019 prize post \cite{Wolfram19} takes a running maximum; the one-step law behind the reset formula is \KNOWN\ (\cite[Lemma~3]{Rowland}; also \cite{Pat,Dib,Nersissian}). Unbounded diagonal periods (Theorem~\ref{fr:unbounded}) is \KNOWN, Gravner--Griffeath \cite[Prop.~3.3]{GG}; it is reproved here only to fix the diagonal-frame notation. Part~IV: the downstream-only proposition, the record-parity theorem, the two-sided criterion, the parent-instance split.
\item[\EXACT] The kernel bounds of Theorem~\ref{du:certified}; the chiral recursion verified on $126$ words; the front data to $d=90\,000$ (E1--E6), with E1--E3 and the (C) half of E4 extended to $8\times10^{5}$ by the third frame of E7, the (U) half of E4 only to $90\,000$; the frozen diagonals and doubling positions to $8\times10^{5}$ (E7); $\tau_d$ exact to $8\times10^{5}$ (Remarks~\ref{fr:concave}, \ref{fr:dyadic}) and per diagonal, with a certificate, to $4\times10^{6}$, so Conjecture~H holds on $[20,4\times10^{6}]$ (E10; also on $[20,4\times10^{7}]$ by exact simulation with two independent simulators, ledger 10.163c2; public elsewhere: $d\le78\,287$ for the settling times, in the table of \cite{Wolfram19}, with $1\,000$ terms in \cite{A363346}, and $d\le200\,000$ for the conjecture, in \cite{Dib}), with bin medians of $M_d/d$ to $4\times10^{6}$ (E9, \MEASURED); the fibre sizes to $k=60$ ($|L_{60}(01)|=103\,220$); the forward trees $F_P$ to $P=16$; the kill depth to $T=120$.
\item[\OPEN] P1 and P2. Conjecture~H ($\tau_d\ge d$ for $d\ge20$, Conjecture~\ref{fr:conj:H}) is also open, and on its own it decides neither: it is a statement about the geometry of the ordered region, and Rule~$60$ has $\tau_d=d+1$ for every $d\ge0$ with centre column constantly $1$, periodic and of density $1$. The statement of Conjecture~H is not ours: it is a question in Wolfram's prize post \cite{Wolfram}, and it is in \cite{Pat} (commit \texttt{08786d1}, on the main branch 28 August 2026) and in \cite{Dib}. Theorem~\ref{fr:thm:Hstar} reduces Conjecture~H to two inequalities on the front, whose margins are measured and wide (Remark~\ref{fr:Cmargin}) and which no proof reaches; every sufficient statement on the front row we have found is equivalent to them, and what a proof needs is a lower bound on the density of ones in the transient interior. Conjecture~H settles no Prize problem (ledger 10.153c2~D). The automaticity route needs an invariant of the chiral Cartier recursion taking infinitely many values; none is known. The kernel criterion of Allouche--Shallit--Yassawi \cite[Thm.~1]{ASY} needs only one residue class modulo $2^{k}$ at each level $k$, with the resulting kernel elements pairwise distinct; in the instances collected there the residues are given in closed form, and no dynamical recursion has yet supplied such a family. The zero-run criterion (Remark~\ref{du:zerorun}; observed to hold at the coin's rate; \MEASURED) and the two-sided criterion of Corollary~\ref{co:cor:twosided} are honest bridges between P1 and P2 without a proof on either side. Conjecture~\ref{fr:universal} (left-half universality) is refuted by the $53\,209$-cell starting word of \cite{Lania}, which takes the other branch after $53\,207$ (\EXACT, our check of that state; ledger 10.235); which finite words take the seed's branch is \OPEN. Whether $L(01)$ contains an interface that is not eventually periodic: one such interface already puts $01$ beyond Theorem~\ref{ex:thm:A} at every $K$ (proof of Proposition~\ref{ex:prop:boundary}), and uncountability of $L(01)$, the hypothesis stated there, is one sufficient condition for it and not a necessary one, since a countable closed shift-invariant set can contain a point that is not eventually periodic. Counting cannot reach uncountability (Corollary~\ref{ex:cor:infimum}), the branching criterion is the only route to it here, and the trap shows that route must succeed despite atoms of the push-forward (Remark~\ref{du:rem:trap}). Random right halves driven by $(01)^\infty$ settle into a stationary stream of positive measured entropy (Remark~\ref{ex:rem:stream}), so the zero-entropy expectation is withdrawn and positive entropy of $L(01)$ is \MEASURED; uncountability remains \OPEN.
\item[Closed] Every bounded-window or bounded-state route to P1 (Section~\ref{sec:record}); the constructive zero-run route (halving law, ledger 10.76; not reproduced in this paper); small linear lifts (Proposition~\ref{du:lift}). The bounded-$f$ kill route is not closed but unsupported by the data (Section~\ref{sec:record}). From the wider literature, also closed: the sublattice linearization that underlies every rigorous Rule~18 result, by Proposition~\ref{du:nosubdyn}; every method resting on a conservation law, by exhaustion over the elementary rules (Remark~\ref{fr:closedfamilies}); the measure-level routes of Section~\ref{sec:p2}, at every trace width, not only at width one; and, for Problem~3, the communication-rank route (Remark~\ref{p3:rank}).
\item[Also closed, by equivalence] The operator-algebraic reformulation. The chain representation of the Cuntz algebra $\mathcal O_2$ (\KNOWN: Kawamura \cite{Kawamura}, Bratteli--Jorgensen \cite{BJ}) attached to a binary word is irreducible exactly when the word is not eventually periodic, so P1 is equivalent to an irreducibility question; but the equivalence is definitional, the only numerical invariant in that theory is the count of distinct suffixes, which is the Myhill--Nerode count whose finiteness \emph{is} eventual periodicity, and $K$-theory is unavailable since $K_*(\mathcal O_2)=0$. Likewise the learning-theoretic route (Remark~\ref{co:windowbound:read}).
\item[Decision points] Two questions now have the property that either answer is informative. The front speed is rational or it is not, and the second answer splits in two (Remark~\ref{fr:tropical}). If it is rational, a universal $(a,\tau)$-blocker would certify it, but that route is closed twice, by its cost (Remark~\ref{fr:rational}) and by its infinite sup-over-environments constant (Remark~\ref{fr:exists}). If it is not, the front's discrepancy from a line is either bounded, as for a Beatty sequence, or unbounded, the third branch, which is the one the growing band of Remark~\ref{fr:tropical} points to (\MEASURED); and the unbounded periods of Theorem~\ref{fr:unbounded} are not the mechanism, since irrationality needs no unbounded system, the dividing line being autonomy rather than dimension (Remark~\ref{fr:tropical}). And the centre column is substitutive or it is not: the recurrence function argues against the primitive substitutivity that family~\textup{(i)} of Proposition~\ref{p2:twofamilies} needs (\MEASURED), which removes one of the two families of technique behind every single-sequence density-$\tfrac12$ theorem we know; the column is consistent with the hypothesis of family~\textup{(ii)}, and what is missing there is any known way to verify that hypothesis for one orbit (Section~\ref{sec:p2}). The cheapest unclaimed step in the paper is neither of these but the near-diagonal
almost-subadditivity of the front (Remark~\ref{fr:exists}): proving
$M_{d+e}\le M_d+M_e+f(d+e)$ for $d\le e\le2d$ with $\sum f(n)/n^{2}<\infty$ would give the existence of
$\lim M_d/d$. We have been aiming at a constant $f=326$, which the measurement supports but which is
strictly harder than necessary: $\sum n^{\alpha-2}$ converges for every $\alpha<1$, so any
$f(n)=O(n^{\alpha})$ with $\alpha<1$ suffices, and Remark~\ref{fr:concave} reduces the whole statement to
a bound on a single index. The coupling route we
proposed in an earlier draft is \textsc{refuted}: the increment is an exact cocycle, so every
subadditivity hypothesis holds with equality.
\end{enumerate}

\begin{remark}[families closed by exhaustion over the elementary rules]\label{fr:closedfamilies}
Three families of technique are unavailable for reasons that are themselves theorems, proved by
exhaustion over the elementary rules. Rule~30 is not number-conserving --- only $204$, $170$, $240$,
$184$ and $226$ are \cite{BFuks} --- and has no additive invariant of second order either, the
Hattori--Takesue conditions leaving only $12$, $14$, $15$, $34$, $35$, $42$, $43$, $51$, $140$, $142$,
$200$ and the trivial cases \cite{Fuks05,HT}. Every soliton, ultradiscrete-integrable, height-function
and hydrodynamic method begins with a conservation law, so this closes them together, and it is why the
Rule~184 identity of Theorem~\ref{co:thm:record} transfers as an identity and not as a method. The
exactly solvable ``Rule~54'' and ``Rule~150'' of the statistical-mechanics literature \cite{BKP,WPG}
are reversible two-step automata rather than these maps. And the strongest measure-level theorem that
does apply, that every one-sided permutative rule is mixing for the uniform measure \cite{Kleveland},
says nothing about either problem.
\end{remark}

\section{Why no measure-level argument can give P2}\label{sec:p2}

Rule~30 is surjective (Hedlund \cite[Thm.~6.6]{Hedlund}) and the balance property behind that --- the
preimage of a length-$k$ cylinder is a disjoint union of exactly $|A|^{r+\ell}$ cylinders of length
$k+r+\ell$ --- is the Blankenship--Rothaus theorem, first published as \cite[Thm.~5.4]{Hedlund}. So the
uniform Bernoulli measure is invariant, and it is tempting to read P2 as a consequence. It is not, and
the obstruction can be stated as a proof rather than as a caution.

\begin{proposition}[the seed is invisible to measure-level hypotheses; \KNOWN]\label{p2:rule90}
No hypothesis of the form ``$F$ is surjective, preserves the uniform measure, is (bi)permutive, is
linear, or asymptotically randomizes every harmonically mixing measure'' can imply that the centre
column of $F$ from a single seed has density $\tfrac12$.
\end{proposition}
\begin{proof}
Rule~90 satisfies all of them: it is bipermutive, surjective, preserves the uniform measure, is a
linear cellular automaton, and therefore asymptotically randomizes every harmonically mixing measure
in the sense of \cite{PY}. Its centre column from the single seed is $1,0,0,0,\dots$, of density $0$.
\end{proof}

Rule~90 thus satisfies strictly more than Rule~30 does and fails the conclusion outright, which is the
sharpest form of the remark that permutivity alone cannot separate the two rules
(\cite[\S4]{Kopra}; \cite[\S1]{Condrey}). Four further points fix the status of this route.

First, the randomization theorems do not merely fail on the initial measure, they never engage: every
one of them quantifies over \emph{linear} automata, and Rule~30 is not linear, so the question of
which measures they randomize does not arise. Pivato and Yassawi do not discuss point masses at all;
what they say is that all known sufficient conditions require the measure to have full support, and a
point mass violates that at every rank. Their \S8 also refutes the conjecture that positive entropy
suffices, so ``the obstruction is that the seed has zero entropy'' would be the wrong diagnosis.
Pivato's survey \cite{PivatoSurvey} contains nothing about single-seed or Dirac initial conditions.

Second, the routes that have worked for a single deterministic orbit rest on extra arithmetic
structure, and they sit at the two ends of one complexity axis (Proposition~\ref{p2:twofamilies}). At the
low-complexity end are the decimation identity $t_{2n}=t_n$ that gives the Thue--Morse density exactly
$\tfrac12$ and the Perron-eigenvector frequencies of primitive substitutions; a finite $k$-kernel belongs
only to the constant-length (automatic) case, since the Perron argument needs primitivity rather than
constant length and so covers primitive substitutive words of non-constant length, the Fibonacci word
for instance. At the high-complexity end is the normality of the Stoneham constants \cite{BM}, through
the hot-spot lemma: its hypothesis is a one-sided bound on block frequencies with no finite $k$-kernel
in it, and self-similarity enters only in verifying that bound (Remark~\ref{p2:hotspot}). Our kernel
figures give a scale, not an exclusion from either end: $A^c_2(2^{21}-1)=138\,672$
(Theorem~\ref{du:certified}; a census at $2^{24}$ terms reported more than $10^{6}$ in \cite{Auto}, not
re-derived here) says that an automaton generating $c$ least-significant digit first would need at
least that many states, not that none exists, since every finite prefix extends to an automatic
sequence (Remark~\ref{du:censuscaveat}). What separates $c$ from the low-complexity end is its
recurrence function, and that separation is \MEASURED.

Third, the measure-rigidity template that replaced symmetry in the quantum setting is unavailable in
principle here, not merely missing. Rule~30 is left-permutive, hence closing \cite[\S3]{Kurka}, and for closing maps the
jointly $(\sigma,F)$-periodic points are dense (Boyle--Kitchens \cite[Thm.~4.4]{BK}); the invariant
simplex is therefore large, and no rigidity statement of the Host--Maass--Mart\'inez kind can hold.

What does constrain a single orbit is small but real: for every finite configuration the width-$2$
traces are never eventually periodic and the orbit has infinitely many limit points
(\cite[Thms.~3.5, 4.7]{Kopra}), and the right Lyapunov exponent is exactly $1$ at \emph{every} point,
seed included, for a left-permutive rule with negative left radius \cite{Kurka}. The left exponent, the
one that carries our front, has no theorem at any point. The measured values of that speed, ours and
Wolfram's two \cite{NKS}, are collected in Remark~\ref{fr:slopebar}, which also gives their precision honestly.
Wolfram's are not the first in print: his own CRYPTO~'85 extended abstract \cite{W85} already gives the left Lyapunov exponent as $0.25$, and Meier and Staffelbach \cite{MS} report that the effect of a
bit change propagates to the left with speed roughly $\tfrac14$, an observation they attribute to an
earlier paper of Wolfram's;
Remark~\ref{fr:rational} says what the literature can and cannot do with such a number either way.

Fourth, and this is the one that cuts against our own programme rather than against the literature.
Every proved statement that a single explicit deterministic sequence has letter density exactly
$\tfrac12$ --- and we have not found an exception --- belongs to one of two families, and the two sit at
opposite ends of a single complexity axis. Our own measurements place the centre column at the
high-complexity end (\MEASURED): outside the first family, and consistent with the hypothesis of the
second, though out of reach of every known way of verifying it.

\begin{proposition}[the two families, at opposite ends of one axis; \KNOWN\ chain, \MEASURED\ placement of $c$]\label{p2:twofamilies}
The published single-orbit density theorems divide as follows.
\begin{enumerate}[nosep]
\item[\textbf{(i)}] \emph{Low complexity, through substitutivity.} Thue--Morse through the decimation
identity $t_{2n}=t_n$; primitive substitutive sequences through the Perron eigenvector; automatic
sequences; Toeplitz words; and the Thue--Morse word in base $3/2$ \cite[Thm.~18]{CERS}. Each requires
the sequence to be substitutive, and the frequency theorem requires \emph{primitivity}, since a
non-primitive substitution need not have letter frequencies at all. Primitivity forces linear
recurrence, hence $p(n)\le Kn$ and $R(n)\le Kn$ for a fixed $K$, where $R$ is the recurrence function.
\item[\textbf{(ii)}] \emph{High complexity, through a block-frequency bound.} The normal numbers:
Champernowne's constant by explicit construction, and the Stoneham constants through the hot-spot
lemma \cite{BM}. The hypothesis is a one-sided upper bound on block frequencies --- for every $y$,
$\liminf_m\limsup_n 2^mA(x,y,n,m)/n<\infty$ --- and nothing else.
\end{enumerate}
The two families sit at opposite ends of a single axis. Family~\textup{(i)} forces $p(n)\le Kn$. Family
\textup{(ii)} forces the reverse: a uniform hot-spot bound $C$ makes every word of length $m$ occur
often enough that $p(m)\ge2^m/C$, so its hypothesis \emph{requires} exponential factor complexity.
\end{proposition}
\noindent Where $c$ sits is a measurement, not part of the proposition: as far as computed it is at the \textup{(ii)} end, not between them (\MEASURED, from the measurements that follow).
\begin{proof}[Measurements]
For \textup{(i)}: on $2^{18}$ terms, the recurrence function satisfies $R(13)\ge105\,901$, so any
linear-recurrence constant for the centre column exceeds $8\,146$; the growth is geometric, with
$R(n)/2^{n}$ tracking $n\log 2$, which is the value for an i.i.d.\ coin, and $p(n)=2^{n}$ exactly for
$n\le16$ on $2^{20}$ terms. A gap between consecutive occurrences found in a prefix is a genuine gap in
the infinite word, so the lower bound on $K$ is not an artefact; what a finite table cannot do is
refute an asymptotic $O(n)$ bound, so this is evidence together with a quantitative constraint --- any
primitive substitution generating this word needs linear-recurrence constant above $8\,146$, which
bounds its alphabet from below --- rather than a proof. For the claim about family~\textup{(ii)}: a uniform bound $C_m\le C$ forces at least $2^m/C$ distinct
words of length $m$, since no word may carry more than $C$ times its fair share of a prefix. Our
measured profile $C_m=\max_w2^m\,\mathrm{count}(w)/N$ at $N=2^{18}$ is
$1.00,1.00,1.01,1.01,1.02,1.04,1.06,1.10,1.13,1.21,1.34,1.58,1.72,2.19,2.75,3.50$ for $m=1,\dots,16$,
tracking an actual coin sequence on the same prefix to within $0.25$ at every $m$, the growth being the
finite-$N$ floor and nothing else. Rule~90 gives $C_m=2^m$ exactly.
\end{proof}

The honest conclusion is therefore not that we fall between the two families but that we sit squarely
at the high-complexity end and cannot use its methods. We are excluded from family~\textup{(i)} by the
recurrence function. We are \emph{not} excluded from family~\textup{(ii)}'s hypothesis --- we are
consistent with it, and the measured profile is the evidence --- but we are excluded from every known
way of \emph{verifying} that hypothesis. That is the precise sense in which P2 has
no available technique, and it is sharper than the caution we gave earlier.

One subtlety is worth spelling out, because it is the obvious objection. Substitutive is strictly
broader than automatic --- automatic means substitutive of constant length --- and a primitive
substitutive word of non-constant length, the Fibonacci word for instance, is aperiodic \emph{and} has
letter frequencies given by the Perron eigenvector. So primitive substitutivity would deliver P1 and
P2 together, and our Part~II's argument against automaticity would not by itself rule it out. What
rules it out is the complexity, not the automaticity: linear recurrence is incompatible with a
recurrence function growing like $n2^{n}$. The exclusion therefore rests on the recurrence function
alone, and the kernel census could not carry it in any case: no prefix census can imply
non-automaticity (Remark~\ref{du:censuscaveat}).

Two further facts close the obvious workarounds. One might hope to prove a density statement at
width~$2$, where Theorem~\ref{ex:thm:zero} does bite, and then descend to the centre column. Rule~90
kills that at every width, not just at width~$1$: counting the rows whose width-$w$ window is not
identically zero, from the single seed, gives
\[
\begin{array}{lrrrrr}
T & w=1 & w=2 & w=4 & w=8 & w=16\\
10^{3} & 1 & 10 & 25 & 50 & 91\\
10^{4} & 1 & 14 & 37 & 78 & 151\\
10^{5} & 1 & 17 & 46 & 99 & 196\\
10^{6} & 1 & 20 & 55 & 120 & 241
\end{array}
\]
which is $O(\log T)$ at each fixed width, so Rule~90's width-$w$ trace is aperiodic for every $w\ge2$
by Theorem~\ref{ex:thm:zero} and has density $0$ for every $w$. No theorem of the form ``rapidly left
expansive and number-like implies the width-$w$ trace has density bounded below'' can be true, at any
width. The same test kills the weakest quantitative form of P1 as well. Reading a row as a binary
expansion, $r_t=\sum_{i\ge0}u^t_i2^{-i-1}$, the first digit of $r_t$ is the centre cell, so
$\limsup r_t-\liminf r_t>\tfrac12$ would force the centre column not to be eventually constant --- the
exact analogue of Vijayaraghavan's remark on $(3/2)^n\bmod1$. Over three thousand steps Rule~30 gives
spread $0.99906$; Rule~90 gives $\liminf=0$, $\limsup=\tfrac12$, spread exactly $\tfrac12$. The
threshold is saturated but not exceeded inside the class, so nothing above it is provable there, and
the conclusion it would have given is in any case now a theorem by a Rule-30-specific argument
\cite[Cor.~5]{Condrey}.

A second single-orbit family exists, and although it concerns linear rules it changes what P2 should be
asked to prove. Kawaharada studies the cumulative count of ones along one orbit of one elementary rule
from a single-site seed, with no measure anywhere, and finds that the limiting density profile is a
\emph{singular function} rather than a constant \cite{Kawa}. That is a proved single-orbit density statement outside both families of
Proposition~\ref{p2:twofamilies} --- no contradiction, since that classification concerns density exactly $\tfrac12$ --- and its conclusion is a negative one: for
those rules the natural density does not converge to a number at all. So P2 is not only unproved but
formulated in a way that fails for the nearest rules where the question has been answered, and any
route to it must explain why Rule~30 should be better behaved than they are.

What P2 would need, then, is a technique that is neither measure-level nor substitutive, and exactly one
such technique exists in print. Gravner and Griffeath prove that a two-dimensional Packard box rule
fills the lattice from a single seed with asymptotic density $\tfrac49$ \cite{GGBox,GGZ2}, and their own
assessment is that exact global asymptotics for a nontrivial local interaction is rare in nonlinear
spatial dynamics. The recipe is: a \emph{finite} alphabet of nested domain types, the density from a
linear recursion over that alphabet, then a refinement argument for uniformity. It needs monotone
solidification, which Rule~30 does not have, and an additive component, which Part~II supplies; but it
breaks at the first step for us, because Theorem~\ref{fr:unbounded} makes the periods of the left wedge
unbounded, so the domain alphabet is infinite, and the centre column lies in the region where no domain
structure is known at all. If a finite domain alphabet is ever found for any part of this orbit, that is
the paper to imitate.

Where the difficulty sits, in one sentence each way. From the front: the stall run is decided by parities of transient cells, and a certificate carried at the front keeps losing them, since a bit of a diagonal it does not carry survives only until the first $1$ of the diagonal to its left --- by the erasure-cone theorem of the referee dialogue (\PROVED\ there and summarised in ledger~10.61, Chat~B, round~5; the argument is not reproduced here) a certificate holding the two front diagonals on a window of rows, the background and boundedly many further diagonals loses at least one row of that window per plateau (Definition~\ref{fr:def:jumps}); on the seed's run, with no further diagonals, the actual loss outruns the refill from the advancing front by $0.46$ to $0.70$ rows per diagonal (\EXACT), but the proved one row per plateau does not, and a net loss in general would need a lower bound on the density of ones in the transient interior (\OPEN). From the column: Section~\ref{co:split} --- after every white centre cell one bit enters from the right, and after every black one the next bit is computed from the column's whole past.

\section{Sources and reproduction}\label{sec:sources}
Ledger \texttt{SESSION\_FINDINGS.md} (\S\S1--10.235, the entries suffixed b, c and c2, and the session-d entries 10.229d--10.243d; entries 10.199--10.229 record the boundary-section, degree-law, influence and ring-defect work of this version, and 10.150c the ledger-to-paper audit) and \texttt{HANDOFF.md}; proofs \texttt{PROOFS\_theorems\_6\_to\_9.md}; the component documents \texttt{rule30\_excluded\_words.tex}, \texttt{duhamel\_decimation\_section.tex}, \texttt{front\_section.tex}, \texttt{collapse\_section.tex}; the automaticity document \cite{Auto} and its independent re-run \texttt{agent\_certificate.md}; the inventory \texttt{rule30\_totality\_paper.md}. Scripts are the \texttt{s3\_*.py}, \texttt{ag1\_*}, \texttt{ag2\_*} files in the working folder, together with the agent working directories \texttt{agent\_reach/} (the trap, the cap witness \texttt{cap/ring\_cycle.txt}, the literature sweep) and \texttt{agent\_topcoef/} (the degree law); the certificates of this version are \texttt{s3\_trap\_repro.py}, \texttt{s3\_cap\_verify.py}, \texttt{s3\_cap\_upper.py}, \texttt{s3\_cap\_plateau.py}, \texttt{s3\_cap\_odd.py}, \texttt{s3\_runcap.py}, \texttt{s3\_gate\_fast.py}, \texttt{s3\_verify\_topcoef.py}, \texttt{s3\_verify\_general.py}, \texttt{s3\_influence.py}, \texttt{s3\_ring\_defect.py} and \texttt{s3\_ring\_defect2.py}, each with its log. This document compiles with Tectonic (XeTeX engine). The searches for other public work behind the citations added in October 2026 are recorded in ledger entries 10.165c2--10.167c2 and 10.235.

\begin{thebibliography}{99}
\bibitem{FS} P.~Flajolet, R.~Sedgewick, \emph{Analytic Combinatorics}, Cambridge University Press, 2009.
\bibitem{SYZS} A.~Szilard, S.~Yu, K.~Zhang, J.~Shallit, Characterizing regular languages with polynomial densities, \emph{MFCS 1992}, LNCS 629, 494--503.
\bibitem{Incitti} R.~Incitti, The growth function of context-free languages, \emph{Theoret.\ Comput.\ Sci.} 255 (2001) 601--605.
\bibitem{Cass02} J.~Cassaigne, Constructing infinite words of intermediate complexity, \emph{Developments in Language Theory} (DLT 2002), LNCS 2450, 173--184.
\bibitem{BG} M.~R.~Bridson, R.~H.~Gilman, Context-free languages of sub-exponential growth, \emph{J.~Comput.\ System Sci.} 64 (2002) 308--310.
\bibitem{DDMP} S.~Donoso, F.~Durand, A.~Maass, S.~Petite, Interplay between finite topological rank minimal Cantor systems, $S$-adic subshifts and their complexity, \emph{Trans.\ Amer.\ Math.\ Soc.} 374 (2021) 3453--3489.
\bibitem{GMachi} R.~I.~Grigorchuk, A.~Mach\`i, An example of an indexed language of intermediate growth, \emph{Theoret.\ Comput.\ Sci.} 215 (1999) 325--327.
\bibitem{BM10} A.~Ya.~Belov, R.~Mikhailov, Free subalgebras of Lie algebras close to nilpotent, \emph{Groups Geom.\ Dyn.} 4 (2010) 15--29; the growth dichotomy for automata algebras is Ufnarovskii, \emph{Mat.\ Zametki} 31 (1982) 465--472.
\bibitem{GMZ} B.~Greenfeld, C.~G.~Moreira, E.~Zelmanov, On the complexity of subshifts and infinite words, arXiv:2408.03403 (2024); and arXiv:2508.16840.
\bibitem{Mitchell} A.~Mitchell, On word complexity and topological entropy of random substitution subshifts, arXiv:2305.04817 (2023).
\bibitem{Pansiot} J.-J.~Pansiot, Complexit\'e des facteurs des mots infinis engendr\'es par morphismes it\'er\'es, \emph{ICALP 1984}, LNCS 172, 380--389.
\bibitem{Jen86} E.~Jen, Global properties of cellular automata, \emph{J.~Stat.~Phys.} 43 (1986) 219--242.
\bibitem{Jen90} E.~Jen, Aperiodicity in one-dimensional cellular automata, \emph{Physica D} 45 (1990) 3--18.
\bibitem{Hedlund} G.~A.~Hedlund, Endomorphisms and automorphisms of the shift dynamical system, \emph{Math.\ Systems Theory} 3 (1969) 320--375.
\bibitem{PY} M.~Pivato, R.~Yassawi, Asymptotic randomization of sofic shifts by linear cellular automata, \emph{Ergodic Theory Dynam.\ Systems} 26 (2006) 1177--1201.
\bibitem{PivatoSurvey} M.~Pivato, The ergodic theory of cellular automata, survey.
\bibitem{BK} M.~Boyle, B.~Kitchens, Periodic points for onto cellular automata, \emph{Indag.\ Math.} 10 (1999) 483--493.
\bibitem{BM} D.~H.~Bailey, M.~Misiurewicz, A strong hot spot theorem, \emph{Proc.\ Amer.\ Math.\ Soc.} 134 (2006) 2495--2501.
\bibitem{Kopra2020} J.~Kopra, Lyapunov exponents of multiplication automata, arXiv:2001.09675 (2020), Thm.~5.5.
\bibitem{CFK} V.~Cyr, J.~Franks, B.~Kra, The spacetime of a shift endomorphism, \emph{Trans.\ Amer.\ Math.\ Soc.} 371 (2019) 461--488, Question~5.7.
\bibitem{CPY} E.~Coven, M.~Pivato, R.~Yassawi, Prevalence of odometers in cellular automata, \emph{Proc.\ Amer.\ Math.\ Soc.} 135 (2007) 815--821, Thm.~4.
\bibitem{GL23} J.~Gravner, X.~Liu, Weak robust periodic solutions in one-dimensional cellular automata, \emph{J.~Cell.\ Autom.} 17 (2023) 251--280.
\bibitem{LN} W.~Li, M.~G.~Nordahl, Transient behavior of cellular automaton rule 110, \emph{Phys.\ Lett.\ A} 166 (1992) 335--339.
\bibitem{Kawa} H.~Kawaharada, Singular functions arising from cellular automata, arXiv:2008.13217 and arXiv:2203.05236.
\bibitem{Kopra2} J.~Kopra, On the trace subshifts of fractional multiplication automata, \emph{Theoret.\ Comput.\ Sci.} 851 (2021) 92--110; arXiv:2005.05112 (2020), Prop.~2.7 (numbering of arXiv v1).
\bibitem{BGrav} J.~M.~Baetens, J.~Gravner, Stability of cellular automata trajectories revisited: branching walks and Lyapunov profiles, \emph{J.~Nonlinear Sci.} 26 (2016) 1329--1367.
\bibitem{PivatoDefect} M.~Pivato, Defect particle kinematics in one-dimensional cellular automata, \emph{Theoret.\ Comput.\ Sci.} 377 (2007) 205--228.
\bibitem{Sutner} K.~Sutner, Almost periodic configurations on linear cellular automata, \emph{Fund.\ Inform.} 58 (2003) 223--240, Thm.~3.2.
\bibitem{GG22} J.~Gravner, D.~Griffeath, The one-dimensional Exactly~1 cellular automaton: replication, periodicity, and chaos from finite seeds, \emph{J.~Stat.\ Phys.} 142 (2011) 168--200.
\bibitem{GGBox} J.~Gravner, D.~Griffeath, Asymptotic densities for Packard Box rules, \emph{Nonlinearity} 22 (2009) 1817--1846.
\bibitem{BT} D.~A.~Barrington, D.~Th\'erien, Finite monoids and the fine structure of $\mathsf{NC}^1$, \emph{J.~ACM} 35 (1988) 941--952.
\bibitem{Bandt} C.~Bandt, Elementary fractal geometry 4: automata-generated topological spaces, \emph{Comm.\ Math.} 33 (2025), arXiv:2312.01486.
\bibitem{Green87} F.~Green, NP-complete problems in cellular automata, \emph{Complex Systems} 1 (1987) 453--474.
\bibitem{Cook} M.~Cook, Universality in elementary cellular automata, \emph{Complex Systems} 15 (2004) 1--40.
\bibitem{NW} T.~Neary, D.~Woods, P-completeness of cellular automaton Rule 110, \emph{ICALP 2006}, LNCS 4051, 132--143.
\bibitem{SS12} M.~Sch\"ule, R.~Stoop, A full computation-relevant topological dynamics classification of elementary cellular automata, \emph{Chaos} 22 (2012) 043143.
\bibitem{Kohlberg} E.~Kohlberg, Invariant half-lines of nonexpansive piecewise-linear transformations, \emph{Math.\ Oper.\ Res.} 5 (1980) 366--372.
\bibitem{KariSurvey} J.~Kari, Theory of cellular automata: a survey, \emph{Theoret.\ Comput.\ Sci.} 334 (2005) 3--33, Open Problem~5.
\bibitem{Meunier} P.-\'E.~Meunier, Unraveling simplicity in elementary cellular automata, arXiv:1406.5306 (2014).
\bibitem{NA} J.~Natal, O.~Al-saadi, Fast simulation of cellular automata by self-composition, arXiv:2409.07065 (2025).
\bibitem{Alman} J.~Alman, Kronecker products, low-depth circuits, and matrix rigidity, \emph{STOC 2021}, 772--785.
\bibitem{Kurka} P.~K\r{u}rka, Languages, equicontinuity and attractors in cellular automata, \emph{Ergodic Theory Dynam.\ Systems} 17 (1997) 417--433; and \emph{Topological and Symbolic Dynamics}, Soc.\ Math.\ France, 2003.
\bibitem{Martinez} G.~J.~Mart\'inez, A note on elementary cellular automata classification, arXiv:1306.5577 (2013), \S4.3; after A.~Wuensche, M.~Lesser, \emph{The Global Dynamics of Cellular Automata}, Addison-Wesley, 1992.
\bibitem{CKMFR} G.~Christol, T.~Kamae, M.~Mend\`es France, G.~Rauzy, Suites alg\'ebriques, automates et substitutions, \emph{Bull.\ Soc.\ Math.\ France} 108 (1980) 401--419.
\bibitem{BCOQ} F.~Baccelli, G.~Cohen, G.~J.~Olsder, J.-P.~Quadrat, \emph{Synchronization and Linearity}, Wiley, 1992.
\bibitem{BGT} V.~D.~Blondel, S.~Gaubert, J.~N.~Tsitsiklis, Approximating the spectral radius of sets of matrices in the max-algebra is NP-hard, \emph{IEEE Trans.\ Automat.\ Control} 45 (2000) 1762--1765.
\bibitem{NT} K.~Nishinari, D.~Takahashi, Analytical properties of ultradiscrete Burgers equation and rule-184 cellular automaton, \emph{J.~Phys.~A} 31 (1998) 5439--5450.
\bibitem{dBE} N.~G.~de Bruijn, P.~Erd\H{o}s, Some linear and some quadratic recursion formulas~II, \emph{Indag.\ Math.} 14 (1952) 152--163.
\bibitem{FR} Z.~F\"uredi, I.~Z.~Ruzsa, Nearly subadditive sequences, \emph{Acta Math.\ Hungar.} 161 (2020) 401--411.
\bibitem{Anashin} V.~Anashin, \emph{Applied Algebraic Dynamics}, de Gruyter Expositions in Mathematics 49, Berlin, 2009.
\bibitem{GGZ2} J.~Gravner, D.~Griffeath, Cellular automaton growth on $\Z^2$: theorems, examples, and problems, \emph{Adv.\ Appl.\ Math.} 21 (1998) 241 ff.
\bibitem{EN} K.~Eloranta, E.~Nummelin, The kink of cellular automaton Rule~18 performs a random walk, \emph{J.~Stat.\ Phys.} 69 (1992) 1131--1136.
\bibitem{RC} A.~Rupe, J.~P.~Crutchfield, Spacetime symmetries, invariant sets, and additive subdynamics of cellular automata, arXiv:1812.11597 (2018).
\bibitem{Fuks25} H.~Fuk\'s, Solving the initial value problem for cellular automata by pattern decomposition, arXiv:2512.24337 (2025), \S2.
\bibitem{Fuks23} H.~Fuk\'s, \emph{Solvable Cellular Automata: Methods and Applications}, Springer, 2023, Ch.~12.
\bibitem{Fuks05} H.~Fuk\'s, Second order additive invariants in elementary cellular automata, arXiv:nlin/0502037 (2005).
\bibitem{HT} T.~Hattori, S.~Takesue, Additive conserved quantities in discrete-time lattice dynamical systems, \emph{Physica D} 49 (1991) 295--322.
\bibitem{BFuks} N.~Boccara, H.~Fuk\'s, Number-conserving cellular automaton rules, \emph{Fund.\ Inform.} 52 (2002) 1--13.
\bibitem{BKP} B.~Bu\v{c}a, K.~Klobas, T.~Prosen, Rule~54: exactly solvable model of nonequilibrium statistical mechanics, \emph{J.~Stat.\ Mech.} (2021) 074001.
\bibitem{WPG} J.~W.~P.~Wilkinson, T.~Prosen, J.~P.~Garrahan, Exact solution of the Rule~150 reversible cellular automaton, arXiv:2110.15085 (2021).
\bibitem{HubnerKim} F.~H\"ubner, S.~W.~P.~Kim, Influence-solvability: a systematic theory of $(1+1)$D solvability and its application to brickwork circuits, arXiv:2606.12538 (2026).
\bibitem{Kleveland} R.~Kleveland, Mixing properties of one-dimensional cellular automata, \emph{Proc.\ Amer.\ Math.\ Soc.} 125 (1997) 1755--1766.
\bibitem{CERS} J.~Cassaigne, B.~Espinoza, M.~Rigo, M.~Stipulanti, Symbols frequencies in the Thue--Morse word in base $3/2$ and related conjectures, arXiv:2602.21895 (2026), Thm.~18.
\bibitem{MOW} O.~Martin, A.~M.~Odlyzko, S.~Wolfram, Algebraic properties of cellular automata, \emph{Comm.\ Math.\ Phys.} 93 (1984) 219--258.
\bibitem{Kopra} J.~Kopra, Rapid left expansivity, a commonality between Wolfram's Rule~30 and powers of $p/q$, \emph{Theoret.\ Comput.\ Sci.} 946 (2023) 113668; arXiv:2202.13809 (2022), under the title \emph{A natural class of cellular automata containing fractional multiplication automata, Rule~30, and others}. Def.~2.4, Thms.~3.5 and 4.7 and Problem~4.8 are cited in the numbering of arXiv v1.
\bibitem{Rowland} E.~S.~Rowland, Local nested structure in Rule 30, \emph{Complex Systems} 16 (2006) 239--258.
\bibitem{GG} J.~Gravner, D.~Griffeath, Robust periodic solutions and evolution from seeds in one-dimensional edge cellular automata, \emph{Theoret.\ Comput.\ Sci.} 466 (2012) 64--86.
\bibitem{Wolfram} S.~Wolfram, Announcing the Rule 30 Prizes, \emph{Stephen Wolfram Writings}, 1 October 2019, \url{https://writings.stephenwolfram.com/2019/10/announcing-the-rule-30-prizes/}; the three problems are also posted at \url{https://rule30prize.org/}.
\bibitem{Eilenberg} S.~Eilenberg, \emph{Automata, Languages and Machines}, Vol.~A, Academic Press, 1974.
\bibitem{AY} B.~Adamczewski, R.~Yassawi, A note on Christol's theorem, arXiv:1906.08703 (2019).
\bibitem{Shallit4} J.~Shallit, Automaticity IV: sequences, sets, and diversity, \emph{J.~Th\'eor.\ Nombres Bordeaux} 8 (1996) 347--367.
\bibitem{Nesme} V.~Nesme, arXiv:2207.13062.
\bibitem{AS} J.-P.~Allouche, J.~Shallit, \emph{Automatic Sequences}, Cambridge University Press, 2003.
\bibitem{KrawczykMullner} E.~Krawczyk, C.~M\"ullner, Automaticity of uniformly recurrent substitutive sequences, arXiv:2111.13134 (2021), Thm.~B (numbering of arXiv v3).
\bibitem{Bridy} A.~Bridy, Automatic sequences and curves over finite fields, \emph{Algebra Number Theory} 11 (2017) 685--712.
\bibitem{LD} B.~Litow, P.~Dumas, Additive cellular automata and algebraic series, \emph{Theoret.\ Comput.\ Sci.} 119 (1993) 345--354.
\bibitem{AvHPS} J.-P.~Allouche, F.~von Haeseler, H.-O.~Peitgen, G.~Skordev, Linear cellular automata, finite automata and Pascal's triangle, \emph{Discrete Appl.\ Math.} 66 (1996) 1--22.
\bibitem{BF} V.~Belitsky, P.~A.~Ferrari, Ballistic annihilation and deterministic surface growth, \emph{J.~Stat.~Phys.} 80 (1995) 517--543.
\bibitem{KS} J.~Krug, H.~Spohn, Universality classes for deterministic surface growth, \emph{Phys.\ Rev.\ A} 38 (1988) 4271--4283.
\bibitem{Q} N.~Berzai, The $0^m1$ family and the quench lemma, working document, 2026 (\texttt{quench\_writeup.md}).
\bibitem{Auto} N.~Berzai (with Claude), The Rule 30 centre column: an exact Duhamel formula and certified bounds on automaticity, working document \texttt{rule30\_automaticity.pdf}, 2026-09-16.
\bibitem{MS} W.~Meier, O.~Staffelbach, Analysis of pseudo random sequences generated by cellular automata, \emph{Advances in Cryptology --- EUROCRYPT '91}, LNCS 547, Springer, 1991, 186--199.
\bibitem{KA} C.~K.~Ko\c{c}, A.~M.~Apohan, Inversion of cellular automata iterations, \emph{IEE Proc.\ Comput.\ Digit.\ Tech.} 144 (1997) 279--284.
\bibitem{Spencer} J.~Spencer, Cellular automata as cryptographic random number generators, arXiv:1306.3546 (2013).
\bibitem{Leon} Le\'on et al., cryptanalysis of cellular-automaton keystream generators, 2003.
\bibitem{W83} S.~Wolfram, Statistical mechanics of cellular automata, \emph{Rev.\ Mod.\ Phys.} 55 (1983) 601--644.
\bibitem{W85} S.~Wolfram, Cryptography with cellular automata (extended abstract), \emph{Advances in Cryptology --- CRYPTO '85}, LNCS 218, Springer, 1986, 429--432.
\bibitem{NKS} S.~Wolfram, \emph{A New Kind of Science}, Wolfram Media, 2002.
\bibitem{Mariot} L.~Mariot, Insights from a decade of cellular automata in cryptography, arXiv:2405.02875 (2024).
\bibitem{CLMR} E.~Chan-L\'opez, A.~Mart\'in-Ruiz, Symmetric nonlinear cellular automata as algebraic references for Rule 30, arXiv:2604.00165 (2026).
\bibitem{Condrey} D.~L.~Condrey, Finite configurations cannot generate a constant trace in Rule 30, arXiv:2609.09431 (2026).
\bibitem{ASY} J.-P.~Allouche, J.~Shallit, R.~Yassawi, How to prove that a sequence is not automatic, \emph{Expo.\ Math.} 40 (2022) 1--22.
\bibitem{RY} E.~Rowland, R.~Yassawi, A characterization of $p$-automatic sequences as columns of linear cellular automata, \emph{Adv.\ Appl.\ Math.} 63 (2015) 68--89.
\bibitem{Rueppel} R.~A.~Rueppel, \emph{Analysis and Design of Stream Ciphers}, Springer, 1986.
\bibitem{Kawamura} K.~Kawamura, Polynomial representations of the Cuntz algebras arising from permutations, arXiv:math/0508273 (2005).
\bibitem{BJ} O.~Bratteli, P.~E.~T.~Jorgensen, Iterated function systems and permutation representations of the Cuntz algebra, \emph{Mem.\ Amer.\ Math.\ Soc.} 139 (1999), no.~663.
\bibitem{Fekete} M.~Fekete, \"Uber die Verteilung der Wurzeln bei gewissen algebraischen Gleichungen mit ganzzahligen Koeffizienten, \emph{Math.\ Z.} 17 (1923) 228--249.
\bibitem{Meinardus} G.~Meinardus, Asymptotische Aussagen \"uber Partitionen, \emph{Math.\ Z.} 59 (1954) 388--398.
\bibitem{IW} R.~Impagliazzo, A.~Wigderson, $\mathrm{P}=\mathrm{BPP}$ if $\mathrm{E}$ requires exponential circuits: derandomizing the XOR lemma, \emph{STOC 1997}, 220--229.
\bibitem{Walnut} H.~Mousavi, Automatic theorem proving in Walnut, arXiv:1603.06017 (2016).
\bibitem{Moore98} C.~Moore, Predicting nonlinear cellular automata quickly by decomposing them into linear ones, \emph{Physica D} 111 (1998) 27--41; arXiv:patt-sol/9701008 (1997), by C.~Moore and T.~Pnin.
\bibitem{FMG} M.~Feder, N.~Merhav, M.~Gutman, Universal prediction of individual sequences, \emph{IEEE Trans.\ Inform.\ Theory} 38 (1992) 1258--1270.
\bibitem{WN} N.~Berzai, The left edge of Rule~30 without Rule~30: the stripe recursion, Wolfram's number $2{,}107{,}985{,}255$, and the forks beyond it, manuscript, 2026; with its code and data at \url{https://github.com/nateberzai-svg/rule30-left-edge}.
\bibitem{Dib} Dibujaron (D.~McPolstra), \texttt{rule30}, public code repository, \url{https://github.com/Dibujaron/rule30}; commit \texttt{94833d21c1efc1e826df}\allowbreak\texttt{3c137777b580817c5e16}, 8 September 2026, file \texttt{blueprint/crystals.md}, item~55. Further commits of the same repository, by short hash: \texttt{0a1cda2} and \texttt{47f8ab5} (the settled configuration, 7--8 September 2026), \texttt{a8dbbc0} (the masking document, 8 September 2026), \texttt{59a9c26} (the period wall, 8 September 2026) and \texttt{e7fb239} (the shield, 9 September 2026); and its Lean lemmas.
\bibitem{Pat} Patto1155, \texttt{rule30-foundry}, public code repository, \url{https://github.com/Patto1155/rule30-foundry}; commit \texttt{04f30f6085210b476e64}\allowbreak\texttt{d742c77dc7c69a075791} (pull request~32), merged 8 September 2026, from a run of about 30 August 2026, files \texttt{docs/experiment-logs/2026-09-08-b1-recovery.md} and \texttt{runs/b1-recovered-2026-09-08/}\allowbreak\texttt{pattern\_map\_walk32.original.json}. Further commits of the same repository, by short hash: \texttt{d612567} (Lemma~A, 19 August 2026; on the main branch 28 August 2026), \texttt{08786d1} (the wedge; on the main branch 28 August 2026) and \texttt{de9c29d} (the period-word sieve, 8 September 2026).
\bibitem{Lania} szymon-lania, public gist with a starting state for the other version of left diagonal 53209 of Rule~30, 21 August 2026, \url{https://gist.github.com/szymon-lania/2535d62468a7912137b5002d900ebdd6}; described in its author's comment on E.~Rowland's video \emph{The Hidden Structure of Rule 30}, \url{https://www.youtube.com/watch?v=HcGjCcLEgN4}.
\bibitem{Cochon} cochon123, \texttt{rule30-prize}, public code repository, \url{https://github.com/cochon123/rule30-prize}; commits \texttt{0ff6930} (10 September 2026) and \texttt{adb0139} (12 September 2026).
\bibitem{WRyan} wryan2986, \texttt{rule30-lab}, public code repository, \url{https://github.com/wryan2986/rule30-lab}; commit \texttt{8960181}, 22 July 2026.
\bibitem{Woffinden} A.~Woffinden, Synchrony with nothing: the Rule~30 Prize Problems, the question beneath them, and a structural definition of randomness, Zenodo, 23 August 2026, \url{https://doi.org/10.5281/zenodo.22065659}.
\bibitem{Ikram} C.~A.~H.~Ikram, Paper~1 --- Non-periodicity problem: finite cones and period-one obstructions for the center-column non-periodicity problem in Rule~30, OSF Preprints, 15 May 2026, \url{https://osf.io/kybmz}.
\bibitem{Nersissian} T.~Nersissian, Diagonal periods and Newton supports of Rules 30, 86 and 135, arXiv:2609.25077 (18 September 2026); and Diagonal bases and diagonal periods of elementary cellular automata, arXiv:2609.25078 (2026).
\bibitem{Brunnbauer} M.~Brunnbauer, Diagonals in elementary cellular automaton 30, web page, 9 December 2019, \url{https://brunni.de/findings30/}.
\bibitem{Wolfram19} S.~Wolfram, code and data of the prize announcement \cite{Wolfram}, 1 October 2019: the code in the post, and the Wolfram Cloud table linked from it, with $78\,288$ diagonals ($d=0$ to $78\,287$).
\bibitem{A363346} P.~Xausa, sequence A363346 in \emph{The On-Line Encyclopedia of Integer Sequences}, 2023: the length of the initial transient on the $n$-th diagonal from the left of Rule~30 started from a single cell; $1\,000$ terms. \url{https://oeis.org/A363346}.
\end{thebibliography}

\appendix

\section{The certificate engine (independent implementation)}
\begin{footnotesize}
\begin{verbatim}
import numpy as np, sys
K=10; N=1<<K; idx=np.arange(N,dtype=np.int64); bits=[(idx>>i)&1 for i in range(K)]
def succ(c,f):
    left=[np.full(N,c,dtype=np.int64)]+bits+[np.full(N,f,dtype=np.int64)]
    out=np.zeros(N,dtype=np.int64)
    for i in range(1,K+1): out|=(left[i-1]^(left[i]|left[i+1]))<<(i-1)
    return out
S0=[succ(0,0),succ(0,1)]; S1=[succ(1,0),succ(1,1)]; u0=(bits[0]==0)
def step(mask,c):
    nxt=np.zeros(N,dtype=bool); tab=S0 if c==0 else S1
    src=np.nonzero(mask)[0]; nxt[tab[0][src]]=True; nxt[tab[1][src]]=True; return nxt
def key(m): return np.packbits(m).tobytes()
def graph(word):
    ALL=np.ones(N,dtype=bool); ids={key(ALL):0}; states=[ALL]; todo=[0]; edges=[]
    while todo:
        s=todo.pop(); fr={():states[s]}
        for ch in word:
            nx={}
            for lab,T in fr.items():
                if ch=='0':
                    for u,sel in ((0,u0),(1,~u0)):
                        Tb=T&sel
                        if Tb.any(): nx[lab+(u,)]=step(Tb,0)
                else: nx[lab]=step(T,1)
            fr=nx
        for lab,T in fr.items():
            if not T.any(): continue
            k=key(T)
            if k not in ids: ids[k]=len(states); states.append(T); todo.append(len(states)-1)
            edges.append((s,ids[k],lab))
    return states,edges
# Tarjan SCC on (states, edges); criterion: for every SCC with an edge,
# #edges (both ends inside) == #vertices.
\end{verbatim}
\end{footnotesize}



\end{document}
